% Generated by scripts/build.py - do not edit.

\documentclass[11pt]{article}

% Shared preamble: lets a single question file or the whole bank compile with XeLaTeX.
\usepackage{amsmath,amssymb,amsthm}
\usepackage{fontspec}
\usepackage{polyglossia}
\setmainlanguage{hebrew}
\setotherlanguage{english}
\newfontfamily\hebrewfont[Script=Hebrew]{David CLM}
\newenvironment{question}{\par\noindent\textbf{שאלה.}\ }{\par}
\newenvironment{hint}{\par\noindent\textbf{רמז.}\ }{\par}
\newenvironment{solution}{\par\noindent\textbf{פתרון.}\ }{\par}


\begin{document}

\title{מאגר שאלות בחינה}\date{}\maketitle

\section*{תשע"ו סמסטר ב מועד ב}

\subsection*{שאלה 1}

\begin{question}
\textbf{הוכיחו או הפריכו:} תהי $f$ פונקציה מונוטונית עולה ממש בקטע $[a,b]$ אז $f$ קמורה.
\end{question}

\begin{hint}
חפשו פונקציה עולה ממש שהגרף שלה "מתעקם" לשני הכיוונים, למשל פולינום אי-זוגי.
\end{hint}

\begin{solution}
\textbf{הטענה אינה נכונה.} מונוטוניות אינה קשורה לקמירות.

דוגמה נגדית: $f(x)=x^3$ בקטע $[a,b]=[-1,1]$. מכיוון ש-$f'(x)=3x^2\ge 0$ ומתאפסת רק בנקודה בודדת $x=0$, הפונקציה עולה ממש בקטע (אפשר גם ישירות: אם $x<y$ אז $y^3-x^3=(y-x)(x^2+xy+y^2)>0$, כי $x^2+xy+y^2=(x+\tfrac y2)^2+\tfrac34y^2>0$ כאשר $x\ne y$).

אולם $f''(x)=6x$, כך ש-$f''<0$ ב-$(-1,0)$ ו-$f''>0$ ב-$(0,1)$. כלומר $f$ קעורה ב-$[-1,0]$ וקמורה ב-$[0,1]$, ולכן אינה קמורה (וגם אינה קעורה) בכל הקטע.

בדיקה ישירה לפי ההגדרה: עבור הנקודות $x_1=-1,\ x_2=0$ ונקודת האמצע $-\tfrac12$:
\[ f\!\left(-\tfrac12\right)=-\tfrac18 \;>\; -\tfrac12=\frac{f(-1)+f(0)}{2}, \]
כלומר המיתר בין $(-1,-1)$ ל-$(0,0)$ נמצא \emph{מתחת} לגרף, בסתירה לקמירות ($f(\tfrac{x_1+x_2}{2})\le \tfrac{f(x_1)+f(x_2)}{2}$). ומנגד, עבור $x_1=0,\ x_2=1$: $f(\tfrac12)=\tfrac18<\tfrac12=\tfrac{f(0)+f(1)}2$, כך שהפונקציה גם אינה קעורה. לכן הדוגמה מפריכה את הטענה בכל אחת מהמוסכמות למילה "קמורה".

(דוגמה נוספת: $f(x)=\sqrt{x}$ ב-$[0,1]$ עולה ממש וקעורה ממש.)
\end{solution}

\subsection*{שאלה 2}

\begin{question}
\textbf{הוכיחו או הפריכו:} הפונקציה הבאה גזירה באפס:
\[ f(x)=\begin{cases} x\sqrt{|x|}\cos(\frac{1}{x}) & x\neq 0\\ 0 & x=0\end{cases}. \]
\end{question}

\begin{hint}
חשבו את הנגזרת לפי ההגדרה, כגבול של מנת ההפרשים, והשתמשו בכך ש"חסומה כפול שואפת לאפס שואפת לאפס".
\end{hint}

\begin{solution}
\textbf{הטענה נכונה.} נחשב לפי הגדרת הנגזרת. עבור $h\neq 0$:
\[ \frac{f(h)-f(0)}{h}=\frac{h\sqrt{|h|}\cos(\frac1h)-0}{h}=\sqrt{|h|}\cos\!\left(\frac1h\right). \]
מתקיים $\left|\sqrt{|h|}\cos(\frac1h)\right|\le \sqrt{|h|}$, ו-$\sqrt{|h|}\to 0$ כאשר $h\to 0$. לפי כלל הסנדוויץ' (או: פונקציה חסומה, $\cos(\frac1h)$, כפול פונקציה השואפת לאפס, $\sqrt{|h|}$, שואפת לאפס):
\[ f'(0)=\lim_{h\to 0}\frac{f(h)-f(0)}{h}=\lim_{h\to0}\sqrt{|h|}\cos\!\left(\frac1h\right)=0. \]
הגבול קיים וסופי, ולכן $f$ גזירה באפס ו-$\boxed{f'(0)=0}$.
\end{solution}

\subsection*{שאלה 3}

\begin{question}
\textbf{הוכיחו או הפריכו:} האינטגרל $\int_2^\infty \frac{dx}{x\ln(x)}dx$ מתכנס.
\end{question}

\begin{hint}
מצאו פונקציה קדומה בעזרת ההצבה $t=\ln x$.
\end{hint}

\begin{solution}
\textbf{הטענה אינה נכונה} -- האינטגרל מתבדר.

הפונקציה $\frac{1}{x\ln x}$ רציפה וחיובית ב-$[2,\infty)$, ולכן האינטגרל הלא אמיתי מוגדר כ-$\lim_{R\to\infty}\int_2^R\frac{dx}{x\ln x}$.
בהצבה $t=\ln x$, $dt=\frac{dx}{x}$:
\[ \int\frac{dx}{x\ln x}=\int\frac{dt}{t}=\ln|t|+C=\ln(\ln x)+C \qquad (x>1). \]
לפי המשפט היסודי (נוסחת ניוטון-לייבניץ):
\[ \int_2^R\frac{dx}{x\ln x}=\ln(\ln R)-\ln(\ln 2)\xrightarrow[R\to\infty]{}\infty, \]
כי $\ln R\to\infty$ ולכן גם $\ln(\ln R)\to\infty$. הגבול אינו סופי, ולכן האינטגרל \textbf{מתבדר}.
\end{solution}

\subsection*{שאלה 4א}

\begin{question}
חשבו את הגבול הבא: $\lim_{x\to 0}\frac{1-\cos(x^2)}{x^2\sin^2(x)}$.
\end{question}

\begin{hint}
השתמשו בפיתוח טיילור של $\cos t$ סביב $0$ עם $t=x^2$, ובגבול היסודי $\frac{\sin x}{x}\to1$.
\end{hint}

\begin{solution}
נכתוב
\[ \frac{1-\cos(x^2)}{x^2\sin^2 x}=\frac{1-\cos(x^2)}{x^4}\cdot\left(\frac{x}{\sin x}\right)^2. \]
לפי פיתוח מקלורין $\cos t=1-\frac{t^2}{2}+o(t^2)$, ועם $t=x^2$: $1-\cos(x^2)=\frac{x^4}{2}+o(x^4)$, ולכן $\frac{1-\cos(x^2)}{x^4}\to\frac12$.
(לחלופין: $1-\cos(x^2)=2\sin^2(\frac{x^2}{2})$, ולכן $\frac{1-\cos(x^2)}{x^4}=\frac12\left(\frac{\sin(x^2/2)}{x^2/2}\right)^2\to\frac12$.)
כמו כן $\frac{x}{\sin x}\to 1$ לפי הגבול היסודי. לפי אריתמטיקה של גבולות:
\[ \lim_{x\to0}\frac{1-\cos(x^2)}{x^2\sin^2 x}=\frac12\cdot 1=\boxed{\frac12}. \]
\end{solution}

\subsection*{שאלה 4ב}

\begin{question}
הראו ש:
\[ \frac15\le \arctan(3)-\arctan(1)\le 1. \]
\end{question}

\begin{hint}
הפעילו את משפט לגרנז' על $\arctan$ בקטע $[1,3]$ וחסמו את הנגזרת.
\end{hint}

\begin{solution}
נגדיר $g(t)=\arctan t$. הפונקציה רציפה ב-$[1,3]$ וגזירה ב-$(1,3)$ עם $g'(t)=\frac{1}{1+t^2}$. לפי משפט לגרנז' קיימת $c\in(1,3)$ כך ש-
\[ \arctan(3)-\arctan(1)=g'(c)(3-1)=\frac{2}{1+c^2}. \]
מכיוון ש-$1<c<3$ מתקיים $2<1+c^2<10$, ולכן
\[ \frac{2}{10}<\frac{2}{1+c^2}<\frac{2}{2}, \qquad\text{כלומר}\qquad \frac15<\arctan(3)-\arctan(1)<1, \]
ובפרט מתקיים אי-השוויון החלש הנדרש. (בדיקה מספרית: $\arctan 3-\arctan 1\approx 0.4636$.)
\end{solution}

\subsection*{שאלה 5}

\begin{question}
\begin{enumerate}
  \item (15 נק') חשבו את האינטגרל הלא מסויים הבא:
  \[ \int\frac{xdx}{(x-1)(x^2+2x+2)}. \]
  \item (5 נק') חשבו את השטח הכלוא מתחת לגרף הפונקציה $f(x)=\frac{x}{(x-1)(x^2+2x+2)}$ בקטע $[-1,0]$.
\end{enumerate}
\end{question}

\begin{hint}
פרקו לשברים חלקיים: $\frac{A}{x-1}+\frac{Bx+C}{x^2+2x+2}$. לגורם הריבועי השלימו לריבוע: $x^2+2x+2=(x+1)^2+1$.
\end{hint}

\begin{hint}
בסעיף ב' בדקו את סימן הפונקציה בקטע לפני שמחשבים את השטח.
\end{hint}

\begin{solution}
\begin{enumerate}
  \item לגורם $x^2+2x+2=(x+1)^2+1$ אין שורשים ממשיים (הדיסקרימיננטה $4-8<0$), ולכן נפרק:
  \[ \frac{x}{(x-1)(x^2+2x+2)}=\frac{A}{x-1}+\frac{Bx+C}{x^2+2x+2}
     \iff x=A(x^2+2x+2)+(Bx+C)(x-1). \]
  הצבת $x=1$: $1=5A$, כלומר $A=\frac15$. השוואת מקדם $x^2$: $0=A+B$, לכן $B=-\frac15$. השוואת המקדם החופשי: $0=2A-C$, לכן $C=\frac25$. כלומר
  \[ \frac{x}{(x-1)(x^2+2x+2)}=\frac{1}{5}\cdot\frac{1}{x-1}-\frac15\cdot\frac{x-2}{x^2+2x+2}. \]
  נכתוב $x-2=\frac12(2x+2)-3$:
  \[ \frac{x-2}{x^2+2x+2}=\frac12\cdot\frac{2x+2}{x^2+2x+2}-\frac{3}{(x+1)^2+1}. \]
  לכן (בעזרת $\int\frac{u'}{u}=\ln|u|$ ו-$\int\frac{dt}{t^2+1}=\arctan t$):
  \[ \boxed{\int\frac{x\,dx}{(x-1)(x^2+2x+2)}=\frac15\ln|x-1|-\frac1{10}\ln(x^2+2x+2)+\frac35\arctan(x+1)+C.} \]
  (בדיקה: גזירת הביטוי מחזירה את האינטגרנד.)

  \item בקטע $[-1,0]$: $x\le 0$, $x-1<0$, $x^2+2x+2>0$, ולכן $f(x)\ge0$ (שווה לאפס רק ב-$x=0$). הפונקציה רציפה בקטע, כך שהשטח הוא $\int_{-1}^0 f(x)\,dx$. נסמן ב-$F$ את הפונקציה הקדומה מסעיף א' ונשתמש בנוסחת ניוטון-לייבניץ:
  \[ F(0)=\frac15\ln1-\frac1{10}\ln2+\frac35\cdot\frac\pi4=-\frac{\ln2}{10}+\frac{3\pi}{20},\qquad
     F(-1)=\frac15\ln2-\frac1{10}\ln1+\frac35\arctan0=\frac{\ln 2}{5}. \]
  לכן
  \[ S=F(0)-F(-1)=\boxed{\frac{3\pi}{20}-\frac{3\ln2}{10}}\approx 0.2633. \]
\end{enumerate}
\end{solution}

\subsection*{שאלה 6א}

\begin{question}
מצאו את תחום ההגדרה, נקודות חיתוך עם הצירים, תחומי עליה וירידה, נקודות קיצון, אסימפטוטות ותחומי קמירות וקעירות של הפונקציה הבאה:
\[ f(x)=(x+2)e^{\frac1x}. \]
\end{question}

\begin{hint}
לאסימפטוטה המשופעת חשבו $\lim\frac{f(x)}{x}$ ואחר כך $\lim (f(x)-x)$; השתמשו בהצבה $t=\frac1x$ ובגבול $\frac{e^t-1}{t}\to1$.
\end{hint}

\begin{solution}
\textbf{תחום הגדרה:} $x\neq0$.

\textbf{חיתוך עם הצירים:} אין חיתוך עם ציר ה-$y$ (כי $0$ אינו בתחום). $f(x)=0\iff x+2=0$ (כי $e^{1/x}>0$), ולכן נקודת החיתוך עם ציר ה-$x$ היא $(-2,0)$. בנוסף $f<0$ עבור $x<-2$ ו-$f>0$ עבור $x>-2,\ x\ne0$.

\textbf{נגזרת ראשונה:}
\[ f'(x)=e^{1/x}+(x+2)e^{1/x}\cdot\left(-\frac1{x^2}\right)=e^{1/x}\,\frac{x^2-x-2}{x^2}=\frac{(x-2)(x+1)}{x^2}e^{1/x}. \]
הסימן נקבע ע"י $(x-2)(x+1)$:
\begin{itemize}
  \item $f'>0$ ב-$(-\infty,-1)$ וב-$(2,\infty)$ -- \textbf{עולה};
  \item $f'<0$ ב-$(-1,0)$ וב-$(0,2)$ -- \textbf{יורדת}.
\end{itemize}
\textbf{קיצון:} ב-$x=-1$ הנגזרת מחליפה סימן מ-$+$ ל-$-$: מקסימום מקומי $f(-1)=e^{-1}=\frac1e$. ב-$x=2$ מ-$-$ ל-$+$: מינימום מקומי $f(2)=4e^{1/2}=4\sqrt e$.

\textbf{אסימפטוטות אנכיות:} כאשר $x\to0^+$, $\frac1x\to+\infty$ ולכן $e^{1/x}\to\infty$ ו-$f(x)\to 2\cdot\infty=+\infty$: $x=0$ אסימפטוטה אנכית (מימין). כאשר $x\to0^-$, $\frac1x\to-\infty$, $e^{1/x}\to0$ ו-$f(x)\to 0$ (אין אסימפטוטה משמאל; לגרף "חור" בראשית).

\textbf{אסימפטוטות משופעות:} $\frac{f(x)}{x}=\left(1+\frac2x\right)e^{1/x}\to1$ כאשר $x\to\pm\infty$. כמו כן
\[ f(x)-x=x\left(e^{1/x}-1\right)+2e^{1/x}=\frac{e^{t}-1}{t}+2e^{t}\Big|_{t=1/x}\;\xrightarrow[t\to0]{}\;1+2=3. \]
לכן $y=x+3$ אסימפטוטה משופעת גם ב-$+\infty$ וגם ב-$-\infty$. (אין אסימפטוטה אופקית.)

\textbf{נגזרת שנייה:} גוזרים את $f'(x)=\left(1-\frac1x-\frac2{x^2}\right)e^{1/x}$:
\[ f''(x)=\left(\frac1{x^2}+\frac4{x^3}\right)e^{1/x}-\frac1{x^2}\left(1-\frac1x-\frac2{x^2}\right)e^{1/x}=\frac{5x+2}{x^4}e^{1/x}. \]
הסימן נקבע ע"י $5x+2$:
\begin{itemize}
  \item $f''<0$ ב-$(-\infty,-\frac25)$ -- \textbf{קעורה} ($\cap$);
  \item $f''>0$ ב-$(-\frac25,0)$ וב-$(0,\infty)$ -- \textbf{קמורה} ($\cup$).
\end{itemize}
נקודת פיתול: $x=-\frac25$, $f(-\frac25)=\frac85e^{-5/2}$.
\end{solution}

\subsection*{שאלה 6ב}

\begin{question}
חשבו את נפח גוף הסיבוב המתקבל מסיבוב גרף הפונקציה:
\[ f(x)=\sin^2(x) \]
סביב ציר ה-$x$ בקטע $[0,\frac\pi2]$.
\end{question}

\begin{hint}
נפח גוף הסיבוב הוא $\pi\int_a^b f^2(x)\,dx$; להורדת החזקה השתמשו ב-$\sin^2x=\frac{1-\cos2x}{2}$.
\end{hint}

\begin{solution}
נפח גוף סיבוב סביב ציר $x$: $V=\pi\int_0^{\pi/2}f^2(x)\,dx=\pi\int_0^{\pi/2}\sin^4x\,dx$. נשתמש בזהויות
\[ \sin^4x=\left(\frac{1-\cos2x}{2}\right)^2=\frac14\left(1-2\cos2x+\cos^22x\right)=\frac14\left(1-2\cos2x+\frac{1+\cos4x}{2}\right)=\frac38-\frac12\cos2x+\frac18\cos4x. \]
לכן
\[ \int_0^{\pi/2}\sin^4x\,dx=\left[\frac38x-\frac14\sin2x+\frac1{32}\sin4x\right]_0^{\pi/2}=\frac{3\pi}{16}, \]
(כי $\sin\pi=\sin2\pi=0$), ומכאן
\[ \boxed{V=\frac{3\pi^2}{16}}. \]
\end{solution}

\section*{תשע"ו סמסטר ב מועד ג}

\subsection*{שאלה 1}

\begin{question}
\textbf{הוכיחו או הפריכו:} למשוואה הבאה יש פתרון ב-$[0,\pi]$:
\[ \ln(x+1)-\sin(x)=1 \]
\end{question}

\begin{hint}
הגדירו $g(x)=\ln(x+1)-\sin(x)-1$ ובדקו את סימנה בקצוות הקטע.
\end{hint}

\begin{solution}
\textbf{הטענה נכונה.} נגדיר $g(x)=\ln(x+1)-\sin x-1$. הפונקציה $g$ רציפה ב-$[0,\pi]$ כסכום של פונקציות רציפות ($\ln(x+1)$ רציפה עבור $x>-1$).

בקצוות:
\[ g(0)=\ln1-\sin0-1=-1<0,\qquad g(\pi)=\ln(\pi+1)-\sin\pi-1=\ln(\pi+1)-1>0, \]
כאשר האי-שוויון האחרון נובע מכך ש-$\pi+1>e$ (כי $\pi+1>4>e$), ולכן $\ln(\pi+1)>\ln e=1$. (מספרית $\ln(\pi+1)\approx1.421$.)

לפי משפט ערך הביניים (משפט בולצאנו) קיימת $c\in(0,\pi)$ שעבורה $g(c)=0$, כלומר $\ln(c+1)-\sin c=1$. לכן למשוואה יש פתרון בקטע.
\end{solution}

\subsection*{שאלה 2}

\begin{question}
\textbf{הוכיחו או הפריכו:} הפונקציה הבאה רציפה באפס:
\[ f(x)=\begin{cases} \frac{\sin(2x)}{\sqrt{x+4}-2} & x\neq 0\\ 5 & x=0\end{cases}. \]
\end{question}

\begin{hint}
הכפילו את המונה והמכנה בצמוד $\sqrt{x+4}+2$.
\end{hint}

\begin{solution}
\textbf{הטענה אינה נכונה.} $f$ רציפה ב-$0$ אם ורק אם $\lim_{x\to0}f(x)=f(0)=5$. נחשב את הגבול (עבור $x\ne0$, $x\ge-4$) בעזרת כפל בצמוד:
\[ \frac{\sin 2x}{\sqrt{x+4}-2}=\frac{\sin 2x\,(\sqrt{x+4}+2)}{(x+4)-4}=\frac{\sin2x}{x}\left(\sqrt{x+4}+2\right)=2\cdot\frac{\sin 2x}{2x}\cdot\left(\sqrt{x+4}+2\right). \]
לפי הגבול היסודי $\frac{\sin2x}{2x}\to1$, ובגלל רציפות השורש $\sqrt{x+4}+2\to4$. לפי אריתמטיקה של גבולות
\[ \lim_{x\to0}f(x)=2\cdot1\cdot4=8\neq5=f(0). \]
לכן $f$ \textbf{אינה רציפה} באפס (יש לה אי-רציפות סליקה; היא הייתה רציפה אילו הגדרנו $f(0)=8$).
\end{solution}

\subsection*{שאלה 3}

\begin{question}
\textbf{הוכיחו או הפריכו:} האינטגרל $\int_0^1\ln^2(x)dx$ מתכנס.
\end{question}

\begin{hint}
מצאו פונקציה קדומה בעזרת אינטגרציה בחלקים (פעמיים), וחשבו את הגבול $\lim_{x\to0^+}x\ln^k x$ בעזרת לופיטל.
\end{hint}

\begin{solution}
\textbf{הטענה נכונה}, וערך האינטגרל הוא $2$.

הפונקציה $\ln^2x$ רציפה ב-$(0,1]$ ואינה חסומה ליד $0$, ולכן זהו אינטגרל לא אמיתי: $\int_0^1\ln^2x\,dx=\lim_{\eps\to0^+}\int_\eps^1\ln^2x\,dx$.

\textbf{פונקציה קדומה} -- אינטגרציה בחלקים עם $u=\ln^2x,\ v'=1$:
\[ \int\ln^2x\,dx=x\ln^2x-\int x\cdot\frac{2\ln x}{x}dx=x\ln^2x-2\int\ln x\,dx=x\ln^2x-2(x\ln x-x)+C. \]
לכן
\[ \int_\eps^1\ln^2x\,dx=\big[x\ln^2x-2x\ln x+2x\big]_\eps^1=2-\left(\eps\ln^2\eps-2\eps\ln\eps+2\eps\right). \]
\textbf{הגבולות:} לפי כלל לופיטל ($\frac{\infty}{\infty}$):
\[ \lim_{\eps\to0^+}\eps\ln\eps=\lim\frac{\ln\eps}{1/\eps}=\lim\frac{1/\eps}{-1/\eps^2}=\lim(-\eps)=0, \]
\[ \lim_{\eps\to0^+}\eps\ln^2\eps=\lim\frac{\ln^2\eps}{1/\eps}=\lim\frac{2\ln\eps/\eps}{-1/\eps^2}=\lim(-2\eps\ln\eps)=0. \]
לכן $\int_0^1\ln^2x\,dx=2-0=2$: האינטגרל \textbf{מתכנס}.

(דרך אחרת: $\ln^2x\le\frac{C}{\sqrt x}$ ליד $0$ כי $\sqrt x\ln^2x\to0$, ו-$\int_0^1\frac{dx}{\sqrt x}$ מתכנס, ולכן לפי מבחן ההשוואה האינטגרל מתכנס.)
\end{solution}

\subsection*{שאלה 4א}

\begin{question}
חשבו את הגבול הבא: $\lim_{x\to0}\frac{\cos(\sin(x))-\cos(x)}{x^4}$.
\end{question}

\begin{hint}
פתחו את $\cos t$ ואת $\sin x$ לפי מקלורין עד סדר $4$ (שימו לב ש-$\sin^2x=x^2-\frac{x^4}{3}+o(x^4)$).
\end{hint}

\begin{solution}
נשתמש בפיתוחי מקלורין:
$\sin x=x-\frac{x^3}{6}+o(x^4)$ ו-$\cos t=1-\frac{t^2}{2}+\frac{t^4}{24}+o(t^4)$.
נחשב:
\[ \sin^2x=\left(x-\frac{x^3}{6}+o(x^4)\right)^2=x^2-\frac{x^4}{3}+o(x^4),\qquad \sin^4x=x^4+o(x^4). \]
מכיוון ש-$\sin x\to0$ ו-$|\sin x|\le|x|$, ניתן להציב $t=\sin x$ (שארית $o(\sin^4x)$ היא $o(x^4)$):
\[ \cos(\sin x)=1-\frac{\sin^2x}{2}+\frac{\sin^4x}{24}+o(x^4)=1-\frac{x^2}{2}+\frac{x^4}{6}+\frac{x^4}{24}+o(x^4)=1-\frac{x^2}{2}+\frac{5x^4}{24}+o(x^4). \]
וכן $\cos x=1-\frac{x^2}{2}+\frac{x^4}{24}+o(x^4)$. לכן
\[ \cos(\sin x)-\cos x=\frac{4x^4}{24}+o(x^4)=\frac{x^4}{6}+o(x^4), \]
ומכאן
\[ \lim_{x\to0}\frac{\cos(\sin x)-\cos x}{x^4}=\boxed{\frac16}. \]
(דרך נוספת: $\cos A-\cos B=-2\sin\frac{A+B}{2}\sin\frac{A-B}{2}$ עם $A=\sin x$, $B=x$; אז $\frac{A+B}{2}\sim x$ ו-$\frac{A-B}{2}=\frac{\sin x-x}{2}\sim-\frac{x^3}{12}$, ומתקבל $-2\cdot x\cdot(-\frac{x^3}{12})=\frac{x^4}{6}$.)
\end{solution}

\subsection*{שאלה 4ב}

\begin{question}
הראו שלכל $0<a<b$ ולכל $p>1$ מתקיים:
\[ pa^{p-1}(b-a)\le b^p-a^p\le pb^{p-1}(b-a) \]
\end{question}

\begin{hint}
הפעילו את משפט לגרנז' על $f(x)=x^p$ בקטע $[a,b]$, והשתמשו במונוטוניות של $x^{p-1}$.
\end{hint}

\begin{solution}
תהי $f(x)=x^p$. היא רציפה ב-$[a,b]$ וגזירה ב-$(a,b)$ (כי $a>0$), עם $f'(x)=px^{p-1}$. לפי משפט לגרנז' קיימת $c\in(a,b)$ כך ש-
\[ b^p-a^p=f'(c)(b-a)=pc^{p-1}(b-a). \]
מכיוון ש-$p>1$, המעריך $p-1>0$ ולכן הפונקציה $x\mapsto x^{p-1}$ עולה ממש ב-$(0,\infty)$. מ-$0<a<c<b$ נובע $a^{p-1}<c^{p-1}<b^{p-1}$. נכפול ב-$p(b-a)>0$:
\[ pa^{p-1}(b-a)<pc^{p-1}(b-a)=b^p-a^p<pb^{p-1}(b-a), \]
ובפרט מתקיים אי-השוויון הנדרש (אף בחוזק).
\end{solution}

\subsection*{שאלה 5}

\begin{question}
\begin{enumerate}
  \item (15 נק') חשבו את האינטגרל הלא מסויים הבא:
  \[ \int\frac{dx}{\cos^3(x)}. \]
  \item (5 נק') חשבו את השטח הכלוא מתחת לגרף הפונקציה $f(x)=\frac{1}{\cos^3(x)}$ בקטע $[0,1]$.
\end{enumerate}
\end{question}

\begin{hint}
כתבו $\frac{1}{\cos^3x}=\frac{1}{\cos x}\cdot\frac1{\cos^2x}$ ובצעו אינטגרציה בחלקים עם $v'=\frac1{\cos^2x}$, $v=\tan x$.
\end{hint}

\begin{hint}
תזדקקו גם ל-$\int\frac{dx}{\cos x}=\ln\left|\frac1{\cos x}+\tan x\right|+C$ (או: הכפילו ב-$\cos x$ והציבו $t=\sin x$).
\end{hint}

\begin{solution}
\begin{enumerate}
  \item נסמן $I=\int\frac{dx}{\cos^3x}$ ונבצע אינטגרציה בחלקים עם $u=\frac1{\cos x}$, $v'=\frac1{\cos^2x}$; אז $u'=\frac{\sin x}{\cos^2x}$, $v=\tan x$:
  \[ I=\frac{\tan x}{\cos x}-\int\frac{\sin x}{\cos^2x}\cdot\frac{\sin x}{\cos x}dx=\frac{\tan x}{\cos x}-\int\frac{1-\cos^2x}{\cos^3x}dx=\frac{\tan x}{\cos x}-I+\int\frac{dx}{\cos x}. \]
  נחשב את $\int\frac{dx}{\cos x}=\int\frac{\cos x\,dx}{1-\sin^2x}$; בהצבה $t=\sin x$:
  \[ \int\frac{dt}{1-t^2}=\frac12\ln\left|\frac{1+t}{1-t}\right|+C=\frac12\ln\left|\frac{1+\sin x}{1-\sin x}\right|+C=\ln\left|\frac{1+\sin x}{\cos x}\right|+C, \]
  (כי $\frac{1+\sin x}{1-\sin x}=\frac{(1+\sin x)^2}{\cos^2x}$). לכן $2I=\frac{\sin x}{\cos^2x}+\ln\left|\frac{1+\sin x}{\cos x}\right|+C$, כלומר
  \[ \boxed{\int\frac{dx}{\cos^3x}=\frac{\sin x}{2\cos^2x}+\frac12\ln\left|\frac{1}{\cos x}+\tan x\right|+C.} \]
  (בדיקה: גזירת התוצאה מחזירה את $\frac1{\cos^3x}$.)

  \item ב-$[0,1]$ מתקיים $0\le x\le1<\frac\pi2$, ולכן $\cos x>0$ ו-$f(x)>0$ רציפה. השטח הוא
  \[ S=\int_0^1\frac{dx}{\cos^3x}=\left[\frac{\sin x}{2\cos^2x}+\frac12\ln\left(\frac{1+\sin x}{\cos x}\right)\right]_0^1
      =\boxed{\frac{\sin1}{2\cos^21}+\frac12\ln\left(\frac{1+\sin1}{\cos1}\right)}\approx2.054, \]
  כי בנקודה $0$ הביטוי שווה ל-$0+\frac12\ln1=0$.
\end{enumerate}
\end{solution}

\subsection*{שאלה 6א}

\begin{question}
שרטטו את התחום הכלוא בין העקומים $x+y=0$ ו-$y=2x-x^2$ ומצאו את שטחו.
\end{question}

\begin{hint}
מצאו את נקודות החיתוך של הישר $y=-x$ עם הפרבולה, וקבעו איזו עקומה נמצאת מעל השנייה.
\end{hint}

\begin{solution}
העקום $x+y=0$ הוא הישר $y=-x$, והעקום $y=2x-x^2$ הוא פרבולה "בוכה" (פתוחה כלפי מטה) עם קודקוד $(1,1)$ החותכת את ציר $x$ ב-$0$ וב-$2$.
\textbf{נקודות חיתוך:} $-x=2x-x^2\iff x^2-3x=0\iff x=0$ או $x=3$, כלומר הנקודות $(0,0)$ ו-$(3,-3)$.
בקטע $(0,3)$: $(2x-x^2)-(-x)=3x-x^2=x(3-x)>0$, כלומר הפרבולה מעל הישר. התחום הוא הקטע שבין הישר (מלמטה) לבין קשת הפרבולה (מלמעלה) עבור $0\le x\le3$. השטח:
\[ S=\int_0^3\big[(2x-x^2)-(-x)\big]dx=\int_0^3(3x-x^2)\,dx=\left[\frac{3x^2}{2}-\frac{x^3}{3}\right]_0^3=\frac{27}{2}-9=\boxed{\frac92}. \]
\end{solution}

\subsection*{שאלה 6ב}

\begin{question}
מצאו את תחום ההגדרה, נקודות חיתוך עם הצירים, תחומי עליה וירידה, נקודות קיצון, אסימפטוטות, תחומי קמירות/קעירות של הפונקציה הבאה:
\[ f(x)=\frac{\sin(x)}{2+\cos(x)}. \]
שרטטו סקיצה של גרף הפונקציה
\end{question}

\begin{hint}
הפונקציה מחזורית עם מחזור $2\pi$ ואי-זוגית, ולכן מספיק לחקור אותה ב-$[0,2\pi]$ (או $[-\pi,\pi]$). שימו לב ש-$2+\cos x\ge1>0$.
\end{hint}

\begin{solution}
\textbf{תחום הגדרה:} $2+\cos x\ge1>0$ לכל $x$, ולכן $f$ מוגדרת (ורציפה וגזירה) בכל $\R$.
\textbf{סימטריה ומחזוריות:} $f(-x)=-f(x)$ (אי-זוגית), ו-$f(x+2\pi)=f(x)$ (מחזורית עם מחזור $2\pi$).

\textbf{חיתוך עם הצירים:} $f(x)=0\iff\sin x=0\iff x=k\pi$, $k\in\Z$. החיתוך עם ציר $y$: $(0,0)$. $f>0$ ב-$(2k\pi,(2k+1)\pi)$ ו-$f<0$ ב-$((2k-1)\pi,2k\pi)$.

\textbf{נגזרת ראשונה:}
\[ f'(x)=\frac{\cos x(2+\cos x)-\sin x\cdot(-\sin x)}{(2+\cos x)^2}=\frac{2\cos x+\cos^2x+\sin^2x}{(2+\cos x)^2}=\frac{2\cos x+1}{(2+\cos x)^2}. \]
$f'(x)=0\iff\cos x=-\frac12\iff x=\pm\frac{2\pi}{3}+2k\pi$. במחזור $[-\pi,\pi]$:
\begin{itemize}
  \item $f'>0$ (\textbf{עולה}) עבור $\cos x>-\frac12$, כלומר ב-$\left(-\frac{2\pi}3+2k\pi,\ \frac{2\pi}3+2k\pi\right)$;
  \item $f'<0$ (\textbf{יורדת}) ב-$\left(\frac{2\pi}3+2k\pi,\ \frac{4\pi}3+2k\pi\right)$.
\end{itemize}
\textbf{קיצון:} מקסימום ב-$x=\frac{2\pi}3+2k\pi$ עם $f=\frac{\sqrt3/2}{2-1/2}=\frac{1}{\sqrt3}$; מינימום ב-$x=-\frac{2\pi}3+2k\pi$ עם $f=-\frac1{\sqrt3}$. אלה גם הקיצון המוחלט: $-\frac1{\sqrt3}\le f(x)\le\frac1{\sqrt3}$ לכל $x$.

\textbf{אסימפטוטות:} אין אסימפטוטות אנכיות ($f$ רציפה בכל $\R$). אין אסימפטוטה אופקית או משופעת: $f$ מחזורית ואינה קבועה, ולכן $\lim_{x\to\pm\infty}f(x)$ אינו קיים (למשל $f(2k\pi)=0$ ו-$f(\frac{2\pi}3+2k\pi)=\frac1{\sqrt3}$), וגם $\frac{f(x)}{x}\to0$ ולכן אסימפטוטה משופעת הייתה חייבת להיות אופקית.

\textbf{נגזרת שנייה:} נגזור את $f'=\frac{2\cos x+1}{(2+\cos x)^2}$:
\[ f''(x)=\frac{-2\sin x(2+\cos x)^2+(2\cos x+1)\cdot2(2+\cos x)\sin x}{(2+\cos x)^4}=\frac{2\sin x\,\big[-(2+\cos x)+2\cos x+1\big]}{(2+\cos x)^3}=\frac{2\sin x(\cos x-1)}{(2+\cos x)^3}. \]
מכיוון ש-$\cos x-1\le0$ (ושווה ל-$0$ רק ב-$x=2k\pi$) והמכנה חיובי, הסימן של $f''$ הפוך לסימן של $\sin x$:
\begin{itemize}
  \item ב-$(2k\pi,(2k+1)\pi)$: $\sin x>0$, לכן $f''<0$ -- \textbf{קעורה} ($\cap$);
  \item ב-$((2k+1)\pi,(2k+2)\pi)$: $\sin x<0$, לכן $f''>0$ -- \textbf{קמורה} ($\cup$).
\end{itemize}
\textbf{נקודות פיתול:} $x=k\pi$ (הנקודות $(k\pi,0)$), כי שם $f''$ מחליפה סימן.

\textbf{סקיצה:} גל מחזורי דמוי סינוס שעובר בראשית, עולה עד למקסימום $\frac1{\sqrt3}\approx0.577$ ב-$x=\frac{2\pi}{3}$, יורד דרך $(\pi,0)$ (נקודת פיתול) עד למינימום $-\frac1{\sqrt3}$ ב-$x=\frac{4\pi}3$, וחוזר ל-$(2\pi,0)$; לעומת $\sin x$ הגרף "מוטה ימינה" -- המקסימום מוזז מ-$\frac\pi2$ ל-$\frac{2\pi}3$. התבנית חוזרת בכל מחזור באורך $2\pi$.
\end{solution}

\section*{תשע"ז סמסטר א מועד א}

\subsection*{שאלה 1}

\begin{question}
השתמשו בנוסחת טיילור כדי לחשב $\sqrt[3]{129}$ עם שגיאה לא גדולה מ-$10^{-3}$.
הסבירו כיצד הערכתם את השגיאה.
\end{question}

\begin{hint}
$129$ קרוב ל-$125=5^3$. פתחו את $f(x)=\sqrt[3]{x}$ סביב $x_0=125$.
\end{hint}

\begin{hint}
העריכו את השארית בצורת לגרנז' $R_n(x)=\frac{f^{(n+1)}(c)}{(n+1)!}(x-x_0)^{n+1}$ עם $125<c<129$, ובדקו איזה $n$ מספיק.
\end{hint}

\begin{solution}
נגדיר $f(x)=\sqrt[3]{x}=x^{1/3}$ ונפתח סביב $x_0=125$ (כי $\sqrt[3]{125}=5$ ידוע במדויק), עם $x=129$, כלומר $x-x_0=4$.

\textbf{נגזרות:}
\[
f'(x)=\frac13 x^{-2/3},\qquad f''(x)=-\frac29 x^{-5/3}.
\]
לכן $f(125)=5$, $f'(125)=\frac13\cdot\frac{1}{25}=\frac{1}{75}$.

\textbf{פולינום טיילור מסדר 1:}
\[
P_1(129)=f(125)+f'(125)\cdot 4=5+\frac{4}{75}=\frac{379}{75}\approx 5.05333.
\]

\textbf{הערכת השגיאה.} לפי משפט טיילור עם שארית לגרנז', קיימת נקודה $c$ בין $125$ ל-$129$ כך ש-
\[
R_1(129)=\frac{f''(c)}{2!}\cdot 4^2=-\frac{2}{9}c^{-5/3}\cdot\frac{16}{2}=-\frac{16}{9\,c^{5/3}}.
\]
מכיוון ש-$c>125$, מתקיים $c^{5/3}>125^{5/3}=5^5=3125$, ולכן
\[
|R_1(129)|<\frac{16}{9\cdot 3125}=\frac{16}{28125}\approx 5.7\cdot 10^{-4}<10^{-3}.
\]
לכן פולינום טיילור מסדר 1 מספיק.

\textbf{תשובה:}
\[
\sqrt[3]{129}\approx 5+\frac{4}{75}=\frac{379}{75}\approx 5.0533,
\]
עם שגיאה קטנה מ-$5.7\cdot10^{-4}<10^{-3}$. (השארית שלילית, כלומר הקירוב גדול מהערך האמיתי; ואכן $\sqrt[3]{129}\approx 5.05277$.)
\end{solution}

\subsection*{שאלה 2א}

\begin{question}
מצאו את הגבול הבא:
$\displaystyle \lim_{x\to 0}\frac{e^x-1}{\sqrt{1+x}-1}$
\end{question}

\begin{solution}
כאשר $x\to0$ המונה והמכנה שואפים ל-$0$ (שתי הפונקציות רציפות ב-$0$ ומתאפסות שם), כלומר גבול מהצורה $\frac00$. המונה והמכנה גזירים בסביבת $0$, ונגזרת המכנה $\frac{1}{2\sqrt{1+x}}\neq0$ שם. לפי כלל לופיטל:
\[
\lim_{x\to0}\frac{e^x-1}{\sqrt{1+x}-1}=\lim_{x\to0}\frac{e^x}{\frac{1}{2\sqrt{1+x}}}=\lim_{x\to0}2\sqrt{1+x}\,e^x=2.
\]
(דרך חלופית: נכפול בצמוד: $\frac{e^x-1}{\sqrt{1+x}-1}=\frac{e^x-1}{x}\cdot(\sqrt{1+x}+1)\to 1\cdot2=2$, לפי הגבול היסודי $\lim_{x\to0}\frac{e^x-1}{x}=1$.)

\textbf{תשובה:} $2$.
\end{solution}

\subsection*{שאלה 2ב}

\begin{question}
מצאו את הגבול הבא:
$\displaystyle \lim_{x\to 0}\frac{\sqrt[3]{8+3x}-2}{\sqrt[4]{16+5x}-2}$
\end{question}

\begin{solution}
בנקודה $x=0$: $\sqrt[3]{8}-2=0$ ו-$\sqrt[4]{16}-2=0$, כלומר גבול מהצורה $\frac00$. שתי הפונקציות גזירות בסביבת $0$:
\[
\left(\sqrt[3]{8+3x}\right)'=\frac13(8+3x)^{-2/3}\cdot3=(8+3x)^{-2/3},\qquad
\left(\sqrt[4]{16+5x}\right)'=\frac54(16+5x)^{-3/4},
\]
ונגזרת המכנה שונה מ-$0$ בסביבת $0$. לפי כלל לופיטל:
\[
\lim_{x\to0}\frac{\sqrt[3]{8+3x}-2}{\sqrt[4]{16+5x}-2}
=\lim_{x\to0}\frac{(8+3x)^{-2/3}}{\frac54(16+5x)^{-3/4}}
=\frac{8^{-2/3}}{\frac54\cdot16^{-3/4}}=\frac{\frac14}{\frac54\cdot\frac18}=\frac{\frac14}{\frac{5}{32}}=\frac85.
\]
\textbf{תשובה:} $\frac85$.
\end{solution}

\subsection*{שאלה 3}

\begin{question}
פונקציה
\[
f(x)=\frac{2-(e^x+e^{-x})\cos(x)}{x^4}
\]
מוגדרת עבור $x\neq0$. האם אפשר להגדיר $f(0)$ כך שהפונקציה תהי רציפה בנקודה $x=0$?
\end{question}

\begin{hint}
יש לבדוק האם $\lim_{x\to0}f(x)$ קיים וסופי.
\end{hint}

\begin{hint}
פתחו את $e^x+e^{-x}$ ואת $\cos x$ לפולינומי מקלורן עד סדר $4$ וכפלו.
\end{hint}

\begin{solution}
אפשר להגדיר את $f(0)$ כך ש-$f$ תהיה רציפה ב-$0$ אם ורק אם הגבול $\lim_{x\to0}f(x)$ קיים וסופי, ואז יש להגדיר את $f(0)$ להיות הגבול.

\textbf{פיתוח מקלורן.} לפי פיתוחי מקלורן של $e^x$ ושל $\cos x$:
\[
e^x+e^{-x}=2\left(1+\frac{x^2}{2}+\frac{x^4}{24}\right)+o(x^5)=2+x^2+\frac{x^4}{12}+o(x^5),
\]
\[
\cos x=1-\frac{x^2}{2}+\frac{x^4}{24}+o(x^5).
\]
נכפול ונשמור איברים עד סדר $4$:
\[
(e^x+e^{-x})\cos x=2+x^2+\frac{x^4}{12}-x^2-\frac{x^4}{2}+\frac{x^4}{12}+o(x^4)
=2+\left(\frac{1}{12}-\frac12+\frac1{12}\right)x^4+o(x^4)=2-\frac{x^4}{3}+o(x^4).
\]
מכאן
\[
2-(e^x+e^{-x})\cos x=\frac{x^4}{3}+o(x^4),
\qquad\text{ולכן}\qquad
f(x)=\frac13+\frac{o(x^4)}{x^4}\xrightarrow[x\to0]{}\frac13.
\]

(דרך חלופית: ארבע הפעלות של כלל לופיטל על $\frac{0}{0}$; הנגזרת הרביעית של המונה ב-$0$ היא $8$ ושל המכנה $24$, ולכן הגבול $\frac{8}{24}=\frac13$.)

\textbf{תשובה:} כן. הגבול קיים ושווה $\frac13$, ולכן אם נגדיר $f(0)=\frac13$ הפונקציה תהיה רציפה ב-$x=0$.
\end{solution}

\subsection*{שאלה 4א}

\begin{question}
חקרו את הפונקציה $y=\ln(e^x+e^{-2x})$ וסרטטו את הגרף שלה.
\underline{צריך למצוא}: תחום הגדרה, אסימפטוטות, תחומי עליה וירידה, נקודות קיצון, תחומי קמירות וקעירות.
\end{question}

\begin{hint}
לאסימפטוטות: כתבו $\ln(e^x+e^{-2x})=x+\ln(1+e^{-3x})=-2x+\ln(1+e^{3x})$.
\end{hint}

\begin{solution}
\textbf{תחום הגדרה:} $e^x+e^{-2x}>0$ לכל $x$, ולכן הפונקציה מוגדרת (ורציפה וגזירה) על כל $\R$.

\textbf{אסימפטוטות:} אין אסימפטוטות אנכיות, כי הפונקציה רציפה על כל $\R$.
עבור $x\to+\infty$:
\[
y=\ln\big(e^x(1+e^{-3x})\big)=x+\ln(1+e^{-3x}),\qquad y-x=\ln(1+e^{-3x})\xrightarrow[x\to+\infty]{}0,
\]
ולכן $y=x$ אסימפטוטה משופעת ב-$+\infty$.
עבור $x\to-\infty$:
\[
y=\ln\big(e^{-2x}(1+e^{3x})\big)=-2x+\ln(1+e^{3x}),\qquad y+2x\xrightarrow[x\to-\infty]{}0,
\]
ולכן $y=-2x$ אסימפטוטה משופעת ב-$-\infty$. אין אסימפטוטות אופקיות. בנוסף $y-x>0$ ו-$y+2x>0$, כלומר הגרף נמצא מעל שתי האסימפטוטות.

\textbf{עליה וירידה:}
\[
y'=\frac{e^x-2e^{-2x}}{e^x+e^{-2x}}.
\]
המכנה חיובי, ו-$e^x-2e^{-2x}>0\iff e^{3x}>2\iff x>\frac{\ln2}{3}$.
לכן $y$ יורדת ב-$\left(-\infty,\frac{\ln 2}{3}\right)$ ועולה ב-$\left(\frac{\ln2}{3},\infty\right)$.

\textbf{קיצון:} ב-$x_0=\frac{\ln2}{3}$ יש מינימום (מוחלט). שם $e^{x_0}=2^{1/3}$, $e^{-2x_0}=2^{-2/3}$, ולכן
\[
y(x_0)=\ln\left(2^{1/3}+2^{-2/3}\right)=\ln\left(3\cdot2^{-2/3}\right)=\ln3-\frac23\ln2\approx0.64.
\]

\textbf{קמירות:} לפי כלל המנה,
\[
y''=\frac{(e^x+4e^{-2x})(e^x+e^{-2x})-(e^x-2e^{-2x})^2}{(e^x+e^{-2x})^2}
=\frac{9e^{-x}}{(e^x+e^{-2x})^2}>0,
\]
(במונה: $e^{2x}+5e^{-x}+4e^{-4x}-e^{2x}+4e^{-x}-4e^{-4x}=9e^{-x}$). לכן הפונקציה קמורה (מחייכת, $\cup$) על כל $\R$, ואין נקודות פיתול.

\textbf{סרטוט:} הגרף יורד מלמעלה לאורך האסימפטוטה $y=-2x$ (מעליה), מגיע למינימום בנקודה $\left(\frac{\ln2}{3},\ \ln3-\frac23\ln2\right)\approx(0.23,\,0.64)$, ועולה לאורך האסימפטוטה $y=x$ (מעליה); כולו קמור. חיתוך עם ציר $y$: $y(0)=\ln2$. אין חיתוך עם ציר $x$ (כי $y\ge y(x_0)>0$).
\end{solution}

\subsection*{שאלה 4ב}

\begin{question}
חשבו אורך העקום $y=\frac25x\sqrt[4]{x}-\frac23\sqrt[4]{x^3}$ בין נקודות חיתוך עם ציר ה-$x$.
\end{question}

\begin{hint}
כתבו $y=\frac25x^{5/4}-\frac23x^{3/4}$ ובדקו ש-$1+(y')^2$ הוא ריבוע שלם.
\end{hint}

\begin{solution}
נכתוב $y=\frac25x^{5/4}-\frac23x^{3/4}$, המוגדרת עבור $x\ge0$.

\textbf{חיתוך עם ציר $x$:} $y=x^{3/4}\left(\frac25x^{1/2}-\frac23\right)=0\iff x=0$ או $\sqrt{x}=\frac53$, כלומר $x=0$ ו-$x=\frac{25}{9}$.

\textbf{נגזרת:}
\[
y'=\frac25\cdot\frac54x^{1/4}-\frac23\cdot\frac34x^{-1/4}=\frac12\left(x^{1/4}-x^{-1/4}\right),
\]
\[
1+(y')^2=1+\frac14\left(x^{1/2}-2+x^{-1/2}\right)=\frac14\left(x^{1/2}+2+x^{-1/2}\right)=\left(\frac{x^{1/4}+x^{-1/4}}{2}\right)^2.
\]
\textbf{אורך העקום} (לפי הנוסחה $L=\int_a^b\sqrt{1+(y')^2}\,dx$; האינטגרל לא אמיתי ב-$0$ אך מתכנס כי $x^{-1/4}$ אינטגרבילית שם):
\[
L=\int_0^{25/9}\frac{x^{1/4}+x^{-1/4}}{2}\,dx=\frac12\left[\frac45x^{5/4}+\frac43x^{3/4}\right]_0^{25/9}.
\]
עבור $x=\frac{25}{9}$: $x^{1/4}=\sqrt{\frac53}$, $x^{5/4}=\frac{25}{9}\sqrt{\frac53}$, $x^{3/4}=\frac53\sqrt{\frac53}$. לכן
\[
L=\frac12\left(\frac{20}{9}+\frac{20}{9}\right)\sqrt{\frac53}=\frac{20}{9}\sqrt{\frac53}=\frac{20\sqrt{15}}{27}\approx2.869.
\]
\end{solution}

\subsection*{שאלה 5א}

\begin{question}
חשבו את השטח החסום על ידי הקווים $y=4x,\ \ y=x,\ \ y=1/x$ בתחום $x\ge0,\ y\ge0$ וציירו את הסקיצה המתאימה.
\end{question}

\begin{hint}
מצאו את נקודות החיתוך של הקווים, וחלקו את האינטגרל לשני קטעים לפי הקו העליון.
\end{hint}

\begin{solution}
\textbf{נקודות חיתוך} (ברביע הראשון): $4x=x\Rightarrow x=0$ (הראשית); $4x=\frac1x\Rightarrow x=\frac12$ (הנקודה $(\frac12,2)$); $x=\frac1x\Rightarrow x=1$ (הנקודה $(1,1)$).

\textbf{סקיצה:} התחום הוא "משולש עקום" עם קודקודים $(0,0)$, $(\frac12,2)$, $(1,1)$: הגבול התחתון הוא הישר $y=x$, והגבול העליון הוא הישר $y=4x$ עבור $0\le x\le\frac12$ וההיפרבולה $y=\frac1x$ עבור $\frac12\le x\le1$.

\textbf{שטח:}
\[
S=\int_0^{1/2}(4x-x)\,dx+\int_{1/2}^1\left(\frac1x-x\right)dx
=\left[\frac{3x^2}{2}\right]_0^{1/2}+\left[\ln x-\frac{x^2}{2}\right]_{1/2}^1
=\frac38+\left(-\frac12-\ln\frac12+\frac18\right)=\frac38+\ln2-\frac38=\ln 2.
\]
\textbf{תשובה:} $S=\ln2$.
\end{solution}

\subsection*{שאלה 5ב}

\begin{question}
סרטטו את העקום $r=2\cos\varphi-1$ (נתון בקואורדינאטות פולריות) וחשבו שטח התחום החסום ע"י העקום.
\end{question}

\begin{hint}
מצאו עבור אילו זוויות $r\ge0$, והשתמשו בנוסחה $S=\frac12\int_\alpha^\beta r^2\,d\varphi$.
\end{hint}

\begin{solution}
\textbf{תחום הזוויות.} $r=2\cos\varphi-1\ge0\iff\cos\varphi\ge\frac12\iff -\frac\pi3\le\varphi\le\frac\pi3$ (בתוך מחזור אחד). בקצוות $\varphi=\pm\frac\pi3$ מתקיים $r=0$ (העקום עובר בראשית), ועבור $\varphi=0$ מתקיים $r=1$ (הנקודה $(1,0)$). העקום סימטרי ביחס לציר $x$ (כי $\cos$ זוגית).

\textbf{סרטוט:} כאשר $\varphi$ עולה מ-$-\frac\pi3$ ל-$\frac\pi3$, הנקודה יוצאת מהראשית בכיוון הזווית $-\frac\pi3$, מגיעה ל-$(1,0)$ וחוזרת לראשית בכיוון הזווית $\frac\pi3$: לולאה סגורה אחת בצורת "טיפה" בחצי המישור $x\ge0$, סימטרית לציר $x$.

\textbf{שטח} (לפי הנוסחה לשטח בקואורדינטות פולריות $S=\frac12\int_\alpha^\beta r^2\,d\varphi$):
\[
(2\cos\varphi-1)^2=4\cos^2\varphi-4\cos\varphi+1=3+2\cos2\varphi-4\cos\varphi,
\]
\[
S=\frac12\int_{-\pi/3}^{\pi/3}\left(3+2\cos2\varphi-4\cos\varphi\right)d\varphi
=\frac12\Big[3\varphi+\sin2\varphi-4\sin\varphi\Big]_{-\pi/3}^{\pi/3}
=\frac12\left(2\pi+\sqrt3-4\sqrt3\right)=\pi-\frac{3\sqrt3}{2}\approx0.544.
\]
\textbf{תשובה:} $S=\pi-\frac{3\sqrt3}{2}$.

\textbf{הערה:} אם מתירים ערכי $r$ שליליים (הנקודה $(r,\varphi)$ עם $r<0$ מסורטטת בכיוון $\varphi+\pi$ במרחק $|r|$), עבור $\frac\pi3<\varphi<\frac{5\pi}3$ מתקבלת לולאה חיצונית גדולה, והעקום כולו הוא לימסון (חילזון פסקל) עם לולאה פנימית — הלולאה שחושבה לעיל. במקרה זה השטח החסום ע"י הלולאה החיצונית הוא
\[
\frac12\int_{\pi/3}^{5\pi/3}(2\cos\varphi-1)^2\,d\varphi=2\pi+\frac{3\sqrt3}{2},
\]
והשטח שבין שתי הלולאות הוא $\left(2\pi+\frac{3\sqrt3}{2}\right)-\left(\pi-\frac{3\sqrt3}{2}\right)=\pi+3\sqrt3$.
\end{solution}

\subsection*{שאלה 6א}

\begin{question}
חשבו את האינטגרל הבא:
$\displaystyle\int\frac{dx}{x^2+x+1}$
\end{question}

\begin{hint}
השלימו לריבוע: $x^2+x+1=\left(x+\frac12\right)^2+\frac34$.
\end{hint}

\begin{solution}
נשלים לריבוע: $x^2+x+1=\left(x+\frac12\right)^2+\frac34$. נציב $t=\frac{2x+1}{\sqrt3}$, כלומר $x+\frac12=\frac{\sqrt3}{2}t$, $dx=\frac{\sqrt3}{2}dt$:
\[
\int\frac{dx}{\left(x+\frac12\right)^2+\frac34}=\int\frac{\frac{\sqrt3}{2}\,dt}{\frac34(t^2+1)}=\frac{2}{\sqrt3}\int\frac{dt}{1+t^2}=\frac{2}{\sqrt3}\arctan t+C.
\]
\textbf{תשובה:}
\[
\int\frac{dx}{x^2+x+1}=\frac{2}{\sqrt3}\arctan\frac{2x+1}{\sqrt3}+C.
\]
(בדיקה: $\frac{d}{dx}\frac{2}{\sqrt3}\arctan\frac{2x+1}{\sqrt3}=\frac{2}{\sqrt3}\cdot\frac{2/\sqrt3}{1+\frac{(2x+1)^2}{3}}=\frac{4}{3+(2x+1)^2}=\frac{1}{x^2+x+1}$.)
\end{solution}

\subsection*{שאלה 6ב}

\begin{question}
חשבו את האינטגרל הבא:
$\displaystyle\int_1^\infty x^2e^{-2x}\,dx$
\end{question}

\begin{hint}
בצעו אינטגרציה בחלקים פעמיים (גוזרים את $x^2$), ואז חשבו את הגבול כאשר הגבול העליון שואף לאינסוף.
\end{hint}

\begin{solution}
\textbf{פונקציה קדומה.} באינטגרציה בחלקים עם $u=x^2$, $dv=e^{-2x}dx$, $v=-\frac12e^{-2x}$:
\[
\int x^2e^{-2x}dx=-\frac{x^2}{2}e^{-2x}+\int xe^{-2x}dx.
\]
שוב בחלקים עם $u=x$:
\[
\int xe^{-2x}dx=-\frac x2e^{-2x}+\frac12\int e^{-2x}dx=-\frac x2e^{-2x}-\frac14e^{-2x}.
\]
לכן
\[
F(x)=\int x^2e^{-2x}dx=-e^{-2x}\left(\frac{x^2}{2}+\frac x2+\frac14\right)+C.
\]
\textbf{האינטגרל הלא אמיתי.} לפי ההגדרה,
\[
\int_1^\infty x^2e^{-2x}dx=\lim_{R\to\infty}\big(F(R)-F(1)\big).
\]
מתקיים $\lim_{R\to\infty}e^{-2R}\left(\frac{R^2}{2}+\frac R2+\frac14\right)=0$ (אקספוננט שולט בפולינום; למשל לפי כלל לופיטל פעמיים על $\frac{R^2}{e^{2R}}$), ו-$F(1)=-e^{-2}\left(\frac12+\frac12+\frac14\right)=-\frac54e^{-2}$. לכן
\[
\int_1^\infty x^2e^{-2x}dx=0+\frac54e^{-2}=\frac{5}{4e^2}\approx0.169.
\]
האינטגרל מתכנס, וערכו $\frac{5}{4e^2}$.
\end{solution}

\section*{תשע"ז סמסטר א מועד ב}

\subsection*{שאלה 1}

\begin{question}
השתמשו בנוסחת טיילור כדי לחשב $\sin 167^\circ$ עם שגיאה לא גדולה מ-$10^{-5}$. הסבירו כיצד הערכתם את השגיאה.
\end{question}

\begin{hint}
$\sin167^\circ=\sin(180^\circ-13^\circ)=\sin13^\circ$. עברו לרדיאנים ופתחו את $\sin x$ סביב $0$.
\end{hint}

\begin{hint}
בפולינום מקלורן של $\sin x$ המקדם של $x^4$ הוא $0$, ולכן $P_3=P_4$ ואפשר להעריך את השארית בעזרת $R_4$.
\end{hint}

\begin{solution}
\textbf{רדוקציה.} $\sin167^\circ=\sin(180^\circ-13^\circ)=\sin13^\circ=\sin x_0$, כאשר
\[
x_0=\frac{13\pi}{180}\approx0.226893 \quad\text{(רדיאנים)}.
\]

\textbf{פולינום מקלורן.} עבור $f(x)=\sin x$: $f(0)=0$, $f'(0)=1$, $f''(0)=0$, $f'''(0)=-1$, $f^{(4)}(0)=0$, ולכן
\[
P_4(x)=P_3(x)=x-\frac{x^3}{6}.
\]

\textbf{הערכת השגיאה.} לפי משפט טיילור עם שארית לגרנז' מסדר $4$, קיים $c$ בין $0$ ל-$x_0$ כך ש-
\[
\sin x_0-P_4(x_0)=R_4(x_0)=\frac{f^{(5)}(c)}{5!}x_0^5=\frac{\cos c}{120}x_0^5 .
\]
מכיוון ש-$|\cos c|\le1$ ו-$0<x_0<0.23$:
\[
|R_4(x_0)|\le\frac{x_0^5}{120}<\frac{0.23^5}{120}\approx\frac{6.44\cdot10^{-4}}{120}\approx5.4\cdot10^{-6}<10^{-5}.
\]
(גם בהערכה גסה יותר $x_0<0.25$ מקבלים $\frac{0.25^5}{120}\approx8.1\cdot10^{-6}<10^{-5}$.)

\textbf{חישוב:}
\[
\sin167^\circ\approx x_0-\frac{x_0^3}{6}\approx0.226893-\frac{0.011680}{6}\approx0.226893-0.001947=0.224946.
\]

\textbf{תשובה:} $\sin167^\circ\approx\frac{13\pi}{180}-\frac16\left(\frac{13\pi}{180}\right)^3\approx0.22495$, עם שגיאה קטנה מ-$10^{-5}$ (הערך המדויק $0.2249511\ldots$).
\end{solution}

\subsection*{שאלה 2א}

\begin{question}
מצאו את הגבול הבא:
$\displaystyle\lim_{x\to0}\frac{\sqrt{\cos x}-\sqrt[3]{\cos x}}{\sin^2x}$
\end{question}

\begin{hint}
השתמשו ב-$\cos x=1-\frac{x^2}{2}+o(x^2)$ ובקירוב $(1+t)^\alpha=1+\alpha t+o(t)$.
\end{hint}

\begin{solution}
זהו גבול מהצורה $\frac00$. לפי פיתוח מקלורן $\cos x=1+t$ עם $t=-\frac{x^2}{2}+o(x^2)\to0$. לפי $(1+t)^\alpha=1+\alpha t+o(t)$:
\[
\sqrt{\cos x}=1-\frac{x^2}{4}+o(x^2),\qquad \sqrt[3]{\cos x}=1-\frac{x^2}{6}+o(x^2),
\]
ולכן
\[
\sqrt{\cos x}-\sqrt[3]{\cos x}=-\frac{x^2}{4}+\frac{x^2}{6}+o(x^2)=-\frac{x^2}{12}+o(x^2).
\]
בנוסף $\sin^2x=x^2+o(x^2)$ (כי $\frac{\sin x}{x}\to1$). לכן
\[
\lim_{x\to0}\frac{\sqrt{\cos x}-\sqrt[3]{\cos x}}{\sin^2x}=\lim_{x\to0}\frac{-\frac{1}{12}+\frac{o(x^2)}{x^2}}{\left(\frac{\sin x}{x}\right)^2}=-\frac1{12}.
\]
(דרך חלופית: נסמן $u=\cos x$; אז $\sin^2x=1-u^2$ ו-$u\to1$, ו-$\lim_{u\to1}\frac{u^{1/2}-u^{1/3}}{1-u^2}$ לפי לופיטל שווה $\frac{\frac12-\frac13}{-2}=-\frac1{12}$.)

\textbf{תשובה:} $-\frac{1}{12}$.
\end{solution}

\subsection*{שאלה 2ב}

\begin{question}
מצאו את הגבול הבא:
$\displaystyle\lim_{x\to0}\left(\frac{1+x}{2+x}\right)^{\frac{1-\sqrt{x}}{1-x}}$
\end{question}

\begin{hint}
בדקו לאן שואפים הבסיס והמעריך בנפרד. האם זהו בכלל ביטוי לא מוגדר?
\end{hint}

\begin{solution}
הביטוי מוגדר רק עבור $x\ge0$ (בגלל $\sqrt x$), ולכן מדובר בגבול חד-צדדי $x\to0^+$.
זהו \emph{לא} ביטוי לא-מוגדר: הבסיס $\frac{1+x}{2+x}\to\frac12>0$ והמעריך $\frac{1-\sqrt x}{1-x}\to1$. נכתוב
\[
\left(\frac{1+x}{2+x}\right)^{\frac{1-\sqrt{x}}{1-x}}=\exp\left(\frac{1-\sqrt{x}}{1-x}\cdot\ln\frac{1+x}{2+x}\right).
\]
לפי אריתמטיקה של גבולות ורציפות $\ln$, המעריך שואף ל-$1\cdot\ln\frac12=\ln\frac12$, ולפי רציפות פונקציית האקספוננט:
\[
\lim_{x\to0^+}\left(\frac{1+x}{2+x}\right)^{\frac{1-\sqrt{x}}{1-x}}=e^{\ln\frac12}=\frac12.
\]
\textbf{תשובה:} $\frac12$.
\end{solution}

\subsection*{שאלה 3א}

\begin{question}
האם קיים ערך $c$ כך ש-
\[
f(x)=\begin{cases}\dfrac1x & x\ge1\\[1ex] 2x+c & x<1\end{cases}
\]
תהיה רציפה ב-$x=1$?
\end{question}

\begin{solution}
$f$ רציפה ב-$x=1$ אם ורק אם $\lim_{x\to1^-}f(x)=\lim_{x\to1^+}f(x)=f(1)$.
\[
f(1)=\frac11=1,\qquad \lim_{x\to1^+}\frac1x=1,\qquad \lim_{x\to1^-}(2x+c)=2+c.
\]
לכן צריך $2+c=1$, כלומר $c=-1$.

\textbf{תשובה:} כן, עבור $c=-1$ (ורק עבורו).
\end{solution}

\subsection*{שאלה 3ב}

\begin{question}
האם למשוואה $\cos(x)-x=0$ יש פתרון בקטע $(-1,1)$, האם הוא יחיד? צטטו במדויק כל משפט שהנכם מסתמכים עליו.
\end{question}

\begin{hint}
הגדירו $g(x)=\cos x-x$ ובדקו את הסימנים של $g(\pm1)$. ליחידות, בדקו את סימן $g'$.
\end{hint}

\begin{solution}
נגדיר $g(x)=\cos x-x$. זו פונקציה רציפה וגזירה על כל $\R$.

\textbf{קיום.} \emph{משפט ערך הביניים (בולצאנו):} אם $g$ רציפה בקטע סגור $[a,b]$ ו-$g(a),g(b)$ בעלי סימנים מנוגדים, אז קיימת $x_0\in(a,b)$ עם $g(x_0)=0$.
כאן $g$ רציפה ב-$[-1,1]$, ו-
\[
g(-1)=\cos1+1>0,\qquad g(1)=\cos1-1<0
\]
(כי $0<\cos1<1$). לכן קיים $x_0\in(-1,1)$ עם $\cos x_0-x_0=0$.

\textbf{יחידות.} $g'(x)=-\sin x-1$. מתקיים $g'(x)\le0$ לכל $x$, ושוויון $g'(x)=0$ רק כאשר $\sin x=-1$, כלומר $x=-\frac\pi2+2\pi k$; אף נקודה כזאת אינה בקטע $(-1,1)$ (כי $\frac\pi2>1$). לכן $g'(x)<0$ לכל $x\in(-1,1)$.
לפי מסקנה ממשפט לגרנז' (\emph{אם $g'<0$ בקטע אז $g$ יורדת ממש בו}), $g$ יורדת ממש ב-$(-1,1)$ ולכן חד-חד-ערכית שם, ויש לה לכל היותר אפס אחד.
(באופן שקול: אם היו שני שורשים $x_1<x_2$, לפי \emph{משפט רול} — $g$ רציפה ב-$[x_1,x_2]$, גזירה ב-$(x_1,x_2)$ ו-$g(x_1)=g(x_2)$ — הייתה קיימת $\xi\in(x_1,x_2)$ עם $g'(\xi)=0$, בסתירה.)

\textbf{תשובה:} כן, למשוואה יש פתרון בקטע $(-1,1)$, והוא יחיד.
\end{solution}

\subsection*{שאלה 4א}

\begin{question}
חקרו את הפונקציה $f(x)=x^2e^{-x}$ ושרטטו את הגרף שלה.
\underline{צריך למצוא}: תחום הגדרה, אסימפטוטות; תחומי עליה וירידה, נקודות קיצון, תחומי קמירות וקעירות.
\end{question}

\begin{solution}
\textbf{תחום הגדרה:} כל $\R$; הפונקציה רציפה וגזירה בכל נקודה, ו-$f(x)\ge0$ עם שוויון רק ב-$x=0$.

\textbf{אסימפטוטות:} אין אנכיות (רציפה על $\R$).
כאשר $x\to+\infty$: $\lim_{x\to\infty}\frac{x^2}{e^x}=0$ (לופיטל פעמיים), ולכן $y=0$ אסימפטוטה אופקית ב-$+\infty$.
כאשר $x\to-\infty$: $f(x)\to+\infty$ ו-$\frac{f(x)}{x}=xe^{-x}\to-\infty$, ולכן אין אסימפטוטה (אופקית או משופעת) ב-$-\infty$.

\textbf{עליה וירידה:}
\[
f'(x)=2xe^{-x}-x^2e^{-x}=x(2-x)e^{-x}.
\]
$f'<0$ ב-$(-\infty,0)$ וב-$(2,\infty)$ (יורדת), ו-$f'>0$ ב-$(0,2)$ (עולה).

\textbf{קיצון:} מינימום (מוחלט) ב-$(0,0)$; מקסימום מקומי ב-$\left(2,\frac4{e^2}\right)\approx(2,0.54)$ (לא מוחלט, כי $f\to\infty$ ב-$-\infty$).

\textbf{קמירות:}
\[
f''(x)=\big((2-2x)-(2x-x^2)\big)e^{-x}=(x^2-4x+2)e^{-x},
\]
שמתאפסת ב-$x=2\pm\sqrt2$. לכן $f$ קמורה ($\cup$) ב-$(-\infty,2-\sqrt2)$ וב-$(2+\sqrt2,\infty)$, וקעורה ($\cap$) ב-$(2-\sqrt2,2+\sqrt2)$. נקודות פיתול ב-$x=2\pm\sqrt2$ (בערך $x\approx0.59$ ו-$x\approx3.41$).

\textbf{סרטוט:} הגרף יורד מ-$+\infty$ (ב-$-\infty$) עד המינימום $(0,0)$ ומשיק לציר $x$ שם, עולה עד המקסימום $(2,4e^{-2})$, ואז יורד ושואף ל-$0$ מלמעלה ($y=0$ אסימפטוטה ב-$+\infty$).
\end{solution}

\subsection*{שאלה 4ב}

\begin{question}
ציירו את הגרף של הפונקציה $r=f(\theta)=\theta^2$ בקואורדינטות קוטביות עבור $0\le\theta\le\pi$ וחשבו את אורך העקום.
\end{question}

\begin{hint}
השתמשו בנוסחה לאורך עקום בקואורדינטות קוטביות $L=\int_\alpha^\beta\sqrt{r^2+(r')^2}\,d\theta$.
\end{hint}

\begin{solution}
\textbf{סרטוט:} $r=\theta^2$ עולה מ-$0$ ל-$\pi^2\approx9.87$ כאשר $\theta$ עולה מ-$0$ ל-$\pi$: ספירלה שיוצאת מהראשית (משיקה לציר $x$ החיובי), עוברת דרך $\left(\theta=\frac\pi2,\ r=\frac{\pi^2}4\approx2.47\right)$ על ציר $y$ החיובי, ומגיעה לנקודה $(-\pi^2,0)$ על ציר $x$ השלילי; כל העקום בחצי המישור העליון $y\ge0$.

\textbf{אורך:} לפי הנוסחה לאורך עקום בקואורדינטות קוטביות, עם $r'=2\theta$:
\[
L=\int_0^\pi\sqrt{r^2+(r')^2}\,d\theta=\int_0^\pi\sqrt{\theta^4+4\theta^2}\,d\theta=\int_0^\pi\theta\sqrt{\theta^2+4}\,d\theta
\]
(כי $\theta\ge0$). בהצבה $u=\theta^2+4$, $du=2\theta\,d\theta$:
\[
L=\frac12\int_4^{\pi^2+4}\sqrt u\,du=\frac13\Big[u^{3/2}\Big]_4^{\pi^2+4}=\frac13\left((\pi^2+4)^{3/2}-8\right)\approx14.55.
\]
\end{solution}

\subsection*{שאלה 5א}

\begin{question}
חשבו את האינטגרל $\displaystyle\int\frac{x^3}{\sqrt{1-x^2}}\,dx$.
\end{question}

\begin{hint}
הציבו $u=1-x^2$ (ואז $x^2=1-u$).
\end{hint}

\begin{solution}
נציב $u=1-x^2$, $du=-2x\,dx$, ו-$x^2=1-u$ (עבור $|x|<1$):
\[
\int\frac{x^3}{\sqrt{1-x^2}}dx=\int\frac{x^2\cdot x\,dx}{\sqrt{1-x^2}}=-\frac12\int\frac{1-u}{\sqrt u}\,du=-\frac12\int\left(u^{-1/2}-u^{1/2}\right)du
=-\sqrt u+\frac13u^{3/2}+C.
\]
\textbf{תשובה:}
\[
\int\frac{x^3}{\sqrt{1-x^2}}dx=-\sqrt{1-x^2}+\frac13(1-x^2)^{3/2}+C=-\frac13(x^2+2)\sqrt{1-x^2}+C.
\]
(בדיקה: גזירת $-\frac13(x^2+2)\sqrt{1-x^2}$ נותנת $-\frac23x\sqrt{1-x^2}+\frac{x(x^2+2)}{3\sqrt{1-x^2}}=\frac{-2x(1-x^2)+x^3+2x}{3\sqrt{1-x^2}}=\frac{x^3}{\sqrt{1-x^2}}$.)
\end{solution}

\subsection*{שאלה 5ב}

\begin{question}
חשבו את האינטגרל הלא-אמיתי $\displaystyle\int_1^\infty\frac{\arctan(x)}{x^2}\,dx$.
\end{question}

\begin{hint}
בצעו אינטגרציה בחלקים עם $u=\arctan x$, $dv=\frac{dx}{x^2}$, ופרקו $\frac{1}{x(1+x^2)}=\frac1x-\frac{x}{1+x^2}$.
\end{hint}

\begin{solution}
\textbf{פונקציה קדומה.} אינטגרציה בחלקים עם $u=\arctan x$, $dv=\frac{dx}{x^2}$, $du=\frac{dx}{1+x^2}$, $v=-\frac1x$:
\[
\int\frac{\arctan x}{x^2}dx=-\frac{\arctan x}{x}+\int\frac{dx}{x(1+x^2)}.
\]
פירוק לשברים חלקיים: $\frac{1}{x(1+x^2)}=\frac1x-\frac{x}{1+x^2}$, ולכן עבור $x>0$
\[
\int\frac{dx}{x(1+x^2)}=\ln x-\frac12\ln(1+x^2)=\ln\frac{x}{\sqrt{1+x^2}}.
\]
\textbf{האינטגרל הלא אמיתי.} לפי ההגדרה, $\int_1^\infty=\lim_{R\to\infty}\int_1^R$:
\[
\int_1^R\frac{\arctan x}{x^2}dx=\left[-\frac{\arctan x}{x}+\ln\frac{x}{\sqrt{1+x^2}}\right]_1^R
=-\frac{\arctan R}{R}+\ln\frac{R}{\sqrt{1+R^2}}+\frac\pi4-\ln\frac1{\sqrt2}.
\]
כאשר $R\to\infty$: $\frac{\arctan R}{R}\to0$ (מונה חסום, מכנה שואף לאינסוף), ו-$\frac{R}{\sqrt{1+R^2}}\to1$ ולכן $\ln\frac{R}{\sqrt{1+R^2}}\to0$ (רציפות $\ln$). לכן
\[
\int_1^\infty\frac{\arctan x}{x^2}dx=\frac\pi4+\frac12\ln2\approx1.132.
\]
\end{solution}

\subsection*{שאלה 6}

\begin{question}
חשבו את נפח גוף הסיבוב המתקבל ע"י סיבוב התחום $D$ סביב ציר ה-$x$, כאשר $D$ התחום בין גרף הפונקציה $f(x)=\dfrac{x}{x^2+6x+8}$ וציר ה-$x$, כאשר $0\le x\le2$.
\end{question}

\begin{hint}
הנפח הוא $V=\pi\int_0^2f^2(x)\,dx$. פרקו קודם את $f$ לשברים חלקיים: $x^2+6x+8=(x+2)(x+4)$.
\end{hint}

\begin{solution}
\textbf{נוסחת הנפח.} לפי הנוסחה לנפח גוף סיבוב סביב ציר $x$:
\[
V=\pi\int_0^2 f^2(x)\,dx=\pi\int_0^2\frac{x^2}{(x^2+6x+8)^2}\,dx.
\]

\textbf{שברים חלקיים.} $x^2+6x+8=(x+2)(x+4)$, ו-
\[
\frac{x}{(x+2)(x+4)}=\frac{A}{x+2}+\frac{B}{x+4},\qquad A=\frac{-2}{-2+4}=-1,\quad B=\frac{-4}{-4+2}=2.
\]
לכן $f(x)=\frac{2}{x+4}-\frac{1}{x+2}$, ו-
\[
f^2(x)=\frac{4}{(x+4)^2}+\frac{1}{(x+2)^2}-\frac{4}{(x+2)(x+4)},\qquad
\frac{4}{(x+2)(x+4)}=\frac{2}{x+2}-\frac{2}{x+4}.
\]

\textbf{אינטגרציה:}
\[
\int_0^2\frac{4\,dx}{(x+4)^2}=\left[-\frac4{x+4}\right]_0^2=-\frac46+1=\frac13,\qquad
\int_0^2\frac{dx}{(x+2)^2}=\left[-\frac1{x+2}\right]_0^2=-\frac14+\frac12=\frac14,
\]
\[
\int_0^2\left(\frac{2}{x+2}-\frac{2}{x+4}\right)dx=\Big[2\ln(x+2)-2\ln(x+4)\Big]_0^2=2\ln\frac46-2\ln\frac24=2\ln\frac43.
\]
לכן
\[
\int_0^2f^2(x)\,dx=\frac13+\frac14-2\ln\frac43=\frac7{12}-2\ln\frac43.
\]

\textbf{תשובה:}
\[
V=\pi\left(\frac{7}{12}-2\ln\frac43\right)\approx0.0250.
\]
(הנפח קטן מאוד, וזה סביר: $0\le f(x)\le f(2)=\frac1{12}$ בקטע.)
\end{solution}

\section*{תשע"ז סמסטר ב מועד א}

\subsection*{שאלה 1}

\begin{question}
חשבו את הנגזרת של הפונקציה $f(x)=\ln(2x)$ בנקודה $x$ ($x>0$), מההגדרה (ללא שימוש בנגזרות ידועות).
\end{question}

\begin{hint}
השתמשו בחוקי הלוגריתם כדי לכתוב את מנת ההפרשים כ-$\frac{\ln(1+\frac hx)}{h}$, ואז בגבול היסודי $\lim_{t\to0}\frac{\ln(1+t)}{t}=1$.
\end{hint}

\begin{solution}
יהי $x>0$. עבור $h\ne0$ עם $|h|<x$ (כך ש-$x+h>0$), לפי חוקי הלוגריתם:
\[ \frac{f(x+h)-f(x)}{h}=\frac{\ln(2x+2h)-\ln(2x)}{h}=\frac{1}{h}\ln\frac{2(x+h)}{2x}=\frac1h\ln\left(1+\frac hx\right). \]
נציב $t=\frac hx$; כאשר $h\to0$ גם $t\to0$ ו-$t\ne0$, ולכן לפי הגבול היסודי $\lim_{t\to0}\frac{\ln(1+t)}{t}=1$ (וכלל ההצבה בגבולות):
\[ \frac1h\ln\left(1+\frac hx\right)=\frac1x\cdot\frac{\ln(1+t)}{t}\xrightarrow[h\to0]{}\frac1x\cdot1. \]
לכן
\[ \boxed{f'(x)=\lim_{h\to0}\frac{f(x+h)-f(x)}{h}=\frac1x}. \]
(שימו לב שהתוצאה זהה לנגזרת של $\ln x$, כי $\ln(2x)=\ln2+\ln x$.)
\end{solution}

\subsection*{שאלה 2}

\begin{question}
האם לפונקציה $f(x)=5-x-x^2-|x+3|$ יש מקסימום בקטע $[-10,10]$? אם כן, מצאו אותו.
\end{question}

\begin{hint}
$f$ רציפה בקטע סגור -- איזה משפט מבטיח קיום מקסימום? כדי למצוא אותו, פתחו את הערך המוחלט ובדקו כל אחד משני התחומים בנפרד.
\end{hint}

\begin{solution}
\textbf{קיום:} $f$ רציפה ב-$[-10,10]$ (סכום של פולינום ושל $-|x+3|$, שהיא רציפה). לפי משפט ויירשטראס, פונקציה רציפה בקטע סגור וחסום מקבלת בו מקסימום (ומינימום). לכן \textbf{כן}, יש מקסימום.

\textbf{מציאתו:} נפתח את הערך המוחלט.
\begin{itemize}
  \item עבור $-3\le x\le10$: $|x+3|=x+3$ ולכן $f(x)=5-x-x^2-x-3=2-2x-x^2=3-(x+1)^2$. זו פרבולה הפוכה עם קודקוד ב-$x=-1\in[-3,10]$, ולכן בתחום זה $f(x)\le3$ ושוויון רק ב-$x=-1$.
  \item עבור $-10\le x<-3$: $|x+3|=-(x+3)$ ולכן $f(x)=5-x-x^2+x+3=8-x^2$. כאן $x^2>9$, ולכן $f(x)<8-9=-1<3$.
\end{itemize}
(לחלופין: $f$ גזירה פרט ל-$x=-3$; $f'(x)=-2-2x$ עבור $x>-3$ מתאפסת ב-$x=-1$, ו-$f'(x)=-2x>0$ עבור $x<-3$; משווים את ערכי $f$ בנקודות החשודות $-10,-3,-1,10$: $f(-10)=-92$, $f(-3)=-1$, $f(-1)=3$, $f(10)=-118$.)

לכן המקסימום של $f$ בקטע הוא $\boxed{f(-1)=3}$, והוא מתקבל בנקודה $x=-1$.
\end{solution}

\subsection*{שאלה 3}

\begin{question}
חשבו:
\[ \int x^3e^{-x^2}\,dx \]
\end{question}

\begin{hint}
הציבו $t=x^2$ (או $t=-x^2$) ואחר כך בצעו אינטגרציה בחלקים.
\end{hint}

\begin{solution}
נציב $t=x^2$, $dt=2x\,dx$. אז $x^3\,dx=x^2\cdot x\,dx=\frac t2\,dt$:
\[ \int x^3e^{-x^2}dx=\frac12\int te^{-t}dt. \]
אינטגרציה בחלקים עם $u=t$, $v'=e^{-t}$ ($u'=1$, $v=-e^{-t}$):
\[ \int te^{-t}dt=-te^{-t}+\int e^{-t}dt=-te^{-t}-e^{-t}+C=-(t+1)e^{-t}+C. \]
נחזור למשתנה $x$:
\[ \boxed{\int x^3e^{-x^2}dx=-\frac12(x^2+1)e^{-x^2}+C.} \]
\textbf{בדיקה:} $\frac{d}{dx}\left[-\frac12(x^2+1)e^{-x^2}\right]=-xe^{-x^2}+\frac12(x^2+1)\cdot2xe^{-x^2}=x^3e^{-x^2}$.
\end{solution}

\subsection*{שאלה 4}

\begin{question}
חשבו:
\[ \int\frac{x^4+1}{x^3-x^2-x+1}\,dx \]
\end{question}

\begin{hint}
המעלה של המונה גדולה ממעלת המכנה: בצעו קודם חילוק פולינומים. פרקו את המכנה: $x^3-x^2-x+1=(x-1)^2(x+1)$.
\end{hint}

\begin{hint}
לשורש הכפול $x=1$ יש שני שברים חלקיים: $\frac{B}{x-1}+\frac{C}{(x-1)^2}$.
\end{hint}

\begin{solution}
\textbf{פירוק המכנה:} $x^3-x^2-x+1=x^2(x-1)-(x-1)=(x-1)(x^2-1)=(x-1)^2(x+1)$.

\textbf{חילוק פולינומים:} המנה היא $x+1$; אכן $(x+1)(x^3-x^2-x+1)=x^4-2x^2+1$, ולכן
\[ x^4+1=(x+1)(x^3-x^2-x+1)+2x^2,\qquad \frac{x^4+1}{x^3-x^2-x+1}=x+1+\frac{2x^2}{(x-1)^2(x+1)}. \]

\textbf{שברים חלקיים:}
\[ \frac{2x^2}{(x-1)^2(x+1)}=\frac{A}{x+1}+\frac{B}{x-1}+\frac{C}{(x-1)^2}\iff 2x^2=A(x-1)^2+B(x-1)(x+1)+C(x+1). \]
הצבת $x=-1$: $2=4A$, כלומר $A=\frac12$. הצבת $x=1$: $2=2C$, כלומר $C=1$. השוואת מקדם $x^2$: $2=A+B$, לכן $B=\frac32$.

\textbf{אינטגרציה:}
\[ \int\frac{x^4+1}{x^3-x^2-x+1}dx=\int\left(x+1+\frac{1/2}{x+1}+\frac{3/2}{x-1}+\frac{1}{(x-1)^2}\right)dx, \]
\[ \boxed{=\frac{x^2}{2}+x+\frac12\ln|x+1|+\frac32\ln|x-1|-\frac{1}{x-1}+C.} \]
(בדיקה: גזירת התוצאה מחזירה את האינטגרנד.)
\end{solution}

\subsection*{שאלה 5}

\begin{question}
נגדיר: $f(x)=\ln(1+x^2)$. חשבו את $f^{(4)}(0)$.
\end{question}

\begin{hint}
במקום לגזור ארבע פעמים, השתמשו בפולינום מקלורין של $\ln(1+t)$ עם $t=x^2$, והשוו מקדמים: המקדם של $x^4$ הוא $\frac{f^{(4)}(0)}{4!}$.
\end{hint}

\begin{solution}
מפיתוח מקלורין $\ln(1+t)=t-\frac{t^2}{2}+o(t^2)$, ועם $t=x^2$:
\[ \ln(1+x^2)=x^2-\frac{x^4}{2}+o(x^4). \]
$f$ גזירה אינסוף פעמים, ולפי יחידות פולינום טיילור (הפולינום היחיד ממעלה $\le4$ המקרב את $f$ עד $o(x^4)$ הוא פולינום טיילור), פולינום טיילור מסדר $4$ של $f$ סביב $0$ הוא $x^2-\frac{x^4}{2}$. מקדם $x^4$ בפולינום טיילור הוא $\frac{f^{(4)}(0)}{4!}$, ולכן
\[ \frac{f^{(4)}(0)}{24}=-\frac12\quad\Longrightarrow\quad\boxed{f^{(4)}(0)=-12}. \]
(בדיקה בגזירה ישירה: $f'(x)=\frac{2x}{1+x^2}$, $f''(x)=\frac{2-2x^2}{(1+x^2)^2}$, $f'''(x)=\frac{4x^3-12x}{(1+x^2)^3}$, $f^{(4)}(x)=\frac{-12x^4+72x^2-12}{(1+x^2)^4}$, ו-$f^{(4)}(0)=-12$.)
\end{solution}

\subsection*{שאלה 6}

\begin{question}
תהי $f$ פונקציה חד חד ערכית וגזירה. נתון: $f(1)=2$, $f(2)=4$, $f(3)=5$, $f(4)=6$, $f'(1)=0$, $f'(2)=2$, $f'(3)=0$, $f'(4)=3$.
חשבו: $\left(f^{-1}\right)'(4)$
\end{question}

\begin{hint}
לפי משפט הנגזרת של הפונקציה ההפוכה, $(f^{-1})'(y)=\frac{1}{f'(f^{-1}(y))}$. מהו $f^{-1}(4)$?
\end{hint}

\begin{solution}
מכיוון ש-$f(2)=4$ ו-$f$ חד-חד-ערכית, $f^{-1}(4)=2$.
לפי משפט הנגזרת של הפונקציה ההפוכה: אם $f$ גזירה ב-$x_0$, $f'(x_0)\ne0$ (ו-$f$ רציפה וחד-חד-ערכית בסביבה), אז $f^{-1}$ גזירה ב-$y_0=f(x_0)$ ו-
\[ (f^{-1})'(y_0)=\frac{1}{f'(x_0)}. \]
כאן $x_0=2$, $y_0=4$, ו-$f'(2)=2\ne0$, ולכן
\[ \boxed{(f^{-1})'(4)=\frac{1}{f'(2)}=\frac12}. \]
(שאר הנתונים -- $f'(4)=3$ למשל -- הם מסיחים: יש להעריך את $f'$ בנקודה $f^{-1}(4)=2$ ולא בנקודה $4$.)
\end{solution}

\subsection*{שאלה 7}

\begin{question}
האם האינטגרל הלא אמיתי
\[ \int_0^1\frac{\sin x}{x^2}\,dx \]
מתכנס?
\end{question}

\begin{hint}
ליד $0$ מתקיים $\sin x\approx x$. השוו (מבחן ההשוואה הגבולי) עם $\frac1x$.
\end{hint}

\begin{solution}
\textbf{האינטגרל מתבדר.}
האינטגרנד רציף ב-$(0,1]$, והבעיה היחידה היא בנקודה $0$. ב-$(0,1]$ מתקיים $\sin x>0$ (כי $0<x\le1<\pi$), כך שהאינטגרנד חיובי ואפשר להשתמש במבחני ההשוואה.

נשווה עם $g(x)=\frac1x>0$:
\[ \lim_{x\to0^+}\frac{\sin x/x^2}{1/x}=\lim_{x\to0^+}\frac{\sin x}{x}=1\in(0,\infty). \]
לפי מבחן ההשוואה הגבולי, $\int_0^1\frac{\sin x}{x^2}dx$ ו-$\int_0^1\frac{dx}{x}$ מתכנסים או מתבדרים יחד. אבל
\[ \int_\eps^1\frac{dx}{x}=-\ln\eps\xrightarrow[\eps\to0^+]{}\infty, \]
כלומר $\int_0^1\frac{dx}{x}$ מתבדר. לכן גם $\int_0^1\frac{\sin x}{x^2}dx$ \textbf{מתבדר}.

(השוואה ישירה: מכיוון ש-$\frac{\sin x}{x}\to1$, קיים $\delta>0$ כך ש-$\frac{\sin x}{x}>\frac12$ ב-$(0,\delta)$, ולכן $\frac{\sin x}{x^2}>\frac{1}{2x}$ שם.)
\end{solution}

\subsection*{שאלה 8}

\begin{question}
יהיו $f(x),g(x)$ פונקציות המוגדרות ורציפות בקרן $[0,+\infty)$. נתון שהאינטגרל הלא אמיתי $\int_0^\infty(f(x)+g(x))\,dx$ מתכנס.
הוכיחו או הפריכו:
האינטגרל הלא אמיתי $\int_0^\infty(f(x)-g(x))\,dx$ מתכנס.
\end{question}

\begin{hint}
נסו פונקציות שסכומן אפס.
\end{hint}

\begin{solution}
\textbf{הטענה אינה נכונה.} דוגמה נגדית: $f(x)=1$, $g(x)=-1$ (שתיהן רציפות ב-$[0,\infty)$).
\begin{itemize}
  \item $f+g\equiv0$, ולכן $\int_0^\infty(f+g)\,dx=\lim_{R\to\infty}\int_0^R0\,dx=0$ -- מתכנס.
  \item $f-g\equiv2$, ולכן $\int_0^R(f-g)\,dx=2R\xrightarrow[R\to\infty]{}\infty$ -- מתבדר.
\end{itemize}
לכן התכנסות $\int_0^\infty(f+g)$ אינה גוררת התכנסות $\int_0^\infty(f-g)$.

(הערה: אם \emph{שני} האינטגרלים $\int_0^\infty f$ ו-$\int_0^\infty g$ היו מתכנסים, אז לפי לינאריות גם $\int_0^\infty(f-g)$ היה מתכנס. כאן נתונה רק התכנסות הסכום.)
\end{solution}

\section*{תשע"ז סמסטר ב מועד ב}

\subsection*{שאלה 1}

\begin{question}
מצאו מספר רציונלי $x=\frac ab$ כך ש-$|\cos(1)-x|<\frac{1}{1000}$.
\end{question}

\begin{hint}
השתמשו בפולינום מקלורין של $\cos$ ובנוסחת השארית של לגרנז', $|R_n(1)|\le\frac{1}{(n+1)!}$. בחרו $n$ כך ש-$(n+1)!>1000$.
\end{hint}

\begin{hint}
שימו לב: בפיתוח של $\cos$ המקדם של $x^{2k+1}$ הוא אפס, ולכן $P_{2k}=P_{2k+1}$ ואפשר להשתמש בשארית מסדר $2k+1$.
\end{hint}

\begin{solution}
פולינום מקלורין של $\cos x$ הוא בעל מקדמים רציונליים, ולכן ערכו ב-$x=1$ רציונלי. נותר לבחור סדר שבו השארית קטנה מ-$\frac1{1000}$.

לפי משפט טיילור עם שארית לגרנז', לכל $n$ קיימת $c$ בין $0$ ל-$1$ כך ש-
\[ \cos1=P_n(1)+R_n(1),\qquad R_n(1)=\frac{\cos^{(n+1)}(c)}{(n+1)!}\cdot1^{n+1}. \]
כל נגזרת של $\cos$ היא $\pm\sin$ או $\pm\cos$, ולכן $|\cos^{(n+1)}(c)|\le1$ ו-$|R_n(1)|\le\frac1{(n+1)!}$.

נבחר $n=7$: $|R_7(1)|\le\frac1{8!}=\frac1{40320}<\frac1{1000}$. מכיוון שמקדם $x^7$ בפיתוח של $\cos$ הוא $0$:
\[ P_7(1)=P_6(1)=1-\frac1{2!}+\frac1{4!}-\frac1{6!}=1-\frac12+\frac1{24}-\frac1{720}=\frac{720-360+30-1}{720}=\frac{389}{720}. \]
לכן
\[ \left|\cos1-\frac{389}{720}\right|\le\frac1{40320}<\frac1{1000}, \]
והמספר הרציונלי $\boxed{x=\frac{389}{720}}$ מקיים את הדרישה. (מספרית $\cos1\approx0.5403023$, $\frac{389}{720}\approx0.5402778$.)

\textbf{הערה:} הקירוב הקודם $1-\frac12+\frac1{24}=\frac{13}{24}$ \emph{אינו} מספיק: החסם שנותן לגרנז' הוא $\frac1{6!}=\frac1{720}>\frac1{1000}$, ואכן $|\cos1-\frac{13}{24}|\approx0.00136>\frac1{1000}$.
\end{solution}

\subsection*{שאלה 2}

\begin{question}
עבור איזה ערך של הפרמטר $A$ הפונקציה $f(x)$ רציפה בנקודה $0$?
\[ f(x)=\begin{cases}\dfrac{1-\cos^2(x)}{2x^2} & x\neq0\\ A & x=0\end{cases} \]
\end{question}

\begin{solution}
$f$ רציפה ב-$0$ אם ורק אם $\lim_{x\to0}f(x)=f(0)=A$. לפי הזהות $1-\cos^2x=\sin^2x$:
\[ \lim_{x\to0}\frac{1-\cos^2x}{2x^2}=\lim_{x\to0}\frac12\left(\frac{\sin x}{x}\right)^2=\frac12\cdot1^2=\frac12, \]
לפי הגבול היסודי $\frac{\sin x}{x}\to1$ ואריתמטיקה של גבולות. לכן $f$ רציפה ב-$0$ אם ורק אם $\boxed{A=\frac12}$.
\end{solution}

\subsection*{שאלה 3}

\begin{question}
חשבו:
\[ \int e^{\sqrt[3]{x}}\,dx \]
\end{question}

\begin{hint}
הציבו $t=\sqrt[3]x$, כלומר $x=t^3$, ואז בצעו אינטגרציה בחלקים פעמיים.
\end{hint}

\begin{solution}
נציב $t=\sqrt[3]x$, כלומר $x=t^3$ ו-$dx=3t^2\,dt$ (ההצבה הפיכה וגזירה על כל $\R$):
\[ \int e^{\sqrt[3]x}dx=3\int t^2e^t\,dt. \]
אינטגרציה בחלקים ($u=t^2$, $v'=e^t$):
\[ \int t^2e^tdt=t^2e^t-2\int te^tdt, \]
ושוב בחלקים ($u=t$, $v'=e^t$): $\int te^tdt=te^t-e^t+C$. לכן
\[ \int t^2e^tdt=t^2e^t-2te^t+2e^t+C=e^t(t^2-2t+2)+C. \]
נחזור ל-$x$:
\[ \boxed{\int e^{\sqrt[3]x}dx=3e^{\sqrt[3]x}\left(\sqrt[3]{x^2}-2\sqrt[3]x+2\right)+C.} \]
\textbf{בדיקה:} $\frac{d}{dt}\left[3e^t(t^2-2t+2)\right]=3e^t(t^2-2t+2+2t-2)=3t^2e^t$, וזה בדיוק האינטגרנד אחרי ההצבה.
\end{solution}

\subsection*{שאלה 4}

\begin{question}
חשבו:
\[ \int\frac{x^2}{\sqrt{4-x^2}}\,dx \]
\end{question}

\begin{hint}
הצבה טריגונומטרית: $x=2\sin t$ עם $t\in(-\frac\pi2,\frac\pi2)$.
\end{hint}

\begin{solution}
האינטגרנד מוגדר עבור $-2<x<2$. נציב $x=2\sin t$, $t\in(-\frac\pi2,\frac\pi2)$ (הצבה הפיכה: $t=\arcsin\frac x2$), $dx=2\cos t\,dt$. בתחום זה $\cos t>0$, ולכן $\sqrt{4-x^2}=\sqrt{4\cos^2t}=2\cos t$. מכאן
\[ \int\frac{x^2}{\sqrt{4-x^2}}dx=\int\frac{4\sin^2t}{2\cos t}\cdot2\cos t\,dt=4\int\sin^2t\,dt=2\int(1-\cos2t)\,dt=2t-\sin2t+C. \]
נחזור ל-$x$: $t=\arcsin\frac x2$, ו-$\sin2t=2\sin t\cos t=2\cdot\frac x2\cdot\frac{\sqrt{4-x^2}}{2}=\frac{x\sqrt{4-x^2}}{2}$. לכן
\[ \boxed{\int\frac{x^2}{\sqrt{4-x^2}}dx=2\arcsin\frac x2-\frac{x\sqrt{4-x^2}}{2}+C.} \]
(בדיקה: גזירת התוצאה נותנת $\frac{2}{\sqrt{4-x^2}}-\frac{\sqrt{4-x^2}}{2}+\frac{x^2}{2\sqrt{4-x^2}}=\frac{4-(4-x^2)+x^2}{2\sqrt{4-x^2}}=\frac{x^2}{\sqrt{4-x^2}}$.)
\end{solution}

\subsection*{שאלה 5}

\begin{question}
מה השטח המקסימלי של משולש שווה שוקיים שהקפו 1?
\end{question}

\begin{hint}
סמנו את אורך השוק $a$ ואת הבסיס $b=1-2a$. הגובה לבסיס הוא $\sqrt{a^2-\frac{b^2}{4}}$. מה התחום המותר של $a$?
\end{hint}

\begin{hint}
נוח יותר למקסם את ריבוע השטח.
\end{hint}

\begin{solution}
נסמן את אורך השוק ב-$a$ ואת אורך הבסיס ב-$b$, עם $2a+b=1$, כלומר $b=1-2a$. כדי שיהיה משולש (לא מנוון) צריך $b>0$ ואי-שוויון המשולש $2a>b$, כלומר $\frac14<a<\frac12$.

הגובה לבסיס (שהוא גם תיכון) הוא, לפי פיתגורס, $h=\sqrt{a^2-\frac{b^2}{4}}$, ולכן
\[ S=\frac12bh=\frac b4\sqrt{4a^2-b^2}=\frac{1-2a}{4}\sqrt{(2a-b)(2a+b)}=\frac{1-2a}{4}\sqrt{4a-1}, \]
כי $2a+b=1$ ו-$2a-b=4a-1$.

נמקסם את $g(a)=16S^2=(1-2a)^2(4a-1)$ ב-$(\frac14,\frac12)$ (מכיוון ש-$S>0$, $S$ ו-$S^2$ מקבלות מקסימום באותה נקודה):
\[ g'(a)=-4(1-2a)(4a-1)+4(1-2a)^2=4(1-2a)\big[(1-2a)-(4a-1)\big]=4(1-2a)(2-6a). \]
בתחום $1-2a>0$, ולכן $g'(a)=0\iff a=\frac13$; $g'>0$ ב-$(\frac14,\frac13)$ ו-$g'<0$ ב-$(\frac13,\frac12)$. לכן $a=\frac13$ היא נקודת מקסימום מוחלט בתחום (בקצוות $g\to0$).

עבור $a=\frac13$: $b=\frac13$ -- המשולש שווה צלעות, ו-
\[ S_{\max}=\frac{1/3}{4}\sqrt{\frac13}=\frac{1}{12\sqrt3}=\boxed{\frac{\sqrt3}{36}}\approx0.0481. \]
\end{solution}

\subsection*{שאלה 6}

\begin{question}
מצאו מינימום ומקסימום מוחלטים של הפונקציה $f(x)=\frac{\sqrt{x^2-1}}{x^2}$.
\end{question}

\begin{hint}
מצאו קודם את תחום ההגדרה, ושימו לב ש-$f$ זוגית. בחנו גם את הגבול ב-$\infty$ ואת ערכי $f$ בקצות התחום.
\end{hint}

\begin{solution}
\textbf{תחום הגדרה:} $x^2-1\ge0$ ו-$x\ne0$, כלומר $|x|\ge1$. $f$ זוגית ($f(-x)=f(x)$), ולכן מספיק לחקור את $x\ge1$.

\textbf{מינימום:} $f(x)\ge0$ בכל התחום, ו-$f(\pm1)=0$. לכן המינימום המוחלט הוא $\boxed{0}$, המתקבל ב-$x=\pm1$.

\textbf{מקסימום:} עבור $x>1$:
\[ f'(x)=\frac{\frac{x}{\sqrt{x^2-1}}\cdot x^2-2x\sqrt{x^2-1}}{x^4}=\frac{x^2-2(x^2-1)}{x^3\sqrt{x^2-1}}=\frac{2-x^2}{x^3\sqrt{x^2-1}}. \]
לכן $f'>0$ ב-$(1,\sqrt2)$ ו-$f'<0$ ב-$(\sqrt2,\infty)$: $f$ עולה ב-$[1,\sqrt2]$ ויורדת ב-$[\sqrt2,\infty)$ (בעזרת רציפות ב-$1$). לכן לכל $x\ge1$ מתקיים $f(x)\le f(\sqrt2)=\frac{\sqrt{2-1}}{2}=\frac12$. (בנוסף $\lim_{x\to\infty}f(x)=\lim\frac{\sqrt{1-1/x^2}}{x}=0$.)
מהזוגיות, המקסימום המוחלט הוא $\boxed{\frac12}$, המתקבל ב-$x=\pm\sqrt2$.
\end{solution}

\subsection*{שאלה 7}

\begin{question}
האם האינטגרל הלא אמיתי
\[ \int_2^\infty\frac{x+7}{\sqrt{x^5-3x}}\,dx \]
מתכנס?
\end{question}

\begin{hint}
עבור $x$ גדול האינטגרנד מתנהג כמו $\frac{x}{x^{5/2}}=\frac{1}{x^{3/2}}$. השתמשו במבחן ההשוואה הגבולי.
\end{hint}

\begin{solution}
\textbf{האינטגרל מתכנס.}
עבור $x\ge2$: $x^5-3x=x(x^4-3)\ge2\cdot13>0$, ולכן האינטגרנד מוגדר, רציף וחיובי ב-$[2,\infty)$, והבעיה היחידה היא בקצה $\infty$.

נשווה עם $g(x)=\frac{1}{x^{3/2}}>0$:
\[ \lim_{x\to\infty}\frac{\frac{x+7}{\sqrt{x^5-3x}}}{x^{-3/2}}=\lim_{x\to\infty}\frac{(x+7)x^{3/2}}{\sqrt{x^5-3x}}=\lim_{x\to\infty}\frac{x^{5/2}\left(1+\frac7x\right)}{x^{5/2}\sqrt{1-\frac{3}{x^4}}}=1\in(0,\infty). \]
לפי מבחן ההשוואה הגבולי, שני האינטגרלים $\int_2^\infty\frac{x+7}{\sqrt{x^5-3x}}dx$ ו-$\int_2^\infty\frac{dx}{x^{3/2}}$ מתכנסים או מתבדרים יחד. האינטגרל $\int_2^\infty\frac{dx}{x^p}$ מתכנס עבור $p>1$, וכאן $p=\frac32$:
\[ \int_2^R x^{-3/2}dx=\left[-2x^{-1/2}\right]_2^R=\frac{2}{\sqrt2}-\frac{2}{\sqrt R}\xrightarrow[R\to\infty]{}\sqrt2. \]
לכן האינטגרל הנתון \textbf{מתכנס}.
\end{solution}

\subsection*{שאלה 8}

\begin{question}
הוכיחו או הפריכו: אם הפונקציה $|f(x)|$ אינטגרבילית לפי רימן בקטע $[a,b]$ אז גם הפונקציה $f(x)$ אינטגרבילית בקטע זה.
\end{question}

\begin{hint}
חפשו פונקציה שמקבלת רק את הערכים $1$ ו-$-1$ באופן "צפוף", בדומה לפונקציית דירכלה.
\end{hint}

\begin{solution}
\textbf{הטענה אינה נכונה.} (הכיוון ההפוך נכון: אם $f$ אינטגרבילית אז גם $|f|$.)

דוגמה נגדית (עבור $a<b$):
\[ f(x)=\begin{cases}1 & x\in\Q\\ -1 & x\notin\Q\end{cases}\qquad x\in[a,b]. \]
אז $|f(x)|=1$ לכל $x$, פונקציה קבועה, ולכן אינטגרבילית ו-$\int_a^b|f|=b-a$.

לעומת זאת $f$ אינה אינטגרבילית לפי רימן: תהי $P$ חלוקה כלשהי של $[a,b]$. בכל תת-קטע $[x_{i-1},x_i]$ (באורך חיובי) יש גם מספר רציונלי וגם מספר אי-רציונלי (צפיפות $\Q$ ו-$\R\setminus\Q$ ב-$\R$), ולכן $M_i=\sup f=1$ ו-$m_i=\inf f=-1$. מכאן
\[ U(f,P)=\sum M_i\Delta x_i=b-a,\qquad L(f,P)=\sum m_i\Delta x_i=-(b-a), \]
כך ש-$U(f,P)-L(f,P)=2(b-a)$ לכל חלוקה, ואינו יכול להיות קטן מ-$\eps$ כרצוננו. לפי קריטריון רימן (או: האינטגרל העליון $b-a$ שונה מהתחתון $-(b-a)$), $f$ \textbf{אינה} אינטגרבילית.
\end{solution}

\section*{תשע"ח סמסטר א מועד א}

\subsection*{שאלה 1א}

\begin{question}
חשבו את הגבול הבא:
גבול הסדרה $\displaystyle \lim_{n\to\infty}\left(\frac{n}{n+1}\right)^{3n}$.
\end{question}

\begin{hint}
הפכו את השבר: $\frac{n}{n+1}=\frac{1}{1+\frac1n}$, והשתמשו בגבול המפורסם $\left(1+\frac1n\right)^n\to e$.
\end{hint}

\begin{solution}
לכל $n$:
\[ \left(\frac{n}{n+1}\right)^{3n}=\frac{1}{\left(\frac{n+1}{n}\right)^{3n}}=\frac{1}{\left[\left(1+\frac1n\right)^{n}\right]^{3}}. \]
ידוע ש-$\left(1+\frac1n\right)^n\xrightarrow[n\to\infty]{} e$. מאריתמטיקה של גבולות (מכפלה של שלוש סדרות מתכנסות, ומנה כשהגבול במכנה $e^3\neq 0$) נקבל
\[ \lim_{n\to\infty}\left(\frac{n}{n+1}\right)^{3n}=\frac{1}{e^3}. \]
\end{solution}

\subsection*{שאלה 1ב}

\begin{question}
חשבו את הגבול הבא:
גבול הפונקציה $\displaystyle \lim_{x\to\pi/2}(\tan x)\cdot\arctan\left(x-\frac{\pi}{2}\right)$.
\end{question}

\begin{hint}
כתבו $\tan x=\frac{1}{\cot x}$ וקבלו ביטוי מהצורה $\frac00$.
\end{hint}

\begin{solution}
עבור $x$ קרוב ל-$\frac\pi2$, $x\neq\frac\pi2$, מתקיים $\tan x=\frac{1}{\cot x}$, ולכן
\[ (\tan x)\cdot\arctan\left(x-\frac{\pi}{2}\right)=\frac{\arctan\left(x-\frac{\pi}{2}\right)}{\cot x}. \]
כאשר $x\to\frac\pi2$: המונה שואף ל-$\arctan 0=0$ והמכנה שואף ל-$\cot\frac\pi2=0$ (שתי הפונקציות רציפות). זהו ביטוי מהצורה $\frac00$, שתי הפונקציות גזירות בסביבה מנוקבת של $\frac\pi2$ ונגזרת המכנה $-\frac{1}{\sin^2x}\neq 0$ שם, ולכן ניתן להפעיל את כלל לופיטל:
\[ \lim_{x\to\pi/2}\frac{\arctan\left(x-\frac{\pi}{2}\right)}{\cot x}
   =\lim_{x\to\pi/2}\frac{\dfrac{1}{1+\left(x-\frac\pi2\right)^2}}{-\dfrac{1}{\sin^2x}}
   =\lim_{x\to\pi/2}\left(-\frac{\sin^2x}{1+\left(x-\frac\pi2\right)^2}\right)=-\frac{1}{1+0}=-1, \]
כאשר הגבול האחרון מחושב מרציפות ($\sin^2\frac\pi2=1$). מכיוון שגבול מנת הנגזרות קיים, לפי כלל לופיטל
\[ \lim_{x\to\pi/2}(\tan x)\cdot\arctan\left(x-\frac{\pi}{2}\right)=-1. \]
(דרך נוספת: $\arctan t\sim t$ כאשר $t\to0$, ו-$(x-\frac\pi2)\tan x=-\frac{(x-\frac\pi2)\cos(x-\frac\pi2)}{\sin(x-\frac\pi2)}\to-1$.)
\end{solution}

\subsection*{שאלה 2}

\begin{question}
מצאו את נקודות אי-הרציפות של הפונקציה
\[ f(x)=\frac{e^{\left(\frac1x-\frac{1}{x+1}\right)}}{e^{\frac1x}-2} \]
וקבעו את סוגיהן.
\end{question}

\begin{hint}
הפונקציה אינה מוגדרת בשלוש נקודות: היכן ש-$x=0$, $x=-1$ ושם $e^{1/x}=2$.
\end{hint}

\begin{hint}
ב-$x\to0^+$ חלקו מונה ומכנה ב-$e^{1/x}$.
\end{hint}

\begin{solution}
\textbf{תחום ההגדרה.} הביטוי מוגדר כאשר $x\neq0$, $x\neq-1$ ו-$e^{1/x}\neq 2$, כלומר $\frac1x\neq\ln2$, כלומר $x\neq\frac1{\ln2}$. בכל נקודה אחרת $f$ היא הרכבה ומנה של פונקציות אלמנטריות רציפות (עם מכנה שונה מאפס), ולכן רציפה. נקודות אי-הרציפות הן לכן $x=0,\ x=-1,\ x=\frac{1}{\ln2}$. נבדוק כל אחת.

\textbf{1. $x=\frac1{\ln2}$.} המונה רציף בנקודה זו וערכו
$e^{\ln2-\frac{1}{\frac1{\ln2}+1}}>0$ (אקספוננט תמיד חיובי). המכנה $e^{1/x}-2$ שואף ל-$0$. כאשר $x\to\left(\frac1{\ln2}\right)^+$ מתקיים $\frac1x<\ln2$, לכן $e^{1/x}-2\to0^-$, ו-
\[ \lim_{x\to(1/\ln2)^+}f(x)=-\infty, \qquad \lim_{x\to(1/\ln2)^-}f(x)=+\infty \]
(משמאל $\frac1x>\ln2$ והמכנה שואף ל-$0^+$). הגבולות החד-צדדיים אינם סופיים, ולכן זו \textbf{נקודת אי-רציפות מסוג שני}.

\textbf{2. $x=-1$.} כאשר $x\to-1$, המכנה רציף ושואף ל-$e^{-1}-2<0$ (מספר שלילי השונה מאפס). נבדוק את המונה:
\begin{itemize}
  \item כאשר $x\to-1^+$: $x+1\to0^+$, לכן $\frac1x-\frac1{x+1}\to-1-\infty=-\infty$, ומכאן $e^{\left(\frac1x-\frac1{x+1}\right)}\to0$. לכן $\lim_{x\to-1^+}f(x)=\frac{0}{e^{-1}-2}=0$.
  \item כאשר $x\to-1^-$: $x+1\to0^-$, לכן $\frac1x-\frac1{x+1}\to+\infty$ והמונה שואף ל-$+\infty$. מכיוון שהמכנה שואף למספר שלילי, $\lim_{x\to-1^-}f(x)=-\infty$.
\end{itemize}
אחד הגבולות החד-צדדיים אינסופי, ולכן זו \textbf{נקודת אי-רציפות מסוג שני}.

\textbf{3. $x=0$.}
\begin{itemize}
  \item כאשר $x\to0^+$: $\frac1x\to+\infty$ ולכן גם המונה וגם המכנה שואפים ל-$\infty$. נחלק מונה ומכנה ב-$e^{1/x}$:
  \[ f(x)=\frac{e^{\frac1x}\cdot e^{-\frac1{x+1}}}{e^{\frac1x}\left(1-2e^{-\frac1x}\right)}=\frac{e^{-\frac{1}{x+1}}}{1-2e^{-\frac1x}}\xrightarrow[x\to0^+]{}\frac{e^{-1}}{1-0}=\frac1e, \]
  כי $e^{-1/x}\to0$ כאשר $\frac1x\to+\infty$. (בפתרון הרשמי הגבול חושב בכלל לופיטל ויצא אותו ערך.)
  \item כאשר $x\to0^-$: $\frac1x\to-\infty$, לכן $e^{1/x}\to0$ והמכנה שואף ל-$-2$; כמו כן $\frac1x-\frac1{x+1}\to-\infty-1=-\infty$, לכן המונה שואף ל-$0$. מכאן $\lim_{x\to0^-}f(x)=\frac{0}{-2}=0$.
\end{itemize}
שני הגבולות החד-צדדיים קיימים וסופיים אך שונים ($\frac1e\neq0$), ולכן זו \textbf{נקודת אי-רציפות מסוג ראשון (קפיצה)}.

\textbf{סיכום:} $x=0$ — אי-רציפות מסוג ראשון (קפיצה); $x=-1$ ו-$x=\frac1{\ln2}$ — אי-רציפות מסוג שני.
\end{solution}

\subsection*{שאלה 3א}

\begin{question}
חשבו את הגבול:
\[ \lim_{x\to0}\left(\frac{\sqrt{1-2x+x^3}-\sqrt[3]{1-3x+x^2}}{x^2}\right) \]
\underline{רמז}: אפשר להשתמש בנוסחאות מקלורין.
\end{question}

\begin{hint}
פתחו כל שורש לפי $(1+t)^p=1+pt+\frac{p(p-1)}{2}t^2+o(t^2)$ עד סדר $x^2$.
\end{hint}

\begin{solution}
נשתמש בפיתוח מקלורין $(1+t)^p=1+pt+\frac{p(p-1)}{2!}t^2+o(t^2)$ כאשר $t\to0$.

עבור השורש הראשון, $t=x^3-2x\to0$ ו-$p=\frac12$:
\begin{align*}
\sqrt{1-2x+x^3}&=1+\frac12(x^3-2x)-\frac18(x^3-2x)^2+o(x^2)\\
&=1-x+\frac{x^3}{2}-\frac18\left(4x^2-4x^4+x^6\right)+o(x^2)=1-x-\frac{x^2}{2}+o(x^2).
\end{align*}
(השתמשנו בכך ש-$t=O(x)$ ולכן $o(t^2)=o(x^2)$, ובכך שהחזקות $x^3,x^4,x^6$ הן $o(x^2)$.)

עבור השורש השני, $t=-3x+x^2$ ו-$p=\frac13$, $\frac{p(p-1)}{2}=\frac{\frac13\cdot(-\frac23)}{2}=-\frac19$:
\begin{align*}
\sqrt[3]{1-3x+x^2}&=1+\frac13(-3x+x^2)-\frac19(-3x+x^2)^2+o(x^2)\\
&=1-x+\frac{x^2}{3}-\frac19\left(9x^2-6x^3+x^4\right)+o(x^2)=1-x-\frac{2x^2}{3}+o(x^2).
\end{align*}
לכן
\[ \frac{\sqrt{1-2x+x^3}-\sqrt[3]{1-3x+x^2}}{x^2}=\frac{\left(-\frac12+\frac23\right)x^2+o(x^2)}{x^2}=\frac16+\frac{o(x^2)}{x^2}\xrightarrow[x\to0]{}\frac16. \]
\[ \boxed{\lim_{x\to0}\frac{\sqrt{1-2x+x^3}-\sqrt[3]{1-3x+x^2}}{x^2}=\frac16} \]
\end{solution}

\subsection*{שאלה 3ב}

\begin{question}
הוכיחו כי $0<x<\arcsin x$ לכל $x\in(0,1]$.
\end{question}

\begin{hint}
התבוננו בפונקציה $g(x)=\arcsin x-x$ ובנגזרת שלה.
\end{hint}

\begin{solution}
אי-השוויון $0<x$ ברור עבור $x\in(0,1]$. נותר להוכיח $x<\arcsin x$.

נגדיר $g(x)=\arcsin x-x$ בקטע $[0,1]$. הפונקציה רציפה ב-$[0,1]$ וגזירה ב-$(0,1)$, עם
\[ g'(x)=\frac{1}{\sqrt{1-x^2}}-1>0 \quad\text{לכל } x\in(0,1), \]
כי $0<\sqrt{1-x^2}<1$ שם.

יהי $b\in(0,1]$. לפי משפט לגרנז' עבור $g$ בקטע $[0,b]$ קיימת נקודה $c\in(0,b)$ כך ש-
\[ g(b)-g(0)=g'(c)\,(b-0). \]
כיוון ש-$g(0)=\arcsin0-0=0$, $g'(c)>0$ ו-$b>0$, נקבל $g(b)>0$, כלומר $\arcsin b>b$.

לכן $0<x<\arcsin x$ לכל $x\in(0,1]$. $\blacksquare$

(הפתרון הרשמי מנוסח בדרך השלילה: מניחים $\arcsin b\le b$, ומלגרנז' מקבלים $\frac{1}{\sqrt{1-c^2}}=\frac{\arcsin b}{b}\le1$, בסתירה ל-$\frac{1}{\sqrt{1-c^2}}>1$. זו אותה הוכחה.)
\end{solution}

\subsection*{שאלה 4}

\begin{question}
חקרו את הפונקציה $f(x)=(x+2)e^{1/x}$ ושרטטו סקיצה של הגרף שלה.
\end{question}

\begin{hint}
לאסימפטוטה המשופעת: $a=\lim\frac{f(x)}{x}$, ואז $b=\lim(f(x)-ax)$. בחישוב $b$ הופיע הגבול $\lim_{x\to\infty}x\left(e^{1/x}-1\right)$.
\end{hint}

\begin{solution}
\textbf{1. תיאור כללי.}
\begin{itemize}
  \item \textbf{תחום הגדרה:} $x\neq0$.
  \item \textbf{נקודות חיתוך עם הצירים:} אין חיתוך עם ציר $y$ כי $0$ אינו בתחום. $f(x)=0\iff x+2=0$ (כי $e^{1/x}>0$), ולכן החיתוך עם ציר $x$ הוא $(-2,0)$.
  \item \textbf{זוגיות:} $f(-1)=e^{-1}$, $f(1)=3e$, ולכן $f$ אינה זוגית ואינה אי-זוגית. היא גם אינה מחזורית.
  \item \textbf{סימן:} כיוון ש-$e^{1/x}>0$, הסימן של $f$ הוא סימן $x+2$: $f(x)>0$ עבור $x>-2$ ($x\neq0$), ו-$f(x)<0$ עבור $x<-2$.
\end{itemize}

\textbf{2. רציפות.} $f$ אלמנטרית ולכן רציפה בכל תחום הגדרתה. בנקודה $x=0$:
\[ \lim_{x\to0^-}(x+2)e^{1/x}=2\cdot0=0,\qquad \lim_{x\to0^+}(x+2)e^{1/x}=2\cdot(+\infty)=+\infty, \]
כי $\frac1x\to\mp\infty$. לכן ב-$x=0$ יש נקודת אי-רציפות מסוג שני.

\textbf{3. אסימפטוטות.}
\begin{itemize}
  \item \textbf{אנכית:} $x=0$ (מימין), כי $\lim_{x\to0^+}f(x)=+\infty$. משמאל הגרף שואף לנקודה $(0,0)$ (שאינה בגרף).
  \item \textbf{משופעת} $y=ax+b$ כאשר $x\to\pm\infty$:
  \[ a=\lim_{x\to\pm\infty}\frac{f(x)}{x}=\lim_{x\to\pm\infty}\left(1+\frac2x\right)e^{1/x}=1\cdot e^0=1. \]
  \[ b=\lim_{x\to\pm\infty}\left(f(x)-x\right)=\lim_{x\to\pm\infty}\left(x\left(e^{1/x}-1\right)+2e^{1/x}\right). \]
  נציב $t=\frac1x\to0$: $x\left(e^{1/x}-1\right)=\frac{e^t-1}{t}\to1$ (גבול מפורסם, או לופיטל: $\frac{e^t}{1}\to1$). כמו כן $2e^{1/x}\to2$. לכן $b=1+2=3$.

  האסימפטוטה המשופעת היא $y=x+3$, גם כאשר $x\to\infty$ וגם כאשר $x\to-\infty$. בפרט $\lim_{x\to\pm\infty}f(x)=\pm\infty$.
\end{itemize}

\textbf{4. עלייה, ירידה וקיצון.}
\[ f'(x)=e^{1/x}+(x+2)e^{1/x}\cdot\left(-\frac{1}{x^2}\right)=e^{1/x}\left(1-\frac{x+2}{x^2}\right)=e^{1/x}\cdot\frac{x^2-x-2}{x^2}=e^{1/x}\cdot\frac{(x-2)(x+1)}{x^2}. \]
כיוון ש-$e^{1/x}>0$ ו-$x^2>0$, סימן $f'$ הוא סימן $(x-2)(x+1)$:
\begin{itemize}
  \item $f'>0$ ב-$(-\infty,-1)$ וב-$(2,\infty)$ — $f$ עולה שם;
  \item $f'<0$ ב-$(-1,0)$ וב-$(0,2)$ — $f$ יורדת שם.
\end{itemize}
לכן ב-$x=-1$ יש מקסימום מקומי, $f(-1)=1\cdot e^{-1}=\frac1e$, וב-$x=2$ יש מינימום מקומי, $f(2)=4e^{1/2}=4\sqrt e$.

\textbf{5. קמירות ונקודות פיתול.} גוזרים את $f'(x)=e^{1/x}\cdot\frac{x^2-x-2}{x^2}=e^{1/x}\left(1-\frac1x-\frac2{x^2}\right)$:
\[ f''(x)=e^{1/x}\left(-\frac1{x^2}\right)\left(1-\frac1x-\frac2{x^2}\right)+e^{1/x}\left(\frac1{x^2}+\frac4{x^3}\right)
=e^{1/x}\cdot\frac{-x^2+x+2+x^2+4x}{x^4}=e^{1/x}\cdot\frac{5x+2}{x^4}. \]
סימן $f''$ הוא סימן $5x+2$. לכן $f$ קמורה (כלפי מעלה, $f''>0$) ב-$\left(-\frac25,0\right)$ וב-$(0,\infty)$, וקעורה (כלפי מטה) ב-$\left(-\infty,-\frac25\right)$. ב-$x=-\frac25$ הנגזרת השנייה מחליפה סימן ו-$f$ רציפה, ולכן זו נקודת פיתול:
\[ \left(-\frac25,\ f\left(-\frac25\right)\right)=\left(-\frac25,\ \frac85e^{-5/2}\right). \]

\textbf{6. סקיצה.} משמאל לציר $y$: כאשר $x\to-\infty$ הגרף נמצא מתחת לאסימפטוטה $y=x+3$ ומתקרב אליה (כי $f(x)-(x+3)=\frac{5}{2x}+O\left(\frac1{x^2}\right)$), הוא עולה וקעור, חותך את ציר $x$ ב-$(-2,0)$, מגיע למקסימום מקומי $\left(-1,\frac1e\right)$, ואחר כך יורד; בנקודת הפיתול $x=-\frac25$ הוא הופך לקמור וממשיך לרדת אל הנקודה $(0,0)$ (שאינה שייכת לגרף). מימין לציר $y$: הגרף יורד מ-$+\infty$ (אסימפטוטה אנכית $x=0$) עד למינימום $(2,4\sqrt e)\approx(2,\,6.6)$, ואחר כך עולה ומתקרב מלמעלה לאסימפטוטה $y=x+3$; בתחום זה הוא קמור.
\end{solution}

\subsection*{שאלה 5א}

\begin{question}
חשבו את האינטגרל $\displaystyle\int\frac{3x-1}{x^2-x+1}\,dx$.
\end{question}

\begin{hint}
כתבו את המונה כצירוף של נגזרת המכנה $2x-1$ וקבוע, והשלימו את המכנה לריבוע.
\end{hint}

\begin{solution}
למכנה $x^2-x+1$ אין שורשים ממשיים (הדיסקרימיננטה $1-4=-3<0$), ולכן $x^2-x+1>0$ לכל $x$. נפרק את המונה לפי נגזרת המכנה $(x^2-x+1)'=2x-1$:
\[ 3x-1=\frac32(2x-1)+\frac12. \]
לכן
\[ \int\frac{3x-1}{x^2-x+1}\,dx=\frac32\int\frac{2x-1}{x^2-x+1}\,dx+\frac12\int\frac{dx}{\left(x-\frac12\right)^2+\frac34}. \]
באינטגרל הראשון נציב $u=x^2-x+1$, $du=(2x-1)\,dx$:
\[ \frac32\int\frac{du}{u}=\frac32\ln|u|=\frac32\ln\left(x^2-x+1\right). \]
בשני נציב $v=x-\frac12$, $dv=dx$, ונשתמש ב-$\int\frac{dv}{v^2+a^2}=\frac1a\arctan\frac va$ עם $a=\frac{\sqrt3}{2}$:
\[ \frac12\int\frac{dv}{v^2+\frac34}=\frac12\cdot\frac{2}{\sqrt3}\arctan\frac{2v}{\sqrt3}=\frac{1}{\sqrt3}\arctan\frac{2x-1}{\sqrt3}. \]
לסיכום
\[ \boxed{\int\frac{3x-1}{x^2-x+1}\,dx=\frac32\ln\left(x^2-x+1\right)+\frac{1}{\sqrt3}\arctan\frac{2x-1}{\sqrt3}+C} \]
(אין צורך בערך מוחלט בתוך ה-$\ln$ כי $x^2-x+1>0$.) בדיקה בגזירה: $\frac32\cdot\frac{2x-1}{x^2-x+1}+\frac1{\sqrt3}\cdot\frac{2/\sqrt3}{1+\frac{(2x-1)^2}{3}}=\frac{3x-\frac32}{x^2-x+1}+\frac{2}{4x^2-4x+4}=\frac{3x-1}{x^2-x+1}$.

\textbf{הערה:} בפתרון הרשמי המקדם של ה-$\arctan$ רשום $\frac{2}{\sqrt3}$; זו טעות (נשכח הגורם $\frac12$), והמקדם הנכון הוא $\frac1{\sqrt3}$.
\end{solution}

\subsection*{שאלה 5ב}

\begin{question}
האם האינטגרל הלא-אמיתי $\displaystyle\int_0^\infty\frac{x}{e^{x^2}}\,dx$ מתכנס? אם כן, יש לחשב את ערכו של האינטגרל. אם לא, יש לנמק מדוע.
\end{question}

\begin{hint}
הציבו $u=x^2$.
\end{hint}

\begin{solution}
לפי ההגדרה, $\int_0^\infty\frac{x}{e^{x^2}}\,dx=\lim_{b\to\infty}\int_0^b xe^{-x^2}\,dx$, אם הגבול קיים. נציב $u=x^2$, $du=2x\,dx$ (כאשר $x$ עובר מ-$0$ ל-$b$, $u$ עובר מ-$0$ ל-$b^2$):
\[ \int_0^b xe^{-x^2}\,dx=\frac12\int_0^{b^2}e^{-u}\,du=\frac12\left[-e^{-u}\right]_0^{b^2}=\frac12\left(1-e^{-b^2}\right). \]
כאשר $b\to\infty$, $e^{-b^2}\to0$, ולכן הגבול קיים וסופי:
\[ \boxed{\int_0^\infty\frac{x}{e^{x^2}}\,dx=\frac12} \]
כלומר האינטגרל מתכנס וערכו $\frac12$.
\end{solution}

\subsection*{שאלה 6א}

\begin{question}
חשבו את נפח גוף הסיבוב המתקבל ע"י סיבוב התחום $D$ סביב ציר ה-$x$, כאשר $D$ התחום בין גרף הפונקציה $f(x)=\tan x$ וציר ה-$x$, כאשר $0\le x\le\frac\pi4$.
\end{question}

\begin{hint}
השתמשו בזהות $\tan^2x=\frac{1}{\cos^2x}-1$.
\end{hint}

\begin{solution}
נפח גוף הסיבוב סביב ציר $x$ הוא $V=\pi\int_a^b f^2(x)\,dx$. לכן, בעזרת $\tan^2x=\frac1{\cos^2x}-1$:
\[ V=\pi\int_0^{\pi/4}\tan^2x\,dx=\pi\int_0^{\pi/4}\left(\frac1{\cos^2x}-1\right)dx=\pi\Big[\tan x-x\Big]_0^{\pi/4}=\pi\left(1-\frac\pi4\right). \]
\[ \boxed{V=\pi-\frac{\pi^2}{4}} \]
\end{solution}

\subsection*{שאלה 6ב}

\begin{question}
נגדיר $r=f(\theta)=\theta^2-1$ בקואורדינטות קטביות. חשבו את אורך העקומה הזאת בין $\theta=1$ ל-$\theta=9$.

(במבחן מצורף איור של העקומה — ספירלה — על רקע רשת קטבית.)
\end{question}

\begin{hint}
אורך עקומה בקואורדינטות קטביות הוא $\int_\alpha^\beta\sqrt{r^2+\left(\frac{dr}{d\theta}\right)^2}\,d\theta$. בדקו שהביטוי מתחת לשורש הוא ריבוע שלם.
\end{hint}

\begin{solution}
אורך עקומה בקואורדינטות קטביות $r=f(\theta)$, $\alpha\le\theta\le\beta$:
\[ L=\int_\alpha^\beta\sqrt{r^2+\left(\frac{dr}{d\theta}\right)^2}\,d\theta. \]
כאן $\frac{dr}{d\theta}=2\theta$, ולכן
\[ r^2+\left(\frac{dr}{d\theta}\right)^2=(\theta^2-1)^2+4\theta^2=\theta^4-2\theta^2+1+4\theta^2=\theta^4+2\theta^2+1=(\theta^2+1)^2. \]
מכיוון ש-$\theta^2+1>0$, $\sqrt{(\theta^2+1)^2}=\theta^2+1$, ו-
\[ L=\int_1^9(\theta^2+1)\,d\theta=\left[\frac{\theta^3}{3}+\theta\right]_1^9=\left(243+9\right)-\left(\frac13+1\right)=252-\frac43=\frac{752}{3}. \]
\[ \boxed{L=\frac{752}{3}=250\tfrac23} \]
\textbf{הערה:} בפתרון הרשמי נכתב $252-\frac13+1=252\frac23$ — טעות סימן בהצבת הגבול התחתון; הערך הנכון הוא $252-\left(\frac13+1\right)=250\frac23$.
\end{solution}

\section*{תשע"ח סמסטר א מועד ב}

\subsection*{שאלה 1א}

\begin{question}
חשבו את הגבול הבא:
גבול הסדרה $\displaystyle\lim_{n\to\infty}\left(\sqrt{(n^2+2)(n^2-4)}-\sqrt{n^4-7}\right)$.
\end{question}

\begin{hint}
הכפילו וחלקו בצמוד $\sqrt{(n^2+2)(n^2-4)}+\sqrt{n^4-7}$.
\end{hint}

\begin{solution}
לכל $n\ge2$ שני השורשים מוגדרים. נכפיל ונחלק בצמוד:
\[ \sqrt{(n^2+2)(n^2-4)}-\sqrt{n^4-7}=\frac{(n^2+2)(n^2-4)-(n^4-7)}{\sqrt{(n^2+2)(n^2-4)}+\sqrt{n^4-7}}. \]
במונה: $(n^2+2)(n^2-4)=n^4-2n^2-8$, ולכן המונה הוא $n^4-2n^2-8-n^4+7=-2n^2-1$. נחלק מונה ומכנה ב-$n^2$:
\[ \frac{-2n^2-1}{\sqrt{(n^2+2)(n^2-4)}+\sqrt{n^4-7}}=\frac{-2-\frac1{n^2}}{\sqrt{\left(1+\frac2{n^2}\right)\left(1-\frac4{n^2}\right)}+\sqrt{1-\frac7{n^4}}}. \]
מאריתמטיקה של גבולות ורציפות השורש, המונה שואף ל-$-2$ והמכנה ל-$1+1=2$. לכן
\[ \boxed{\lim_{n\to\infty}\left(\sqrt{(n^2+2)(n^2-4)}-\sqrt{n^4-7}\right)=-1} \]
\end{solution}

\subsection*{שאלה 1ב}

\begin{question}
חשבו את הגבול הבא:
גבול הפונקציה $\displaystyle\lim_{x\to1^-}\frac{\frac\pi2-\arcsin x}{\sqrt[3]{1-x^2}}$.
\end{question}

\begin{hint}
זהו ביטוי מהצורה $\frac00$; הפעילו את כלל לופיטל ופשטו את מנת הנגזרות.
\end{hint}

\begin{solution}
כאשר $x\to1^-$: $\arcsin x\to\frac\pi2$ (רציפות), ולכן המונה שואף ל-$0$; גם המכנה $\sqrt[3]{1-x^2}\to0$. זהו ביטוי מהצורה $\frac00$. בקטע $(0,1)$ שתי הפונקציות גזירות:
\[ \left(\frac\pi2-\arcsin x\right)'=-\frac{1}{\sqrt{1-x^2}},\qquad \left((1-x^2)^{1/3}\right)'=\frac13(1-x^2)^{-2/3}\cdot(-2x)=-\frac{2x}{3(1-x^2)^{2/3}}\neq0. \]
מנת הנגזרות:
\[ \frac{-\frac{1}{(1-x^2)^{1/2}}}{-\frac{2x}{3(1-x^2)^{2/3}}}=\frac{3(1-x^2)^{2/3}}{2x(1-x^2)^{1/2}}=\frac{3(1-x^2)^{1/6}}{2x}\xrightarrow[x\to1^-]{}\frac{3\cdot0}{2}=0. \]
גבול מנת הנגזרות קיים, ולכן לפי כלל לופיטל
\[ \boxed{\lim_{x\to1^-}\frac{\frac\pi2-\arcsin x}{\sqrt[3]{1-x^2}}=0} \]
(אינטואיציה: $\frac\pi2-\arcsin x=\arccos x\approx\sqrt{2(1-x)}$, שהוא "גדול פחות" מ-$\sqrt[3]{1-x^2}\approx\sqrt[3]{2(1-x)}$.)
\end{solution}

\subsection*{שאלה 2}

\begin{question}
מצאו את נקודות אי-הרציפות של הפונקציה
\[ f(x)=\frac{x\ln|x|}{x^2-1} \]
וקבעו את סוגיהן.
\end{question}

\begin{hint}
בנקודות $x=\pm1$ גם המונה מתאפס — השתמשו בגבול $\lim_{t\to0}\frac{\ln(1+t)}{t}=1$ או בלופיטל. ב-$x=0$ השתמשו בכך ש-$x\ln|x|\to0$.
\end{hint}

\begin{solution}
\textbf{תחום ההגדרה:} $\ln|x|$ מוגדר עבור $x\neq0$, והמכנה מתאפס עבור $x=\pm1$. לכן $f$ מוגדרת עבור $x\notin\{0,1,-1\}$, ושם היא רציפה כמנה של פונקציות רציפות עם מכנה שונה מאפס. נקודות אי-הרציפות הן $x=0,\,1,\,-1$.

\textbf{$x=0$:} ידוע ש-$\lim_{x\to0}x\ln|x|=0$ (למשל בלופיטל: $\lim_{x\to0^+}\frac{\ln x}{1/x}=\lim_{x\to0^+}\frac{1/x}{-1/x^2}=\lim_{x\to0^+}(-x)=0$, ומשמאל זה אותו דבר כי $x\ln|x|$ אי-זוגית). המכנה שואף ל-$-1$. לכן
\[ \lim_{x\to0}f(x)=\frac{0}{-1}=0. \]
הגבול קיים וסופי אך $f$ אינה מוגדרת ב-$0$: \textbf{אי-רציפות סליקה}.

\textbf{$x=1$:} בסביבת $1$, $|x|=x$, והביטוי מהצורה $\frac00$. נכתוב
\[ f(x)=\frac{x}{x+1}\cdot\frac{\ln x}{x-1}. \]
נציב $t=x-1\to0$: $\frac{\ln x}{x-1}=\frac{\ln(1+t)}{t}\to1$ (גבול מפורסם). לכן
\[ \lim_{x\to1}f(x)=\frac{1}{2}\cdot1=\frac12. \]
\textbf{אי-רציפות סליקה}.

\textbf{$x=-1$:} נשים לב ש-$f(-x)=\frac{-x\ln|x|}{x^2-1}=-f(x)$, כלומר $f$ אי-זוגית. לכן
\[ \lim_{x\to-1}f(x)=-\lim_{x\to1}f(x)=-\frac12. \]
(או ישירות: בסביבת $-1$, $f(x)=\frac{x}{x-1}\cdot\frac{\ln(-x)}{x+1}$, ועם $t=-x-1\to0$: $\frac{\ln(-x)}{x+1}=\frac{\ln(1+t)}{-t}\to-1$, ו-$\frac{x}{x-1}\to\frac12$.)
\textbf{אי-רציפות סליקה}.

\textbf{סיכום:} שלוש נקודות אי-הרציפות $x=0,1,-1$ כולן סליקות; ניתן להגדיר $f(0)=0$, $f(1)=\frac12$, $f(-1)=-\frac12$ ולקבל פונקציה רציפה על כל $\R$.
\end{solution}

\subsection*{שאלה 3א}

\begin{question}
תהי $f(x)=x^2\ln(1+x)$. חשבו את $f^{(11)}(0)$.

\underline{רמז}: אפשר להשתמש בנוסחאות מקלורין.
\end{question}

\begin{hint}
המקדם של $x^{11}$ בפולינום מקלורין של $f$ שווה ל-$\frac{f^{(11)}(0)}{11!}$. איזה איבר בפיתוח של $\ln(1+x)$ תורם ל-$x^{11}$ אחרי הכפלה ב-$x^2$?
\end{hint}

\begin{solution}
פיתוח מקלורין של $\ln(1+x)$:
\[ \ln(1+x)=x-\frac{x^2}{2}+\frac{x^3}{3}-\dots+(-1)^{n-1}\frac{x^n}{n}+o(x^n). \]
נכפיל ב-$x^2$ (עם $n=9$):
\[ f(x)=x^2\ln(1+x)=x^3-\frac{x^4}{2}+\frac{x^5}{3}-\dots+(-1)^{8}\frac{x^{11}}{9}+o(x^{11}). \]
זהו פולינום ממעלה $\le 11$ ועוד $o(x^{11})$, ולכן לפי יחידות פולינום טיילור הוא פולינום מקלורין מסדר $11$ של $f$ ($f$ גזירה אינסוף פעמים בסביבת $0$). המקדם של $x^{11}$ בפולינום מקלורין הוא $\frac{f^{(11)}(0)}{11!}$, ולכן
\[ \frac{f^{(11)}(0)}{11!}=\frac19\quad\Longrightarrow\quad f^{(11)}(0)=\frac{11!}{9}=\frac{39916800}{9}=4435200. \]
\[ \boxed{f^{(11)}(0)=\frac{11!}{9}=4435200} \]
\end{solution}

\subsection*{שאלה 3ב}

\begin{question}
הוכיחו כי יש לפולינום $x^3-4x^2+7x+13$ שורש ממשי אחד ויחיד.
\end{question}

\begin{hint}
קיום — משפט ערך הביניים. יחידות — הראו שהנגזרת חיובית תמיד.
\end{hint}

\begin{solution}
יהי $p(x)=x^3-4x^2+7x+13$.

\textbf{קיום:} $p$ רציף (פולינום). $p(-2)=-8-16-14+13=-25<0$ ו-$p(0)=13>0$. לפי משפט ערך הביניים קיים $c\in(-2,0)$ עם $p(c)=0$.

\textbf{יחידות:} $p'(x)=3x^2-8x+7$. הדיסקרימיננטה היא $64-84=-20<0$ והמקדם המוביל חיובי, ולכן $p'(x)>0$ לכל $x\in\R$. מכאן $p$ עולה ממש על $\R$ (מסקנה ממשפט לגרנז'), ולכן אינה יכולה לקבל את הערך $0$ פעמיים.

(לחלופין: אם היו שני שורשים $a<b$, לפי משפט רול היה $c\in(a,b)$ עם $p'(c)=0$, סתירה.)

לכן לפולינום שורש ממשי אחד ויחיד (הנמצא בקטע $(-2,0)$; בקירוב $c\approx-1.054$). $\blacksquare$
\end{solution}

\subsection*{שאלה 4}

\begin{question}
חקרו את הפונקציה $f(x)=x^2e^{-x^2}$ ושרטטו סקיצה של הגרף שלה.
\end{question}

\begin{hint}
הפונקציה זוגית — מספיק לחקור אותה עבור $x\ge0$.
\end{hint}

\begin{solution}
\textbf{1. תיאור כללי.}
\begin{itemize}
  \item \textbf{תחום הגדרה:} כל $\R$.
  \item \textbf{חיתוך עם הצירים:} $f(x)=0\iff x=0$ (כי $e^{-x^2}>0$). נקודת החיתוך היחידה היא $(0,0)$.
  \item \textbf{זוגיות:} $f(-x)=(-x)^2e^{-(-x)^2}=f(x)$, כלומר $f$ זוגית והגרף סימטרי ביחס לציר $y$. היא אינה מחזורית.
  \item \textbf{סימן:} $f(x)\ge0$ לכל $x$, ושוויון רק ב-$x=0$.
\end{itemize}

\textbf{2. רציפות:} $f$ אלמנטרית ומוגדרת על כל $\R$, ולכן רציפה בכל $\R$ (ואין אסימפטוטות אנכיות).

\textbf{3. התנהגות באינסוף ואסימפטוטות:}
\[ \lim_{x\to\pm\infty}x^2e^{-x^2}=\lim_{t\to\infty}\frac{t}{e^t}=0 \]
(הצבנו $t=x^2$; הגבול $\frac{t}{e^t}\to0$ — גבול מוכר, או לופיטל: $\frac{1}{e^t}\to0$). לכן $y=0$ אסימפטוטה אופקית ב-$\pm\infty$.

\textbf{4. עלייה, ירידה וקיצון:}
\[ f'(x)=2xe^{-x^2}+x^2e^{-x^2}(-2x)=2xe^{-x^2}(1-x^2)=-2x(x-1)(x+1)e^{-x^2}. \]
נקודות קריטיות: $x=0,\pm1$. סימן $f'$ הוא סימן $x(1-x^2)$:
\begin{itemize}
  \item ב-$(-\infty,-1)$: $f'>0$, $f$ עולה;
  \item ב-$(-1,0)$: $f'<0$, $f$ יורדת;
  \item ב-$(0,1)$: $f'>0$, $f$ עולה;
  \item ב-$(1,\infty)$: $f'<0$, $f$ יורדת.
\end{itemize}
לכן: ב-$x=\pm1$ מקסימום מקומי (וגם מוחלט), $f(\pm1)=e^{-1}=\frac1e$; ב-$x=0$ מינימום מקומי (וגם מוחלט), $f(0)=0$. טווח הפונקציה הוא $\left[0,\frac1e\right]$ (מרציפות, ממשפט ערך הביניים, ומכך ש-$f\to0$ באינסוף).

\textbf{5. קמירות ונקודות פיתול:}
\[ f''(x)=\left(2(x-x^3)e^{-x^2}\right)'=2(1-3x^2)e^{-x^2}+2(x-x^3)(-2x)e^{-x^2}=2e^{-x^2}\left(2x^4-5x^2+1\right). \]
$f''(x)=0\iff 2x^4-5x^2+1=0\iff x^2=\frac{5\pm\sqrt{17}}{4}$. נסמן
\[ x_1=\sqrt{\tfrac{5-\sqrt{17}}{4}}\approx0.468,\qquad x_2=\sqrt{\tfrac{5+\sqrt{17}}{4}}\approx1.510. \]
הפולינום $2s^2-5s+1$ (עם $s=x^2$) חיובי מחוץ לשורשיו ושלילי ביניהם, לכן:
\begin{itemize}
  \item $f''>0$ (קמורה, "כלפי מעלה") עבור $|x|<x_1$ ועבור $|x|>x_2$;
  \item $f''<0$ (קעורה, "כלפי מטה") עבור $x_1<|x|<x_2$.
\end{itemize}
ארבע נקודות פיתול: $x=\pm x_1$ עם $f(\pm x_1)=\frac{5-\sqrt{17}}{4}e^{-\frac{5-\sqrt{17}}{4}}\approx0.176$, ו-$x=\pm x_2$ עם $f(\pm x_2)=\frac{5+\sqrt{17}}{4}e^{-\frac{5+\sqrt{17}}{4}}\approx0.233$.

\textbf{6. סקיצה:} גרף בצורת "גבנון כפול" סימטרי ביחס לציר $y$: יוצא מ-$y=0$ ב-$-\infty$ (מעל הציר), עולה למקסימום $\left(-1,\frac1e\right)$, יורד למינימום $(0,0)$ (שם הגרף משיק לציר $x$), עולה שוב למקסימום $\left(1,\frac1e\right)$ ויורד ושואף ל-$0$ כאשר $x\to\infty$. הקמירות מתחלפת בנקודות $\pm0.468$ ו-$\pm1.510$.
\end{solution}

\subsection*{שאלה 5א}

\begin{question}
חשבו את האינטגרל $\displaystyle\int(4-x^2)^{-3/2}\,dx$. \quad [\underline{רמז}: אפשר להציב $x=2\sin\theta$.]
\end{question}

\begin{solution}
האינטגרנד מוגדר עבור $|x|<2$. נציב $x=2\sin\theta$ עם $\theta\in\left(-\frac\pi2,\frac\pi2\right)$; אז $dx=2\cos\theta\,d\theta$ ו-$\cos\theta>0$, ולכן
\[ 4-x^2=4-4\sin^2\theta=4\cos^2\theta,\qquad (4-x^2)^{3/2}=8\cos^3\theta. \]
\[ \int(4-x^2)^{-3/2}\,dx=\int\frac{2\cos\theta}{8\cos^3\theta}\,d\theta=\frac14\int\frac{d\theta}{\cos^2\theta}=\frac14\tan\theta+C. \]
נחזור למשתנה $x$: $\sin\theta=\frac x2$, $\cos\theta=\sqrt{1-\frac{x^2}{4}}=\frac{\sqrt{4-x^2}}{2}$, ולכן $\tan\theta=\frac{x}{\sqrt{4-x^2}}$.
\[ \boxed{\int(4-x^2)^{-3/2}\,dx=\frac{x}{4\sqrt{4-x^2}}+C} \]
בדיקה: $\left(\frac{x}{4\sqrt{4-x^2}}\right)'=\frac14\cdot\frac{\sqrt{4-x^2}+\frac{x^2}{\sqrt{4-x^2}}}{4-x^2}=\frac14\cdot\frac{4}{(4-x^2)^{3/2}}=(4-x^2)^{-3/2}$.
\end{solution}

\subsection*{שאלה 5ב}

\begin{question}
מצאו את השטח הכלוא בתוך הלולאה הנוצרת על-ידי הגרף של הפונקציה $r=\tan\frac\theta2$ בקואורדינטות קטביות.

(במבחן מצורף איור: העקומה עוברת בראשית, יוצרת לולאה קטנה מעל הראשית, ושני ענפיה נמשכים למעלה ולצדדים.)
\end{question}

\begin{hint}
מצאו עבור אילו ערכים של $\theta$ העקומה חוזרת לאותה נקודה: השוו את הנקודות המתקבלות עבור $\theta=\frac\pi2$ ו-$\theta=-\frac\pi2$.
\end{hint}

\begin{hint}
שטח בקואורדינטות קטביות: $S=\frac12\int_\alpha^\beta r^2\,d\theta$, והשתמשו ב-$\tan^2u=\frac1{\cos^2u}-1$.
\end{hint}

\begin{solution}
העקומה $r=\tan\frac\theta2$ מוגדרת עבור $\theta\in(-\pi,\pi)$, ו-$r\to\pm\infty$ כאשר $\theta\to\pm\pi$ (אלה שני הענפים). הנקודה הקרטזית היא $(r\cos\theta,\ r\sin\theta)$.
\begin{itemize}
  \item $\theta=0$: $r=0$ — הראשית.
  \item $\theta=\frac\pi2$: $r=\tan\frac\pi4=1$, הנקודה $(0,1)$.
  \item $\theta=-\frac\pi2$: $r=\tan\left(-\frac\pi4\right)=-1$, הנקודה $(-1\cdot\cos(-\frac\pi2),\,-1\cdot\sin(-\frac\pi2))=(0,1)$.
\end{itemize}
כלומר העקומה עוברת בנקודה $(0,1)$ פעמיים, ב-$\theta=\pm\frac\pi2$: היא חותכת את עצמה שם. כאשר $\theta$ עובר מ-$-\frac\pi2$ ל-$\frac\pi2$ העקומה יוצאת מ-$(0,1)$, עוברת דרך הראשית וחוזרת ל-$(0,1)$ — זו הלולאה. (עבור $\theta\in(-\frac\pi2,0)$, $r<0$ והנקודות נמצאות ברביע השני; עבור $\theta\in(0,\frac\pi2)$ ברביע הראשון; הלולאה סימטרית ביחס לציר $y$.)

השטח הכלוא בלולאה, לפי $S=\frac12\int_\alpha^\beta r^2\,d\theta$ (הנוסחה תקפה גם כאשר $r<0$, כי $r^2$ הוא ריבוע המרחק מהראשית):
\[ S=\frac12\int_{-\pi/2}^{\pi/2}\tan^2\frac\theta2\,d\theta=\int_0^{\pi/2}\tan^2\frac\theta2\,d\theta \]
(האינטגרנד זוגי). עם $\tan^2u=\frac1{\cos^2u}-1$:
\[ S=\int_0^{\pi/2}\left(\frac{1}{\cos^2\frac\theta2}-1\right)d\theta=\left[2\tan\frac\theta2-\theta\right]_0^{\pi/2}=2\tan\frac\pi4-\frac\pi2=2-\frac\pi2. \]
\[ \boxed{S=2-\frac\pi2\approx0.429} \]
\end{solution}

\subsection*{שאלה 6א}

\begin{question}
האם האינטגרל הלא-אמיתי $\displaystyle\int_1^\infty\frac{e^x}{e^{e^x}}\,dx$ מתכנס? אם כן, יש לחשב את ערכו של האינטגרל. אם לא, יש לנמק מדוע.
\end{question}

\begin{hint}
הציבו $u=e^x$.
\end{hint}

\begin{solution}
לפי ההגדרה $\int_1^\infty\frac{e^x}{e^{e^x}}dx=\lim_{b\to\infty}\int_1^b e^xe^{-e^x}dx$. נציב $u=e^x$, $du=e^x\,dx$; כאשר $x$ עובר מ-$1$ ל-$b$, $u$ עובר מ-$e$ ל-$e^b$:
\[ \int_1^b e^xe^{-e^x}\,dx=\int_e^{e^b}e^{-u}\,du=\left[-e^{-u}\right]_e^{e^b}=e^{-e}-e^{-e^b}. \]
כאשר $b\to\infty$: $e^b\to\infty$ ולכן $e^{-e^b}\to0$. הגבול קיים וסופי, כלומר האינטגרל מתכנס:
\[ \boxed{\int_1^\infty\frac{e^x}{e^{e^x}}\,dx=e^{-e}} \]
\end{solution}

\subsection*{שאלה 6ב}

\begin{question}
חשבו את נפח גוף הסיבוב המתקבל ע"י סיבוב התחום $D$ סביב ציר ה-$x$, כאשר $D$ התחום בין גרף הפונקציה $f(x)=\frac{\arctan x}{\sqrt{1+x^2}}$ וציר ה-$x$, כאשר $0\le x\le1$.
\end{question}

\begin{hint}
שימו לב ש-$\frac{1}{1+x^2}$ היא הנגזרת של $\arctan x$, והציבו $u=\arctan x$.
\end{hint}

\begin{solution}
$V=\pi\int_0^1 f^2(x)\,dx=\pi\int_0^1\frac{\arctan^2x}{1+x^2}\,dx$. נציב $u=\arctan x$, $du=\frac{dx}{1+x^2}$; כאשר $x$ עובר מ-$0$ ל-$1$, $u$ עובר מ-$0$ ל-$\frac\pi4$:
\[ V=\pi\int_0^{\pi/4}u^2\,du=\pi\left[\frac{u^3}{3}\right]_0^{\pi/4}=\pi\cdot\frac{1}{3}\cdot\frac{\pi^3}{64}. \]
\[ \boxed{V=\frac{\pi^4}{192}} \]
\end{solution}

\section*{תשע"ח סמסטר א מועד ג}

\subsection*{שאלה 1א}

\begin{question}
חשבו את הגבול הבא:
גבול הסדרה $\displaystyle\lim_{n\to\infty}\left(\frac{(n+4)!\cdot n}{(n+5)!-(n+2)!}\right)$.
\end{question}

\begin{hint}
הוציאו $(n+2)!$ כגורם משותף במונה ובמכנה.
\end{hint}

\begin{solution}
נשתמש ב-$(n+4)!=(n+2)!\,(n+3)(n+4)$ וב-$(n+5)!=(n+2)!\,(n+3)(n+4)(n+5)$. נצמצם ב-$(n+2)!\neq0$:
\[ \frac{(n+4)!\cdot n}{(n+5)!-(n+2)!}=\frac{n(n+3)(n+4)}{(n+3)(n+4)(n+5)-1}. \]
במונה ובמכנה פולינומים ממעלה $3$ עם מקדם מוביל $1$. נחלק ב-$n^3$:
\[ =\frac{\left(1+\frac3n\right)\left(1+\frac4n\right)}{\left(1+\frac3n\right)\left(1+\frac4n\right)\left(1+\frac5n\right)-\frac1{n^3}}\xrightarrow[n\to\infty]{}\frac{1}{1-0}=1 \]
לפי אריתמטיקה של גבולות.
\[ \boxed{\lim_{n\to\infty}\frac{(n+4)!\cdot n}{(n+5)!-(n+2)!}=1} \]
\end{solution}

\subsection*{שאלה 1ב}

\begin{question}
חשבו את הגבול הבא:
גבול הפונקציה $\displaystyle\lim_{x\to0}\tan x\ln(x^2)$.
\end{question}

\begin{hint}
כתבו $\tan x\ln(x^2)=\frac{\tan x}{x}\cdot 2x\ln|x|$.
\end{hint}

\begin{solution}
עבור $0<|x|<\frac\pi2$: $\ln(x^2)=2\ln|x|$, ולכן
\[ \tan x\ln(x^2)=\frac{\tan x}{x}\cdot 2x\ln|x|. \]
ידוע ש-$\lim_{x\to0}\frac{\tan x}{x}=1$ (גבול מפורסם). כמו כן $\lim_{x\to0}x\ln|x|=0$: עבור $x\to0^+$ לפי לופיטל ($\frac\infty\infty$)
\[ \lim_{x\to0^+}\frac{\ln x}{1/x}=\lim_{x\to0^+}\frac{1/x}{-1/x^2}=\lim_{x\to0^+}(-x)=0, \]
ועבור $x\to0^-$ זה נובע מכך ש-$x\ln|x|$ אי-זוגית. לכן
\[ \boxed{\lim_{x\to0}\tan x\ln(x^2)=1\cdot2\cdot0=0} \]
\end{solution}

\subsection*{שאלה 2}

\begin{question}
מצאו את נקודות אי-הרציפות של הפונקציה
\[ f(x)=\frac{x^2-x}{e^{x^3-x}-1} \]
וקבעו את סוגיהן.
\end{question}

\begin{hint}
המכנה מתאפס כאשר $x^3-x=0$. השתמשו בגבול $\lim_{t\to0}\frac{e^t-1}{t}=1$ עם $t=x^3-x$.
\end{hint}

\begin{solution}
\textbf{תחום ההגדרה:} $e^{x^3-x}-1=0\iff x^3-x=0\iff x(x-1)(x+1)=0$. לכן $f$ מוגדרת ורציפה (מנה של פונקציות רציפות) לכל $x\notin\{0,1,-1\}$, ואלה נקודות אי-הרציפות.

נסמן $t=t(x)=x^3-x$. עבור $x\notin\{0,1,-1\}$:
\[ f(x)=\frac{x^2-x}{x^3-x}\cdot\frac{x^3-x}{e^{x^3-x}-1}=\frac{x(x-1)}{x(x-1)(x+1)}\cdot\frac{t}{e^t-1}=\frac{1}{x+1}\cdot\frac{t}{e^t-1}. \]
כאשר $x\to0$ או $x\to\pm1$, $t\to0$ ו-$t\neq0$, ולכן לפי הגבול המפורסם ומשפט הגבול של הרכבה $\frac{t}{e^t-1}\to1$.

\textbf{$x=0$:} $\lim_{x\to0}f(x)=\frac{1}{0+1}\cdot1=1$. הגבול קיים וסופי — \textbf{אי-רציפות סליקה}.

\textbf{$x=1$:} $\lim_{x\to1}f(x)=\frac{1}{2}\cdot1=\frac12$ — \textbf{אי-רציפות סליקה}.

\textbf{$x=-1$:} כאן המונה שואף ל-$(-1)^2-(-1)=2\neq0$ והמכנה שואף ל-$0$. נבדוק סימנים: $x^3-x=x(x-1)(x+1)$, ובסביבת $-1$, $x(x-1)\approx2>0$. לכן
\begin{itemize}
  \item עבור $x\to-1^+$: $x+1>0$, $t\to0^+$, $e^t-1\to0^+$ ולכן $f(x)\to+\infty$;
  \item עבור $x\to-1^-$: $x+1<0$, $t\to0^-$, $e^t-1\to0^-$ ולכן $f(x)\to-\infty$.
\end{itemize}
הגבולות החד-צדדיים אינסופיים — \textbf{אי-רציפות מסוג שני}.

\textbf{סיכום:} $x=0$ ו-$x=1$ — נקודות אי-רציפות סליקות (עם גבולות $1$ ו-$\frac12$ בהתאמה); $x=-1$ — אי-רציפות מסוג שני.
\end{solution}

\subsection*{שאלה 3א}

\begin{question}
חשבו את הגבול:
\[ \lim_{x\to0}\left(\frac{1-\cos(x^3)}{x^2\ln(1+x^2)\sin^2x}\right) \]
\underline{רמז}: אפשר להשתמש בנוסחאות מקלורין.
\end{question}

\begin{hint}
$1-\cos u=\frac{u^2}{2}+o(u^2)$, $\ln(1+u)=u+o(u)$, $\sin x=x+o(x)$. מה סדר הגודל של המונה ושל המכנה?
\end{hint}

\begin{solution}
פיתוחי מקלורין: $\cos u=1-\frac{u^2}{2}+o(u^2)$ עם $u=x^3$ נותן
\[ 1-\cos(x^3)=\frac{x^6}{2}+o(x^6). \]
כמו כן $\ln(1+x^2)=x^2+o(x^2)$ ו-$\sin x=x+o(x)$, ולכן $\sin^2x=x^2+o(x^2)$. המכנה:
\[ x^2\ln(1+x^2)\sin^2x=x^2\left(x^2+o(x^2)\right)\left(x^2+o(x^2)\right)=x^6+o(x^6). \]
לכן
\[ \frac{1-\cos(x^3)}{x^2\ln(1+x^2)\sin^2x}=\frac{\frac12+\frac{o(x^6)}{x^6}}{1+\frac{o(x^6)}{x^6}}\xrightarrow[x\to0]{}\frac12. \]
(באופן שקול: $\frac{1-\cos(x^3)}{x^6}\to\frac12$, $\frac{\ln(1+x^2)}{x^2}\to1$, $\frac{\sin^2x}{x^2}\to1$, ומפרקים את הביטוי למכפלה.)
\[ \boxed{\lim_{x\to0}\frac{1-\cos(x^3)}{x^2\ln(1+x^2)\sin^2x}=\frac12} \]
\end{solution}

\subsection*{שאלה 3ב}

\begin{question}
תהי $f$ פונקציה גזירה לכל מספר ממשי $x$, ונניח שיש בדיוק $k$ פתרונות למשוואה $f'(x)=0$. הוכיחו שיש לכל היותר $k+1$ פתרונות למשוואה $f(x)=0$.
\end{question}

\begin{hint}
בין כל שני שורשים של $f$ יש שורש של $f'$ (משפט רול). הוכיחו בדרך השלילה.
\end{hint}

\begin{solution}
נניח בשלילה שלמשוואה $f(x)=0$ יש לפחות $k+2$ פתרונות שונים, ונסמן $k+2$ מהם $x_1<x_2<\dots<x_{k+2}$.

לכל $i=1,\dots,k+1$, הפונקציה $f$ רציפה בקטע $[x_i,x_{i+1}]$ וגזירה בקטע $(x_i,x_{i+1})$ (כי היא גזירה בכל $\R$, ולכן גם רציפה), ו-$f(x_i)=f(x_{i+1})=0$. לפי \textbf{משפט רול} קיימת נקודה $c_i\in(x_i,x_{i+1})$ כך ש-$f'(c_i)=0$.

הקטעים $(x_1,x_2),(x_2,x_3),\dots,(x_{k+1},x_{k+2})$ זרים בזוגות, ולכן $c_1<c_2<\dots<c_{k+1}$ הן $k+1$ נקודות שונות שבהן $f'=0$. זו סתירה להנחה שלמשוואה $f'(x)=0$ יש בדיוק $k$ פתרונות.

לכן למשוואה $f(x)=0$ יש לכל היותר $k+1$ פתרונות. $\blacksquare$
\end{solution}

\subsection*{שאלה 4}

\begin{question}
חקרו את הפונקציה $f(x)=\sqrt{\frac{x^3}{x-1}}$ ושרטטו סקיצה של הגרף שלה. אין צורך לחשב את $f''(x)$ בשאלה זו, ואין צורך לעסוק בקמירות, קעירות ונקודות פיתול.
\end{question}

\begin{hint}
תחום ההגדרה: היכן ש-$\frac{x^3}{x-1}\ge0$. שימו לב שהוא מורכב משני חלקים נפרדים.
\end{hint}

\begin{hint}
באסימפטוטה המשופעת שימו לב ש-$\sqrt{x^2}=|x|$, ולכן השיפועים ב-$+\infty$ וב-$-\infty$ שונים.
\end{hint}

\begin{solution}
\textbf{1. תיאור כללי.}
\begin{itemize}
  \item \textbf{תחום הגדרה:} צריך $x\neq1$ ו-$\frac{x^3}{x-1}\ge0$. סימן $x^3$ הוא סימן $x$, לכן המנה אי-שלילית כאשר $x\le0$ (מונה $\le0$, מכנה $<0$) או $x>1$ (שניהם חיוביים). התחום: $(-\infty,0]\cup(1,\infty)$.
  \item \textbf{חיתוך עם הצירים:} $f(x)=0\iff x=0$. נקודת החיתוך היחידה (עם שני הצירים) היא $(0,0)$.
  \item \textbf{זוגיות/מחזוריות:} התחום אינו סימטרי ביחס ל-$0$, ולכן $f$ אינה זוגית ואינה אי-זוגית; היא אינה מחזורית.
  \item \textbf{סימן:} $f(x)\ge0$ בכל התחום, ושוויון רק ב-$x=0$.
\end{itemize}

\textbf{2. רציפות:} $f$ היא הרכבה של פונקציות אלמנטריות, ולכן רציפה בכל תחום הגדרתה (ב-$x=0$ — רציפה משמאל). הנקודה $x=1$ אינה בתחום, וכאשר $x\to1^+$: $x^3\to1$, $x-1\to0^+$, ולכן $f(x)\to+\infty$.

\textbf{3. אסימפטוטות.}
\begin{itemize}
  \item \textbf{אנכית:} $x=1$ (מימין).
  \item \textbf{כאשר $x\to+\infty$:} $f(x)=\sqrt{x^2\cdot\frac{x}{x-1}}=x\sqrt{\frac{x}{x-1}}$ (כי $x>0$).
  \[ a=\lim_{x\to\infty}\frac{f(x)}{x}=\lim_{x\to\infty}\sqrt{\frac{x}{x-1}}=1, \]
  \[ b=\lim_{x\to\infty}\left(f(x)-x\right)=\lim_{x\to\infty}x\left(\sqrt{\tfrac{x}{x-1}}-1\right)=\lim_{x\to\infty}\frac{x\left(\frac{x}{x-1}-1\right)}{\sqrt{\frac{x}{x-1}}+1}=\lim_{x\to\infty}\frac{\frac{x}{x-1}}{\sqrt{\frac{x}{x-1}}+1}=\frac{1}{2}. \]
  אסימפטוטה: $y=x+\frac12$.
  \item \textbf{כאשר $x\to-\infty$:} כאן $\sqrt{x^2}=|x|=-x$, ולכן $f(x)=-x\sqrt{\frac{x}{x-1}}$.
  \[ a=\lim_{x\to-\infty}\frac{f(x)}{x}=-\lim_{x\to-\infty}\sqrt{\frac{x}{x-1}}=-1, \]
  \[ b=\lim_{x\to-\infty}\left(f(x)+x\right)=\lim_{x\to-\infty}(-x)\left(\sqrt{\tfrac{x}{x-1}}-1\right)=-\lim_{x\to-\infty}\frac{\frac{x}{x-1}}{\sqrt{\frac{x}{x-1}}+1}=-\frac12. \]
  אסימפטוטה: $y=-x-\frac12$.
\end{itemize}

\textbf{4. עלייה, ירידה וקיצון.} נסמן $g(x)=\frac{x^3}{x-1}$. אז
\[ g'(x)=\frac{3x^2(x-1)-x^3}{(x-1)^2}=\frac{x^2(2x-3)}{(x-1)^2}, \]
ובנקודות שבהן $g(x)>0$ (כלומר $x<0$ או $x>1$), לפי כלל השרשרת
\[ f'(x)=\frac{g'(x)}{2\sqrt{g(x)}}=\frac{x^2(2x-3)}{2(x-1)^2\sqrt{\frac{x^3}{x-1}}}. \]
סימן $f'$ הוא סימן $2x-3$ (שאר הגורמים חיוביים):
\begin{itemize}
  \item ב-$(-\infty,0)$: $2x-3<0$ ולכן $f$ יורדת (ממש). כיוון ש-$f$ רציפה ב-$0$ משמאל, היא יורדת בכל $(-\infty,0]$.
  \item ב-$\left(1,\frac32\right)$: $f'<0$, $f$ יורדת.
  \item ב-$\left(\frac32,\infty\right)$: $f'>0$, $f$ עולה.
\end{itemize}
לכן:
\begin{itemize}
  \item ב-$x=\frac32$ מינימום מקומי: $f\left(\frac32\right)=\sqrt{\frac{27/8}{1/2}}=\sqrt{\frac{27}{4}}=\frac{3\sqrt3}{2}\approx2.6$.
  \item ב-$x=0$ (קצה התחום משמאל) מינימום מקומי, ואף מוחלט, $f(0)=0$. הנגזרת החד-צדדית שם: $\frac{f(x)-f(0)}{x}=\frac{|x|^{3/2}}{x\sqrt{1-x}}\to0$ כאשר $x\to0^-$, כלומר הגרף מגיע ל-$(0,0)$ באופן אופקי (משיק לציר $x$).
\end{itemize}

\textbf{5. סקיצה.} \emph{ענף שמאלי} ($x\le0$): הגרף יורד מ-$+\infty$, קרוב מעל לאסימפטוטה $y=-x-\frac12$ כאשר $x\to-\infty$, עד הנקודה $(0,0)$ שבה הוא מגיע משיק לציר $x$ ונעצר. \emph{ענף ימני} ($x>1$): הגרף יורד מ-$+\infty$ (צמוד לאסימפטוטה האנכית $x=1$) עד המינימום $\left(\frac32,\frac{3\sqrt3}{2}\right)$, ואחר כך עולה ומתקרב (מלמעלה) לאסימפטוטה $y=x+\frac12$. בקטע $(0,1]$ אין גרף.
\end{solution}

\subsection*{שאלה 5א}

\begin{question}
חשבו את האינטגרל $\displaystyle\int\frac{x^4+4x^3+4x^2+8x+1}{(x-1)(x+2)^2}\,dx$.
\end{question}

\begin{hint}
מעלת המונה גדולה ממעלת המכנה — בצעו קודם חילוק פולינומים, ואז פרקו לשברים חלקיים מהצורה $\frac{A}{x-1}+\frac{B}{x+2}+\frac{C}{(x+2)^2}$.
\end{hint}

\begin{solution}
המכנה: $(x-1)(x+2)^2=(x-1)(x^2+4x+4)=x^3+3x^2-4$. מעלת המונה ($4$) גדולה ממעלת המכנה ($3$), לכן נחלק פולינומים:
\[ (x+1)(x^3+3x^2-4)=x^4+4x^3+3x^2-4x-4, \]
ולכן
\[ x^4+4x^3+4x^2+8x+1=(x+1)(x^3+3x^2-4)+(x^2+12x+5). \]
כלומר
\[ \frac{x^4+4x^3+4x^2+8x+1}{(x-1)(x+2)^2}=x+1+\frac{x^2+12x+5}{(x-1)(x+2)^2}. \]
פירוק לשברים חלקיים:
\[ \frac{x^2+12x+5}{(x-1)(x+2)^2}=\frac{A}{x-1}+\frac{B}{x+2}+\frac{C}{(x+2)^2}, \]
\[ x^2+12x+5=A(x+2)^2+B(x-1)(x+2)+C(x-1). \]
\begin{itemize}
  \item $x=1$: $18=9A$, כלומר $A=2$.
  \item $x=-2$: $4-24+5=-15=-3C$, כלומר $C=5$.
  \item השוואת מקדמי $x^2$: $1=A+B$, כלומר $B=-1$.
\end{itemize}
(בדיקה, מקדם חופשי: $4A-2B-C=8+2-5=5$, כנדרש.) לכן
\[ \int\frac{x^4+4x^3+4x^2+8x+1}{(x-1)(x+2)^2}\,dx=\int\left(x+1+\frac{2}{x-1}-\frac{1}{x+2}+\frac{5}{(x+2)^2}\right)dx \]
\[ \boxed{=\frac{x^2}{2}+x+2\ln|x-1|-\ln|x+2|-\frac{5}{x+2}+C} \]
\end{solution}

\subsection*{שאלה 5ב}

\begin{question}
חשבו את אורך העקומה $f(x)=x^{3/2}$ בין $x=0$ ל-$x=9$.
\end{question}

\begin{hint}
אורך עקומה: $L=\int_a^b\sqrt{1+(f'(x))^2}\,dx$.
\end{hint}

\begin{solution}
$f'(x)=\frac32x^{1/2}$, ולכן $1+(f'(x))^2=1+\frac94x$. אורך העקומה:
\[ L=\int_0^9\sqrt{1+\frac94x}\,dx. \]
נציב $u=1+\frac94x$, $du=\frac94dx$; כאשר $x$ עובר מ-$0$ ל-$9$, $u$ עובר מ-$1$ ל-$1+\frac{81}{4}=\frac{85}{4}$:
\[ L=\frac49\int_1^{85/4}u^{1/2}\,du=\frac49\cdot\frac23\left[u^{3/2}\right]_1^{85/4}=\frac{8}{27}\left(\left(\frac{85}{4}\right)^{3/2}-1\right)=\frac{8}{27}\left(\frac{85\sqrt{85}}{8}-1\right). \]
\[ \boxed{L=\frac{85\sqrt{85}-8}{27}\approx28.73} \]
\end{solution}

\subsection*{שאלה 6א}

\begin{question}
האם האינטגרל הלא-אמיתי $\displaystyle\int_0^\infty\frac{1}{e^x+e^{-x}}\,dx$ מתכנס? אם כן, יש לחשב את ערכו של האינטגרל. אם לא, יש לנמק מדוע.
\end{question}

\begin{hint}
כפלו מונה ומכנה ב-$e^x$ והציבו $u=e^x$.
\end{hint}

\begin{solution}
נכפול מונה ומכנה ב-$e^x>0$: $\frac{1}{e^x+e^{-x}}=\frac{e^x}{e^{2x}+1}$. לפי ההגדרה,
\[ \int_0^\infty\frac{dx}{e^x+e^{-x}}=\lim_{b\to\infty}\int_0^b\frac{e^x}{1+(e^x)^2}\,dx. \]
נציב $u=e^x$, $du=e^x\,dx$; $u$ עובר מ-$1$ ל-$e^b$:
\[ \int_0^b\frac{e^x\,dx}{1+e^{2x}}=\int_1^{e^b}\frac{du}{1+u^2}=\arctan(e^b)-\arctan1=\arctan(e^b)-\frac\pi4. \]
כאשר $b\to\infty$, $e^b\to\infty$ ו-$\arctan(e^b)\to\frac\pi2$. הגבול קיים וסופי, ולכן האינטגרל מתכנס:
\[ \boxed{\int_0^\infty\frac{dx}{e^x+e^{-x}}=\frac\pi2-\frac\pi4=\frac\pi4} \]
\end{solution}

\subsection*{שאלה 6ב}

\begin{question}
חשבו את נפח גוף הסיבוב המתקבל ע"י סיבוב התחום $D$ סביב ציר ה-$x$, כאשר $D$ התחום בין גרף הפונקציה $f(x)=\sqrt{\arcsin x}$ וציר ה-$x$, כאשר $0\le x\le1$.
\end{question}

\begin{hint}
יש לחשב $\int\arcsin x\,dx$ — השתמשו באינטגרציה בחלקים עם $u=\arcsin x$, $v'=1$.
\end{hint}

\begin{solution}
$V=\pi\int_0^1f^2(x)\,dx=\pi\int_0^1\arcsin x\,dx$. נחשב בחלקים, עם $u=\arcsin x$, $v'=1$, $u'=\frac{1}{\sqrt{1-x^2}}$, $v=x$:
\[ \int\arcsin x\,dx=x\arcsin x-\int\frac{x}{\sqrt{1-x^2}}\,dx=x\arcsin x+\sqrt{1-x^2}+C \]
(האינטגרל האחרון — בהצבה $t=1-x^2$, $dt=-2x\,dx$: $\int\frac{x\,dx}{\sqrt{1-x^2}}=-\frac12\int t^{-1/2}dt=-\sqrt{1-x^2}$.) הפונקציה $x\arcsin x+\sqrt{1-x^2}$ רציפה ב-$[0,1]$ (האינטגרנד $\arcsin x$ רציף בקטע הסגור), ולכן לפי המשפט היסודי
\[ \int_0^1\arcsin x\,dx=\left[x\arcsin x+\sqrt{1-x^2}\right]_0^1=\left(\frac\pi2+0\right)-(0+1)=\frac\pi2-1. \]
\[ \boxed{V=\pi\left(\frac\pi2-1\right)=\frac{\pi^2}{2}-\pi} \]
\end{solution}

\section*{תשע"ח סמסטר ב מועד א}

\subsection*{שאלה 1א}

\begin{question}
חשבו את הנגזרת של הפונקציה $f(x)=\cos x$ בנקודה $x$ לפי הגדרת הנגזרת (תשובה לפי טבלת הנגזרות לא תתקבל).
\end{question}

\begin{hint}
השתמשו בזהות $\cos\alpha-\cos\beta=-2\sin\frac{\alpha+\beta}{2}\sin\frac{\alpha-\beta}{2}$ ובגבול היסודי $\lim_{t\to0}\frac{\sin t}{t}=1$.
\end{hint}

\begin{solution}
לפי ההגדרה $f'(x)=\lim_{h\to0}\frac{f(x+h)-f(x)}{h}$. לפי הזהות $\cos\alpha-\cos\beta=-2\sin\frac{\alpha+\beta}{2}\sin\frac{\alpha-\beta}{2}$:
\[
\frac{\cos(x+h)-\cos x}{h}=\frac{-2\sin\left(x+\frac h2\right)\sin\frac h2}{h}=-\sin\left(x+\frac h2\right)\cdot\frac{\sin\frac h2}{\frac h2}.
\]
כאשר $h\to0$: $\sin\left(x+\frac h2\right)\to\sin x$ (רציפות הסינוס) ו-$\frac{\sin(h/2)}{h/2}\to1$ (הגבול היסודי). לפי אריתמטיקה של גבולות:
\[
(\cos x)'=\lim_{h\to0}\frac{\cos(x+h)-\cos x}{h}=-\sin x.
\]
\end{solution}

\subsection*{שאלה 1ב}

\begin{question}
מצאו את כל האסימפטוטות של הפונקציה $f(x)=\frac{\ln(2x)}{\ln(x)}$.
\end{question}

\begin{hint}
כתבו $\ln(2x)=\ln2+\ln x$.
\end{hint}

\begin{solution}
\textbf{תחום הגדרה:} $x>0$ (בשביל $\ln x$ ו-$\ln 2x$) ו-$\ln x\neq0$, כלומר $x\in(0,1)\cup(1,\infty)$. בתחום זה
\[
f(x)=\frac{\ln2+\ln x}{\ln x}=1+\frac{\ln2}{\ln x}.
\]
\textbf{אנכיות:} כאשר $x\to1^+$, $\ln x\to0^+$ ולכן $f(x)\to+\infty$; כאשר $x\to1^-$, $\ln x\to0^-$ ולכן $f(x)\to-\infty$. לכן $x=1$ אסימפטוטה אנכית.
כאשר $x\to0^+$: $\ln x\to-\infty$ ולכן $f(x)\to1$ — גבול סופי, ולכן $x=0$ \emph{אינה} אסימפטוטה אנכית.

\textbf{אופקיות/משופעות:} הפונקציה מוגדרת רק עבור $x>0$, ולכן בודקים רק $x\to+\infty$: $\ln x\to\infty$ ולכן $f(x)\to1$. לכן $y=1$ אסימפטוטה אופקית ב-$+\infty$ (ואין משופעת).

\textbf{תשובה:} אסימפטוטה אנכית $x=1$ ואסימפטוטה אופקית $y=1$ (כאשר $x\to+\infty$).
\end{solution}

\subsection*{שאלה 1ג}

\begin{question}
עבור איזה ערכים של $a$ לגרף הפונקציה $f(x)=x^4-ax^2$ אין נקודות פיתול?
\end{question}

\begin{hint}
נקודת פיתול היא נקודה שבה $f''$ מחליפה סימן. חשבו את $f''$ ובדקו לאילו ערכים של $a$ היא מחליפה סימן.
\end{hint}

\begin{solution}
$f''(x)=12x^2-2a$. נקודת פיתול היא נקודה שבה $f''$ מחליפה סימן.
\begin{itemize}
  \item אם $a>0$: $f''(x)=0\iff x=\pm\sqrt{a/6}$, ו-$f''$ (פרבולה) מחליפה סימן בכל אחת מהנקודות האלה, לכן יש שתי נקודות פיתול.
  \item אם $a=0$: $f''(x)=12x^2\ge0$, מתאפסת רק ב-$x=0$ בלי להחליף סימן — אין נקודות פיתול ($f=x^4$ קמורה).
  \item אם $a<0$: $f''(x)=12x^2-2a>0$ לכל $x$ — אין נקודות פיתול.
\end{itemize}
\textbf{תשובה:} $a\le0$.
\end{solution}

\subsection*{שאלה 2א}

\begin{question}
לאילו ערכים של פרמטר $a$ למשוואה $x^4-a\cdot x^3+27=0$ אין שורשים ממשיים?
\end{question}

\begin{hint}
$x=0$ אינו שורש. בודדו את הפרמטר: $a=\frac{x^4+27}{x^3}=x+\frac{27}{x^3}$, ומצאו את טווח הערכים של הפונקציה $g(x)=x+\frac{27}{x^3}$.
\end{hint}

\begin{solution}
$x=0$ אינו שורש (מתקבל $27\neq0$). עבור $x\neq0$:
\[
x^4-ax^3+27=0\iff a=\frac{x^4+27}{x^3}=x+\frac{27}{x^3}=:g(x).
\]
לכן למשוואה יש שורש ממשי אם ורק אם $a$ שייך לתמונה של $g$ על $\R\setminus\{0\}$. נמצא את התמונה.

\textbf{עבור $x>0$:} $g'(x)=1-\frac{81}{x^4}$, שמתאפסת ב-$x=3$; $g'<0$ ב-$(0,3)$ ו-$g'>0$ ב-$(3,\infty)$. לכן $x=3$ מינימום מוחלט על $(0,\infty)$, עם $g(3)=3+1=4$. בנוסף $g(x)\to+\infty$ כאשר $x\to0^+$ וכאשר $x\to+\infty$. מכיוון ש-$g$ רציפה, לפי משפט ערך הביניים התמונה של $(0,\infty)$ היא $[4,\infty)$.

\textbf{עבור $x<0$:} $g$ פונקציה אי-זוגית ($g(-x)=-g(x)$), ולכן התמונה של $(-\infty,0)$ היא $(-\infty,-4]$.

לכן התמונה של $g$ היא $(-\infty,-4]\cup[4,\infty)$, ולמשוואה יש שורש ממשי אם ורק אם $|a|\ge4$. (בדיקה: עבור $a=4$, $x=3$: $81-108+27=0$.)

\textbf{תשובה:} אין שורשים ממשיים אם ורק אם $-4<a<4$.
\end{solution}

\subsection*{שאלה 2ב}

\begin{question}
חשבו את נפח הגוף שמתקבל על ידי סיבוב סביב ציר ה-$x$ של התחום המישורי החסום על ידי הקו $y=\frac{x}{\sqrt{x^2+1}}$ וציר ה-$x$ בקטע $[0,1]$. ציירו את הסקיצה המתאימה.
\end{question}

\begin{solution}
\textbf{סקיצה:} $y=\frac{x}{\sqrt{x^2+1}}$ עולה מ-$0$ (ב-$x=0$) עד $\frac1{\sqrt2}\approx0.71$ (ב-$x=1$), ואי-שלילית בקטע. התחום הוא השטח שבין הגרף לציר $x$ עבור $0\le x\le1$, וסיבובו סביב ציר $x$ נותן גוף דמוי "גביע".

\textbf{נפח:} לפי הנוסחה לנפח גוף סיבוב $V=\pi\int_a^by^2\,dx$:
\[
V=\pi\int_0^1\frac{x^2}{x^2+1}dx=\pi\int_0^1\left(1-\frac{1}{x^2+1}\right)dx=\pi\Big[x-\arctan x\Big]_0^1=\pi\left(1-\frac\pi4\right).
\]
\textbf{תשובה:} $V=\pi\left(1-\frac\pi4\right)\approx0.674$.
\end{solution}

\subsection*{שאלה 3א}

\begin{question}
חשבו: $\displaystyle\int_{-1}^15x^4\big(\ln(x+2)+\sin x\big)dx$
\end{question}

\begin{hint}
$5x^4\sin x$ היא פונקציה אי-זוגית, ולכן האינטגרל שלה על $[-1,1]$ מתאפס. לחלק השני, אינטגרציה בחלקים עם $u=\ln(x+2)$, $dv=5x^4dx$.
\end{hint}

\begin{solution}
\textbf{החלק עם הסינוס.} $g(x)=5x^4\sin x$ אי-זוגית ($g(-x)=-g(x)$) ורציפה, ולכן $\int_{-1}^1 5x^4\sin x\,dx=0$.

\textbf{החלק עם הלוגריתם.} אינטגרציה בחלקים עם $u=\ln(x+2)$, $dv=5x^4dx$, $du=\frac{dx}{x+2}$, $v=x^5$:
\[
\int_{-1}^15x^4\ln(x+2)\,dx=\Big[x^5\ln(x+2)\Big]_{-1}^1-\int_{-1}^1\frac{x^5}{x+2}dx=\ln3-\int_{-1}^1\frac{x^5}{x+2}dx
\]
(כי ב-$x=-1$: $(-1)^5\ln1=0$). חילוק פולינומים:
\[
\frac{x^5}{x+2}=x^4-2x^3+4x^2-8x+16-\frac{32}{x+2}.
\]
על הקטע הסימטרי $[-1,1]$ האינטגרלים של החזקות האי-זוגיות מתאפסים, ולכן
\[
\int_{-1}^1\frac{x^5}{x+2}dx=\int_{-1}^1\left(x^4+4x^2+16\right)dx-32\int_{-1}^1\frac{dx}{x+2}
=\frac25+\frac83+32-32\ln3=\frac{526}{15}-32\ln3.
\]
\textbf{תשובה:}
\[
\int_{-1}^15x^4\big(\ln(x+2)+\sin x\big)dx=\ln3-\frac{526}{15}+32\ln3=33\ln3-\frac{526}{15}\approx1.1875.
\]
\end{solution}

\subsection*{שאלה 3ב}

\begin{question}
רשמו פולינום מקלורן מסדר 2 עבור הפונקציה $y(x)$ אשר מוגדרת בצורה סתומה על ידי המשוואה $y\cdot\sin x-x\cdot\cos x+y=2$.
\end{question}

\begin{hint}
הציבו $x=0$ כדי למצוא את $y(0)$, ואז גזרו את המשוואה (גזירה סתומה) פעמיים והציבו $x=0$.
\end{hint}

\begin{solution}
\textbf{ערך ב-$0$:} בהצבה $x=0$: $y(0)\cdot0-0+y(0)=2$, ולכן $y(0)=2$.

\textbf{נגזרת ראשונה} (גזירה סתומה של שני האגפים לפי $x$, כאשר $y=y(x)$):
\[
y'\sin x+y\cos x-\cos x+x\sin x+y'=0.\tag{*}
\]
בהצבה $x=0$: $y(0)-1+y'(0)=0$, ולכן $y'(0)=1-2=-1$.

\textbf{נגזרת שנייה} (גזירה של (*)):
\[
y''\sin x+y'\cos x+y'\cos x-y\sin x+\sin x+\sin x+x\cos x+y''=0.
\]
בהצבה $x=0$: $2y'(0)+y''(0)=0$, ולכן $y''(0)=2$.

\textbf{פולינום מקלורן מסדר 2:}
\[
P_2(x)=y(0)+y'(0)x+\frac{y''(0)}{2}x^2=2-x+x^2.
\]
(בדיקה: מהמשוואה $y=\frac{2+x\cos x}{1+\sin x}$, ופיתוח ישיר נותן אותו פולינום.)
\end{solution}

\subsection*{שאלה 4א}

\begin{question}
$\displaystyle\lim_{x\to3}\frac{\sqrt{x^2-2x+6}-\sqrt{x^2+2x-6}}{x^2-4x+3}$.
\end{question}

\begin{hint}
הכפילו בצמוד של המונה.
\end{hint}

\begin{solution}
ב-$x=3$ המונה והמכנה מתאפסים ($\sqrt9-\sqrt9=0$, $9-12+3=0$). נכפול בצמוד:
\[
\sqrt{x^2-2x+6}-\sqrt{x^2+2x-6}=\frac{(x^2-2x+6)-(x^2+2x-6)}{\sqrt{x^2-2x+6}+\sqrt{x^2+2x-6}}=\frac{-4(x-3)}{\sqrt{x^2-2x+6}+\sqrt{x^2+2x-6}},
\]
ו-$x^2-4x+3=(x-1)(x-3)$. לכן עבור $x\neq3$ (בסביבת $3$):
\[
\frac{\sqrt{x^2-2x+6}-\sqrt{x^2+2x-6}}{x^2-4x+3}=\frac{-4}{(x-1)\left(\sqrt{x^2-2x+6}+\sqrt{x^2+2x-6}\right)}\xrightarrow[x\to3]{}\frac{-4}{2\cdot(3+3)}=-\frac13.
\]
\textbf{תשובה:} $-\frac13$.
\end{solution}

\subsection*{שאלה 4ב}

\begin{question}
האם הפונקציה
\[
f(x)=\begin{cases}(1+\sin2x)^{\frac{2}{3x}}+\arctan x\cdot\cos\frac1x & x\neq0\\ e^{4/3} & x=0\end{cases}
\]
רציפה בנקודה $x=0$?
\end{question}

\begin{hint}
טפלו בכל מחובר בנפרד: הראשון הוא מהצורה $1^\infty$; השני הוא "אפסה כפול חסומה".
\end{hint}

\begin{solution}
יש לבדוק האם $\lim_{x\to0}f(x)=f(0)=e^{4/3}$.

\textbf{המחובר הראשון.} בסביבת $0$ מתקיים $1+\sin2x>0$, ולכן
\[
(1+\sin2x)^{\frac{2}{3x}}=\exp\left(\frac{2}{3x}\ln(1+\sin2x)\right).
\]
נחשב את המעריך:
\[
\frac{2\ln(1+\sin2x)}{3x}=\frac23\cdot\frac{\ln(1+\sin2x)}{\sin2x}\cdot\frac{\sin2x}{2x}\cdot2\xrightarrow[x\to0]{}\frac23\cdot1\cdot1\cdot2=\frac43,
\]
לפי הגבולות היסודיים $\lim_{t\to0}\frac{\ln(1+t)}{t}=1$ (עם $t=\sin2x\to0$) ו-$\lim_{t\to0}\frac{\sin t}{t}=1$. לפי רציפות האקספוננט, המחובר הראשון שואף ל-$e^{4/3}$.

\textbf{המחובר השני.} $\arctan x\to0$ כאשר $x\to0$, ו-$\left|\cos\frac1x\right|\le1$. מכפלה של פונקציה השואפת ל-$0$ בפונקציה חסומה שואפת ל-$0$, ולכן $\arctan x\cdot\cos\frac1x\to0$.

לכן $\lim_{x\to0}f(x)=e^{4/3}+0=e^{4/3}=f(0)$.

\textbf{תשובה:} כן, $f$ רציפה ב-$x=0$.
\end{solution}

\subsection*{שאלה 4ג}

\begin{question}
האם האינטגרל הלא אמיתי $\displaystyle\int_0^{+\infty}\frac{x\cdot e^{-x}}{(x+1)^4}dx$ מתכנס?
\end{question}

\begin{hint}
השוו ל-$e^{-x}$.
\end{hint}

\begin{solution}
האינטגרנד $g(x)=\frac{xe^{-x}}{(x+1)^4}$ רציף ואי-שלילי על $[0,\infty)$, ולכן האינטגרל לא אמיתי רק בגלל הגבול האינסופי, ואפשר להשתמש במבחן ההשוואה.
עבור $x\ge0$: $0\le x<x+1\le(x+1)^4$, ולכן $\frac{x}{(x+1)^4}\le1$ ו-
\[
0\le g(x)\le e^{-x}.
\]
האינטגרל $\int_0^\infty e^{-x}dx=\lim_{R\to\infty}\left(1-e^{-R}\right)=1$ מתכנס. לפי \emph{מבחן ההשוואה} לאינטגרלים לא אמיתיים של פונקציות אי-שליליות, גם $\int_0^\infty g(x)\,dx$ מתכנס (וערכו בין $0$ ל-$1$).

\textbf{תשובה:} כן, האינטגרל מתכנס.
\end{solution}

\section*{תשע"ח סמסטר ב מועד ב}

\subsection*{שאלה 1א}

\begin{question}
חשבו את הנגזרת של הפונקציה $f(x)=\sin2x$ בנקודה $x$ לפי הגדרת הנגזרת (תשובה לפי טבלת הנגזרות לא תתקבל).
\end{question}

\begin{hint}
השתמשו בזהות $\sin\alpha-\sin\beta=2\cos\frac{\alpha+\beta}{2}\sin\frac{\alpha-\beta}{2}$.
\end{hint}

\begin{solution}
לפי ההגדרה ולפי הזהות $\sin\alpha-\sin\beta=2\cos\frac{\alpha+\beta}{2}\sin\frac{\alpha-\beta}{2}$:
\[
\frac{\sin(2x+2h)-\sin2x}{h}=\frac{2\cos(2x+h)\sin h}{h}=2\cos(2x+h)\cdot\frac{\sin h}{h}.
\]
כאשר $h\to0$: $\cos(2x+h)\to\cos2x$ (רציפות הקוסינוס) ו-$\frac{\sin h}{h}\to1$ (הגבול היסודי). לכן
\[
(\sin2x)'=\lim_{h\to0}\frac{\sin2(x+h)-\sin2x}{h}=2\cos2x.
\]
\end{solution}

\subsection*{שאלה 1ב}

\begin{question}
מצאו את כל הנקודות האי רציפות של הפונקציה ותמיינו אותם
\[
f(x)=\frac{1}{25-5^{1/(x-2)}}.
\]
\end{question}

\begin{hint}
הנקודות החשודות הן $x=2$ (שם $\frac1{x-2}$ לא מוגדר) והנקודה שבה המכנה מתאפס. ב-$x=2$ חשבו גבולות חד-צדדיים בנפרד.
\end{hint}

\begin{solution}
\textbf{תחום הגדרה:} צריך $x\neq2$ וגם $5^{1/(x-2)}\neq25$, כלומר $\frac1{x-2}\neq2$, כלומר $x\neq\frac52$. בכל נקודה אחרת $f$ היא הרכבה ומנה של פונקציות רציפות, ולכן רציפה. נקודות אי-הרציפות הן $x=2$ ו-$x=\frac52$.

\textbf{הנקודה $x=2$:}
\begin{itemize}
  \item $x\to2^+$: $\frac1{x-2}\to+\infty$, ולכן $5^{1/(x-2)}\to+\infty$ ו-$f(x)\to0$.
  \item $x\to2^-$: $\frac1{x-2}\to-\infty$, ולכן $5^{1/(x-2)}\to0$ ו-$f(x)\to\frac1{25}$.
\end{itemize}
שני הגבולות החד-צדדיים קיימים, סופיים ושונים: \textbf{אי-רציפות מסוג ראשון (קפיצה)}.

\textbf{הנקודה $x=\frac52$:} המכנה $25-5^{1/(x-2)}$ שואף ל-$0$. עבור $2<x<\frac52$: $\frac1{x-2}>2$, ולכן $5^{1/(x-2)}>25$ והמכנה שלילי: $f(x)\to-\infty$ כאשר $x\to\frac52^-$. עבור $x>\frac52$ המכנה חיובי: $f(x)\to+\infty$ כאשר $x\to\frac52^+$. הגבולות החד-צדדיים אינסופיים: \textbf{אי-רציפות מסוג שני} (ו-$x=\frac52$ אסימפטוטה אנכית).
\end{solution}

\subsection*{שאלה 1ג}

\begin{question}
מהו שיפוע המשיק לקו $y=\int_0^x\frac{dt}{1+t^4}$ בנקודה עליו $x=1$.
\end{question}

\begin{solution}
הפונקציה $\frac1{1+t^4}$ רציפה על $\R$, ולכן לפי \emph{המשפט היסודי של החשבון האינפיניטסימלי} הפונקציה $y(x)=\int_0^x\frac{dt}{1+t^4}$ גזירה ו-$y'(x)=\frac1{1+x^4}$. שיפוע המשיק ב-$x=1$:
\[
y'(1)=\frac1{1+1}=\frac12.
\]
\end{solution}

\subsection*{שאלה 2א}

\begin{question}
הוכיחו שלכל $x>0$ מתקיים $3\sqrt[3]{x^2}-2x\le1$.
\end{question}

\begin{hint}
מצאו את המקסימום של $g(x)=3x^{2/3}-2x$ בקטע $(0,\infty)$.
\end{hint}

\begin{solution}
נגדיר $g(x)=3x^{2/3}-2x$ עבור $x>0$; היא גזירה שם ו-
\[
g'(x)=2x^{-1/3}-2=2\left(\frac{1}{\sqrt[3]x}-1\right).
\]
$g'(x)>0$ עבור $0<x<1$ ו-$g'(x)<0$ עבור $x>1$. לכן $g$ עולה ב-$(0,1]$ ויורדת ב-$[1,\infty)$ (מסקנה ממשפט לגרנז'), ולכן $x=1$ נקודת מקסימום מוחלט של $g$ ב-$(0,\infty)$:
\[
g(x)\le g(1)=3-2=1\qquad\text{לכל }x>0.
\]
זהו בדיוק $3\sqrt[3]{x^2}-2x\le1$, עם שוויון רק ב-$x=1$. $\blacksquare$
\end{solution}

\subsection*{שאלה 2ב}

\begin{question}
חשבו את השטח הכלוא בין הקווים $y=x\cdot e^x$, $y=x-1$, $x=0$, $x=2$. ציירו את הסקיצה המתאימה.
\end{question}

\begin{hint}
בדקו איזה מהגרפים נמצא מעל השני בקטע $[0,2]$. לאינטגרל של $xe^x$ השתמשו באינטגרציה בחלקים.
\end{hint}

\begin{solution}
\textbf{מי מעל מי.} עבור $x\in[0,2]$: $e^x\ge1$ ו-$x\ge0$, ולכן $xe^x\ge x>x-1$. כלומר הגרף של $y=xe^x$ מעל הישר $y=x-1$ בכל הקטע (ואין נקודות חיתוך בקטע).

\textbf{סקיצה:} הישר $y=x-1$ עובר מ-$(0,-1)$ ל-$(2,1)$; העקום $y=xe^x$ עולה מ-$(0,0)$ ל-$(2,2e^2)\approx(2,14.8)$. התחום חסום משמאל ע"י $x=0$, מימין ע"י $x=2$, מלמטה ע"י הישר ומלמעלה ע"י העקום.

\textbf{שטח:}
\[
S=\int_0^2\big(xe^x-(x-1)\big)dx.
\]
באינטגרציה בחלקים ($u=x$, $dv=e^xdx$): $\int xe^xdx=xe^x-e^x=(x-1)e^x$. לכן
\[
S=\Big[(x-1)e^x\Big]_0^2-\left[\frac{x^2}{2}-x\right]_0^2=\big(e^2-(-1)\big)-(2-2)=e^2+1.
\]
\textbf{תשובה:} $S=e^2+1\approx8.39$.
\end{solution}

\subsection*{שאלה 3א}

\begin{question}
חשבו:
\begin{enumerate}
  \item[(א)] (8 נק') $\displaystyle\int\frac{x^3+2x^2-4}{x^3-2x^2}dx$
  \item[(ב)] (7 נק') $\displaystyle\int\frac{x}{\sqrt{x+1}}dx$
\end{enumerate}
\end{question}

\begin{hint}
ב-(א) חלקו פולינומים ואז פרקו לשברים חלקיים: $x^3-2x^2=x^2(x-2)$. ב-(ב) הציבו $u=x+1$.
\end{hint}

\begin{solution}
\begin{enumerate}
  \item[(א)] \textbf{חילוק פולינומים:} $x^3+2x^2-4=(x^3-2x^2)+4x^2-4$, ולכן
  \[
  \frac{x^3+2x^2-4}{x^3-2x^2}=1+\frac{4x^2-4}{x^2(x-2)}.
  \]
  \textbf{שברים חלקיים:}
  \[
  \frac{4x^2-4}{x^2(x-2)}=\frac Ax+\frac B{x^2}+\frac C{x-2}\iff 4x^2-4=Ax(x-2)+B(x-2)+Cx^2.
  \]
  הצבת $x=2$: $12=4C\Rightarrow C=3$. הצבת $x=0$: $-4=-2B\Rightarrow B=2$. השוואת מקדמי $x^2$: $A+C=4\Rightarrow A=1$.
  (בדיקה: $x^2-2x+2x-4+3x^2=4x^2-4$.) לכן
  \[
  \int\frac{x^3+2x^2-4}{x^3-2x^2}dx=\int\left(1+\frac1x+\frac2{x^2}+\frac3{x-2}\right)dx=x+\ln|x|-\frac2x+3\ln|x-2|+C.
  \]

  \item[(ב)] נציב $u=x+1$, $du=dx$, $x=u-1$:
  \[
  \int\frac{x}{\sqrt{x+1}}dx=\int\frac{u-1}{\sqrt u}du=\int\left(u^{1/2}-u^{-1/2}\right)du=\frac23u^{3/2}-2u^{1/2}+C.
  \]
  \textbf{תשובה:} $\displaystyle\int\frac{x}{\sqrt{x+1}}dx=\frac23(x+1)^{3/2}-2\sqrt{x+1}+C=\frac23(x-2)\sqrt{x+1}+C$.
\end{enumerate}
\end{solution}

\subsection*{שאלה 3ב}

\begin{question}
רשמו פולינום מקלורן מסדר 2 עבור הפונקציה $f(x)=\cos(2x)-\arctan(x+1)$.
\end{question}

\begin{solution}
פולינום מקלורן מסדר 2: $P_2(x)=f(0)+f'(0)x+\frac{f''(0)}{2}x^2$.
\[
f(0)=\cos0-\arctan1=1-\frac\pi4.
\]
\[
f'(x)=-2\sin2x-\frac{1}{1+(x+1)^2},\qquad f'(0)=0-\frac12=-\frac12.
\]
\[
f''(x)=-4\cos2x+\frac{2(x+1)}{\big(1+(x+1)^2\big)^2},\qquad f''(0)=-4+\frac{2}{4}=-\frac72.
\]
\textbf{תשובה:}
\[
P_2(x)=1-\frac\pi4-\frac12x-\frac74x^2.
\]
\end{solution}

\subsection*{שאלה 4א}

\begin{question}
לאילו ערכים של של הפרמטרים $a$ ו-$b$ הפונקציה
\[
f(x)=\begin{cases}\dfrac{\tan(2x)-x}{x+\sin(2x)} & -\frac\pi4<x<0\\[1ex] ax+b & 0\le x\le1\\[0.5ex] x^{1/(x-1)} & x>1\end{cases}
\]
תהיה רציפה בתחום $x>-\frac\pi4$?
\end{question}

\begin{hint}
הפונקציה רציפה בתוך כל אחד משלושת הקטעים; יש לבדוק רק את נקודות התפר $x=0$ ו-$x=1$. ב-$x=0$ חלקו מונה ומכנה ב-$x$; ב-$x=1$ זהו גבול מהצורה $1^\infty$.
\end{hint}

\begin{solution}
\textbf{רציפות בתוך הקטעים.} ב-$\left(-\frac\pi4,0\right)$: $\tan2x$ מוגדרת ורציפה (כי $2x\in\left(-\frac\pi2,0\right)$), והמכנה $x+\sin2x$ שלילי ממש (שני המחוברים שליליים), ולכן $f$ רציפה שם כמנה של רציפות. ב-$(0,1)$ $f$ פולינום. ב-$(1,\infty)$: $x^{1/(x-1)}=e^{\frac{\ln x}{x-1}}$ רציפה. נותר לבדוק את $x=0$ ו-$x=1$.

\textbf{הנקודה $x=0$.} $f(0)=b$, והגבול מימין הוא $b$. משמאל, נחלק מונה ומכנה ב-$x$:
\[
\lim_{x\to0^-}\frac{\tan2x-x}{x+\sin2x}=\lim_{x\to0^-}\frac{2\cdot\frac{\tan2x}{2x}-1}{1+2\cdot\frac{\sin2x}{2x}}=\frac{2-1}{1+2}=\frac13,
\]
לפי הגבולות היסודיים $\frac{\sin t}{t}\to1$, $\frac{\tan t}{t}\to1$. לכן צריך $b=\frac13$.

\textbf{הנקודה $x=1$.} $f(1)=a+b$, והגבול משמאל הוא $a+b$. מימין:
\[
\lim_{x\to1^+}x^{1/(x-1)}=\lim_{x\to1^+}e^{\frac{\ln x}{x-1}}=e^1=e,
\]
כי $\lim_{x\to1}\frac{\ln x}{x-1}=1$ (למשל לפי לופיטל: $\frac{1/x}{1}\to1$, או לפי $\ln(1+t)/t\to1$ עם $t=x-1$), ולפי רציפות האקספוננט. לכן צריך $a+b=e$, כלומר $a=e-\frac13$.

\textbf{תשובה:} $a=e-\frac13$, $b=\frac13$.
\end{solution}

\subsection*{שאלה 4ב}

\begin{question}
חשבו $\displaystyle\int_1^{+\infty}\frac{\ln x}{x^4}dx$.
\end{question}

\begin{hint}
אינטגרציה בחלקים עם $u=\ln x$, $dv=x^{-4}dx$.
\end{hint}

\begin{solution}
\textbf{פונקציה קדומה.} אינטגרציה בחלקים עם $u=\ln x$, $dv=x^{-4}dx$, $du=\frac{dx}x$, $v=-\frac1{3x^3}$:
\[
\int\frac{\ln x}{x^4}dx=-\frac{\ln x}{3x^3}+\frac13\int x^{-4}dx=-\frac{\ln x}{3x^3}-\frac1{9x^3}+C.
\]
\textbf{האינטגרל הלא אמיתי:}
\[
\int_1^{\infty}\frac{\ln x}{x^4}dx=\lim_{R\to\infty}\left[-\frac{\ln x}{3x^3}-\frac{1}{9x^3}\right]_1^R
=\lim_{R\to\infty}\left(-\frac{\ln R}{3R^3}-\frac1{9R^3}\right)+\frac19=\frac19,
\]
כי $\lim_{R\to\infty}\frac{\ln R}{R^3}=0$ (לפי לופיטל: $\frac{1/R}{3R^2}=\frac1{3R^3}\to0$).

\textbf{תשובה:} האינטגרל מתכנס ושווה $\frac19$.
\end{solution}

\section*{תשע"ט סמסטר א מועד א}

\subsection*{שאלה 1א}

\begin{question}
מצאו את הגבול הבא:
$\displaystyle\lim_{x\to0}\big(\cos x\big)^{\frac{1}{\sin^2 2x}}$
\end{question}

\begin{hint}
כתבו $(\cos x)^{1/\sin^22x}=e^{\frac{\ln\cos x}{\sin^22x}}$ וחשבו את גבול המעריך.
\end{hint}

\begin{solution}
זהו גבול מהצורה $1^\infty$. בסביבה מנוקבת קטנה של $0$ מתקיים $\cos x>0$ ו-$\sin2x\neq0$, ולכן
\[ (\cos x)^{\frac1{\sin^22x}}=e^{\frac{\ln\cos x}{\sin^22x}}. \]
נחשב את גבול המעריך. נכתוב
\[ \frac{\ln\cos x}{\sin^22x}=\frac{\ln\cos x}{x^2}\cdot\frac{(2x)^2}{\sin^22x}\cdot\frac14. \]
הגורם האמצעי שואף ל-$1$ לפי $\lim_{t\to0}\frac{\sin t}{t}=1$. עבור הגורם הראשון (צורה $\frac00$) נשתמש בכלל לופיטל:
\[ \lim_{x\to0}\frac{\ln\cos x}{x^2}=\lim_{x\to0}\frac{-\tan x}{2x}=-\frac12. \]
לכן גבול המעריך הוא $-\frac12\cdot1\cdot\frac14=-\frac18$, ומרציפות פונקציית האקספוננט
\[ \lim_{x\to0}(\cos x)^{\frac1{\sin^22x}}=e^{-1/8}. \]
\end{solution}

\subsection*{שאלה 1ב}

\begin{question}
מצאו את הגבול הבא:
$\displaystyle\lim_{x\to\infty}\frac{\sin\left(\frac{1}{\sqrt x}\right)}{\sqrt{x+1}-\sqrt x}$
\end{question}

\begin{hint}
הכפילו בצמוד $\sqrt{x+1}+\sqrt x$ והשתמשו ב-$\lim_{t\to0}\frac{\sin t}{t}=1$ עם $t=\frac1{\sqrt x}$.
\end{hint}

\begin{solution}
נכפיל ונחלק בצמוד:
\[ \frac{\sin\frac1{\sqrt x}}{\sqrt{x+1}-\sqrt x}=\sin\!\left(\tfrac1{\sqrt x}\right)\left(\sqrt{x+1}+\sqrt x\right)
=\frac{\sin\frac1{\sqrt x}}{\frac1{\sqrt x}}\cdot\frac{\sqrt{x+1}+\sqrt x}{\sqrt x}
=\frac{\sin\frac1{\sqrt x}}{\frac1{\sqrt x}}\cdot\left(\sqrt{1+\tfrac1x}+1\right). \]
כאשר $x\to\infty$, $t=\frac1{\sqrt x}\to0^+$ ולכן הגורם הראשון שואף ל-$1$ (גבול יסודי והרכבת גבולות), והגורם השני שואף ל-$1+1=2$. לכן הגבול הוא $\boxed{2}$.
\end{solution}

\subsection*{שאלה 2}

\begin{question}
כדי לחשב $a=\sqrt[3]{29}$ באופן מקורב מציגים
\[ a=3\cdot\sqrt[3]{29/27}=3\cdot\sqrt[3]{1+2/27} \]
ומשתמשים בפולינום מקלורין ממעלה 2 של פונקציה $f(x)=3\cdot(1+x)^{1/3}$. העריכו את שגיאת הקירוב בעזרת נוסחת טיילור.
\end{question}

\begin{hint}
השארית בצורת לגרנז' היא $R_2(x)=\frac{f'''(c)}{3!}x^3$ עבור $c$ בין $0$ ל-$x$. חסמו את $|f'''(c)|$ עבור $c\in(0,\frac2{27})$.
\end{hint}

\begin{solution}
\textbf{פולינום מקלורין.} נגזור את $f(x)=3(1+x)^{1/3}$:
\[ f'(x)=(1+x)^{-2/3},\qquad f''(x)=-\frac23(1+x)^{-5/3},\qquad f'''(x)=\frac{10}{9}(1+x)^{-8/3}. \]
לכן $f(0)=3$, $f'(0)=1$, $f''(0)=-\frac23$ ו-
\[ P_2(x)=3+x-\frac13x^2. \]
\textbf{הקירוב.} עבור $x=\frac2{27}$:
\[ a\approx P_2\!\left(\tfrac2{27}\right)=3+\frac2{27}-\frac13\cdot\frac{4}{729}=3+\frac{162}{2187}-\frac{4}{2187}=\frac{6719}{2187}\approx3.072245. \]
\textbf{הערכת השגיאה.} לפי נוסחת טיילור עם שארית לגרנז', קיים $c\in(0,\frac2{27})$ כך ש-
\[ a-P_2\!\left(\tfrac2{27}\right)=R_2\!\left(\tfrac2{27}\right)=\frac{f'''(c)}{3!}\left(\frac2{27}\right)^3=\frac{10}{54}(1+c)^{-8/3}\cdot\frac{8}{19683}. \]
מכיוון ש-$c>0$ מתקיים $0<(1+c)^{-8/3}<1$, ולכן
\[ 0<R_2<\frac{10}{54}\cdot\frac{8}{19683}=\frac{40}{531441}\approx7.53\cdot10^{-5}. \]
כלומר השגיאה קטנה מ-$7.6\cdot10^{-5}$ (ובפרט מ-$10^{-4}$), והקירוב הוא קירוב מלמטה (השארית חיובית). לשם השוואה, $\sqrt[3]{29}=3.0723168\ldots$, והשגיאה בפועל היא כ-$7.17\cdot10^{-5}$.
\end{solution}

\subsection*{שאלה 3א}

\begin{question}
תהיה פונקציה $f(x)$ מוגדרת בסביבת נקודה $x_0$. השלימו את ההגדרה הבאה:

\textbf{מספר $a$ הוא גבול של $f(x)$ כאשר $x$ שואף ל-$x_0$ אם לכל \ldots}
\end{question}

\begin{solution}
(הגדרת קושי, $\eps$--$\delta$) מספר $a$ הוא גבול של $f(x)$ כאשר $x$ שואף ל-$x_0$ אם לכל $\eps>0$ קיים $\delta>0$ כך שלכל $x$ המקיים $0<|x-x_0|<\delta$ מתקיים $|f(x)-a|<\eps$.

(הגדרה שקולה לפי היינה: לכל סדרה $x_n\to x_0$ עם $x_n\neq x_0$ מתקיים $f(x_n)\to a$.)
\end{solution}

\subsection*{שאלה 3ב}

\begin{question}
הסבירו מדוע לפונקציה $f(x)=\cos\left(\dfrac{x+1}{x-1}\right)$ אין גבול כאשר $x$ שואף ל-$1$.
\end{question}

\begin{hint}
השתמשו בהגדרת הגבול לפי היינה: מצאו שתי סדרות השואפות ל-$1$ שעליהן ערכי הפונקציה שואפים לגבולות שונים.
\end{hint}

\begin{solution}
נשתמש בהגדרת היינה. אם היה קיים $L=\lim_{x\to1}f(x)$, אז לכל סדרה $x_n\to1$, $x_n\ne1$, היה מתקיים $f(x_n)\to L$.

נשים לב שאם $\frac{x+1}{x-1}=y$ אז $x=\frac{y+1}{y-1}$. נבחר
\[ x_n=\frac{2\pi n+1}{2\pi n-1},\qquad z_n=\frac{(2n+1)\pi+1}{(2n+1)\pi-1}. \]
שתי הסדרות שואפות ל-$1$ (מחלקים מונה ומכנה ב-$n$) ושונות מ-$1$. מתקיים $\frac{x_n+1}{x_n-1}=2\pi n$ ו-$\frac{z_n+1}{z_n-1}=(2n+1)\pi$, ולכן
\[ f(x_n)=\cos(2\pi n)=1\to1,\qquad f(z_n)=\cos\big((2n+1)\pi\big)=-1\to-1. \]
קיבלנו שתי סדרות השואפות ל-$1$ שעליהן ערכי הפונקציה מתכנסים לגבולות שונים, בסתירה ליחידות הגבול. לכן ל-$f$ אין גבול ב-$1$.

(באופן אינטואיטיבי: כאשר $x\to1^\pm$ הביטוי $\frac{x+1}{x-1}\to\pm\infty$, והקוסינוס מתנודד בין $-1$ ל-$1$ אינסוף פעמים.)
\end{solution}

\subsection*{שאלה 4}

\begin{question}
חקרו את הפונקציה $y=(x+2)e^{1/x}$ ושרטטו את הגרף שלה.

\underline{צריך לחקור}: תחום הגדרה, רציפות, אסימפטוטות, תחומי עליה וירידה, נקודות קיצון, תחומי קמירות וקעירות.
\end{question}

\begin{hint}
לאסימפטוטה המשופעת: $m=\lim\frac{y}{x}$, $n=\lim(y-mx)$; בחישוב $n$ הציבו $t=\frac1x$ והשתמשו ב-$\lim_{t\to0}\frac{e^t-1}{t}=1$.
\end{hint}

\begin{hint}
בדקו בנפרד את הגבולות החד-צדדיים ב-$x=0$: כאשר $x\to0^-$, $e^{1/x}\to0$.
\end{hint}

\begin{solution}
\textbf{תחום הגדרה ורציפות.} $x\neq0$. בתחום הפונקציה רציפה (מכפלה והרכבה של פונקציות רציפות). נקודת חיתוך עם ציר $x$: $x=-2$; $y>0$ עבור $x>-2$, $x\ne0$, ו-$y<0$ עבור $x<-2$.

\textbf{אסימפטוטות אנכיות.} כאשר $x\to0^+$: $\frac1x\to+\infty$, $e^{1/x}\to+\infty$ ו-$x+2\to2$, לכן $y\to+\infty$: הישר $x=0$ אסימפטוטה אנכית (מימין). כאשר $x\to0^-$: $e^{1/x}\to0$ ולכן $y\to0$ – משמאל אין אסימפטוטה (הגרף "נכנס" לראשית; ב-$x=0$ אי-רציפות עיקרית).

\textbf{אסימפטוטות משופעות.} $m=\lim_{x\to\pm\infty}\frac{(x+2)e^{1/x}}{x}=\lim\left(1+\frac2x\right)e^{1/x}=1$. כעת
\[ y-x=x\big(e^{1/x}-1\big)+2e^{1/x}=\frac{e^{1/x}-1}{1/x}+2e^{1/x}\xrightarrow[x\to\pm\infty]{}1+2=3, \]
לפי $\lim_{t\to0}\frac{e^t-1}{t}=1$ עם $t=\frac1x\to0$. לכן $y=x+3$ אסימפטוטה משופעת בשני הכיוונים.

\textbf{עליה וירידה.}
\[ y'=e^{1/x}+(x+2)e^{1/x}\cdot\left(-\frac1{x^2}\right)=e^{1/x}\,\frac{x^2-x-2}{x^2}=e^{1/x}\,\frac{(x-2)(x+1)}{x^2}. \]
סימן $y'$ כסימן $(x-2)(x+1)$: $y'>0$ ב-$(-\infty,-1)$ וב-$(2,\infty)$ – הפונקציה עולה; $y'<0$ ב-$(-1,0)$ וב-$(0,2)$ – הפונקציה יורדת.

\textbf{נקודות קיצון.} ב-$x=-1$ הנגזרת מחליפה סימן מ-$+$ ל-$-$: מקסימום מקומי $y(-1)=e^{-1}=\frac1e$. ב-$x=2$ מ-$-$ ל-$+$: מינימום מקומי $y(2)=4\sqrt e\approx6.59$.

\textbf{קמירות וקעירות.} חישוב ישיר (גזירת $e^{1/x}(1-\frac1x-\frac2{x^2})$) נותן
\[ y''=e^{1/x}\left(-\frac1{x^2}\right)\left(1-\frac1x-\frac2{x^2}\right)+e^{1/x}\left(\frac1{x^2}+\frac4{x^3}\right)=e^{1/x}\,\frac{5x+2}{x^4}. \]
סימן $y''$ כסימן $5x+2$: $y''<0$ (קעורה, $\cap$) ב-$(-\infty,-\frac25)$, ו-$y''>0$ (קמורה, $\cup$) ב-$(-\frac25,0)$ וב-$(0,\infty)$. נקודת פיתול: $x=-\frac25$, $y\left(-\frac25\right)=\frac85e^{-5/2}\approx0.131$.

\textbf{סקיצה.} מהפיתוח $y=x+3+\frac{5}{2x}+o\!\left(\frac1x\right)$ נובע שהגרף נמצא מתחת לאסימפטוטה כאשר $x\to-\infty$ ומעליה כאשר $x\to+\infty$. משמאל: הגרף מתקרב מלמטה לישר $y=x+3$, עולה וחוצה את ציר $x$ ב-$x=-2$, מגיע למקסימום $(-1,\frac1e)$, יורד דרך נקודת הפיתול $(-\frac25,\frac85e^{-5/2})$ ושואף ל-$0$ כאשר $x\to0^-$. מימין ל-$0$ הגרף יורד מ-$+\infty$ (אסימפטוטה $x=0$) עד המינימום $(2,4\sqrt e)$, ואז עולה ומתקרב לאסימפטוטה $y=x+3$ כאשר $x\to+\infty$.
\end{solution}

\subsection*{שאלה 5א}

\begin{question}
חשבו את האינטגרל הבא:
$\displaystyle\int e^{2x}\sin(3x)\,dx$
\end{question}

\begin{hint}
בצעו אינטגרציה בחלקים פעמיים; תקבלו משוואה שבה האינטגרל המבוקש מופיע בשני האגפים.
\end{hint}

\begin{solution}
נסמן $I=\int e^{2x}\sin3x\,dx$. אינטגרציה בחלקים עם $u=\sin3x$, $dv=e^{2x}dx$ ($v=\frac12e^{2x}$):
\[ I=\frac12e^{2x}\sin3x-\frac32\int e^{2x}\cos3x\,dx. \]
שוב בחלקים, $u=\cos3x$, $dv=e^{2x}dx$:
\[ \int e^{2x}\cos3x\,dx=\frac12e^{2x}\cos3x+\frac32\int e^{2x}\sin3x\,dx=\frac12e^{2x}\cos3x+\frac32I. \]
לכן
\[ I=\frac12e^{2x}\sin3x-\frac34e^{2x}\cos3x-\frac94I\;\Longrightarrow\;\frac{13}{4}I=\frac{e^{2x}}{4}\big(2\sin3x-3\cos3x\big), \]
ולבסוף
\[ \int e^{2x}\sin3x\,dx=\frac{e^{2x}\big(2\sin3x-3\cos3x\big)}{13}+C. \]
(בדיקה: גזירת התוצאה נותנת $\frac{e^{2x}}{13}\big(4\sin3x-6\cos3x+6\cos3x+9\sin3x\big)=e^{2x}\sin3x$.)
\end{solution}

\subsection*{שאלה 5ב}

\begin{question}
חשבו את האינטגרל הבא:
$\displaystyle\int_0^1\frac{x^4}{x^2+2x+1}\,dx$
\end{question}

\begin{hint}
$x^2+2x+1=(x+1)^2$; הציבו $u=x+1$ (או בצעו חילוק פולינומים).
\end{hint}

\begin{solution}
$x^2+2x+1=(x+1)^2$. נציב $u=x+1$, $x=u-1$, $u\in[1,2]$:
\[ \int_0^1\frac{x^4}{(x+1)^2}\,dx=\int_1^2\frac{(u-1)^4}{u^2}\,du=\int_1^2\frac{u^4-4u^3+6u^2-4u+1}{u^2}\,du=\int_1^2\left(u^2-4u+6-\frac4u+\frac1{u^2}\right)du. \]
\[ =\left[\frac{u^3}3-2u^2+6u-4\ln u-\frac1u\right]_1^2=\left(\frac83-8+12-4\ln2-\frac12\right)-\left(\frac13-2+6-1\right)=\left(\frac{37}{6}-4\ln2\right)-\frac{10}{3}=\frac{17}{6}-4\ln2. \]
\textbf{תשובה:} $\dfrac{17}{6}-4\ln2\approx0.0607$.
\end{solution}

\subsection*{שאלה 6}

\begin{question}
האם האינטגרל הלא אמיתי הבא $\displaystyle\int_0^1\frac{\ln x}{\sqrt x}\,dx$ מתכנס? אם כן, חשבו את האינטגרל.
\end{question}

\begin{hint}
הנקודה הבעייתית היא $x=0$. חשבו $\int_\eps^1$ באינטגרציה בחלקים ($u=\ln x$, $dv=x^{-1/2}dx$) והעבירו גבול $\eps\to0^+$.
\end{hint}

\begin{solution}
הפונקציה $\frac{\ln x}{\sqrt x}$ רציפה ב-$(0,1]$ ולא חסומה ליד $0$ (שואפת ל-$-\infty$), ולכן זהו אינטגרל לא אמיתי מסוג שני, המוגדר כ-$\lim_{\eps\to0^+}\int_\eps^1\frac{\ln x}{\sqrt x}\,dx$.

באינטגרציה בחלקים, $u=\ln x$, $dv=x^{-1/2}dx$, $du=\frac{dx}x$, $v=2\sqrt x$:
\[ \int_\eps^1\frac{\ln x}{\sqrt x}\,dx=\Big[2\sqrt x\ln x\Big]_\eps^1-\int_\eps^1\frac{2\sqrt x}{x}\,dx=-2\sqrt\eps\ln\eps-\Big[4\sqrt x\Big]_\eps^1=-2\sqrt\eps\ln\eps-4+4\sqrt\eps. \]
כעת $\lim_{\eps\to0^+}\sqrt\eps\ln\eps=\lim_{\eps\to0^+}\frac{\ln\eps}{\eps^{-1/2}}$, ולפי לופיטל (צורה $\frac{-\infty}{\infty}$):
\[ \lim_{\eps\to0^+}\frac{1/\eps}{-\frac12\eps^{-3/2}}=\lim_{\eps\to0^+}\left(-2\sqrt\eps\right)=0. \]
גם $4\sqrt\eps\to0$. לכן הגבול קיים וסופי: האינטגרל \textbf{מתכנס}, ו-
\[ \int_0^1\frac{\ln x}{\sqrt x}\,dx=-4. \]
\end{solution}

\section*{תשע"ט סמסטר א מועד ב}

\subsection*{שאלה 1א}

\begin{question}
מצאו את הגבול הבא:
$\displaystyle\lim_{x\to0}\frac{\cos x-\sqrt{1-2x^2}}{\sin^2x}$
\end{question}

\begin{hint}
השתמשו בפיתוחי מקלורן $\cos x=1-\frac{x^2}2+o(x^2)$, $\sqrt{1+t}=1+\frac t2+o(t)$, ו-$\sin^2x\sim x^2$.
\end{hint}

\begin{solution}
לפי פיתוחי מקלורן:
\[ \cos x=1-\frac{x^2}{2}+o(x^2),\qquad \sqrt{1-2x^2}=1+\frac12(-2x^2)+o(x^2)=1-x^2+o(x^2). \]
לכן המונה הוא $\cos x-\sqrt{1-2x^2}=\frac{x^2}{2}+o(x^2)$. כמו כן $\frac{\sin^2x}{x^2}\to1$. מכאן
\[ \lim_{x\to0}\frac{\cos x-\sqrt{1-2x^2}}{\sin^2x}=\lim_{x\to0}\frac{\frac12+\frac{o(x^2)}{x^2}}{\frac{\sin^2x}{x^2}}=\frac12. \]
(אפשר גם בלופיטל: אחרי החלפת $\sin^2x$ ב-$x^2$, $\lim\frac{-\sin x+\frac{2x}{\sqrt{1-2x^2}}}{2x}=\frac{-1+2}{2}=\frac12$.)
\end{solution}

\subsection*{שאלה 1ב}

\begin{question}
מצאו את הגבול הבא:
$\displaystyle\lim_{x\to0}\left(\frac{1-4x^2}{1-3x^2}\right)^{\frac{2}{x^2}}$

רמז: $f^g=e^{g\ln f}$.
\end{question}

\begin{hint}
אחרי המעבר לצורה $e^{g\ln f}$, כתבו $\ln\frac{1-4x^2}{1-3x^2}=\ln(1-4x^2)-\ln(1-3x^2)$ והשתמשו בגבול היסודי $\lim_{t\to0}\frac{\ln(1+t)}{t}=1$.
\end{hint}

\begin{solution}
זהו גבול מהצורה $1^\infty$. בסביבה של $0$ הבסיס חיובי, ולפי הרמז
\[ \left(\frac{1-4x^2}{1-3x^2}\right)^{\frac2{x^2}}=\exp\!\left(\frac{2}{x^2}\ln\frac{1-4x^2}{1-3x^2}\right)=\exp\!\left(2\cdot\frac{\ln(1-4x^2)-\ln(1-3x^2)}{x^2}\right). \]
לפי הגבול היסודי $\lim_{t\to0}\frac{\ln(1+t)}{t}=1$:
\[ \frac{\ln(1-4x^2)}{x^2}=-4\cdot\frac{\ln(1-4x^2)}{-4x^2}\to-4,\qquad \frac{\ln(1-3x^2)}{x^2}\to-3. \]
לכן המעריך שואף ל-$2(-4+3)=-2$, ומרציפות האקספוננט הגבול הוא $e^{-2}$.
\end{solution}

\subsection*{שאלה 2}

\begin{question}
תהי $f(x)=\ln x$. רשמו פולינום טיילור ממעלה 3 של פונקציה $f$ בסביבת נקודה $x_0=1$. השתמשו בפולינום זה כדי לחשב קרוב ל-$\ln(1.1)$. העריכו את שגיאת הקרוב הזה בעזרת נוסחת טיילור.
\end{question}

\begin{hint}
השארית בצורת לגרנז' היא $R_3(x)=\frac{f^{(4)}(c)}{4!}(x-1)^4$ עם $c$ בין $1$ ל-$x$.
\end{hint}

\begin{solution}
\textbf{הפולינום.} $f(x)=\ln x$, $f'(x)=\frac1x$, $f''(x)=-\frac1{x^2}$, $f'''(x)=\frac2{x^3}$, $f^{(4)}(x)=-\frac6{x^4}$. בנקודה $1$: $f(1)=0$, $f'(1)=1$, $f''(1)=-1$, $f'''(1)=2$. לכן
\[ P_3(x)=(x-1)-\frac{(x-1)^2}{2}+\frac{(x-1)^3}{3}. \]
\textbf{הקירוב.} עבור $x=1.1$, $x-1=0.1$:
\[ \ln(1.1)\approx P_3(1.1)=0.1-\frac{0.01}{2}+\frac{0.001}{3}=0.1-0.005+0.000\overline{3}=0.095\overline{3}. \]
\textbf{הערכת השגיאה.} לפי נוסחת טיילור עם שארית לגרנז', קיים $c\in(1,1.1)$ כך ש-
\[ R_3(1.1)=\frac{f^{(4)}(c)}{4!}(0.1)^4=-\frac{6}{24c^4}\cdot10^{-4}=-\frac{10^{-4}}{4c^4}. \]
מכיוון ש-$c>1$, $c^4>1$ ולכן
\[ |R_3(1.1)|<\frac{10^{-4}}{4}=2.5\cdot10^{-5}. \]
כמו כן $R_3<0$, כלומר הקירוב $0.09533$ גדול מעט מהערך האמיתי. (לשם השוואה $\ln1.1=0.0953102\ldots$, והשגיאה בפועל כ-$2.3\cdot10^{-5}$.)
\end{solution}

\subsection*{שאלה 3א}

\begin{question}
האם פונקציה
\[ f(x)=\begin{cases}\dfrac{1}{\tan x}-\dfrac1x & x\neq0\\[2mm] 0 & x=0\end{cases} \]
גזירה בנקודה $x=0$? הסבירו.
\end{question}

\begin{hint}
חשבו את גבול מנת ההפרשים $\frac{f(h)-f(0)}{h}=\frac{h-\tan h}{h^2\tan h}$, למשל בעזרת פיתוח $\tan h=h+\frac{h^3}{3}+o(h^3)$.
\end{hint}

\begin{solution}
(הפונקציה מוגדרת בסביבה מנוקבת $0<|x|<\frac\pi2$ של $0$, ושם $\tan x\ne0$.) לפי הגדרת הנגזרת:
\[ f'(0)=\lim_{h\to0}\frac{f(h)-f(0)}{h}=\lim_{h\to0}\frac{\frac1{\tan h}-\frac1h}{h}=\lim_{h\to0}\frac{h-\tan h}{h^2\tan h}. \]
לפי פיתוח מקלורן $\tan h=h+\frac{h^3}{3}+o(h^3)$, ולכן $h-\tan h=-\frac{h^3}{3}+o(h^3)$. כמו כן $h^2\tan h=h^3\cdot\frac{\tan h}{h}$ ו-$\frac{\tan h}{h}\to1$. לכן
\[ f'(0)=\lim_{h\to0}\frac{-\frac13+\frac{o(h^3)}{h^3}}{\frac{\tan h}{h}}=-\frac13. \]
(אפשר גם בלופיטל: $\lim\frac{h-\tan h}{h^3}=\lim\frac{1-\frac1{\cos^2h}}{3h^2}=\lim\frac{-\tan^2h}{3h^2}=-\frac13$.)

\textbf{תשובה:} כן, $f$ גזירה ב-$0$ ו-$f'(0)=-\frac13$. (בפרט $f$ רציפה ב-$0$: $\lim_{x\to0}f(x)=\lim\frac{x-\tan x}{x\tan x}=0=f(0)$.)
\end{solution}

\subsection*{שאלה 3ב}

\begin{question}
הוכיחו או הפריכו: אם פונקציה $f(x)$ גזירה בנקודה $x_0$ היא רציפה ב-$x_0$.
\end{question}

\begin{hint}
כתבו $f(x)-f(x_0)=\frac{f(x)-f(x_0)}{x-x_0}\cdot(x-x_0)$.
\end{hint}

\begin{solution}
\textbf{הטענה נכונה.} נניח ש-$f$ גזירה ב-$x_0$, כלומר הגבול $f'(x_0)=\lim_{x\to x_0}\frac{f(x)-f(x_0)}{x-x_0}$ קיים וסופי. לכל $x\ne x_0$ בסביבת $x_0$:
\[ f(x)-f(x_0)=\frac{f(x)-f(x_0)}{x-x_0}\cdot(x-x_0). \]
לפי אריתמטיקה של גבולות (מכפלת שני גבולות סופיים):
\[ \lim_{x\to x_0}\big(f(x)-f(x_0)\big)=f'(x_0)\cdot0=0, \]
כלומר $\lim_{x\to x_0}f(x)=f(x_0)$, ו-$f$ רציפה ב-$x_0$. $\blacksquare$

(הכיוון ההפוך אינו נכון: $|x|$ רציפה ב-$0$ אך אינה גזירה שם.)
\end{solution}

\subsection*{שאלה 4}

\begin{question}
חקרו את הפונקציה $f(x)=\dfrac{x^3}{x^2-4}$ ושרטטו את הגרף שלה.

\underline{צריך לחקור}: תחום הגדרה, רציפות, אסימפטוטות, תחומי עליה וירידה, נקודות קיצון, תחומי קמירות וקעירות.
\end{question}

\begin{hint}
שימו לב שהפונקציה אי-זוגית. חילוק פולינומים נותן $f(x)=x+\frac{4x}{x^2-4}$.
\end{hint}

\begin{solution}
\textbf{תחום הגדרה ורציפות.} $x\ne\pm2$. בתחום $f$ רציפה (פונקציה רציונלית). $f(-x)=-f(x)$, כלומר $f$ אי-זוגית והגרף סימטרי ביחס לראשית. חיתוך עם הצירים: רק $(0,0)$.

\textbf{אסימפטוטות.} אנכיות: ב-$x=\pm2$ המונה $\pm8\ne0$ והמכנה מתאפס, לכן:
\[ \lim_{x\to2^\pm}f(x)=\pm\infty,\qquad \lim_{x\to-2^\pm}f(x)=\pm\infty \]
(לדוגמה ליד $-2$: המונה $\approx-8$, והמכנה $x^2-4$ שלילי עבור $x\to-2^+$ וחיובי עבור $x\to-2^-$). לכן $x=2$ ו-$x=-2$ אסימפטוטות אנכיות. משופעת: לפי חילוק פולינומים
\[ f(x)=x+\frac{4x}{x^2-4},\qquad \frac{4x}{x^2-4}\xrightarrow[x\to\pm\infty]{}0, \]
ולכן $y=x$ אסימפטוטה משופעת בשני הכיוונים.

\textbf{עליה וירידה.}
\[ f'(x)=\frac{3x^2(x^2-4)-x^3\cdot2x}{(x^2-4)^2}=\frac{x^2(x^2-12)}{(x^2-4)^2}. \]
נקודות חשודות: $x=0$, $x=\pm2\sqrt3$. $f'>0$ עבור $|x|>2\sqrt3$ – $f$ עולה ב-$(-\infty,-2\sqrt3)$ וב-$(2\sqrt3,\infty)$. $f'<0$ עבור $|x|<2\sqrt3$, $x\ne0,\pm2$ – $f$ יורדת ב-$(-2\sqrt3,-2)$, ב-$(-2,2)$ (כולל $x=0$, שבה $f'=0$ אך הסימן לא מתחלף) וב-$(2,2\sqrt3)$.

\textbf{קיצון.} $x=-2\sqrt3$: מקסימום מקומי, $f(-2\sqrt3)=\frac{-24\sqrt3}{8}=-3\sqrt3$. $x=2\sqrt3$: מינימום מקומי, $f(2\sqrt3)=3\sqrt3$. ב-$x=0$ אין קיצון.

\textbf{קמירות וקעירות.}
\[ f''(x)=\frac{8x(x^2+12)}{(x^2-4)^3}. \]
מכיוון ש-$x^2+12>0$, סימן $f''$ כסימן $\frac{x}{(x^2-4)^3}$, כלומר כסימן $x(x^2-4)$:
\begin{itemize}
  \item $x<-2$: שלילי – קעורה ($\cap$);
  \item $-2<x<0$: חיובי – קמורה ($\cup$);
  \item $0<x<2$: שלילי – קעורה ($\cap$);
  \item $x>2$: חיובי – קמורה ($\cup$).
\end{itemize}
נקודת פיתול: $(0,0)$ (שם $f''$ מחליפה סימן ו-$f$ מוגדרת). ב-$x=\pm2$ $f$ אינה מוגדרת.

\textbf{סקיצה.} ב-$(-\infty,-2)$: הגרף עולה מתחת לאסימפטוטה $y=x$ עד המקסימום $(-2\sqrt3,-3\sqrt3)\approx(-3.46,-5.2)$ ואז יורד ל-$-\infty$ ליד $x=-2$. ב-$(-2,2)$: יורד מ-$+\infty$ דרך נקודת הפיתול $(0,0)$ (עם משיק אופקי) אל $-\infty$. ב-$(2,\infty)$: יורד מ-$+\infty$ עד המינימום $(2\sqrt3,3\sqrt3)\approx(3.46,5.2)$ ואז עולה ומתקרב מלמעלה לאסימפטוטה $y=x$.
\end{solution}

\subsection*{שאלה 5}

\begin{question}
חשבו את האינטגרל הבא:
\[ \int_0^\pi x^2\sin(3x)\,dx \]
\end{question}

\begin{hint}
בצעו אינטגרציה בחלקים פעמיים, בכל פעם גוזרים את החזקה של $x$.
\end{hint}

\begin{solution}
נבצע אינטגרציה בחלקים עם $u=x^2$, $dv=\sin3x\,dx$ ($du=2x\,dx$, $v=-\frac13\cos3x$):
\[ \int_0^\pi x^2\sin3x\,dx=\left[-\frac{x^2\cos3x}{3}\right]_0^\pi+\frac23\int_0^\pi x\cos3x\,dx. \]
מכיוון ש-$\cos3\pi=-1$: $\left[-\frac{x^2\cos3x}{3}\right]_0^\pi=\frac{\pi^2}{3}$.

שוב בחלקים, $u=x$, $dv=\cos3x\,dx$ ($v=\frac13\sin3x$):
\[ \int_0^\pi x\cos3x\,dx=\left[\frac{x\sin3x}{3}\right]_0^\pi-\frac13\int_0^\pi\sin3x\,dx=0-\frac13\left[-\frac{\cos3x}{3}\right]_0^\pi=-\frac13\cdot\frac{1+1}{3}=-\frac29, \]
(כי $\sin3\pi=0$ ו-$-\cos3\pi+\cos0=2$). לכן
\[ \int_0^\pi x^2\sin3x\,dx=\frac{\pi^2}{3}+\frac23\cdot\left(-\frac29\right)=\frac{\pi^2}{3}-\frac{4}{27}\approx3.142. \]
\end{solution}

\subsection*{שאלה 6א}

\begin{question}
הוכיחו כי למשוואה $e^{2x}-x-5=0$ יש בדיוק שני פתרונות ממשיים.
\end{question}

\begin{hint}
בדקו את המונוטוניות של $g(x)=e^{2x}-x-5$ בעזרת הנגזרת, ואת ערכה בנקודת המינימום. קיום – לפי משפט ערך הביניים, יחידות בכל קטע – לפי מונוטוניות.
\end{hint}

\begin{solution}
נגדיר $g(x)=e^{2x}-x-5$, רציפה וגזירה בכל $\R$, עם $g'(x)=2e^{2x}-1$. $g'(x)=0\iff e^{2x}=\frac12\iff x_0=-\frac{\ln2}{2}$. עבור $x<x_0$: $g'<0$, $g$ יורדת ממש; עבור $x>x_0$: $g'>0$, $g$ עולה ממש. ערך המינימום:
\[ g(x_0)=\frac12+\frac{\ln2}{2}-5\approx-4.15<0. \]
\textbf{קיום:} $\lim_{x\to-\infty}g(x)=+\infty$ (כי $e^{2x}\to0$ ו-$-x\to+\infty$), ולמשל $g(-6)=e^{-12}+1>0$. כמו כן $g(2)=e^4-7>0$. לפי משפט ערך הביניים יש שורש ב-$(-6,x_0)$ ושורש ב-$(x_0,2)$.

\textbf{יחידות:} בכל אחד מהקטעים $(-\infty,x_0]$ ו-$[x_0,\infty)$ הפונקציה מונוטונית ממש ולכן חח"ע, כך שיש בכל אחד לכל היותר שורש אחד (ו-$x_0$ עצמו אינו שורש). לכן למשוואה בדיוק שני פתרונות ממשיים.
\end{solution}

\subsection*{שאלה 6ב}

\begin{question}
נתונה עקומה בקואורדינאטות קוטביות $r=1+\cos\varphi$ $(0\le\varphi\le2\pi)$. חשבו את האורך העקומה.
\end{question}

\begin{hint}
אורך עקומה קוטבית הוא $L=\int_\alpha^\beta\sqrt{r^2+(r')^2}\,d\varphi$, והשתמשו ב-$1+\cos\varphi=2\cos^2\frac\varphi2$.
\end{hint}

\begin{solution}
העקומה (קרדיואידה) היא $x=r\cos\varphi$, $y=r\sin\varphi$, ונוסחת האורך בקואורדינטות קוטביות היא
\[ L=\int_0^{2\pi}\sqrt{r^2+\left(\frac{dr}{d\varphi}\right)^2}\,d\varphi. \]
כאן $r'=-\sin\varphi$ ולכן
\[ r^2+r'^2=(1+\cos\varphi)^2+\sin^2\varphi=2+2\cos\varphi=4\cos^2\frac\varphi2. \]
מכאן $\sqrt{r^2+r'^2}=2\left|\cos\frac\varphi2\right|$. עבור $\varphi\in[0,2\pi]$, $\frac\varphi2\in[0,\pi]$; ה-$\cos$ אי-שלילי ב-$[0,\pi]$ (לגבי $\varphi$) ואי-חיובי ב-$[\pi,2\pi]$. מהסימטריה:
\[ L=2\int_0^{\pi}2\cos\frac\varphi2\,d\varphi=4\left[2\sin\frac\varphi2\right]_0^\pi=8. \]
\textbf{תשובה:} אורך העקומה הוא $8$.
\end{solution}

\section*{תשע"ט סמסטר א מועד מיוחד}

\subsection*{שאלה 1א}

\begin{question}
מצאו את הגבול הבא:
$\displaystyle \lim_{x\to\infty}\left(1+2x-4x^2+x^4\right)^{1/\ln(2x^2+1)}$.
\end{question}

\begin{hint}
כתבו את הביטוי בצורה $e^{\ln(\cdots)/\ln(2x^2+1)}$ והוציאו את החזקה הגבוהה ביותר מתוך כל לוגריתם.
\end{hint}

\begin{solution}
נסמן $p(x)=1+2x-4x^2+x^4$. עבור $x$ גדול מספיק $p(x)>0$ (כי $p(x)\to\infty$), ולכן
\[ p(x)^{1/\ln(2x^2+1)} = \exp\left(\frac{\ln p(x)}{\ln(2x^2+1)}\right). \]
נוציא את החזקה המובילה מכל לוגריתם:
\[ \ln p(x) = \ln\left(x^4\left(1+\tfrac{2}{x^3}-\tfrac{4}{x^2}+\tfrac{1}{x^4}\right)\right) = 4\ln x + \ln\left(1+\tfrac{2}{x^3}-\tfrac{4}{x^2}+\tfrac{1}{x^4}\right), \]
\[ \ln(2x^2+1) = 2\ln x + \ln\left(2+\tfrac{1}{x^2}\right). \]
כאשר $x\to\infty$, הביטויים $\ln\left(1+\tfrac{2}{x^3}-\tfrac{4}{x^2}+\tfrac{1}{x^4}\right)\to \ln 1 = 0$ ו-$\ln\left(2+\tfrac1{x^2}\right)\to\ln 2$ (רציפות הלוגריתם), ואילו $\ln x\to\infty$. נחלק מונה ומכנה ב-$\ln x$:
\[ \frac{\ln p(x)}{\ln(2x^2+1)} = \frac{4 + \frac{\ln\left(1+\frac{2}{x^3}-\frac{4}{x^2}+\frac{1}{x^4}\right)}{\ln x}}{2+\frac{\ln\left(2+\frac{1}{x^2}\right)}{\ln x}} \xrightarrow[x\to\infty]{} \frac{4+0}{2+0}=2. \]
(אפשר גם להשתמש בכלל לופיטל במקרה $\frac{\infty}{\infty}$: $\frac{p'(x)/p(x)}{4x/(2x^2+1)}=\frac{(4x^3-8x+2)(2x^2+1)}{4x\,p(x)}\to\frac{8}{4}=2$.)

מרציפות פונקציית האקספוננט נקבל
\[ \lim_{x\to\infty}\left(1+2x-4x^2+x^4\right)^{1/\ln(2x^2+1)} = e^{2}. \]
\textbf{תשובה:} $e^2$.
\end{solution}

\subsection*{שאלה 1ב}

\begin{question}
מצאו את הגבול הבא:
$\displaystyle \lim_{n\to\infty}\left(\sqrt{(n^2+1)(n^2-4)}-\sqrt{n^4-9}\right)$.
\end{question}

\begin{hint}
הכפילו וחלקו בצמוד $\sqrt{(n^2+1)(n^2-4)}+\sqrt{n^4-9}$.
\end{hint}

\begin{solution}
זהו ביטוי מהצורה $\infty-\infty$. נכפיל ונחלק בצמוד (עבור $n\ge 2$ שני השורשים מוגדרים וחיוביים):
\[ \sqrt{(n^2+1)(n^2-4)}-\sqrt{n^4-9} = \frac{(n^2+1)(n^2-4)-(n^4-9)}{\sqrt{(n^2+1)(n^2-4)}+\sqrt{n^4-9}}. \]
במונה: $(n^2+1)(n^2-4) = n^4-3n^2-4$, ולכן המונה הוא $n^4-3n^2-4-n^4+9 = -3n^2+5$. נחלק מונה ומכנה ב-$n^2$:
\[ \frac{-3+\frac{5}{n^2}}{\sqrt{\left(1+\frac{1}{n^2}\right)\left(1-\frac{4}{n^2}\right)}+\sqrt{1-\frac{9}{n^4}}} \xrightarrow[n\to\infty]{} \frac{-3}{1+1} = -\frac{3}{2}, \]
לפי אריתמטיקה של גבולות ורציפות השורש.

\textbf{תשובה:} $-\dfrac32$.
\end{solution}

\subsection*{שאלה 2}

\begin{question}
חשבו את האינטגרל הבא $\displaystyle\int_0^{\pi/4} x\tan^2 x\,dx$.
\end{question}

\begin{hint}
השתמשו בזהות $\tan^2 x = \frac{1}{\cos^2 x}-1$ ופצלו לשני אינטגרלים.
\end{hint}

\begin{hint}
את $\int x\cdot\frac{1}{\cos^2 x}\,dx$ חשבו באינטגרציה בחלקים, עם $u=x$ ו-$v'=\frac{1}{\cos^2x}$.
\end{hint}

\begin{solution}
לפי הזהות $1+\tan^2 x = \frac{1}{\cos^2 x}$:
\[ \int_0^{\pi/4} x\tan^2 x\,dx = \int_0^{\pi/4} \frac{x}{\cos^2 x}\,dx - \int_0^{\pi/4} x\,dx. \]
(הפונקציות רציפות בקטע $[0,\pi/4]$, ולכן כל האינטגרלים קיימים.)

\textbf{האינטגרל הראשון} — אינטגרציה בחלקים עם $u=x$, $v'=\frac{1}{\cos^2x}$, כלומר $u'=1$, $v=\tan x$:
\[ \int_0^{\pi/4}\frac{x}{\cos^2x}\,dx = \Big[x\tan x\Big]_0^{\pi/4} - \int_0^{\pi/4}\tan x\,dx. \]
מתקיים $\int \tan x\,dx = \int\frac{\sin x}{\cos x}dx = -\ln|\cos x|+C$ (הצבה $t=\cos x$), ולכן
\[ \int_0^{\pi/4}\frac{x}{\cos^2x}\,dx = \frac{\pi}{4}\cdot 1 - 0 + \Big[\ln\cos x\Big]_0^{\pi/4} = \frac{\pi}{4} + \ln\frac{\sqrt2}{2} - \ln 1 = \frac{\pi}{4} - \frac{\ln 2}{2}. \]

\textbf{האינטגרל השני:} $\displaystyle\int_0^{\pi/4}x\,dx = \frac{x^2}{2}\Big|_0^{\pi/4} = \frac{\pi^2}{32}$.

לסיכום:
\[ \int_0^{\pi/4} x\tan^2 x\,dx = \frac{\pi}{4} - \frac{\ln 2}{2} - \frac{\pi^2}{32} \approx 0.1304. \]
\end{solution}

\subsection*{שאלה 3א}

\begin{question}
רשמו פולינום מקלורן ממעלה 6 עבור פונקציה $y=\sin\left(\sin(x^2)\right)$.
\end{question}

\begin{hint}
הציבו $u=\sin(x^2)=x^2-\frac{x^6}{6}+o(x^6)$ בפיתוח $\sin u = u-\frac{u^3}{6}+o(u^3)$.
\end{hint}

\begin{solution}
נשתמש בפיתוחי מקלורן הידועים $\sin t = t - \frac{t^3}{6} + o(t^4)$ (כאשר $t\to0$).

ראשית, עם $t=x^2$: $\sin(x^2) = x^2 - \frac{x^6}{6} + o(x^8)$.

נסמן $u=\sin(x^2)$; אז $u\to 0$ כאשר $x\to0$ ו-$u = O(x^2)$. לפי הפיתוח $\sin u = u - \frac{u^3}{6} + o(u^4)$, ומכיוון ש-$o(u^4)=o(x^8)$:
\[ u^3 = \left(x^2-\tfrac{x^6}{6}+o(x^8)\right)^3 = x^6 + o(x^6), \]
ולכן
\[ \sin(\sin(x^2)) = x^2 - \frac{x^6}{6} - \frac{x^6}{6} + o(x^6) = x^2 - \frac{x^6}{3} + o(x^6). \]
מיחידות פולינום טיילור (פולינום ממעלה $\le 6$ המקיים $f(x)-P(x)=o(x^6)$ הוא פולינום טיילור), פולינום מקלורן ממעלה 6 הוא
\[ P_6(x) = x^2 - \frac{x^6}{3}. \]
\end{solution}

\subsection*{שאלה 3ב}

\begin{question}
עבור איזה ערך של $a$ פונקציה
\[ f(x)=\begin{cases} \dfrac{1}{\ln x}-\dfrac{1}{x-1} & x\neq 1\\[2mm] a & x=1\end{cases} \]
רציפה בתחום הגדרתה?
\end{question}

\begin{hint}
יש לחשב $\lim_{x\to1}\left(\frac{1}{\ln x}-\frac{1}{x-1}\right)$: הביאו למכנה משותף וקבלו מקרה $\frac00$.
\end{hint}

\begin{solution}
תחום ההגדרה של $f$ הוא $x>0$ (נדרש $\ln x$ מוגדר; בנקודה $x=1$ הפונקציה מוגדרת להיות $a$). בכל נקודה $x>0$, $x\ne1$, הפונקציה רציפה כהרכבה, הפרש ומנה של פונקציות רציפות עם מכנים שונים מאפס ($\ln x\ne0$ ו-$x-1\ne0$). לכן $f$ רציפה בתחום הגדרתה אם ורק אם היא רציפה ב-$x=1$, כלומר אם ורק אם $a=\lim_{x\to1}f(x)$.

נחשב: $\displaystyle \frac{1}{\ln x}-\frac{1}{x-1} = \frac{x-1-\ln x}{(x-1)\ln x}$, מקרה $\frac{0}{0}$. לפי כלל לופיטל (המונה והמכנה גזירים בסביבה של 1):
\[ \lim_{x\to1}\frac{x-1-\ln x}{(x-1)\ln x} = \lim_{x\to1}\frac{1-\frac1x}{\ln x+\frac{x-1}{x}} = \lim_{x\to1}\frac{x-1}{x\ln x + x-1}, \]
שוב מקרה $\frac00$, ולפי לופיטל שוב:
\[ \lim_{x\to1}\frac{1}{\ln x + 1 + 1} = \frac{1}{2}. \]
(לחלופין, עם $x=1+h$: $\ln(1+h)=h-\frac{h^2}{2}+o(h^2)$, והמונה הוא $\frac{h^2}{2}+o(h^2)$ והמכנה $h^2+o(h^2)$.)

\textbf{תשובה:} $a=\dfrac12$.
\end{solution}

\subsection*{שאלה 4}

\begin{question}
חקרו את הפונקציה $f(x)=\dfrac{e^{2x-1}}{x^2}$ וסרטטו את הגרף שלה.
צריך למצוא: תחום הגדרה, אסימפטוטות, תחומי עליה וירידה, נקודות קיצון, נקודות חיתוך של הגרף עם הצירים, תחומי קמירות וקעירות.
\end{question}

\begin{hint}
$f'(x)=\frac{2e^{2x-1}(x-1)}{x^3}$. שימו לב שהסימן של $x^3$ משתנה ב-$x=0$.
\end{hint}

\begin{solution}
\textbf{תחום הגדרה:} $x\neq 0$, כלומר $(-\infty,0)\cup(0,\infty)$.

\textbf{חיתוך עם הצירים:} $x=0$ אינו בתחום, ולכן אין חיתוך עם ציר $y$. מכיוון ש-$e^{2x-1}>0$ ו-$x^2>0$, מתקיים $f(x)>0$ לכל $x$ בתחום, ולכן אין חיתוך עם ציר $x$.

\textbf{אסימפטוטות:}
\begin{itemize}
  \item \emph{אנכית:} $\lim_{x\to0^\pm}\frac{e^{2x-1}}{x^2} = \frac{e^{-1}}{0^+} = +\infty$, לכן $x=0$ אסימפטוטה אנכית (הגרף עולה ל-$+\infty$ משני הצדדים).
  \item \emph{ב-$-\infty$:} $e^{2x-1}\to0$ ו-$\frac1{x^2}\to0$, לכן $\lim_{x\to-\infty}f(x)=0$, ו-$y=0$ אסימפטוטה אופקית משמאל.
  \item \emph{ב-$+\infty$:} $\frac{f(x)}{x}=\frac{e^{2x-1}}{x^3}\to\infty$ (אקספוננט גובר על כל חזקה, למשל לפי לופיטל שלוש פעמים), ולכן אין אסימפטוטה אופקית או משופעת בצד ימין, ו-$f(x)\to+\infty$.
\end{itemize}

\textbf{נגזרת ראשונה:} לפי כלל המנה,
\[ f'(x) = \frac{2e^{2x-1}x^2 - e^{2x-1}\cdot 2x}{x^4} = \frac{2e^{2x-1}(x-1)}{x^3}. \]
$f'(x)=0 \iff x=1$. הסימן של $f'$ נקבע ע"י $\frac{x-1}{x^3}$:
\begin{itemize}
  \item $x<0$: $x-1<0$, $x^3<0$ ולכן $f'>0$ — $f$ \textbf{עולה} ב-$(-\infty,0)$.
  \item $0<x<1$: $x-1<0$, $x^3>0$ ולכן $f'<0$ — $f$ \textbf{יורדת} ב-$(0,1)$.
  \item $x>1$: $f'>0$ — $f$ \textbf{עולה} ב-$(1,\infty)$.
\end{itemize}
לכן ב-$x=1$ יש \textbf{מינימום מקומי}: $f(1)=\frac{e^{1}}{1}=e$, כלומר הנקודה $(1,e)$. אין נקודות קיצון נוספות ($x=0$ אינו בתחום).

\textbf{נגזרת שנייה:} נגזור את $f'(x)=2e^{2x-1}\cdot\frac{x-1}{x^3}$:
\[ \left(\frac{x-1}{x^3}\right)' = \frac{x^3-3x^2(x-1)}{x^6} = \frac{-2x+3}{x^4}, \]
\[ f''(x) = 2e^{2x-1}\left(\frac{2(x-1)}{x^3}+\frac{3-2x}{x^4}\right) = \frac{2e^{2x-1}\left(2x^2-4x+3\right)}{x^4}. \]
לטרינום $2x^2-4x+3$ דיסקרימיננטה $16-24<0$ ומקדם מוביל חיובי, ולכן הוא חיובי לכל $x$. גם $e^{2x-1}>0$ ו-$x^4>0$. לכן $f''(x)>0$ לכל $x\ne0$:
$f$ \textbf{קמורה} (קמורה כלפי מעלה, $\cup$) בכל אחד מהקטעים $(-\infty,0)$ ו-$(0,\infty)$, ואין נקודות פיתול.

\textbf{סרטוט:} משמאל הגרף מתחיל קרוב לציר $x$ (מעליו), עולה וקמור ושואף ל-$+\infty$ כאשר $x\to0^-$. מימין ל-$0$ הגרף יורד מ-$+\infty$ עד המינימום $(1,e)$ ואז עולה במהירות אקספוננציאלית ל-$+\infty$. כל הגרף מעל ציר $x$.
\end{solution}

\subsection*{שאלה 5}

\begin{question}
משוואה של העקום בקואורדינטות קוטביות היא: $r=2(1-\sin\varphi)$.
סרטטו את העקום וחשבו שטח התחום החסום ע"י העקום.
\end{question}

\begin{hint}
שטח תחום בקואורדינטות קוטביות: $S=\frac12\int_\alpha^\beta r^2(\varphi)\,d\varphi$. כאן $r\ge0$ לכל $\varphi$, ולכן העקום נסגר כאשר $\varphi$ עובר על $[0,2\pi]$.
\end{hint}

\begin{solution}
\textbf{סרטוט:} מתקיים $0\le r=2(1-\sin\varphi)\le 4$ לכל $\varphi$, ו-$r$ מחזורית עם מחזור $2\pi$, ולכן העקום סגור ומתקבל כאשר $\varphi\in[0,2\pi]$. כמה ערכים:
\[ \varphi=0:\ r=2;\quad \varphi=\tfrac{\pi}{2}:\ r=0;\quad \varphi=\pi:\ r=2;\quad \varphi=\tfrac{3\pi}{2}:\ r=4. \]
כלומר העקום עובר בנקודות $(2,0)$, $(0,0)$, $(-2,0)$, $(0,-4)$ (בקואורדינטות קרטזיות). מכיוון ש-$r(\pi-\varphi)=r(\varphi)$, העקום סימטרי ביחס לציר $y$. זוהי \textbf{קרדיואידה} ("לב") שהחוד שלה בראשית (כאשר $\varphi=\frac\pi2$) והיא "פתוחה" כלפי מטה: הנקודה הרחוקה ביותר היא $(0,-4)$.

\textbf{שטח:} לפי הנוסחה לשטח בקואורדינטות קוטביות, $S=\frac12\int_0^{2\pi}r^2\,d\varphi$:
\[ S = \frac12\int_0^{2\pi}4(1-\sin\varphi)^2\,d\varphi = 2\int_0^{2\pi}\left(1-2\sin\varphi+\sin^2\varphi\right)d\varphi. \]
נחשב כל איבר בנפרד:
\[ \int_0^{2\pi}1\,d\varphi=2\pi,\qquad \int_0^{2\pi}\sin\varphi\,d\varphi = \Big[-\cos\varphi\Big]_0^{2\pi}=0, \]
\[ \int_0^{2\pi}\sin^2\varphi\,d\varphi = \int_0^{2\pi}\frac{1-\cos2\varphi}{2}\,d\varphi = \Big[\frac{\varphi}{2}-\frac{\sin2\varphi}{4}\Big]_0^{2\pi} = \pi. \]
לכן
\[ S = 2\left(2\pi - 0 + \pi\right) = 6\pi. \]
\textbf{תשובה:} שטח התחום הוא $6\pi$.
\end{solution}

\subsection*{שאלה 6א}

\begin{question}
האם האינטגרל הלא אמיתי הבא $\displaystyle\int_{-\infty}^{-2}\frac{dx}{x^2-2x}$ מתכנס? אם כן, מצאו אותו.
\end{question}

\begin{hint}
פרקו לשברים חלקיים: $\frac{1}{x(x-2)}=\frac12\left(\frac{1}{x-2}-\frac1x\right)$.
\end{hint}

\begin{solution}
בקטע $(-\infty,-2]$ המכנה $x^2-2x=x(x-2)>0$ (שני הגורמים שליליים), ולכן הפונקציה רציפה וחיובית שם, והבעיה היחידה היא הגבול $-\infty$. לפי ההגדרה,
\[ \int_{-\infty}^{-2}\frac{dx}{x^2-2x} = \lim_{R\to-\infty}\int_R^{-2}\frac{dx}{x(x-2)}. \]
פירוק לשברים חלקיים: $\frac{1}{x(x-2)}=\frac{A}{x}+\frac{B}{x-2}$, ומ-$1=A(x-2)+Bx$ נקבל ($x=0$) $A=-\frac12$ ו-($x=2$) $B=\frac12$. לכן
\[ \int\frac{dx}{x(x-2)} = \frac12\ln|x-2|-\frac12\ln|x| + C = \frac12\ln\left|\frac{x-2}{x}\right| + C. \]
לפיכך
\[ \int_R^{-2}\frac{dx}{x(x-2)} = \frac12\ln\frac{-4}{-2} - \frac12\ln\left|\frac{R-2}{R}\right| = \frac12\ln2-\frac12\ln\left|1-\frac2R\right|. \]
כאשר $R\to-\infty$, $1-\frac2R\to1$ ומרציפות הלוגריתם $\ln\left|1-\frac2R\right|\to0$. לכן הגבול קיים וסופי:
\[ \int_{-\infty}^{-2}\frac{dx}{x^2-2x} = \frac{\ln2}{2}. \]
\textbf{תשובה:} האינטגרל מתכנס וערכו $\frac12\ln2$.
(אפשר גם לראות את ההתכנסות ממבחן ההשוואה הגבולי עם $\frac1{x^2}$, אך החישוב הישיר נותן גם את הערך.)
\end{solution}

\subsection*{שאלה 6ב}

\begin{question}
חשבו את האינטגרל הבא: $\displaystyle\int_{\ln2}^{\ln3}\frac{dx}{e^x-e^{-x}}$.
\end{question}

\begin{hint}
הכפילו מונה ומכנה ב-$e^x$ והציבו $t=e^x$.
\end{hint}

\begin{solution}
בקטע $[\ln2,\ln3]$ מתקיים $e^x>e^{-x}$, ולכן האינטגרנד רציף והאינטגרל אמיתי. נכפיל מונה ומכנה ב-$e^x$:
\[ \int_{\ln2}^{\ln3}\frac{dx}{e^x-e^{-x}} = \int_{\ln2}^{\ln3}\frac{e^x\,dx}{e^{2x}-1}. \]
נציב $t=e^x$, $dt=e^x\,dx$; הגבולות: $x=\ln2\mapsto t=2$, $x=\ln3\mapsto t=3$:
\[ = \int_2^3\frac{dt}{t^2-1} = \int_2^3\frac12\left(\frac{1}{t-1}-\frac{1}{t+1}\right)dt = \frac12\Big[\ln\frac{t-1}{t+1}\Big]_2^3 = \frac12\left(\ln\frac12-\ln\frac13\right) = \frac12\ln\frac32. \]
\textbf{תשובה:} $\dfrac12\ln\dfrac32\approx0.2027$.
\end{solution}

\section*{תשע"ט סמסטר ב מועד א}

\subsection*{שאלה 1א}

\begin{question}
חשבו את הנגזרת של הפונקציה $f(x)=\sqrt{5x+3}$ בנקודה $x$ לפי הגדרת הנגזרת (תשובה לפי טבלת הנגזרות לא תתקבל).
\end{question}

\begin{hint}
הכפילו את המונה והמכנה בצמוד $\sqrt{5(x+h)+3}+\sqrt{5x+3}$.
\end{hint}

\begin{solution}
הפונקציה מוגדרת עבור $x\ge -\tfrac35$, ונחשב את הנגזרת בנקודה $x>-\tfrac35$. לפי ההגדרה, בעזרת הכפלה בצמוד:
\begin{align*}
f'(x)&=\lim_{h\to0}\frac{\sqrt{5(x+h)+3}-\sqrt{5x+3}}{h}
=\lim_{h\to0}\frac{\big(5(x+h)+3\big)-(5x+3)}{h\left(\sqrt{5x+5h+3}+\sqrt{5x+3}\right)}\\
&=\lim_{h\to0}\frac{5}{\sqrt{5x+5h+3}+\sqrt{5x+3}}=\frac{5}{2\sqrt{5x+3}},
\end{align*}
כאשר במעבר האחרון השתמשנו ברציפות פונקציית השורש ($\sqrt{5x+5h+3}\to\sqrt{5x+3}>0$). לכן
\[ f'(x)=\frac{5}{2\sqrt{5x+3}},\qquad x>-\tfrac35. \]
(בנקודה $x=-\tfrac35$ המנה $\frac{\sqrt{5h}}{h}\to+\infty$, ולכן שם הפונקציה אינה גזירה.)
\end{solution}

\subsection*{שאלה 1ב}

\begin{question}
מצאו את כל האסימפטוטות של הפונקציה: $f(x)=\dfrac{x^2+x-6}{|x|-2}$.
\end{question}

\begin{hint}
פרקו את המונה $x^2+x-6=(x+3)(x-2)$ וטפלו בנפרד במקרים $x>0$ ו-$x<0$.
\end{hint}

\begin{solution}
תחום ההגדרה: $|x|\neq2$, כלומר $x\neq\pm2$. נפרק $x^2+x-6=(x+3)(x-2)$.

\textbf{עבור $x\ge0$, $x\ne 2$:} $|x|-2=x-2$ ולכן $f(x)=\dfrac{(x+3)(x-2)}{x-2}=x+3$.

\textbf{עבור $x<0$, $x\ne-2$:} $|x|-2=-x-2=-(x+2)$ ולכן
\[ f(x)=-\frac{(x+3)(x-2)}{x+2}=-\frac{x^2+x-6}{x+2}=-x+1+\frac{4}{x+2}, \]
(חילוק פולינומים: $x^2+x-6=(x+2)(x-1)-4$).

\textbf{אסימפטוטות אנכיות:} ב-$x=2$: $\lim_{x\to2}f(x)=\lim_{x\to2}(x+3)=5$ סופי, ולכן זו נקודת אי-רציפות סליקה ואין שם אסימפטוטה. ב-$x=-2$: $\lim_{x\to-2^\pm}\left(-x+1+\frac{4}{x+2}\right)=\pm\infty$, ולכן $x=-2$ אסימפטוטה אנכית.

\textbf{אסימפטוטות משופעות/אופקיות:} כאשר $x\to+\infty$ מתקיים $f(x)=x+3$ בדיוק, ולכן $y=x+3$ אסימפטוטה משופעת ב-$+\infty$. כאשר $x\to-\infty$:
\[ f(x)-(-x+1)=\frac{4}{x+2}\xrightarrow[x\to-\infty]{}0, \]
ולכן $y=-x+1$ אסימפטוטה משופעת ב-$-\infty$. (אין אסימפטוטות אופקיות.)

\textbf{תשובה:} $x=-2$ (אנכית), $y=x+3$ (ב-$+\infty$), $y=-x+1$ (ב-$-\infty$).
\end{solution}

\subsection*{שאלה 1ג}

\begin{question}
מהו שיפוע המשיק לקו $y=\int_0^{x+1}e^{-2t^2}\,dt$ בנקודה עליו $x=0$.
\end{question}

\begin{hint}
השתמשו במשפט היסודי של החשבון האינטגרלי יחד עם כלל השרשרת.
\end{hint}

\begin{solution}
נסמן $F(u)=\int_0^u e^{-2t^2}\,dt$. הפונקציה $e^{-2t^2}$ רציפה, ולכן לפי המשפט היסודי $F'(u)=e^{-2u^2}$. מכאן $y(x)=F(x+1)$ ולפי כלל השרשרת
\[ y'(x)=F'(x+1)\cdot 1=e^{-2(x+1)^2}. \]
שיפוע המשיק בנקודה $x=0$ הוא $y'(0)=e^{-2}=\dfrac{1}{e^2}$.
\end{solution}

\subsection*{שאלה 2א}

\begin{question}
הוכיחו כי למשוואה $x^2-\ln x=0$ אין פתרונות ממשיים.
\end{question}

\begin{hint}
מצאו את ערך המינימום של $g(x)=x^2-\ln x$ בתחום $x>0$ והראו שהוא חיובי.
\end{hint}

\begin{solution}
נגדיר $g(x)=x^2-\ln x$, המוגדרת וגזירה בתחום $x>0$ (לכן אין למשוואה פתרונות עם $x\le0$). נגזור:
\[ g'(x)=2x-\frac1x=\frac{2x^2-1}{x}. \]
עבור $x>0$: $g'(x)<0$ כאשר $0<x<\frac{1}{\sqrt2}$ ו-$g'(x)>0$ כאשר $x>\frac1{\sqrt2}$. לכן $g$ יורדת ממש ב-$(0,\frac1{\sqrt2}]$ ועולה ממש ב-$[\frac1{\sqrt2},\infty)$, ו-$x=\frac1{\sqrt2}$ נקודת מינימום מוחלט של $g$ בתחום:
\[ g\!\left(\tfrac1{\sqrt2}\right)=\frac12-\ln\frac1{\sqrt2}=\frac12+\frac{\ln2}{2}>0. \]
לכן $g(x)\ge\frac{1+\ln 2}{2}>0$ לכל $x>0$, ובפרט $g(x)\ne0$: למשוואה אין פתרונות ממשיים.
\end{solution}

\subsection*{שאלה 2ב}

\begin{question}
חשבו את השטח הכלוא בין הקווים $y=x\cdot e^x$, $y=x-1$, $x=0$, $x=2$. ציירו את הסקיצה המתאימה.
\end{question}

\begin{hint}
בדקו איזה מהגרפים נמצא מעל השני בקטע $[0,2]$, ואת $\int xe^x\,dx$ חשבו באינטגרציה בחלקים.
\end{hint}

\begin{solution}
לכל $x\in[0,2]$ מתקיים $e^x\ge1$ ולכן $xe^x\ge x>x-1$. כלומר בקטע הגרף של $y=xe^x$ נמצא מעל הישר $y=x-1$ (הסקיצה: הישר $y=x-1$ עובר דרך $(0,-1)$ ו-$(2,1)$, והעקומה $y=xe^x$ עוברת דרך $(0,0)$ ו-$(2,2e^2)$; התחום חסום משמאל ומימין בישרים $x=0$, $x=2$). השטח:
\[ S=\int_0^2\big(xe^x-(x-1)\big)\,dx. \]
באינטגרציה בחלקים ($u=x$, $dv=e^xdx$): $\int xe^x\,dx=xe^x-\int e^x\,dx=(x-1)e^x+C$. לכן
\[ \int_0^2 xe^x\,dx=\big[(x-1)e^x\big]_0^2=e^2-(-1)=e^2+1, \qquad \int_0^2(x-1)\,dx=\Big[\tfrac{x^2}{2}-x\Big]_0^2=0. \]
\textbf{תשובה:} $S=e^2+1$.
\end{solution}

\subsection*{שאלה 3א}

\begin{question}
חשבו:
\begin{enumerate}
  \item[א.] (8 נק') $\displaystyle\int\frac{(x+1)^2}{x(x+2)^2}\,dx$
  \item[ב.] (7 נק') $\displaystyle\int_{-\pi/4}^{\pi/4}\left(\frac{1}{\cos^2x\cdot(2+\tan x)^2}+x\cos^3x\right)dx$
\end{enumerate}
\end{question}

\begin{hint}
בסעיף א' פרקו לשברים חלקיים מהצורה $\frac{A}{x}+\frac{B}{x+2}+\frac{C}{(x+2)^2}$.
\end{hint}

\begin{hint}
בסעיף ב' האינטגרל של פונקציה אי-זוגית על קטע סימטרי הוא $0$; בחלק השני הציבו $u=\tan x$.
\end{hint}

\begin{solution}
\begin{enumerate}
  \item[א.] מעלת המונה (2) קטנה ממעלת המכנה (3), ולכן נפרק ישירות לשברים חלקיים:
  \[ \frac{(x+1)^2}{x(x+2)^2}=\frac{A}{x}+\frac{B}{x+2}+\frac{C}{(x+2)^2}
  \;\Longrightarrow\; (x+1)^2=A(x+2)^2+Bx(x+2)+Cx. \]
  הצבת $x=0$: $1=4A$, כלומר $A=\frac14$. הצבת $x=-2$: $1=-2C$, כלומר $C=-\frac12$. השוואת מקדמי $x^2$: $1=A+B$, ולכן $B=\frac34$. לכן
  \[ \int\frac{(x+1)^2}{x(x+2)^2}\,dx=\frac14\ln|x|+\frac34\ln|x+2|+\frac{1}{2(x+2)}+C. \]
  \item[ב.] נפצל לשני אינטגרלים. הפונקציה $x\cos^3x$ אי-זוגית (מכפלה של $x$ האי-זוגית ב-$\cos^3x$ הזוגית) ורציפה, ולכן האינטגרל שלה על הקטע הסימטרי $[-\frac\pi4,\frac\pi4]$ שווה $0$.

  באינטגרל הראשון נציב $u=\tan x$, $du=\frac{dx}{\cos^2x}$; כאשר $x$ עובר מ-$-\frac\pi4$ ל-$\frac\pi4$, $u$ עובר מ-$-1$ ל-$1$ (ובקטע זה $2+\tan x\ge1>0$):
  \[ \int_{-\pi/4}^{\pi/4}\frac{dx}{\cos^2x\,(2+\tan x)^2}=\int_{-1}^{1}\frac{du}{(2+u)^2}=\left[-\frac{1}{2+u}\right]_{-1}^{1}=-\frac13+1=\frac23. \]
  \textbf{תשובה:} $\frac23$.
\end{enumerate}
\end{solution}

\subsection*{שאלה 3ב}

\begin{question}
רשמו פולינום מקלורן מסדר 2 עבור הפונקציה $f(x)=\cos(2x)-\arctan(x+1)$.
\end{question}

\begin{hint}
פתחו כל מחובר בנפרד: ל-$\cos(2x)$ השתמשו בפיתוח הידוע של $\cos t$, ול-$\arctan(x+1)$ חשבו את $g(0),g'(0),g''(0)$.
\end{hint}

\begin{solution}
פולינום מקלורן מסדר 2: $P_2(x)=f(0)+f'(0)x+\frac{f''(0)}{2}x^2$. נחשב בנפרד לכל מחובר.

עבור $\cos(2x)$: לפי הפיתוח $\cos t=1-\frac{t^2}{2}+o(t^2)$ נקבל $\cos 2x=1-2x^2+o(x^2)$.

עבור $g(x)=\arctan(x+1)$:
\[ g(0)=\arctan1=\frac\pi4,\quad g'(x)=\frac{1}{1+(x+1)^2}\Rightarrow g'(0)=\frac12,\quad g''(x)=-\frac{2(x+1)}{\big(1+(x+1)^2\big)^2}\Rightarrow g''(0)=-\frac24=-\frac12. \]
לכן $\arctan(x+1)=\frac\pi4+\frac x2-\frac{x^2}{4}+o(x^2)$. בחיסור:
\[ P_2(x)=\left(1-\frac\pi4\right)-\frac12x-\frac74x^2. \]
\end{solution}

\subsection*{שאלה 4א}

\begin{question}
לאילו ערכים של של הפרמטרים $a$ ו-$b$ הפונקציה
\[ f(x)=\begin{cases}\dfrac{1-\cos x}{x\ln(1+x)} & -1<x<0\\[2mm] ax+b & 0\le x\le1\\[1mm] x^{1/(x-1)} & x>1\end{cases} \]
תהיה רציפה בתחום $x>-1$?
\end{question}

\begin{hint}
בכל אחד מהקטעים הפתוחים הפונקציה רציפה כהרכבה של פונקציות אלמנטריות; יש לבדוק רק את $x=0$ ואת $x=1$.
\end{hint}

\begin{hint}
את הגבול $\lim_{x\to1^+}x^{1/(x-1)}$ חשבו בעזרת $x^{1/(x-1)}=e^{\frac{\ln x}{x-1}}$.
\end{hint}

\begin{solution}
בקטעים $(-1,0)$, $(0,1)$, $(1,\infty)$ הפונקציה רציפה (מנה/הרכבה של פונקציות רציפות, והמכנה $x\ln(1+x)\ne0$ ב-$(-1,0)$; ב-$(1,\infty)$: $x^{1/(x-1)}=e^{\ln x/(x-1)}$). נותר לבדוק את נקודות התפר.

\textbf{הנקודה $x=0$:} $f(0)=b=\lim_{x\to0^+}f(x)$. משמאל, בעזרת הגבולות הידועים $\lim_{x\to0}\frac{1-\cos x}{x^2}=\frac12$ ו-$\lim_{x\to0}\frac{\ln(1+x)}{x}=1$:
\[ \lim_{x\to0^-}\frac{1-\cos x}{x\ln(1+x)}=\lim_{x\to0^-}\frac{1-\cos x}{x^2}\cdot\frac{x}{\ln(1+x)}=\frac12\cdot1=\frac12. \]
לכן רציפות ב-$0$ שקולה ל-$b=\frac12$.

\textbf{הנקודה $x=1$:} $f(1)=a+b=\lim_{x\to1^-}f(x)$. מימין:
\[ \lim_{x\to1^+}x^{1/(x-1)}=\lim_{x\to1^+}e^{\frac{\ln x}{x-1}}. \]
הגבול $\lim_{x\to1}\frac{\ln x}{x-1}$ הוא מהצורה $\frac00$, ולפי לופיטל $=\lim_{x\to1}\frac{1/x}{1}=1$ (זו גם הנגזרת של $\ln$ ב-$1$). מרציפות האקספוננט הגבול הוא $e^1=e$. לכן רציפות ב-$1$ שקולה ל-$a+b=e$.

\textbf{תשובה:} $b=\frac12$, $a=e-\frac12$.
\end{solution}

\subsection*{שאלה 4ב}

\begin{question}
חשבו $\displaystyle\int_1^{+\infty}\frac{\ln x}{x^4}\,dx$.
\end{question}

\begin{hint}
בצעו אינטגרציה בחלקים עם $u=\ln x$, $dv=x^{-4}dx$ ואחר כך השאיפו את הגבול העליון לאינסוף.
\end{hint}

\begin{solution}
לפי הגדרת האינטגרל הלא אמיתי, $\int_1^\infty=\lim_{R\to\infty}\int_1^R$. באינטגרציה בחלקים עם $u=\ln x$, $dv=x^{-4}dx$ ($du=\frac{dx}{x}$, $v=-\frac{1}{3x^3}$):
\[ \int_1^R\frac{\ln x}{x^4}\,dx=\left[-\frac{\ln x}{3x^3}\right]_1^R+\frac13\int_1^R\frac{dx}{x^4}=-\frac{\ln R}{3R^3}+\frac13\left[-\frac{1}{3x^3}\right]_1^R=-\frac{\ln R}{3R^3}+\frac19\left(1-\frac{1}{R^3}\right). \]
כאשר $R\to\infty$: $\frac{\ln R}{R^3}\to0$ (לפי לופיטל: $\lim\frac{1/R}{3R^2}=0$) ו-$\frac1{R^3}\to0$. לכן האינטגרל מתכנס ו-
\[ \int_1^{+\infty}\frac{\ln x}{x^4}\,dx=\frac19. \]
\end{solution}

\section*{תשע"ט סמסטר ב מועד ב}

\subsection*{שאלה 1א}

\begin{question}
חשבו את הנגזרת של הפונקציה $f(x)=\sin(3x+4)$ בנקודה $x$ לפי הגדרת הנגזרת (תשובה לפי טבלת הנגזרות לא תתקבל).
\end{question}

\begin{hint}
השתמשו בזהות $\sin\alpha-\sin\beta=2\cos\frac{\alpha+\beta}{2}\sin\frac{\alpha-\beta}{2}$ ובגבול $\lim_{t\to0}\frac{\sin t}{t}=1$.
\end{hint}

\begin{solution}
לפי ההגדרה ובעזרת הזהות $\sin\alpha-\sin\beta=2\cos\frac{\alpha+\beta}{2}\sin\frac{\alpha-\beta}{2}$:
\begin{align*}
f'(x)&=\lim_{h\to0}\frac{\sin(3x+3h+4)-\sin(3x+4)}{h}
=\lim_{h\to0}\frac{2\cos\!\left(3x+4+\frac{3h}{2}\right)\sin\frac{3h}{2}}{h}\\
&=\lim_{h\to0}3\cos\!\left(3x+4+\tfrac{3h}{2}\right)\cdot\frac{\sin\frac{3h}{2}}{\frac{3h}{2}}=3\cos(3x+4)\cdot1,
\end{align*}
כאשר השתמשנו ברציפות $\cos$ ובגבול היסודי $\lim_{t\to0}\frac{\sin t}{t}=1$. לכן $f'(x)=3\cos(3x+4)$.
\end{solution}

\subsection*{שאלה 1ב}

\begin{question}
מצאו את כל האסימפטוטות של הפונקציה: $f(x)=\dfrac{x\sqrt{x^2-1}}{2x^2-1}$.
\end{question}

\begin{hint}
שימו לב לתחום ההגדרה $|x|\ge1$ ולכך ש-$\sqrt{x^2}=|x|$.
\end{hint}

\begin{solution}
\textbf{תחום הגדרה:} צריך $x^2-1\ge0$, כלומר $|x|\ge1$, וגם $2x^2-1\ne0$, כלומר $x\ne\pm\frac1{\sqrt2}$ – אך נקודות אלה ממילא אינן בתחום. לכן התחום הוא $(-\infty,-1]\cup[1,\infty)$.

\textbf{אסימפטוטות אנכיות:} בתחום המכנה מקיים $2x^2-1\ge1>0$, ולכן $f$ רציפה בכל תחומה (גם בקצוות $x=\pm1$, שם $f(\pm1)=0$). אין נקודה שבה $f$ שואפת לאינסוף, ולכן אין אסימפטוטות אנכיות.

\textbf{אסימפטוטות אופקיות:} מכיוון ש-$\sqrt{x^2-1}=|x|\sqrt{1-\frac1{x^2}}$,
\[ f(x)=\frac{x|x|\sqrt{1-\frac1{x^2}}}{x^2\left(2-\frac1{x^2}\right)}=\frac{|x|}{x}\cdot\frac{\sqrt{1-\frac1{x^2}}}{2-\frac1{x^2}}. \]
כאשר $x\to+\infty$: $\frac{|x|}{x}=1$ ולכן $f(x)\to\frac12$. כאשר $x\to-\infty$: $\frac{|x|}{x}=-1$ ולכן $f(x)\to-\frac12$. מכיוון שהגבולות סופיים, אין אסימפטוטות משופעות נוספות.

\textbf{תשובה:} $y=\frac12$ (ב-$+\infty$) ו-$y=-\frac12$ (ב-$-\infty$); אין אסימפטוטות אנכיות.
\end{solution}

\subsection*{שאלה 1ג}

\begin{question}
עבור אילו ערכים של $a$ יש לקו $y=x^2+a\sin x$ נקודות פיתול?
\end{question}

\begin{hint}
בדקו מתי $y''$ מחליפה סימן.
\end{hint}

\begin{solution}
$y'=2x+a\cos x$, $y''=2-a\sin x$. נקודת פיתול היא נקודה שבה $y''$ מחליפה סימן.
\begin{itemize}
  \item אם $|a|<2$: $|a\sin x|\le|a|<2$ ולכן $y''>0$ לכל $x$ – אין נקודות פיתול.
  \item אם $a=2$: $y''=2(1-\sin x)\ge0$, ומתאפסת רק בנקודות מבודדות; הסימן לא מתחלף – אין פיתול. באופן דומה עבור $a=-2$: $y''=2(1+\sin x)\ge0$.
  \item אם $|a|>2$: המשוואה $\sin x=\frac2a$ עם $0<\left|\frac2a\right|<1$ יש לה פתרונות $x_0$, ובהם $\cos x_0\ne0$. לכן $(y'')'(x_0)=-a\cos x_0\ne0$, כלומר $y''$ עוברת דרך $0$ בצורה מונוטונית ממש ומחליפה סימן ב-$x_0$ – נקודת פיתול.
\end{itemize}
\textbf{תשובה:} $|a|>2$.
\end{solution}

\subsection*{שאלה 2א}

\begin{question}
הוכיחו כי למשוואה $\dfrac{1}{x+1}+\dfrac1x+\dfrac{1}{x-1}=1$ יש לפחות שני פתרונות.
\end{question}

\begin{hint}
הגדירו $g(x)=\frac{1}{x+1}+\frac1x+\frac1{x-1}-1$ ובדקו את הגבולות החד-צדדיים של $g$ בקצוות הקטעים $(-1,0)$ ו-$(0,1)$. השתמשו במשפט ערך הביניים.
\end{hint}

\begin{solution}
נגדיר $g(x)=\frac{1}{x+1}+\frac1x+\frac1{x-1}-1$, הרציפה בכל קטע שאינו מכיל את $-1,0,1$.

\textbf{הקטע $(-1,0)$:} כאשר $x\to-1^+$, המחובר $\frac1{x+1}\to+\infty$ ושאר המחוברים חסומים, ולכן $g(x)\to+\infty$. כאשר $x\to0^-$, $\frac1x\to-\infty$ ולכן $g(x)\to-\infty$. מכאן קיימות נקודות $-1<p<q<0$ עם $g(p)>0$ ו-$g(q)<0$. $g$ רציפה ב-$[p,q]$, ולפי משפט ערך הביניים קיימת $x_1\in(p,q)$ עם $g(x_1)=0$.

\textbf{הקטע $(0,1)$:} כאשר $x\to0^+$, $\frac1x\to+\infty$ ולכן $g\to+\infty$; כאשר $x\to1^-$, $\frac1{x-1}\to-\infty$ ולכן $g\to-\infty$. באותו אופן, לפי משפט ערך הביניים קיים $x_2\in(0,1)$ עם $g(x_2)=0$.

הקטעים זרים, ולכן $x_1\ne x_2$ ולמשוואה לפחות שני פתרונות. (באותו אופן יש פתרון שלישי ב-$(1,\infty)$: $g\to+\infty$ כאשר $x\to1^+$ ו-$g\to-1$ כאשר $x\to\infty$.)
\end{solution}

\subsection*{שאלה 2ב}

\begin{question}
חשבו את נפחו של גוף הסיבוב של תחום מישורי:
\[ 0\le x\le2,\quad 0\le y\le\sqrt{\frac{x}{x+2}} \]
סביב ציר ה-$x$. ציירו את הסקיצה המתאימה.
\end{question}

\begin{hint}
נפח גוף סיבוב סביב ציר $x$ הוא $V=\pi\int_a^b y^2\,dx$.
\end{hint}

\begin{solution}
הסקיצה: הפונקציה $y=\sqrt{\frac{x}{x+2}}$ עולה מ-$y(0)=0$ ל-$y(2)=\sqrt{\frac12}$; התחום הוא השטח שבין הגרף לציר $x$ בקטע $[0,2]$, והגוף נוצר מסיבובו סביב ציר $x$. לפי נוסחת נפח גוף סיבוב:
\[ V=\pi\int_0^2\frac{x}{x+2}\,dx=\pi\int_0^2\left(1-\frac{2}{x+2}\right)dx=\pi\Big[x-2\ln(x+2)\Big]_0^2=\pi\big(2-2\ln4+2\ln2\big). \]
\textbf{תשובה:} $V=\pi(2-2\ln2)=2\pi(1-\ln2)\approx1.928$.
\end{solution}

\subsection*{שאלה 3א}

\begin{question}
חשבו:
\begin{enumerate}
  \item[א.] (7 נק') $\displaystyle\int\frac{x^2-8}{x^3+4x^2}\,dx$
  \item[ב.] (8 נק') $\displaystyle\int_0^1\ln(1+\sqrt x)\,dx$
\end{enumerate}
\end{question}

\begin{hint}
בסעיף א' פרקו $x^3+4x^2=x^2(x+4)$ ולשברים חלקיים. בסעיף ב' הציבו $t=\sqrt x$ ואחר כך בצעו אינטגרציה בחלקים.
\end{hint}

\begin{solution}
\begin{enumerate}
  \item[א.] $x^3+4x^2=x^2(x+4)$, ומעלת המונה קטנה ממעלת המכנה. נפרק:
  \[ \frac{x^2-8}{x^2(x+4)}=\frac Ax+\frac B{x^2}+\frac C{x+4}\;\Longrightarrow\;x^2-8=Ax(x+4)+B(x+4)+Cx^2. \]
  הצבת $x=0$: $-8=4B$, $B=-2$. הצבת $x=-4$: $8=16C$, $C=\frac12$. מקדמי $x^2$: $1=A+C$, $A=\frac12$. לכן
  \[ \int\frac{x^2-8}{x^3+4x^2}\,dx=\frac12\ln|x|+\frac2x+\frac12\ln|x+4|+C. \]
  \item[ב.] נציב $x=t^2$, $dx=2t\,dt$, $t\in[0,1]$:
  \[ \int_0^1\ln(1+\sqrt x)\,dx=\int_0^1 2t\ln(1+t)\,dt. \]
  באינטגרציה בחלקים ($u=\ln(1+t)$, $dv=2t\,dt$, $v=t^2$):
  \[ =\Big[t^2\ln(1+t)\Big]_0^1-\int_0^1\frac{t^2}{1+t}\,dt=\ln2-\int_0^1\left(t-1+\frac1{1+t}\right)dt=\ln2-\left(\frac12-1+\ln2\right)=\frac12. \]
  \textbf{תשובה:} $\frac12$.
\end{enumerate}
\end{solution}

\subsection*{שאלה 3ב}

\begin{question}
רשמו פולינום טיילור מסדר 2 עבור הפונקציה
\[ f(x)=\int_0^x\frac{1-t}{1+t^2}\,dt \]
סביב הנקודה $x=1$.
\end{question}

\begin{hint}
לפי המשפט היסודי $f'(x)=\frac{1-x}{1+x^2}$; את $f(1)$ חשבו ישירות מהאינטגרל.
\end{hint}

\begin{solution}
פולינום טיילור מסדר 2 סביב $1$: $P_2(x)=f(1)+f'(1)(x-1)+\frac{f''(1)}{2}(x-1)^2$.

הפונקציה $\frac{1-t}{1+t^2}$ רציפה בכל $\R$, ולכן לפי המשפט היסודי של החשבון האינטגרלי $f'(x)=\frac{1-x}{1+x^2}$, ומכאן $f'(1)=0$. נגזור שוב:
\[ f''(x)=\frac{-(1+x^2)-(1-x)\cdot2x}{(1+x^2)^2}=\frac{x^2-2x-1}{(1+x^2)^2},\qquad f''(1)=\frac{-2}{4}=-\frac12. \]
ולבסוף
\[ f(1)=\int_0^1\frac{dt}{1+t^2}-\int_0^1\frac{t\,dt}{1+t^2}=\Big[\arctan t\Big]_0^1-\Big[\tfrac12\ln(1+t^2)\Big]_0^1=\frac\pi4-\frac{\ln2}{2}. \]
\textbf{תשובה:}
\[ P_2(x)=\frac\pi4-\frac{\ln2}{2}-\frac14(x-1)^2. \]
\end{solution}

\subsection*{שאלה 4א}

\begin{question}
לאילו ערכים של של הפרמטרים $a$ ו-$b$ הפונקציה
\[ f(x)=\begin{cases}\dfrac{\sin(3x)}{\ln(1-x)} & x<0\\[2mm] ax+b & 0\le x\le1\\[1mm] \big(\cos(x-1)\big)^{1/(x-1)} & x>1\end{cases} \]
תהיה רציפה לכל $x$ ממשי?
\end{question}

\begin{hint}
יש לבדוק רציפות רק בנקודות התפר $x=0$ ו-$x=1$. את הגבול מימין ל-$1$ כתבו בצורה $e^{\frac{\ln\cos(x-1)}{x-1}}$.
\end{hint}

\begin{solution}
עבור $x<0$ מתקיים $1-x>1$ ולכן $\ln(1-x)>0$, כך שבקטע $(-\infty,0)$ הפונקציה רציפה כמנה של פונקציות רציפות עם מכנה שונה מאפס. ב-$(0,1)$ היא פולינום. עבור $x>1$ הביטוי $\big(\cos(x-1)\big)^{1/(x-1)}=e^{\frac{\ln\cos(x-1)}{x-1}}$ מוגדר ורציף כאשר $\cos(x-1)>0$, ובפרט בקטע $1<x<1+\frac\pi2$, שהוא הסביבה הרלוונטית לרציפות ב-$1$. (הערה: עבור $x>1$ שבהם $\cos(x-1)\le0$ הביטוי אינו מוגדר במובן הממשי, ולכן את הדרישה "רציפה לכל $x$" יש להבין כרציפות בתחום ההגדרה.) נבדוק את נקודות התפר.

\textbf{הנקודה $x=0$:} $f(0)=b=\lim_{x\to0^+}f(x)$. משמאל:
\[ \lim_{x\to0^-}\frac{\sin3x}{\ln(1-x)}=\lim_{x\to0^-}\frac{\sin3x}{3x}\cdot\frac{-x}{\ln(1-x)}\cdot(-3)=1\cdot1\cdot(-3)=-3, \]
לפי $\lim_{t\to0}\frac{\sin t}{t}=1$ ו-$\lim_{t\to0}\frac{\ln(1+t)}{t}=1$ (עם $t=-x$). לכן $b=-3$.

\textbf{הנקודה $x=1$:} $f(1)=a+b=\lim_{x\to1^-}f(x)$. מימין, נסמן $u=x-1\to0^+$:
\[ \big(\cos u\big)^{1/u}=e^{\frac{\ln\cos u}{u}},\qquad \lim_{u\to0^+}\frac{\ln\cos u}{u}\overset{\text{לופיטל}}{=}\lim_{u\to0^+}\frac{-\tan u}{1}=0. \]
מרציפות האקספוננט הגבול הוא $e^0=1$. לכן $a+b=1$.

\textbf{תשובה:} $b=-3$, $a=4$.
\end{solution}

\subsection*{שאלה 4ב}

\begin{question}
חשבו $\displaystyle\int_{0.25}^{+\infty}\frac{dx}{(4x+1)\sqrt x}$.
\end{question}

\begin{hint}
הציבו $t=\sqrt x$.
\end{hint}

\begin{solution}
נציב $t=\sqrt x$, $x=t^2$, $dx=2t\,dt$; כאשר $x=0.25$, $t=\frac12$, וכאשר $x\to\infty$, $t\to\infty$:
\[ \int_{0.25}^{R}\frac{dx}{(4x+1)\sqrt x}=\int_{1/2}^{\sqrt R}\frac{2t\,dt}{(4t^2+1)t}=\int_{1/2}^{\sqrt R}\frac{2\,dt}{1+(2t)^2}=\Big[\arctan(2t)\Big]_{1/2}^{\sqrt R}=\arctan(2\sqrt R)-\frac\pi4. \]
כאשר $R\to\infty$, $\arctan(2\sqrt R)\to\frac\pi2$. לכן האינטגרל מתכנס ו-
\[ \int_{0.25}^{+\infty}\frac{dx}{(4x+1)\sqrt x}=\frac\pi2-\frac\pi4=\frac\pi4. \]
\end{solution}

\section*{תש"פ סמסטר א בוחן}

\subsection*{שאלה 1}

\begin{question}
הגבול $\displaystyle\lim_{x\to 0}\frac{\sin ax}{\sin bx}\,\frac{x}{\ln(1+cx)}$ שווה ל-
\begin{enumerate}
  \item $\frac{ac}{b}$
  \item $\frac{a}{bc}$
  \item $\frac{bc}{a}$
  \item $\frac{b}{ac}$
  \item הגבול לא קיים
\end{enumerate}
\end{question}

\begin{hint}
פרקו את הביטוי למכפלה של גבולות יסודיים: $\frac{\sin t}{t}\to 1$ ו-$\frac{\ln(1+t)}{t}\to 1$ כאשר $t\to 0$.
\end{hint}

\begin{solution}
\textbf{התשובה הנכונה: (ב)} $\frac{a}{bc}$ (בהנחה הסמויה בשאלה ש-$a,b,c\neq 0$, אחרת הביטוי אינו מוגדר או שהגבול שונה).

נכפול ונחלק כך שיופיעו גבולות יסודיים. עבור $x$ קרוב ל-0 ושונה ממנו:
\[
\frac{\sin ax}{\sin bx}\cdot\frac{x}{\ln(1+cx)}
=\frac{\sin ax}{ax}\cdot\frac{bx}{\sin bx}\cdot\frac{cx}{\ln(1+cx)}\cdot\frac{ax\cdot x}{bx\cdot cx}
=\frac{\sin ax}{ax}\cdot\frac{bx}{\sin bx}\cdot\frac{cx}{\ln(1+cx)}\cdot\frac{a}{bc}.
\]
כאשר $x\to 0$ גם $ax,bx,cx\to 0$, ולכן לפי הגבול היסודי $\lim_{t\to 0}\frac{\sin t}{t}=1$ (והצבה $t=ax$, $t=bx$) מתקיים
$\frac{\sin ax}{ax}\to 1$ ו-$\frac{bx}{\sin bx}\to 1$, ולפי הגבול היסודי $\lim_{t\to 0}\frac{\ln(1+t)}{t}=1$ (מדף הנוסחאות, $\lim_{t\to0}\frac{\log_a(1+t)}{t}=\frac{1}{\ln a}$ עם בסיס $e$) מתקיים $\frac{cx}{\ln(1+cx)}\to 1$.
לפי אריתמטיקה של גבולות:
\[
\lim_{x\to 0}\frac{\sin ax}{\sin bx}\,\frac{x}{\ln(1+cx)}=1\cdot 1\cdot 1\cdot\frac{a}{bc}=\boxed{\frac{a}{bc}}.
\]

\textbf{למה האחרות שגויות:} (א), (ג), (ד) הן צירופים אחרים של $a,b,c$, ששונים מ-$\frac{a}{bc}$ באופן כללי: למשל עבור $a=2,\ b=3,\ c=5$ הגבול הוא $\frac{2}{15}$, בעוד (א) נותנת $\frac{10}{3}$, (ג) נותנת $\frac{15}{2}$ ו-(ד) נותנת $\frac{3}{10}$. (ה) שגויה כי הראינו שהגבול קיים.
\end{solution}

\subsection*{שאלה 2}

\begin{question}
הגבול $\displaystyle\lim_{x\to\infty}\left(\frac{x^2-1}{x^2+3x}\right)^{\frac{x}{3}}$ שווה ל-
\begin{enumerate}
  \item $e^3$
  \item $e$
  \item $e^{\frac13}$
  \item $e^{-1}$
  \item $\infty$
\end{enumerate}
\end{question}

\begin{hint}
זהו גבול מהצורה $1^\infty$. כתבו את הבסיס כ-$1+t(x)$ עם $t(x)\to 0$ והשתמשו ב-$\lim_{y\to\infty}\left(1+\frac1y\right)^y=e$.
\end{hint}

\begin{solution}
\textbf{התשובה הנכונה: (ד)} $e^{-1}$.

הבסיס שואף ל-1 (חלוקה ב-$x^2$: $\frac{1-1/x^2}{1+3/x}\to 1$) והמעריך שואף ל-$\infty$, כלומר זהו גבול מהצורה $1^\infty$. נכתוב
\[
\frac{x^2-1}{x^2+3x}=1+\frac{x^2-1-x^2-3x}{x^2+3x}=1+t(x),\qquad t(x)=\frac{-3x-1}{x^2+3x}\xrightarrow[x\to\infty]{}0 .
\]
עבור $x$ גדול $t(x)\neq 0$ ולכן
\[
\left(1+t(x)\right)^{\frac x3}=\Bigl[\left(1+t(x)\right)^{\frac{1}{t(x)}}\Bigr]^{t(x)\cdot\frac{x}{3}} .
\]
לפי הגבול היסודי $\lim_{s\to 0}(1+s)^{1/s}=e$ (הצבה $s=t(x)\to 0$) הבסיס הפנימי שואף ל-$e$. המעריך:
\[
t(x)\cdot\frac x3=\frac{-3x-1}{x^2+3x}\cdot\frac{x}{3}=\frac{-3x^2-x}{3x^2+9x}=\frac{-3-\frac1x}{3+\frac9x}\xrightarrow[x\to\infty]{}-1 .
\]
מכיוון ש-$u^v=e^{v\ln u}$ ופונקציות $\ln$ ו-$\exp$ רציפות, אם $u\to e$ ו-$v\to -1$ אז $u^v\to e^{-1}$. לכן
\[
\lim_{x\to\infty}\left(\frac{x^2-1}{x^2+3x}\right)^{\frac{x}{3}}=\boxed{e^{-1}} .
\]

\textbf{למה האחרות שגויות:} (א), (ב), (ג) מתקבלות משגיאה בסימן או במקדם של המעריך (למשל התעלמות מכך ש-$t(x)$ שלילי). (ה) שגויה: הבסיס קטן מ-1 עבור $x$ גדול (המונה קטן מהמכנה), ולכן החזקה אף חסומה ע"י 1 ואינה יכולה לשאוף ל-$\infty$.
\end{solution}

\subsection*{שאלה 3}

\begin{question}
הנגזרת של $e^{2x}\ln\sqrt{1+x}$ שווה ל-
\begin{enumerate}
  \item $e^{2x}\ln\sqrt{1+x}+e^{2x}\frac{1}{2(1+x)}$
  \item $2e^{2x}\ln\sqrt{1+x}+e^{2x}\frac{1}{2(1+x)}$
  \item $2e^{2x}\ln\sqrt{1+x}+e^{2x}\frac{1}{(1+x)}$
  \item $e^{2x}\ln\sqrt{1+x}+e^{2x}\frac{1}{(1+x)}$
  \item $2e^{2x}\sqrt{1+x}+2e^{2x}\frac{1}{(1+x)}$
\end{enumerate}
\end{question}

\begin{solution}
\textbf{התשובה הנכונה: (ב)}.

לפי כלל המכפלה $(fg)'=f'g+fg'$ עם $f(x)=e^{2x}$, $g(x)=\ln\sqrt{1+x}$ (מוגדרת עבור $x>-1$).

לפי כלל השרשרת: $f'(x)=e^{2x}\cdot 2=2e^{2x}$.

עבור $g$ נשתמש בכך ש-$\ln\sqrt{1+x}=\frac12\ln(1+x)$, ולכן $g'(x)=\frac12\cdot\frac{1}{1+x}=\frac{1}{2(1+x)}$.
(לחלופין, בכלל השרשרת: $g'(x)=\frac{1}{\sqrt{1+x}}\cdot\frac{1}{2\sqrt{1+x}}=\frac{1}{2(1+x)}$.)

לכן
\[
\left(e^{2x}\ln\sqrt{1+x}\right)'=\boxed{2e^{2x}\ln\sqrt{1+x}+e^{2x}\frac{1}{2(1+x)}} .
\]

\textbf{למה האחרות שגויות:} (א) ו-(ד) שוכחות את הגורם 2 מנגזרת המעריך $2x$ (כלל השרשרת). (ג) ו-(ד) שוכחות את הגורם $\frac12$ שנובע מהשורש. (ה) מחליפה את $\ln\sqrt{1+x}$ ב-$\sqrt{1+x}$ ואת המקדם של האיבר השני ב-2, וזו אינה הנגזרת.
\end{solution}

\subsection*{שאלה 4}

\begin{question}
הנגזרת $y'(0)$ של הפונקציה הסתומה $y(x)$ המוגדרת על ידי המשוואה
$\cos\left(\frac{\pi}{2}+x+y\right)-y=0$ בסביבה של הנקודה $(0,0)$ שווה ל-
\begin{enumerate}
  \item $\pi/2$
  \item $-\frac12$
  \item $1+\pi/2$
  \item $\frac12$
  \item $0$
\end{enumerate}
\end{question}

\begin{hint}
גזרו את שני אגפי המשוואה לפי $x$, כאשר $y=y(x)$ (כלל השרשרת), ואז הציבו $x=0,\ y=0$.
\end{hint}

\begin{solution}
\textbf{התשובה הנכונה: (ב)} $-\frac12$.

ראשית, הנקודה $(0,0)$ מקיימת את המשוואה: $\cos\frac{\pi}{2}-0=0$.
נגזור את שני אגפי המשוואה $\cos\left(\frac{\pi}{2}+x+y(x)\right)-y(x)=0$ לפי $x$, תוך שימוש בכלל השרשרת:
\[
-\sin\left(\frac{\pi}{2}+x+y\right)\cdot\left(1+y'\right)-y'=0 .
\]
נציב $x=0$, $y(0)=0$: $\sin\frac{\pi}{2}=1$, ולכן
\[
-(1+y'(0))-y'(0)=0\quad\Longrightarrow\quad -1-2y'(0)=0\quad\Longrightarrow\quad y'(0)=\boxed{-\frac12}.
\]
(בדיקה: לפי הזהות $\cos(\frac{\pi}{2}+\theta)=-\sin\theta$ המשוואה היא $-\sin(x+y)=y$; גזירה נותנת $-\cos(x+y)(1+y')=y'$, ובנקודה $(0,0)$ שוב $y'(0)=-\frac12$.)

\textbf{למה האחרות שגויות:} (ד) מתקבלת משגיאת סימן בנגזרת של $\cos$. (ה) מתקבלת אם שוכחים לגזור את $y$ שבתוך הקוסינוס או את האיבר $-y$. (א) ו-(ג) מתקבלות מבלבול בין ערך הזווית $\frac{\pi}{2}$ לבין הנגזרת; הנגזרת היא ערך יחיד שחישבנו, $-\frac12$.
\end{solution}

\subsection*{שאלה 5}

\begin{question}
הפונקציה $f(x)$ רציפה בנקודה $x=x_0$. איזו טענה נכונה ?
\begin{enumerate}
  \item $f(x)$ גזירה בנקודה $x=x_0$
  \item קיים גבול $\displaystyle\lim_{x\to x_0}\frac{f(x)-f(x_0)}{x-x_0}$
  \item קיים גבול $\displaystyle\lim_{x\to x_0}\left(f(x)-f(x_0)\right)$
  \item קיים גבול $\displaystyle\lim_{h\to 0}\frac{f(x_0+h)-f(x_0)}{h}$
  \item $f'(x)$ רציפה בנקודה $x=x_0$
\end{enumerate}
\end{question}

\begin{hint}
רציפות אינה גוררת גזירות. חשבו על $f(x)=|x|$ בנקודה $x_0=0$.
\end{hint}

\begin{solution}
\textbf{התשובה הנכונה: (ג)}.

לפי ההגדרה, $f$ רציפה ב-$x_0$ אם ורק אם $\lim_{x\to x_0}f(x)=f(x_0)$. לפי אריתמטיקה של גבולות (גבול של קבוע הוא הקבוע עצמו) זה שקול לכך ש-
\[
\lim_{x\to x_0}\bigl(f(x)-f(x_0)\bigr)=f(x_0)-f(x_0)=0 ,
\]
ובפרט הגבול בטענה (ג) קיים (ושווה ל-0).

\textbf{למה האחרות שגויות:} נתבונן ב-$f(x)=|x|$ וב-$x_0=0$. הפונקציה רציפה ב-0, אבל
\[
\frac{f(0+h)-f(0)}{h}=\frac{|h|}{h}=\begin{cases}1 & h>0\\ -1 & h<0\end{cases}
\]
ולכן הגבולות החד-צדדיים של מנת ההפרשים הם $1$ ו-$-1$ והגבול אינו קיים. מכאן:
\begin{itemize}
  \item (ב) ו-(ד) שגויות: אלה שתי צורות כתיבה של הגדרת הנגזרת $f'(x_0)$, והראינו שהגבול לא קיים.
  \item (א) שגויה: גזירות ב-$x_0$ פירושה בדיוק קיום הגבול ב-(ב)/(ד).
  \item (ה) שגויה: $f'(0)$ כלל אינה מוגדרת, ולכן $f'$ אינה רציפה ב-0. (יתרה מזו, גם פונקציה גזירה יכולה להיות בעלת נגזרת לא רציפה, למשל $x^2\sin\frac1x$ עם $f(0)=0$.)
\end{itemize}
\end{solution}

\section*{תש"פ סמסטר א מועד א}

\subsection*{חלק א', שאלה 1}

\begin{question}
חקרו את הפונקציה $y=\frac{e^{2x}}{x}$ לפי הסעיפים הבאים:
\begin{enumerate}
  \item תחום הגדרה, נקודות חיתוך עם הצירים, זוגיות, אי-זוגיות
  \item רציפות, נקודות אי-רציפות
  \item אסימפטוטות
  \item תחומי עליה וירידה, נקודות קיצון
  \item תחומי קמירות וקעירות, נקודות פיתול
  \item תיאור גרפי.
\end{enumerate}
\end{question}

\begin{hint}
$y'=\frac{e^{2x}(2x-1)}{x^2}$ ו-$y''=\frac{2e^{2x}(2x^2-2x+1)}{x^3}$. בדקו את הסימן של כל גורם.
\end{hint}

\begin{solution}
\begin{enumerate}
  \item \textbf{תחום הגדרה:} $x\neq0$, כלומר $(-\infty,0)\cup(0,\infty)$.

  \textbf{חיתוך עם הצירים:} $x=0$ לא בתחום — אין חיתוך עם ציר $y$. $e^{2x}>0$ תמיד, ולכן $y\ne0$ — אין חיתוך עם ציר $x$. (הסימן: $y<0$ עבור $x<0$ ו-$y>0$ עבור $x>0$.)

  \textbf{זוגיות:} $y(-x)=\frac{e^{-2x}}{-x}$. למשל $y(1)=e^2$, $y(-1)=-e^{-2}$, ולכן $y(-1)\ne y(1)$ וגם $y(-1)\ne -y(1)$. הפונקציה \textbf{אינה זוגית ואינה אי-זוגית}.

  \item הפונקציה היא מנה של פונקציות רציפות ($e^{2x}$ ו-$x$) ולכן רציפה בכל נקודה שבה המכנה שונה מאפס, כלומר בכל תחום הגדרתה. הנקודה היחידה שבעייתית היא $x=0$ (שבה הפונקציה אינה מוגדרת):
  \[ \lim_{x\to0^-}\frac{e^{2x}}{x}=\frac{1}{0^-}=-\infty,\qquad \lim_{x\to0^+}\frac{e^{2x}}{x}=\frac{1}{0^+}=+\infty, \]
  ולכן $x=0$ היא נקודת אי-רציפות מסוג שני (עיקרית).

  \item \textbf{אסימפטוטה אנכית:} $x=0$ (לפי הגבולות בסעיף הקודם).

  \textbf{ב-$-\infty$:} $e^{2x}\to0$ ו-$\frac1x\to0$, לכן $\lim_{x\to-\infty}y=0$ — אסימפטוטה אופקית $y=0$ משמאל (הגרף מתחת לציר).

  \textbf{ב-$+\infty$:} $\frac{y}{x}=\frac{e^{2x}}{x^2}\to\infty$ (לפי לופיטל פעמיים: $\lim\frac{e^{2x}}{x^2}=\lim\frac{2e^{2x}}{2x}=\lim\frac{4e^{2x}}{2}=\infty$), לכן אין אסימפטוטה אופקית או משופעת מימין.

  \item לפי כלל המנה:
  \[ y'=\frac{2e^{2x}\cdot x - e^{2x}}{x^2} = \frac{e^{2x}(2x-1)}{x^2}. \]
  מכיוון ש-$e^{2x}>0$ ו-$x^2>0$, הסימן של $y'$ הוא הסימן של $2x-1$:
  \begin{itemize}
    \item $y'<0$ עבור $x<0$ ועבור $0<x<\frac12$: הפונקציה \textbf{יורדת} ב-$(-\infty,0)$ וב-$(0,\frac12)$.
    \item $y'>0$ עבור $x>\frac12$: הפונקציה \textbf{עולה} ב-$(\frac12,\infty)$.
  \end{itemize}
  ב-$x=\frac12$ הנגזרת מחליפה סימן מ-$-$ ל-$+$, ולכן זו נקודת \textbf{מינימום מקומי}: $y(\tfrac12)=\frac{e}{1/2}=2e$, כלומר $(\frac12,2e)$. אין נקודות קיצון נוספות.

  \item נגזור את $y'=e^{2x}\cdot\frac{2x-1}{x^2}$:
  \[ \left(\frac{2x-1}{x^2}\right)'=\frac{2x^2-2x(2x-1)}{x^4}=\frac{2-2x}{x^3}, \]
  \[ y''=e^{2x}\left(\frac{2(2x-1)}{x^2}+\frac{2-2x}{x^3}\right)=\frac{e^{2x}\left(4x^2-2x+2-2x\right)}{x^3}=\frac{2e^{2x}\left(2x^2-2x+1\right)}{x^3}. \]
  לטרינום $2x^2-2x+1$ דיסקרימיננטה $4-8<0$ ומקדם מוביל חיובי, ולכן הוא חיובי תמיד. מכאן הסימן של $y''$ הוא הסימן של $x^3$, כלומר של $x$:
  \begin{itemize}
    \item $x<0$: $y''<0$ — הפונקציה \textbf{קעורה} ($\cap$) ב-$(-\infty,0)$.
    \item $x>0$: $y''>0$ — הפונקציה \textbf{קמורה} ($\cup$) ב-$(0,\infty)$.
  \end{itemize}
  סימן $y''$ מתחלף רק ב-$x=0$, שאינה בתחום ההגדרה, ולכן \textbf{אין נקודות פיתול}.

  \item \textbf{תיאור גרפי:} משמאל לציר $y$ הגרף מתחת לציר $x$: מתחיל קרוב ל-$y=0$ ב-$-\infty$, יורד (וקעור) ושואף ל-$-\infty$ כאשר $x\to0^-$. מימין לציר $y$ הגרף מעל ציר $x$: יורד מ-$+\infty$ (ליד $x=0^+$) עד נקודת המינימום $(\frac12,2e)\approx(0.5,\,5.44)$, ואז עולה (בקמירות) ל-$+\infty$ במהירות אקספוננציאלית.
\end{enumerate}
\end{solution}

\subsection*{חלק א', שאלה 2}

\begin{question}
מצאו את הערך הגדול ביותר ואת הערך הקטן ביותר של הפונקציה $y=\sqrt[3]{x}(1+2x)$ בקטע $[-8,1]$.
\end{question}

\begin{hint}
המועמדים הם קצות הקטע, נקודות שבהן $y'=0$, ונקודות שבהן $y$ אינה גזירה (שימו לב ל-$x=0$).
\end{hint}

\begin{solution}
הפונקציה $y=x^{1/3}+2x^{4/3}$ רציפה בקטע הסגור $[-8,1]$, ולכן לפי \textbf{משפט ויירשטראס} היא מקבלת בו מקסימום ומינימום. נקודות אלה הן בקצות הקטע או בנקודות פנימיות קריטיות (שבהן $y'=0$ או $y'$ לא קיימת), לפי משפט פרמה.

\textbf{נגזרת} (עבור $x\ne0$):
\[ y' = \frac13x^{-2/3}+\frac83x^{1/3} = \frac{1+8x}{3\sqrt[3]{x^2}}. \]
\begin{itemize}
  \item $y'=0\iff x=-\frac18\in(-8,1)$.
  \item ב-$x=0$ הנגזרת אינה קיימת: $\frac{y(h)-y(0)}{h}=\frac{h^{1/3}(1+2h)}{h}=\frac{1+2h}{h^{2/3}}\to+\infty$. לכן $x=0$ נקודה קריטית.
\end{itemize}

\textbf{ערכי הפונקציה במועמדים:}
\begin{align*}
y(-8) &= \sqrt[3]{-8}\,(1-16) = (-2)(-15) = 30,\\
y\left(-\tfrac18\right) &= \sqrt[3]{-\tfrac18}\left(1-\tfrac14\right) = -\tfrac12\cdot\tfrac34 = -\tfrac38,\\
y(0) &= 0,\\
y(1) &= 1\cdot 3 = 3.
\end{align*}

\textbf{תשובה:} הערך הגדול ביותר הוא $30$ (מתקבל ב-$x=-8$), והערך הקטן ביותר הוא $-\frac38$ (מתקבל ב-$x=-\frac18$).
\end{solution}

\subsection*{חלק א', שאלה 3}

\begin{question}
מצאו את האינטגרל
\[ \int\left(\frac{x^2-x-4}{(x-1)(x^2+2x+1)}+e^{2x}x^2\right)dx. \]
\end{question}

\begin{hint}
שימו לב ש-$x^2+2x+1=(x+1)^2$, ופרקו לשברים חלקיים מהצורה $\frac{A}{x-1}+\frac{B}{x+1}+\frac{C}{(x+1)^2}$.
\end{hint}

\begin{hint}
את $\int x^2e^{2x}dx$ חשבו באינטגרציה בחלקים פעמיים (גוזרים את הפולינום).
\end{hint}

\begin{solution}
לפי לינאריות האינטגרל נחשב כל מחובר בנפרד.

\textbf{המחובר הראשון.} $x^2+2x+1=(x+1)^2$. המונה ממעלה 2 והמכנה ממעלה 3, ולכן נפרק לשברים חלקיים:
\[ \frac{x^2-x-4}{(x-1)(x+1)^2} = \frac{A}{x-1}+\frac{B}{x+1}+\frac{C}{(x+1)^2}, \]
כלומר $x^2-x-4 = A(x+1)^2+B(x-1)(x+1)+C(x-1)$.
\begin{itemize}
  \item $x=1$: $1-1-4=4A$, ולכן $A=-1$.
  \item $x=-1$: $1+1-4=-2C$, ולכן $C=1$.
  \item השוואת מקדמי $x^2$: $1=A+B$, ולכן $B=2$.
\end{itemize}
(בדיקה, מקדם חופשי: $A-B-C=-1-2-1=-4$ $\checkmark$.) לכן
\[ \int\frac{x^2-x-4}{(x-1)(x+1)^2}dx = -\ln|x-1| + 2\ln|x+1| - \frac{1}{x+1}+C_1. \]

\textbf{המחובר השני.} אינטגרציה בחלקים עם $u=x^2$, $v'=e^{2x}$ ($v=\frac12e^{2x}$):
\[ \int x^2e^{2x}dx = \frac{x^2}{2}e^{2x} - \int xe^{2x}dx. \]
שוב בחלקים עם $u=x$, $v'=e^{2x}$:
\[ \int xe^{2x}dx = \frac{x}{2}e^{2x}-\frac12\int e^{2x}dx = \frac x2e^{2x}-\frac14e^{2x}+C_2. \]
לכן
\[ \int x^2e^{2x}dx = e^{2x}\left(\frac{x^2}{2}-\frac{x}{2}+\frac14\right)+C_3. \]

\textbf{תשובה:}
\[ \int\left(\frac{x^2-x-4}{(x-1)(x^2+2x+1)}+e^{2x}x^2\right)dx = -\ln|x-1|+2\ln|x+1|-\frac{1}{x+1}+e^{2x}\left(\frac{x^2}{2}-\frac{x}{2}+\frac14\right)+C. \]
(בדיקה: גזירת התוצאה מחזירה את האינטגרנד.)
\end{solution}

\subsection*{חלק ב', שאלה 1}

\begin{question}
תהי $f(x)$ רציפה בקטע $[a,b]$. מה מהבאים חייב להיות נכון?
\begin{enumerate}
  \item[(א)] $f(x)$ לא גזירה בקטע $(a,b)$
  \item[(ב)] $f(x)$ מקיימת את התנאי המשפט לגרנז' בקטע $[a,b]$
  \item[(ג)] $f(x)$ מקיימת את התנאי המשפט רול בקטע $[a,b]$
  \item[(ד)] $f(x)$ גזירה בקטע $(a,b)$
  \item[(ה)] קיים $M$ כך ש-$f(x)<M$ לכל $x$ ששייך ל-$[a,b]$
\end{enumerate}
\end{question}

\begin{hint}
איזה משפט מבטיח משהו לכל פונקציה רציפה בקטע סגור?
\end{hint}

\begin{solution}
\textbf{התשובה הנכונה: (ה).}

לפי \textbf{משפט ויירשטראס הראשון}, פונקציה רציפה בקטע סגור $[a,b]$ חסומה בו. לכן קיים $K$ כך ש-$f(x)\le K$ לכל $x\in[a,b]$, ואז $M=K+1$ מקיים $f(x)<M$ לכל $x\in[a,b]$.

שאר התשובות אינן חייבות להתקיים:
\begin{itemize}
  \item (א) לא נכון: למשל $f(x)=x$ רציפה וגזירה ב-$(a,b)$.
  \item (ב) ו-(ד) לא נכונות: משפט לגרנז' דורש גם גזירות ב-$(a,b)$, ופונקציה רציפה אינה חייבת להיות גזירה. למשל בקטע $[-1,1]$ הפונקציה $f(x)=|x|$ רציפה אך אינה גזירה ב-$0$.
  \item (ג) לא נכון: משפט רול דורש גם גזירות וגם $f(a)=f(b)$; למשל $f(x)=x$ רציפה אך $f(a)\ne f(b)$.
\end{itemize}
\end{solution}

\subsection*{חלק ב', שאלה 2}

\begin{question}
תהי $f(x)$ גזירה בקטע $(a,b)$. מה מהבאים חייב להיות נכון?
\begin{enumerate}
  \item[(א)] $f'(x)$ רציפה ב-$(a,b)$
  \item[(ב)] $f(x)$ מונוטונית ב-$(a,b)$
  \item[(ג)] $f(x)+\sin(x)$ רציפה ב-$(a,b)$
  \item[(ד)] $f(x)$ מונוטונית עולה ב-$(a,b)$
  \item[(ה)] קיים $c\in(a,b)$ כך ש-$f'(c)=0$
\end{enumerate}
\end{question}

\begin{hint}
מה גזירות בנקודה גוררת?
\end{hint}

\begin{solution}
\textbf{התשובה הנכונה: (ג).}

פונקציה גזירה בנקודה רציפה בה, ולכן $f$ רציפה ב-$(a,b)$. גם $\sin x$ רציפה, וסכום של פונקציות רציפות רציף; לכן $f(x)+\sin(x)$ רציפה ב-$(a,b)$.

שאר התשובות אינן חייבות להתקיים:
\begin{itemize}
  \item (א) לא נכון: הנגזרת של פונקציה גזירה אינה חייבת להיות רציפה. הדוגמה הקלאסית: $f(x)=x^2\sin\frac1x$ עבור $x\neq0$ ו-$f(0)=0$, בקטע $(-1,1)$. היא גזירה בכל נקודה ($f'(0)=\lim_{h\to0}h\sin\frac1h=0$), אך $f'(x)=2x\sin\frac1x-\cos\frac1x$ עבור $x\ne0$, ול-$f'$ אין גבול ב-$0$.
  \item (ב), (ד) לא נכונות: למשל $f(x)=x^2$ בקטע $(-1,1)$ גזירה אך אינה מונוטונית.
  \item (ה) לא נכון: למשל $f(x)=x$ גזירה ו-$f'(x)=1\ne0$ לכל $x$. (משפט רול דורש גם רציפות בקטע סגור ו-$f(a)=f(b)$.)
\end{itemize}
\end{solution}

\subsection*{חלק ב', שאלה 3}

\begin{question}
הנגזרת $y'(0)$ של הפונקציה הסתומה $y(x)$ המוגדרת על ידי המשוואה $\cos\left(\frac{\pi}{2}+x+y\right)-y=0$ בסביבה של הנקודה $(0,0)$ שווה ל-
\begin{enumerate}
  \item[(א)] $\frac12$
  \item[(ב)] $0$
  \item[(ג)] $\frac{\pi}{2}$
  \item[(ד)] $1+\frac{\pi}{2}$
  \item[(ה)] $-\frac12$
\end{enumerate}
\end{question}

\begin{hint}
גזרו את שני אגפי המשוואה לפי $x$, כאשר $y=y(x)$ (כלל השרשרת), והציבו $x=0$, $y=0$.
\end{hint}

\begin{solution}
\textbf{התשובה הנכונה: (ה).}

ראשית, הנקודה $(0,0)$ מקיימת את המשוואה: $\cos\frac{\pi}{2}-0=0$. נגזור את שני אגפי המשוואה לפי $x$, כאשר $y=y(x)$, לפי כלל השרשרת:
\[ -\sin\left(\frac{\pi}{2}+x+y\right)\cdot\left(1+y'\right) - y' = 0. \]
נציב $x=0$, $y(0)=0$; אז $\sin\frac{\pi}{2}=1$:
\[ -(1+y'(0)) - y'(0) = 0 \quad\Longrightarrow\quad -1-2y'(0)=0 \quad\Longrightarrow\quad y'(0) = -\frac12. \]
(לחלופין: $\cos\left(\frac\pi2+t\right)=-\sin t$, ולכן המשוואה היא $\sin(x+y)+y=0$; גזירה נותנת $\cos(x+y)(1+y')+y'=0$, וב-$(0,0)$: $1+2y'=0$.)
\end{solution}

\subsection*{חלק ב', שאלה 4}

\begin{question}
אם $f(x)=\frac{1-e^{-\frac{x^2}{2}}}{x}$, אזי בסביבה של $0$ מתקיים
\begin{enumerate}
  \item[(א)] $f(x)=\frac{x}{2}+\frac{x^3}{4}+o(x^3)$
  \item[(ב)] $f(x)=\frac{x}{2}+\frac{x^3}{8}+o(x^3)$
  \item[(ג)] $f(x)=\frac{x}{2}-\frac{x^3}{8}+o(x^3)$
  \item[(ד)] $f(x)=\frac{x}{2}-\frac{x^3}{2}+o(x^3)$
  \item[(ה)] $f(x)=\frac{x}{2}-\frac{x^3}{4}+o(x^3)$
\end{enumerate}
\end{question}

\begin{hint}
הציבו $t=-\frac{x^2}{2}$ בפיתוח $e^t=1+t+\frac{t^2}{2}+o(t^2)$.
\end{hint}

\begin{solution}
\textbf{התשובה הנכונה: (ג).}

לפי פיתוח מקלורן $e^t=1+t+\frac{t^2}{2}+o(t^2)$ כאשר $t\to0$, עם $t=-\frac{x^2}{2}$ (ואז $o(t^2)=o(x^4)$):
\[ e^{-x^2/2} = 1-\frac{x^2}{2}+\frac{x^4}{8}+o(x^4). \]
לכן
\[ 1-e^{-x^2/2} = \frac{x^2}{2}-\frac{x^4}{8}+o(x^4), \]
ובחלוקה ב-$x$ (עבור $x\neq0$; $\frac{o(x^4)}{x}=o(x^3)$):
\[ f(x)=\frac{x}{2}-\frac{x^3}{8}+o(x^3). \]
\end{solution}

\subsection*{חלק ב', שאלה 5}

\begin{question}
השטח החסום על ידי העקומות $y(x)=5x-x^2$ ו-$y(x)=2x$ שווה ל-
\begin{enumerate}
  \item[(א)] $\frac92$
  \item[(ב)] $1$
  \item[(ג)] $\frac{45}{2}$
  \item[(ד)] $\frac{25}{6}$
  \item[(ה)] $\frac{27}{2}$
\end{enumerate}
\end{question}

\begin{hint}
מצאו את נקודות החיתוך של העקומות, וקבעו איזו מהן עליונה בין נקודות החיתוך.
\end{hint}

\begin{solution}
\textbf{התשובה הנכונה: (א).}

נקודות חיתוך: $5x-x^2=2x\iff x^2-3x=0\iff x=0$ או $x=3$.
בקטע $(0,3)$ מתקיים $(5x-x^2)-2x = 3x-x^2 = x(3-x)>0$, כלומר הפרבולה מעל הישר. לכן
\[ S=\int_0^3\left[(5x-x^2)-2x\right]dx = \int_0^3(3x-x^2)\,dx = \left[\frac{3x^2}{2}-\frac{x^3}{3}\right]_0^3 = \frac{27}{2}-9 = \frac92. \]
\end{solution}

\section*{תש"פ סמסטר א מועד ב}

\subsection*{חלק א', שאלה 1}

\begin{question}
חקרו את הפונקציה $y=\frac{x^3-1}{4x^2}$ לפי הסעיפים הבאים:
\begin{enumerate}
  \item תחום הגדרה, נקודות חיתוך עם הצירים, זוגיות, אי-זוגיות
  \item רציפות, נקודות אי-רציפות
  \item אסימפטוטות
  \item תחומי עליה וירידה, נקודות קיצון
  \item תחומי קמירות וקעירות, נקודות פיתול
  \item תיאור גרפי.
\end{enumerate}
\end{question}

\begin{hint}
כדאי לכתוב את הפונקציה בצורה $y=\frac{x}{4}-\frac{1}{4x^2}$: ממנה רואים מיד את האסימפטוטה המשופעת, והגזירה פשוטה.
\end{hint}

\begin{solution}
נסמן $f(x)=\frac{x^3-1}{4x^2}=\frac{x}{4}-\frac{1}{4x^2}$.
\begin{enumerate}
  \item \textbf{תחום הגדרה:} המכנה מתאפס רק ב-$x=0$, ולכן $D=\{x\in\R : x\neq 0\}$.

  \textbf{חיתוך עם הצירים:} $x=0$ אינו בתחום, ולכן אין חיתוך עם ציר $y$. חיתוך עם ציר $x$: $x^3-1=0 \iff x=1$, כלומר הנקודה $(1,0)$.

  \textbf{זוגיות:} $f(-x)=\frac{-x^3-1}{4x^2}$. למשל $f(1)=0$ ואילו $f(-1)=-\frac12$, כך ש-$f(-1)\neq f(1)$ וגם $f(-1)\neq -f(1)$. לכן הפונקציה אינה זוגית ואינה אי-זוגית.

  \item $f$ היא מנה של פולינומים, ולכן רציפה בכל תחום הגדרתה. הנקודה היחידה שבה $f$ אינה מוגדרת היא $x=0$. שם המונה שואף ל-$-1$ והמכנה $4x^2\to 0^+$, ולכן
  \[ \lim_{x\to 0^-} f(x)=\lim_{x\to 0^+} f(x)=-\infty , \]
  כלומר ב-$x=0$ יש נקודת אי-רציפות מסוג שני (אינסופית).

  \item \textbf{אסימפטוטה אנכית:} מהסעיף הקודם, $x=0$ היא אסימפטוטה אנכית (מ-2 הצדדים הפונקציה שואפת ל-$-\infty$). אין אסימפטוטות אנכיות נוספות כי $f$ רציפה בכל נקודה אחרת.

  \textbf{אסימפטוטה משופעת:}
  \[ m=\lim_{x\to\pm\infty}\frac{f(x)}{x}=\lim_{x\to\pm\infty}\frac{x^3-1}{4x^3}=\frac14,\qquad
     n=\lim_{x\to\pm\infty}\Big(f(x)-\frac{x}{4}\Big)=\lim_{x\to\pm\infty}\Big(-\frac{1}{4x^2}\Big)=0 . \]
  לכן $y=\frac{x}{4}$ היא אסימפטוטה משופעת גם ב-$+\infty$ וגם ב-$-\infty$. מאחר ש-$f(x)-\frac{x}{4}=-\frac{1}{4x^2}<0$, הגרף נמצא מתחת לאסימפטוטה. אין אסימפטוטה אופקית (כי $f(x)\to\pm\infty$ כאשר $x\to\pm\infty$).

  \item גוזרים:
  \[ f'(x)=\frac14+\frac{1}{2x^3}=\frac{x^3+2}{4x^3}. \]
  $f'(x)=0 \iff x^3=-2 \iff x=-\sqrt[3]{2}$. טבלת סימנים (המונה $x^3+2$ חיובי עבור $x>-\sqrt[3]{2}$, המכנה $4x^3$ חיובי עבור $x>0$):
  \begin{itemize}
    \item $x<-\sqrt[3]{2}$: מונה שלילי, מכנה שלילי, $f'>0$ --- $f$ עולה.
    \item $-\sqrt[3]{2}<x<0$: מונה חיובי, מכנה שלילי, $f'<0$ --- $f$ יורדת.
    \item $x>0$: מונה חיובי, מכנה חיובי, $f'>0$ --- $f$ עולה.
  \end{itemize}
  לכן $f$ עולה ב-$(-\infty,-\sqrt[3]{2})$ וב-$(0,\infty)$, ויורדת ב-$(-\sqrt[3]{2},0)$.
  ב-$x=-\sqrt[3]{2}$ הנגזרת עוברת מחיובית לשלילית, ולכן זו נקודת מקסימום מקומי:
  \[ f(-\sqrt[3]{2})=\frac{-2-1}{4\sqrt[3]{4}}=-\frac{3}{4\sqrt[3]{4}}\approx -0.47 . \]
  אין נקודות קיצון נוספות (ב-$x=0$ הפונקציה אינה מוגדרת).

  \item
  \[ f''(x)=\Big(\frac14+\frac12x^{-3}\Big)'=-\frac{3}{2x^4}. \]
  לכל $x\neq 0$ מתקיים $f''(x)<0$. לכן $f$ קעורה ($\cap$) בכל אחד מהקטעים $(-\infty,0)$ ו-$(0,\infty)$, ואינה קמורה ($\cup$) באף קטע. הסימן של $f''$ לא מתחלף בתוך תחום ההגדרה, ולכן אין נקודות פיתול.

  \item \textbf{תיאור גרפי:} בצד שמאל הגרף נמצא מתחת לאסימפטוטה $y=\frac{x}{4}$ וצמוד אליה כש-$x\to-\infty$. הוא עולה עד למקסימום $\big(-\sqrt[3]{2},-\frac{3}{4\sqrt[3]{4}}\big)\approx(-1.26,-0.47)$, ואז יורד ל-$-\infty$ כש-$x\to0^-$. בצד ימין הוא עולה מ-$-\infty$ (כש-$x\to0^+$), חותך את ציר $x$ ב-$(1,0)$, ומתקרב מלמטה לאסימפטוטה $y=\frac{x}{4}$ כש-$x\to\infty$. הגרף קעור בכל מקום.
\end{enumerate}
\end{solution}

\subsection*{חלק א', שאלה 2}

\begin{question}
חשבו את הגבול
\[ \lim_{x\to 0}\left(\frac{\sqrt{x+9}-3}{\sqrt{(x+2)^2+5}-x-3}+(\cos(x))^{\frac{1}{\sin^2(2x)}}\right) \]
\end{question}

\begin{hint}
חשבו כל מחובר לחוד. במחובר הראשון הכפילו מונה ומכנה בצמודים (כדי להיפטר מהשורשים).
\end{hint}

\begin{hint}
במחובר השני כתבו $(\cos x)^{\frac{1}{\sin^2(2x)}}=e^{\frac{\ln\cos x}{\sin^2(2x)}}$, וחשבו את גבול המעריך.
\end{hint}

\begin{solution}
נחשב כל מחובר לחוד. אם לשניהם יש גבול סופי, גבול הסכום שווה לסכום הגבולות (אריתמטיקה של גבולות).

\textbf{המחובר הראשון.} עבור $x=0$ מתקבל $\frac{0}{0}$. נכפיל בצמודים. במונה:
\[ \sqrt{x+9}-3=\frac{(x+9)-9}{\sqrt{x+9}+3}=\frac{x}{\sqrt{x+9}+3}. \]
במכנה, $(x+2)^2+5=x^2+4x+9$, ולכן
\[ \sqrt{x^2+4x+9}-(x+3)=\frac{x^2+4x+9-(x+3)^2}{\sqrt{x^2+4x+9}+x+3}=\frac{-2x}{\sqrt{x^2+4x+9}+x+3}. \]
לכן, עבור $x\neq0$ קטן,
\[ \frac{\sqrt{x+9}-3}{\sqrt{(x+2)^2+5}-x-3}=\frac{x}{\sqrt{x+9}+3}\cdot\frac{\sqrt{x^2+4x+9}+x+3}{-2x}
=-\frac{\sqrt{x^2+4x+9}+x+3}{2(\sqrt{x+9}+3)}\xrightarrow[x\to0]{}-\frac{3+3}{2\cdot 6}=-\frac12 , \]
כאשר השתמשנו ברציפות השורש.

\textbf{המחובר השני.} ליד $0$ מתקיים $\cos x>0$, ולכן
\[ (\cos x)^{\frac{1}{\sin^2(2x)}}=\exp\Big(\frac{\ln(\cos x)}{\sin^2(2x)}\Big). \]
נחשב את גבול המעריך בעזרת פיתוחי טיילור סביב $0$:
$\cos x=1-\frac{x^2}{2}+o(x^2)$, ולכן $\ln(\cos x)=\ln\big(1+(-\frac{x^2}{2}+o(x^2))\big)=-\frac{x^2}{2}+o(x^2)$.
וגם $\sin(2x)=2x+o(x)$, ולכן $\sin^2(2x)=4x^2+o(x^2)$. מכאן
\[ \lim_{x\to0}\frac{\ln(\cos x)}{\sin^2(2x)}=\lim_{x\to0}\frac{-\frac{x^2}{2}+o(x^2)}{4x^2+o(x^2)}=-\frac18 . \]
(אפשר לקבל את זה גם בלופיטל: $\frac{\ln\cos x}{\sin^2 2x}$ הוא מהצורה $\frac00$, ויחס הנגזרות $\frac{-\tan x}{4\sin(2x)\cos(2x)}=\frac{-\tan x}{2\sin(4x)}\to -\frac{1}{8}$, כי $\frac{\tan x}{x}\to1$ ו-$\frac{\sin 4x}{4x}\to 1$.)
פונקציית האקספוננט רציפה, ולכן המחובר השני שואף ל-$e^{-1/8}$.

\textbf{תשובה:}
\[ \lim_{x\to 0}\left(\frac{\sqrt{x+9}-3}{\sqrt{(x+2)^2+5}-x-3}+(\cos x)^{\frac{1}{\sin^2(2x)}}\right)=\boxed{e^{-1/8}-\frac12}\approx 0.3825 . \]
\end{solution}

\subsection*{חלק א', שאלה 3}

\begin{question}
מה השטח הכלוא בין העקומות $y=\frac{4}{x}$ ו-$y=5-x$?
מה הנפח המתקבל מסיבוב השטח הזה סביב ציר $x$?
\end{question}

\begin{hint}
מצאו קודם את נקודות החיתוך של העקומות, ובדקו איזו מהן נמצאת מעל השנייה בין נקודות החיתוך.
\end{hint}

\begin{hint}
בסיבוב סביב ציר $x$ מתקבל גוף עם חור (שיטת הדסקיות/הטבעות): $V=\pi\int_a^b\big(y_{\text{עליונה}}^2-y_{\text{תחתונה}}^2\big)\,dx$.
\end{hint}

\begin{solution}
\textbf{נקודות חיתוך:} $\frac4x=5-x \iff 4=5x-x^2 \iff x^2-5x+4=0 \iff (x-1)(x-4)=0$, כלומר $x=1$ ו-$x=4$ (הנקודות $(1,4)$ ו-$(4,1)$).

\textbf{מי מעל מי:} בקטע $(1,4)$, למשל ב-$x=2$: $5-2=3>2=\frac42$. ובאופן כללי, עבור $x>0$,
$5-x-\frac4x=\frac{-(x^2-5x+4)}{x}=\frac{-(x-1)(x-4)}{x}>0$ כאשר $1<x<4$. לכן הישר נמצא מעל ההיפרבולה, ושתי הפונקציות חיוביות בקטע.

\textbf{השטח:} לפי הנוסחה לשטח בין שני גרפים,
\[ S=\int_1^4\Big(5-x-\frac4x\Big)dx=\Big[5x-\frac{x^2}{2}-4\ln x\Big]_1^4
=\Big(20-8-4\ln4\Big)-\Big(5-\frac12\Big)=\frac{15}{2}-8\ln2 . \]
כלומר $\boxed{S=\frac{15}{2}-8\ln 2\approx 1.955}$.

\textbf{הנפח:} בסיבוב סביב ציר $x$, כל חתך אנכי בנקודה $x\in[1,4]$ הוא טבעת עם רדיוס חיצוני $5-x$ ורדיוס פנימי $\frac4x$ (שני הרדיוסים חיוביים). לכן
\[ V=\pi\int_1^4\Big((5-x)^2-\frac{16}{x^2}\Big)dx=\pi\Big[-\frac{(5-x)^3}{3}+\frac{16}{x}\Big]_1^4
=\pi\Big[\Big(-\frac13+4\Big)-\Big(-\frac{64}{3}+16\Big)\Big]=\pi\Big(\frac{11}{3}+\frac{16}{3}\Big)=9\pi . \]
כלומר $\boxed{V=9\pi}$.
\end{solution}

\subsection*{חלק ב', שאלה 1}

\begin{question}
תהי $f(x)$ פונקציה כך ש-$f(x)+\sin(x)$ גזירה בקטע הפתוח $(a,b)$. מה מהבאים חייב להיות נכון?
\begin{enumerate}
  \item $f(x)+\cos(x)$ גזירה בקטע הפתוח $(a,b)$.
  \item $f(x)+\sin(x)$ רציפה בקטע הסגור $[a,b]$.
  \item $f(x)$ רציפה בקטע הסגור $[a,b]$.
  \item $f(x)$ חסומה בקטע הפתוח $(a,b)$.
  \item $f(x)+\sin(x)$ חסומה בקטע הפתוח $(a,b)$.
\end{enumerate}
\end{question}

\begin{hint}
כתבו את $f+\cos$ כסכום של $f+\sin$ ופונקציה גזירה ידועה.
\end{hint}

\begin{solution}
\textbf{התשובה הנכונה: א.}

נסמן $g(x)=f(x)+\sin x$, גזירה ב-$(a,b)$. אז
\[ f(x)+\cos x=g(x)-\sin x+\cos x , \]
וזה סכום של פונקציות גזירות ב-$(a,b)$ ($\sin$ ו-$\cos$ גזירות בכל $\R$), ולכן גזיר ב-$(a,b)$. כלומר א' נכונה תמיד.

שאר הטענות אינן חייבות להתקיים. דוגמה נגדית: $f(x)=\frac{1}{x-a}-\sin x$ ב-$(a,b)$. אז $f(x)+\sin x=\frac{1}{x-a}$ גזירה ב-$(a,b)$, אבל:
\begin{itemize}
  \item $f+\sin$ ו-$f$ אינן מוגדרות (ובוודאי אינן רציפות) ב-$x=a$ --- ב' וג' נופלות. (גם אם נגדיר להן ערך כלשהו ב-$a$, הן אינן חסומות ליד $a$ ולכן אינן רציפות בה.)
  \item $\frac{1}{x-a}\to+\infty$ כאשר $x\to a^+$, ולכן $f+\sin$ אינה חסומה ב-$(a,b)$ --- ה' נופלת; וגם $f=\frac{1}{x-a}-\sin x$ אינה חסומה (כי $\sin$ חסומה) --- ד' נופלת.
\end{itemize}
\end{solution}

\subsection*{חלק ב', שאלה 2}

\begin{question}
הגרף הנתון הוא גרף בקואורדינטות פולריות של הפונקציה:

(בשרטוט: שתי לולאות סגורות ושוות, בצורת עיגול, אחת מעל ציר $x$ ואחת מתחתיו, הנוגעות זו בזו בראשית הצירים; הלולאות סימטריות סביב ציר $y$ ונמתחות לאורכו.)
\begin{enumerate}
  \item $r(\phi)=\sin^2\phi$
  \item $r(\phi)=1+\cos\phi$
  \item $r(\phi)=1-\cos\phi$
  \item $r(\phi)=\cos^2\phi$
  \item $r(\phi)=\cos\phi$
\end{enumerate}
\end{question}

\begin{hint}
בדקו עבור כל פונקציה באילו זוויות $r=0$ (שם העקומה עוברת בראשית) ובאילו זוויות $r$ מקסימלי (הנקודות הרחוקות ביותר מהראשית).
\end{hint}

\begin{solution}
\textbf{התשובה הנכונה: א.}

בגרף העקומה עוברת בראשית בכיוון ציר $x$ (שם שתי הלולאות נוגעות), והנקודות הרחוקות ביותר מהראשית נמצאות על ציר $y$, למעלה ולמטה, באותו מרחק. כמו כן יש שתי לולאות.

\begin{itemize}
  \item $r=\sin^2\phi$: מתקיים $r\ge0$ לכל $\phi$, $r=0$ בדיוק עבור $\phi=0,\pi$ (כיוון ציר $x$), ו-$r=1$ מקסימלי עבור $\phi=\frac{\pi}{2},\frac{3\pi}{2}$, כלומר בנקודות $(0,1)$ ו-$(0,-1)$. כש-$\phi$ עובר מ-$0$ ל-$\pi$ נוצרת לולאה בחצי העליון, וכש-$\phi$ עובר מ-$\pi$ ל-$2\pi$ נוצרת לולאה בחצי התחתון. זה בדיוק הגרף הנתון.
  \item $r=1\pm\cos\phi$: קרדיואידות --- לולאה אחת בלבד (עם "שקע" בראשית), סימטרית סביב ציר $x$. לא מתאים.
  \item $r=\cos^2\phi$: שתי לולאות, אבל לאורך ציר $x$ (מקסימום ב-$\phi=0,\pi$). לא מתאים.
  \item $r=\cos\phi$: מעגל אחד, $x^2+y^2=x$, שמרכזו $(\frac12,0)$. לא מתאים.
\end{itemize}
\end{solution}

\subsection*{חלק ב', שאלה 3}

\begin{question}
הנגזרת $y'(0)$ של הפונקציה הסתומה $y(x)$ המוגדרת על ידי המשוואה $\sin(x-y)-y=0$ בסביבה של הנקודה $(0,0)$ שווה ל-
\begin{enumerate}
  \item $\frac12$
  \item $-\frac12$
  \item $\frac{\pi}{2}$
  \item $1+\frac{\pi}{2}$
  \item $0$
\end{enumerate}
\end{question}

\begin{solution}
\textbf{התשובה הנכונה: א.}

הנקודה $(0,0)$ מקיימת את המשוואה: $\sin(0)-0=0$. נגזור את שני האגפים לפי $x$, כאשר $y=y(x)$ (כלל השרשרת):
\[ \cos(x-y)\cdot(1-y')-y'=0 . \]
נציב $x=0,\ y=0$: $\cos 0\cdot(1-y'(0))-y'(0)=0$, כלומר $1-2y'(0)=0$, ולכן $y'(0)=\frac12$.

(באופן כללי: $y'=\frac{\cos(x-y)}{1+\cos(x-y)}$, וב-$(0,0)$ זה $\frac{1}{2}$.)
\end{solution}

\subsection*{חלק ב', שאלה 4}

\begin{question}
אם $f(x)=\sin^2(x)$, אזי בסביבה של $0$ מתקיים
\begin{enumerate}
  \item $f(x)=x^2-\frac{x^4}{3}+o(x^4)$
  \item $f(x)=x^2+\frac{x^4}{3}+o(x^4)$
  \item $f(x)=x^2-\frac{x^4}{6}+o(x^4)$
  \item $f(x)=x^2+\frac{x^4}{6}+o(x^4)$
  \item $f(x)=x^2+\frac{x^4}{36}+o(x^4)$
\end{enumerate}
\end{question}

\begin{hint}
השתמשו בזהות $\sin^2x=\frac{1-\cos(2x)}{2}$ ובפיתוח המוכר של $\cos$.
\end{hint}

\begin{solution}
\textbf{התשובה הנכונה: א.}

לפי הזהות $\sin^2x=\frac{1-\cos 2x}{2}$ ופיתוח מקלורן $\cos t=1-\frac{t^2}{2}+\frac{t^4}{24}+o(t^4)$ עם $t=2x$:
\[ \cos 2x=1-2x^2+\frac{16x^4}{24}+o(x^4)=1-2x^2+\frac{2x^4}{3}+o(x^4), \]
ולכן
\[ \sin^2x=\frac{1-\cos2x}{2}=x^2-\frac{x^4}{3}+o(x^4). \]
(בדיקה בדרך אחרת: $\sin x=x-\frac{x^3}{6}+o(x^4)$, ולכן $\sin^2x=x^2-2\cdot x\cdot\frac{x^3}{6}+o(x^4)=x^2-\frac{x^4}{3}+o(x^4)$.)
\end{solution}

\subsection*{חלק ב', שאלה 5}

\begin{question}
אם $f(x)$ פונקציה ממשית שנגזרתה היא $f'(x)=(x-1)^2(x+2)(x-5)$ אזי:
\begin{enumerate}
  \item ל-$f(x)$ יש מקסימום ב-$x=-2$ ומינימום ב-$x=5$, ואין נק' קיצון ב-$x=1$.
  \item ל-$f(x)$ יש מקסימום ב-$x=5$ ומינימום ב-$x=-2$, ואין נק' קיצון ב-$x=1$.
  \item ל-$f(x)$ יש מקסימום ב-$x=-2$ וב-$x=1$ ומינימום ב-$x=5$.
  \item ל-$f(x)$ יש מקסימום ב-$x=5$ ומינימום ב-$x=-2$ וב-$x=1$.
  \item ל-$f(x)$ אין נקודות קיצון.
\end{enumerate}
\end{question}

\begin{hint}
בדקו את סימן הנגזרת משני צידי כל נקודה קריטית. שימו לב שהגורם $(x-1)^2$ אינו מחליף סימן.
\end{hint}

\begin{solution}
\textbf{התשובה הנכונה: א.}

$f$ גזירה בכל $\R$, ולכן נקודות קיצון אפשריות רק היכן ש-$f'=0$ (משפט פרמה): $x=-2,\,1,\,5$. נבדוק את סימן $f'$. הגורם $(x-1)^2\ge0$ ואינו מחליף סימן, ולכן הסימן נקבע על ידי $(x+2)(x-5)$:
\begin{itemize}
  \item $x<-2$: $(x+2)(x-5)>0$, ולכן $f'>0$ --- $f$ עולה.
  \item $-2<x<5$, $x\neq 1$: $(x+2)(x-5)<0$, ולכן $f'<0$ --- $f$ יורדת.
  \item $x>5$: $f'>0$ --- $f$ עולה.
\end{itemize}
לכן ב-$x=-2$ הנגזרת עוברת מחיובית לשלילית --- מקסימום מקומי; ב-$x=5$ היא עוברת משלילית לחיובית --- מינימום מקומי; וב-$x=1$ הנגזרת שלילית משני הצדדים, כך ש-$f$ יורדת ממש בסביבת $1$ ואין שם קיצון.
\end{solution}

\section*{תש"פ סמסטר א מועד ג}

\subsection*{חלק א', שאלה 1}

\begin{question}
חקרו את הפונקציה $y=\frac{2x-1}{(x-1)^2}$ לפי הסעיפים הבאים:
\begin{enumerate}
  \item תחום ההגדרה של הפונקציה, נקודות חיתוך עם צירים, זוגיות, אי-זוגיות.
  \item רציפות, נקודות אי-רציפות.
  \item אסימפטוטות
  \item תחום העלייה וירידה, נקודות קיצון.
  \item תחומי קמירות וקעירות, נקודות פיתול.
  \item תיאור גרפי.
\end{enumerate}
\end{question}

\begin{hint}
בגזירה לפי כלל המנה כדאי לצמצם גורם $(x-1)$ אחד מהמונה והמכנה לפני שמפשטים.
\end{hint}

\begin{solution}
נסמן $f(x)=\frac{2x-1}{(x-1)^2}$.
\begin{enumerate}
  \item \textbf{תחום הגדרה:} $D=\{x\in\R: x\neq1\}$.

  \textbf{חיתוך עם הצירים:} $f(0)=\frac{-1}{1}=-1$, כלומר $(0,-1)$ על ציר $y$. $f(x)=0\iff 2x-1=0\iff x=\frac12$, כלומר $(\frac12,0)$ על ציר $x$.

  \textbf{זוגיות:} תחום ההגדרה אינו סימטרי ביחס ל-$0$ ($-1\in D$ אבל $1\notin D$), ולכן הפונקציה אינה זוגית ואינה אי-זוגית.

  \item $f$ מנה של פולינומים, ולכן רציפה בכל תחום הגדרתה. ב-$x=1$: המונה שואף ל-$1>0$ והמכנה $(x-1)^2\to0^+$, ולכן $\lim_{x\to1^\pm}f(x)=+\infty$ --- נקודת אי-רציפות מסוג שני (אינסופית).

  \item \textbf{אנכית:} $x=1$ (משני הצדדים $f\to+\infty$).
  \textbf{אופקית:} $\lim_{x\to\pm\infty}\frac{2x-1}{(x-1)^2}=\lim_{x\to\pm\infty}\frac{\frac2x-\frac1{x^2}}{(1-\frac1x)^2}=0$, ולכן $y=0$ אסימפטוטה אופקית ב-$\pm\infty$ (ולכן אין אסימפטוטה משופעת).

  \item לפי כלל המנה:
  \[ f'(x)=\frac{2(x-1)^2-(2x-1)\cdot2(x-1)}{(x-1)^4}=\frac{2(x-1)-2(2x-1)}{(x-1)^3}=\frac{-2x}{(x-1)^3}. \]
  $f'(x)=0\iff x=0$. סימנים:
  \begin{itemize}
    \item $x<0$: מונה $-2x>0$, מכנה $<0$, ולכן $f'<0$ --- יורדת.
    \item $0<x<1$: מונה $<0$, מכנה $<0$, ולכן $f'>0$ --- עולה.
    \item $x>1$: מונה $<0$, מכנה $>0$, ולכן $f'<0$ --- יורדת.
  \end{itemize}
  לכן $f$ יורדת ב-$(-\infty,0)$ וב-$(1,\infty)$ ועולה ב-$(0,1)$. ב-$x=0$ הנגזרת עוברת משלילית לחיובית: \textbf{מינימום מקומי} $(0,-1)$. (זה גם מינימום מוחלט: עבור $x<\frac12$ הפונקציה יורדת ל-$-1$ ואז עולה, ועבור $x\ge\frac12$, $f\ge0$.)

  \item נגזור את $f'(x)=-2x(x-1)^{-3}$:
  \[ f''(x)=-2(x-1)^{-3}+6x(x-1)^{-4}=\frac{-2(x-1)+6x}{(x-1)^4}=\frac{2(2x+1)}{(x-1)^4}. \]
  המכנה חיובי לכל $x\neq1$, ולכן הסימן הוא של $2x+1$:
  $f''<0$ עבור $x<-\frac12$ --- $f$ קעורה ($\cap$) ב-$(-\infty,-\frac12)$;
  $f''>0$ עבור $-\frac12<x<1$ ו-$x>1$ --- $f$ קמורה ($\cup$) ב-$(-\frac12,1)$ וב-$(1,\infty)$.
  ב-$x=-\frac12$ הסימן מתחלף והפונקציה רציפה, ולכן זו \textbf{נקודת פיתול}: $f(-\frac12)=\frac{-2}{9/4}=-\frac89$, כלומר $(-\frac12,-\frac89)$.

  \item \textbf{תיאור גרפי:} משמאל הגרף מתחיל מתחת לציר $x$, קרוב לאסימפטוטה $y=0$ כש-$x\to-\infty$, יורד (קעור) עד נקודת הפיתול $(-\frac12,-\frac89)$, ממשיך לרדת (קמור) עד המינימום $(0,-1)$, עולה וחותך את ציר $x$ ב-$(\frac12,0)$ ושואף ל-$+\infty$ כש-$x\to1^-$. מימין ל-$x=1$ הוא יורד מ-$+\infty$ ומתקרב מלמעלה ל-$y=0$ כש-$x\to\infty$ (קמור).
\end{enumerate}
\end{solution}

\subsection*{חלק א', שאלה 2}

\begin{question}
נתונה עקומה בצורה פולארית: $r=a(1+\cos\varphi)$, כאשר $a$ פרמטר ממשי וחיובי, $0\le\varphi\le2\pi$.
חשב את \textbf{אורך} העקומה. \textbf{חשבו את השטח} החסום על ידי העקומה
\end{question}

\begin{hint}
אורך עקומה פולרית: $L=\int_\alpha^\beta\sqrt{r^2+(r')^2}\,d\varphi$. שטח: $S=\frac12\int_\alpha^\beta r^2\,d\varphi$.
\end{hint}

\begin{hint}
השתמשו בזהות $1+\cos\varphi=2\cos^2\frac{\varphi}{2}$, ושימו לב לסימן של $\cos\frac\varphi2$ בקטע.
\end{hint}

\begin{solution}
העקומה (קרדיואידה) היא $r(\varphi)=a(1+\cos\varphi)\ge0$, ו-$r'(\varphi)=-a\sin\varphi$.

\textbf{אורך:} לפי נוסחת אורך עקומה בקואורדינטות פולריות,
\[ r^2+(r')^2=a^2\big(1+2\cos\varphi+\cos^2\varphi+\sin^2\varphi\big)=2a^2(1+\cos\varphi)=4a^2\cos^2\frac{\varphi}{2}. \]
לכן
\[ L=\int_0^{2\pi}\sqrt{r^2+(r')^2}\,d\varphi=\int_0^{2\pi}2a\Big|\cos\frac{\varphi}{2}\Big|\,d\varphi . \]
עבור $0\le\varphi\le\pi$, $\cos\frac\varphi2\ge0$, ועבור $\pi\le\varphi\le2\pi$, $\cos\frac\varphi2\le0$. מסימטריה (או ישירות):
\[ L=2\int_0^{\pi}2a\cos\frac{\varphi}{2}\,d\varphi=4a\Big[2\sin\frac\varphi2\Big]_0^{\pi}=8a . \]
כלומר $\boxed{L=8a}$.

\textbf{שטח:} לפי נוסחת השטח בקואורדינטות פולריות,
\[ S=\frac12\int_0^{2\pi}a^2(1+\cos\varphi)^2d\varphi=\frac{a^2}{2}\int_0^{2\pi}\Big(1+2\cos\varphi+\frac{1+\cos2\varphi}{2}\Big)d\varphi
=\frac{a^2}{2}\Big[\frac32\varphi+2\sin\varphi+\frac{\sin2\varphi}{4}\Big]_0^{2\pi}=\frac{a^2}{2}\cdot3\pi . \]
כלומר $\boxed{S=\frac{3\pi a^2}{2}}$.
\end{solution}

\subsection*{חלק א', שאלה 3}

\begin{question}
חשב את הגבול הבא:
\[ \lim_{x\to0^+}\left((x^2+x+1))^{\frac{1}{x^2}}+x\cdot\ln x\right) \]
\end{question}

\begin{hint}
חשבו כל מחובר בנפרד. במחובר הראשון כתבו $(x^2+x+1)^{1/x^2}=e^{\frac{\ln(1+x+x^2)}{x^2}}$.
\end{hint}

\begin{solution}
(בנוסח המקורי מופיעה סוגר מיותרת; המחובר הראשון הוא $(x^2+x+1)^{1/x^2}$.)

\textbf{המחובר השני:} $x\ln x=\frac{\ln x}{1/x}$, מהצורה $\frac{-\infty}{\infty}$. לפי כלל לופיטל:
\[ \lim_{x\to0^+}\frac{\ln x}{1/x}=\lim_{x\to0^+}\frac{1/x}{-1/x^2}=\lim_{x\to0^+}(-x)=0 . \]

\textbf{המחובר הראשון:} עבור $x>0$, $x^2+x+1>1$, ולכן
\[ (x^2+x+1)^{\frac{1}{x^2}}=\exp\Big(\frac{\ln(1+x+x^2)}{x^2}\Big). \]
המעריך: $\frac{\ln(1+x+x^2)}{x^2}=\frac{\ln(1+x+x^2)}{x+x^2}\cdot\frac{x+x^2}{x^2}=\frac{\ln(1+t)}{t}\Big|_{t=x+x^2}\cdot\Big(\frac1x+1\Big)$.
כאשר $x\to0^+$: $t=x+x^2\to0^+$ ולכן $\frac{\ln(1+t)}{t}\to1$, ואילו $\frac1x+1\to+\infty$. לכן המעריך שואף ל-$+\infty$, ומכאן
\[ \lim_{x\to0^+}(x^2+x+1)^{\frac{1}{x^2}}=+\infty . \]
(אינטואיטיבית: זה גבול מהצורה $1^\infty$, אבל הבסיס מתקרב ל-$1$ "לאט" --- כמו $1+x$ --- בעוד שהמעריך גדל כמו $\frac{1}{x^2}$, ולכן $\big(1+x\big)^{1/x^2}\approx e^{1/x}\to\infty$.)

\textbf{סיכום:} סכום של פונקציה השואפת ל-$+\infty$ ופונקציה השואפת לגבול סופי ($0$) שואף ל-$+\infty$. לכן
\[ \lim_{x\to0^+}\left((x^2+x+1)^{\frac{1}{x^2}}+x\ln x\right)=\boxed{+\infty} . \]
\end{solution}

\subsection*{חלק ב', שאלה 1}

\begin{question}
נתונה $h(x)=\begin{cases}x^2+a & x<-1\\ x^3-8 & x\ge-1\end{cases}$. לאיזה ערך של הפרמטר $a$ הפונקציה היא רציפה לכל $x$ ממשי,
\begin{enumerate}
  \item $-1$
  \item $-9$
  \item $-10$
  \item $-8$
  \item $-2$
\end{enumerate}
\end{question}

\begin{solution}
\textbf{התשובה הנכונה: ג ($a=-10$).}

בכל אחד מהקטעים $x<-1$ ו-$x>-1$ הפונקציה פולינום ולכן רציפה. נותר לבדוק את $x=-1$:
\[ h(-1)=\lim_{x\to-1^+}h(x)=(-1)^3-8=-9,\qquad \lim_{x\to-1^-}h(x)=(-1)^2+a=1+a . \]
רציפות ב-$-1$ מתקיימת אם ורק אם $1+a=-9$, כלומר $a=-10$.
\end{solution}

\subsection*{חלק ב', שאלה 2}

\begin{question}
תהי $f(x)$ פונקציה גזירה בקטע פתוח $(a,b)$ וידוע כי נגזרתה חיובית, כלומר $f'(x)>0$. מה מהבאים חייב להיות נכון?
\begin{enumerate}
  \item $f(x)$ פונקציה חסומה בקטע פתוח $(a,b)$.
  \item $f(x)$ פונקציה המקיימת $f(x)>0$ בקטע פתוח $(a,b)$.
  \item $f(x)$ פונקציה רציפה בקטע סגור $[a,b]$.
  \item $f(x)$ פונקציה מונוטונית עולה בקטע פתוח $(a,b)$.
  \item $f(x)$ פונקציה מונוטונית עולה בקטע סגור $[a,b]$.
\end{enumerate}
\end{question}

\begin{hint}
מה אומר משפט לגרנז' על $f(x_2)-f(x_1)$ כאשר $a<x_1<x_2<b$?
\end{hint}

\begin{solution}
\textbf{התשובה הנכונה: ד.}

יהיו $a<x_1<x_2<b$. $f$ גזירה ולכן רציפה ב-$[x_1,x_2]$ וגזירה ב-$(x_1,x_2)$, ולפי משפט לגרנז' קיימת $c\in(x_1,x_2)$ כך ש-
\[ f(x_2)-f(x_1)=f'(c)(x_2-x_1)>0 . \]
לכן $f$ עולה (ממש) ב-$(a,b)$.

דוגמאות נגדיות לשאר הטענות, למשל בקטע $(0,1)$:
\begin{itemize}
  \item $f(x)=-\frac1x$: $f'(x)=\frac1{x^2}>0$, אבל $f$ אינה חסומה ב-$(0,1)$ (א' נופלת), אינה מוגדרת ב-$0$ ובפרט אינה רציפה ב-$[0,1]$ (ג' נופלת) ואינה מונוטונית בקטע הסגור (ה' נופלת, $f$ כלל לא מוגדרת שם).
  \item $f(x)=x-5$: $f'=1>0$ אבל $f<0$ בכל $(0,1)$ (ב' נופלת).
\end{itemize}
\end{solution}

\subsection*{חלק ב', שאלה 3}

\begin{question}
נתונה פונקציה $F(x)=\int_0^x\sqrt{t^2+1}\,dt$. חשבו את $F'(\sqrt2)$.
\begin{enumerate}
  \item $\sqrt3$
  \item $\frac{1}{\sqrt3}$
  \item $\sqrt2$
  \item $3$
  \item $\sqrt2+1$
\end{enumerate}
\end{question}

\begin{solution}
\textbf{התשובה הנכונה: א.}

הפונקציה $t\mapsto\sqrt{t^2+1}$ רציפה בכל $\R$, ולכן לפי המשפט היסודי של החשבון האינטגרלי $F$ גזירה ו-$F'(x)=\sqrt{x^2+1}$. לכן
\[ F'(\sqrt2)=\sqrt{2+1}=\sqrt3 . \]
\end{solution}

\subsection*{חלק ב', שאלה 4}

\begin{question}
מצאו את הערך הקטן ביותר בתחום הגדרה של הפונקציה: $f(x)=x^4+4x$
\begin{enumerate}
  \item $-4$
  \item $1$
  \item $-3$
  \item $-5$
  \item $-2$
\end{enumerate}
\end{question}

\begin{solution}
\textbf{התשובה הנכונה: ג ($-3$).}

$f$ פולינום, מוגדר וגזיר בכל $\R$. $f'(x)=4x^3+4=4(x+1)(x^2-x+1)$, והגורם $x^2-x+1$ חיובי תמיד (דיסקרימיננטה שלילית). לכן $f'<0$ עבור $x<-1$ ו-$f'>0$ עבור $x>-1$: $f$ יורדת ב-$(-\infty,-1]$ ועולה ב-$[-1,\infty)$. מכאן ש-$x=-1$ היא נקודת מינימום מוחלט, והערך הקטן ביותר הוא
\[ f(-1)=1-4=-3 . \]
\end{solution}

\subsection*{חלק ב', שאלה 5}

\begin{question}
תהי $f(x)$ פונקציה גזירה \textbf{פעמיים} בקטע פתוח $(a,b)$. מה מהבאים חייב להיות נכון?
\begin{enumerate}
  \item $f'(x)$ רציפה בקטע פתוח $(a,b)$.
  \item $f(x)$ בעלת נקודת פיתול בקטע פתוח $(a,b)$.
  \item $f(x)$ פונקציה מונוטונית בקטע פתוח $(a,b)$.
  \item $f''(x)$ פונקציה רציפה בקטע פתוח $(a,b)$.
  \item $f(x)$ פונקציה חסומה בקטע פתוח $(a,b)$.
\end{enumerate}
\end{question}

\begin{hint}
פונקציה גזירה בנקודה היא רציפה בה. הפעילו זאת על $f'$.
\end{hint}

\begin{solution}
\textbf{התשובה הנכונה: א.}

$f$ גזירה פעמיים ב-$(a,b)$, כלומר $f'$ גזירה בכל נקודה של $(a,b)$. פונקציה גזירה בנקודה רציפה בה, ולכן $f'$ רציפה ב-$(a,b)$.

דוגמאות נגדיות לשאר הטענות, בקטע $(0,1)$ (אלא אם צוין אחרת):
\begin{itemize}
  \item $f(x)=x^2$ (ב-$(0,1)$): $f''=2>0$, אין נקודת פיתול (ב' נופלת).
  \item $f(x)=\left(x-\frac12\right)^2$: לא מונוטונית ב-$(0,1)$ (ג' נופלת).
  \item $f(x)=x^4\sin\frac1x$ ($f(0)=0$) בקטע $(-1,1)$: גזירה פעמיים, אבל $f''(x)=12x^2\sin\frac1x-6x\cos\frac1x-\sin\frac1x$ עבור $x\neq0$ ו-$f''(0)=0$, ו-$f''$ אינה רציפה ב-$0$ (ד' נופלת).
  \item $f(x)=\frac1x$ ב-$(0,1)$: גזירה פעמיים אך אינה חסומה (ה' נופלת).
\end{itemize}
\end{solution}

\section*{תש"פ סמסטר ב מועד ב}

\subsection*{שאלה 1א}

\begin{question}
חשבו את הנגזרת של הפונקציה $f(x)=\sin(2x+1)$ בנקודה $x$ לפי הגדרת הנגזרת (תשובה לפי טבלת הנגזרות לא תתקבל).
\end{question}

\begin{hint}
השתמשו בזהות $\sin\alpha-\sin\beta=2\sin\frac{\alpha-\beta}{2}\cos\frac{\alpha+\beta}{2}$ ובגבול $\lim_{t\to0}\frac{\sin t}{t}=1$.
\end{hint}

\begin{solution}
לפי ההגדרה, $f'(x)=\lim_{h\to0}\frac{f(x+h)-f(x)}{h}$. לפי הזהות לסכום והפרש פונקציות,
\[ \sin(2x+2h+1)-\sin(2x+1)=2\sin(h)\cos(2x+1+h). \]
לכן
\[ f'(x)=\lim_{h\to0}\frac{2\sin h\cos(2x+1+h)}{h}=2\lim_{h\to0}\frac{\sin h}{h}\cdot\lim_{h\to0}\cos(2x+1+h)=2\cdot1\cdot\cos(2x+1), \]
כאשר השתמשנו בגבול היסודי $\frac{\sin h}{h}\to1$ וברציפות $\cos$. כלומר $\boxed{f'(x)=2\cos(2x+1)}$.
\end{solution}

\subsection*{שאלה 1ב}

\begin{question}
חשבו את הגבולות הבאים:
\begin{enumerate}
  \item (5 נק') $\displaystyle\lim_{x\to\infty}\frac{x+2\sin x}{3x-\cos x}$,
  \item (5 נק') $\displaystyle\lim_{x\to+\infty}\left(\sqrt{x+\sqrt{x}}-\sqrt{x-\sqrt{x}}\right)$
\end{enumerate}
\end{question}

\begin{hint}
בגבול הראשון חלקו מונה ומכנה ב-$x$ והשתמשו בכך ש-$\sin x,\cos x$ חסומות.
\end{hint}

\begin{hint}
בגבול השני הכפילו וחלקו בצמוד.
\end{hint}

\begin{solution}
\begin{enumerate}
  \item נחלק מונה ומכנה ב-$x$ (עבור $x>0$):
  \[ \frac{x+2\sin x}{3x-\cos x}=\frac{1+2\frac{\sin x}{x}}{3-\frac{\cos x}{x}} . \]
  מאחר ש-$|\sin x|,|\cos x|\le1$ ו-$\frac1x\to0$, לפי "חסומה כפול שואפת לאפס" $\frac{\sin x}{x}\to0$ ו-$\frac{\cos x}{x}\to0$. לכן הגבול הוא $\boxed{\frac13}$.
  \item נכפיל ונחלק בצמוד (עבור $x>1$ שני השורשים מוגדרים):
  \[ \sqrt{x+\sqrt x}-\sqrt{x-\sqrt x}=\frac{(x+\sqrt x)-(x-\sqrt x)}{\sqrt{x+\sqrt x}+\sqrt{x-\sqrt x}}=\frac{2\sqrt x}{\sqrt{x+\sqrt x}+\sqrt{x-\sqrt x}}
  =\frac{2}{\sqrt{1+\frac{1}{\sqrt x}}+\sqrt{1-\frac{1}{\sqrt x}}} , \]
  כאשר חילקנו מונה ומכנה ב-$\sqrt x$. כאשר $x\to+\infty$, $\frac{1}{\sqrt x}\to0$, ולכן (רציפות השורש) הגבול הוא $\frac{2}{1+1}=\boxed{1}$.
\end{enumerate}
\end{solution}

\subsection*{שאלה 1ג}

\begin{question}
לאילו ערכים של הפרמטר $a$ לגרף של פונקציה $f(x)=\frac{x^2-9}{x+a}$ אין אסימפטוטה אנכית?
\end{question}

\begin{hint}
מתי השורש של המכנה הוא גם שורש של המונה?
\end{hint}

\begin{solution}
הפונקציה רציפה בכל נקודה שבה $x\neq -a$, ולכן אסימפטוטה אנכית יכולה להיות רק ב-$x=-a$. ב-$x\to-a$ המכנה שואף ל-$0$, ולכן:
\begin{itemize}
  \item אם המונה אינו מתאפס ב-$-a$, כלומר $a^2-9\neq0$, אז $|f(x)|\to\infty$ כאשר $x\to-a$, ו-$x=-a$ היא אסימפטוטה אנכית.
  \item אם $a^2=9$, כלומר $a=3$ או $a=-3$: עבור $a=3$, $f(x)=\frac{(x-3)(x+3)}{x+3}=x-3$ לכל $x\neq-3$, ולכן $\lim_{x\to-3}f(x)=-6$ סופי (נקודת אי-רציפות סליקה), ואין אסימפטוטה אנכית. באופן דומה עבור $a=-3$, $f(x)=x+3$ לכל $x\neq3$ ו-$\lim_{x\to3}f(x)=6$.
\end{itemize}
לכן אין אסימפטוטה אנכית בדיוק כאשר $\boxed{a=3 \text{ או } a=-3}$.
\end{solution}

\subsection*{שאלה 2א}

\begin{question}
עבור אילו ערכים של הפרמטר $a$ למשוואה $(x-1)^2e^x-a=0$ יש פתרון ממשי אחד?
\end{question}

\begin{hint}
חקרו את הפונקציה $g(x)=(x-1)^2e^x$ (תחומי מונוטוניות, קיצון, גבולות ב-$\pm\infty$), וספרו כמה פעמים הישר האופקי $y=a$ חותך את הגרף.
\end{hint}

\begin{solution}
המשוואה שקולה ל-$g(x)=a$ עבור $g(x)=(x-1)^2e^x$. נחקור את $g$, שהיא גזירה בכל $\R$:
\[ g'(x)=2(x-1)e^x+(x-1)^2e^x=e^x(x-1)(x+1). \]
$e^x>0$, ולכן $g'>0$ ב-$(-\infty,-1)$, $g'<0$ ב-$(-1,1)$ ו-$g'>0$ ב-$(1,\infty)$. כלומר $g$ עולה ממש ב-$(-\infty,-1]$, יורדת ממש ב-$[-1,1]$ ועולה ממש ב-$[1,\infty)$, עם מקסימום מקומי $g(-1)=\frac{4}{e}$ ומינימום מקומי $g(1)=0$.

גבולות: $\lim_{x\to\infty}g(x)=\infty$; וכאשר $x\to-\infty$, $(x-1)^2e^x=\frac{(x-1)^2}{e^{-x}}\to0$ (האקספוננט גובר על פולינום, למשל לפי לופיטל פעמיים). כמו כן $g(x)\ge0$ לכל $x$, עם שוויון רק ב-$x=1$.

לפי משפט ערך הביניים ומונוטוניות ממש בכל קטע, בכל קטע מונוטוניות כל ערך בטווח של הקטע מתקבל בדיוק פעם אחת. טווחי הקטעים: $(-\infty,-1]\mapsto(0,\frac4e]$, $[-1,1]\mapsto[0,\frac4e]$, $[1,\infty)\mapsto[0,\infty)$. לכן מספר הפתרונות:
\begin{itemize}
  \item $a<0$: אין פתרונות ($g\ge0$).
  \item $a=0$: פתרון יחיד, $x=1$.
  \item $0<a<\frac4e$: שלושה פתרונות (אחד בכל קטע מונוטוניות).
  \item $a=\frac4e$: שני פתרונות ($x=-1$, ועוד אחד ב-$(1,\infty)$).
  \item $a>\frac4e$: פתרון יחיד (ב-$(1,\infty)$; בשני הקטעים האחרים $g\le\frac4e<a$).
\end{itemize}
\textbf{תשובה:} למשוואה פתרון ממשי יחיד עבור $\boxed{a=0 \text{ או } a>\frac{4}{e}}$.
\end{solution}

\subsection*{שאלה 2ב}

\begin{question}
מצאו את נפח הגוף הנוצר ע"י סיבוב סביב ציר ה-$x$ של הצורה החסומה על ידי הקווים $y=xe^{-x}$, $y=0$, $x=2$. שרטט את הסקיצה המתאימה.
\end{question}

\begin{hint}
הגבולות של האינטגרל הם $x=0$ (שם $xe^{-x}=0$) ו-$x=2$. את $\int x^2e^{-2x}dx$ מחשבים באינטגרציה בחלקים פעמיים.
\end{hint}

\begin{solution}
\textbf{סקיצה:} $y=xe^{-x}$ עוברת בראשית, חיובית עבור $x>0$, עולה עד למקסימום $(1,\frac1e)$ ($y'=(1-x)e^{-x}$) ואחר כך יורדת, וב-$x=2$ ערכה $2e^{-2}\approx0.27$. הצורה החסומה היא האזור שבין הגרף לבין ציר $x$, עבור $0\le x\le 2$ (הגרף חותך את $y=0$ ב-$x=0$).

\textbf{הנפח:} לפי נוסחת נפח גוף סיבוב סביב ציר $x$,
\[ V=\pi\int_0^2\big(xe^{-x}\big)^2dx=\pi\int_0^2x^2e^{-2x}dx . \]
נחשב באינטגרציה בחלקים פעמיים:
\[ \int x^2e^{-2x}dx=-\frac{x^2}{2}e^{-2x}+\int xe^{-2x}dx=-\frac{x^2}{2}e^{-2x}-\frac{x}{2}e^{-2x}+\frac12\int e^{-2x}dx
=-e^{-2x}\Big(\frac{x^2}{2}+\frac{x}{2}+\frac14\Big)+C . \]
לכן
\[ V=\pi\Big[-e^{-2x}\Big(\frac{x^2}{2}+\frac{x}{2}+\frac14\Big)\Big]_0^2=\pi\Big(-\frac{13}{4}e^{-4}+\frac14\Big)=\boxed{\frac{\pi}{4}\big(1-13e^{-4}\big)}\approx0.4357 . \]
\end{solution}

\subsection*{שאלה 3א}

\begin{question}
חשבו:
\begin{enumerate}
  \item (7 נק') $\displaystyle\int\frac{5x+2}{x^4+x^3+2x^2}\,dx$
  \item (8 נק') $\displaystyle\int_0^1\ln(1+\sqrt{x})\,dx$
\end{enumerate}
\end{question}

\begin{hint}
באינטגרל הראשון פרקו לשברים חלקיים: $x^4+x^3+2x^2=x^2(x^2+x+2)$, והגורם הריבועי אי-פריק.
\end{hint}

\begin{hint}
באינטגרל השני הציבו $t=\sqrt x$ ואז בצעו אינטגרציה בחלקים.
\end{hint}

\begin{solution}
\begin{enumerate}
  \item המכנה: $x^4+x^3+2x^2=x^2(x^2+x+2)$, ו-$x^2+x+2$ אי-פריק (הדיסקרימיננטה $1-8<0$). נפרק לשברים חלקיים:
  \[ \frac{5x+2}{x^2(x^2+x+2)}=\frac{A}{x}+\frac{B}{x^2}+\frac{Cx+D}{x^2+x+2}, \]
  כלומר $5x+2=Ax(x^2+x+2)+B(x^2+x+2)+(Cx+D)x^2$. השוואת מקדמים: מקדם חופשי $2B=2\Rightarrow B=1$; מקדם $x$: $2A+B=5\Rightarrow A=2$; מקדם $x^3$: $A+C=0\Rightarrow C=-2$; מקדם $x^2$: $A+B+D=0\Rightarrow D=-3$. לכן
  \[ \frac{5x+2}{x^4+x^3+2x^2}=\frac{2}{x}+\frac{1}{x^2}-\frac{2x+3}{x^2+x+2}. \]
  את המחובר האחרון נפרק: $2x+3=(2x+1)+2$, ו-$x^2+x+2=\big(x+\frac12\big)^2+\frac74$. לכן
  \[ \int\frac{2x+3}{x^2+x+2}dx=\ln(x^2+x+2)+2\int\frac{dx}{(x+\frac12)^2+\frac74}=\ln(x^2+x+2)+\frac{4}{\sqrt7}\arctan\frac{2x+1}{\sqrt7}. \]
  \textbf{תשובה:}
  \[ \int\frac{5x+2}{x^4+x^3+2x^2}dx=2\ln|x|-\frac1x-\ln(x^2+x+2)-\frac{4}{\sqrt7}\arctan\frac{2x+1}{\sqrt7}+C . \]
  \item נציב $t=\sqrt x$, כלומר $x=t^2$, $dx=2t\,dt$, והגבולות $0\to0$, $1\to1$:
  \[ \int_0^1\ln(1+\sqrt x)\,dx=\int_0^12t\ln(1+t)\,dt . \]
  אינטגרציה בחלקים עם $u=\ln(1+t)$, $v'=2t$ ($v=t^2$):
  \[ =\Big[t^2\ln(1+t)\Big]_0^1-\int_0^1\frac{t^2}{1+t}dt=\ln2-\int_0^1\Big(t-1+\frac{1}{1+t}\Big)dt=\ln2-\Big(\frac12-1+\ln2\Big)=\boxed{\frac12} . \]
  (השתמשנו בחילוק פולינומים $\frac{t^2}{1+t}=t-1+\frac{1}{1+t}$.)
\end{enumerate}
\end{solution}

\subsection*{שאלה 3ב}

\begin{question}
רשמו את פולינום מקלורן מסדר $n=3$ עבור הפונקציה:
\[ f(x)=\sin(\sin(x)) \]
\end{question}

\begin{hint}
הרכיבו את הפיתוח של $\sin$ בתוך עצמו.
\end{hint}

\begin{solution}
לפי פיתוח מקלורן, $\sin u=u-\frac{u^3}{6}+o(u^3)$ ו-$\sin x=x-\frac{x^3}{6}+o(x^3)$. נציב $u=\sin x$; מאחר ש-$u=x+o(x)$, מתקיים $u^3=x^3+o(x^3)$ ו-$o(u^3)=o(x^3)$. לכן
\[ \sin(\sin x)=\Big(x-\frac{x^3}{6}\Big)-\frac{x^3}{6}+o(x^3)=x-\frac{x^3}{3}+o(x^3). \]
מיחידות פולינום טיילור, פולינום מקלורן מסדר 3 הוא
\[ \boxed{P_3(x)=x-\frac{x^3}{3}} . \]
(בדיקה ישירה: $f'(x)=\cos(\sin x)\cos x$, $f'(0)=1$; $f$ אי-זוגית ולכן $f(0)=f''(0)=0$; וחישוב נותן $f'''(0)=-2$, ואכן $\frac{-2}{3!}=-\frac13$.)
\end{solution}

\subsection*{שאלה 4א}

\begin{question}
לאילו ערכים של הפרמטרים $a$ ו-$b$ הפונקציה
\[ f(x)=\begin{cases} \dfrac{\sin(3x)}{\ln(1-x)} & x<0\\[2mm] ax+b & 0\le x\le1\\[2mm] (\cos(x-1))^{1/(x-1)} & 1<x<\dfrac{\pi}{2}+1\end{cases} \]
תהיה רציפה לכל $x$ ממשי קטן מ-$\frac{\pi}{2}+1$?
\end{question}

\begin{hint}
דרשו שהגבולות החד-צדדיים בנקודות התפר $x=0$ ו-$x=1$ יהיו שווים לערך הפונקציה. בגבול ב-$x=1^+$ (מהצורה $1^\infty$) עברו ל-$e^{\ln(\cdot)}$.
\end{hint}

\begin{solution}
בכל אחד מהקטעים הפתוחים $(-\infty,0)$, $(0,1)$, $(1,\frac\pi2+1)$ הפונקציה רציפה כהרכבה ומנה של פונקציות אלמנטריות: עבור $x<0$, $1-x>1$ ולכן $\ln(1-x)>0$ (המכנה לא מתאפס); עבור $1<x<\frac\pi2+1$, $0<x-1<\frac\pi2$ ולכן $\cos(x-1)>0$ והחזקה מוגדרת. נותר לבדוק את נקודות התפר.

\textbf{ב-$x=0$:} $f(0)=b=\lim_{x\to0^+}f(x)$. משמאל, בעזרת $\frac{\sin 3x}{3x}\to1$ ו-$\frac{\ln(1-x)}{-x}\to1$:
\[ \lim_{x\to0^-}\frac{\sin3x}{\ln(1-x)}=\lim_{x\to0^-}\frac{\sin3x}{3x}\cdot\frac{-x}{\ln(1-x)}\cdot\frac{3x}{-x}=1\cdot1\cdot(-3)=-3 . \]
לכן רציפות ב-$0$ דורשת $b=-3$.

\textbf{ב-$x=1$:} $f(1)=a+b=\lim_{x\to1^-}f(x)$. מימין נסמן $t=x-1\to0^+$:
\[ (\cos t)^{1/t}=\exp\Big(\frac{\ln\cos t}{t}\Big),\qquad \lim_{t\to0^+}\frac{\ln\cos t}{t}\overset{\text{לופיטל}}{=}\lim_{t\to0^+}\frac{-\tan t}{1}=0 , \]
(הגבול מהצורה $\frac00$, והנגזרות קיימות), ולכן מרציפות האקספוננט $\lim_{x\to1^+}f(x)=e^0=1$. רציפות ב-$1$ דורשת $a+b=1$, כלומר $a=4$.

\textbf{תשובה:} $\boxed{a=4,\ b=-3}$.
\end{solution}

\subsection*{שאלה 4ב}

\begin{question}
מצאו את משוואת המשיק עבור הפונקציה $y(x)$ הנתונה בצורה סתומה $e^{x\cdot y}+y^3=9$ בנקודה $(0,2)$.
\end{question}

\begin{hint}
גזרו את שני אגפי המשוואה לפי $x$, כאשר $y=y(x)$, והציבו את הנקודה.
\end{hint}

\begin{solution}
הנקודה על העקומה: $e^{0}+2^3=1+8=9$. נגזור את המשוואה לפי $x$ (גזירה סתומה, $y=y(x)$):
\[ e^{xy}\,(y+xy')+3y^2y'=0 . \]
נציב $x=0$, $y=2$: $1\cdot(2+0)+12y'=0$, ולכן $y'(0)=-\frac16$. משוואת המשיק:
\[ \boxed{y=2-\frac{x}{6}} . \]
\end{solution}

\subsection*{שאלה 4ג}

\begin{question}
חשבו $\displaystyle\int_{0.25}^{+\infty}\frac{dx}{(4x+1)\sqrt{x}}$.
\end{question}

\begin{hint}
הציבו $t=\sqrt x$.
\end{hint}

\begin{solution}
זהו אינטגרל לא אמיתי (גבול עליון אינסופי). נציב $t=\sqrt x$, $x=t^2$, $dx=2t\,dt$; כש-$x=0.25$, $t=\frac12$, וכש-$x\to\infty$, $t\to\infty$:
\[ \int_{0.25}^{R}\frac{dx}{(4x+1)\sqrt x}=\int_{1/2}^{\sqrt R}\frac{2t\,dt}{(4t^2+1)t}=\int_{1/2}^{\sqrt R}\frac{2\,dt}{1+(2t)^2}=\Big[\arctan(2t)\Big]_{1/2}^{\sqrt R}=\arctan(2\sqrt R)-\frac{\pi}{4}. \]
כאשר $R\to\infty$, $\arctan(2\sqrt R)\to\frac\pi2$. לכן האינטגרל מתכנס ו-
\[ \int_{0.25}^{+\infty}\frac{dx}{(4x+1)\sqrt{x}}=\frac\pi2-\frac\pi4=\boxed{\frac{\pi}{4}} . \]
\end{solution}

\section*{תשפ"ב סמסטר א מועד א}

\subsection*{שאלה 1}

\begin{question}
נגדיר פונקציה על ידי
\[ f(x)=\begin{cases} x^3 \sin(\frac{5}{x}) & x\neq 0\\ 0 & x=0 \end{cases} \]
\begin{enumerate}
  \item (6 נק') האם $f$ רציפה בנקודה $x=0$ ? נמקו.
  \item (6 נק') האם $f$ גזירה בנקודה $x=0$ ? נמקו.
  \item (5 נק') האם הנגזרת $f'(x)$ רציפה בנקודה $x=0$ ? נמקו.
\end{enumerate}
\end{question}

\begin{hint}
"אפסה כפול חסומה": $\left|\sin(\frac{5}{x})\right|\le 1$ לכל $x\neq 0$.
\end{hint}

\begin{hint}
בסעיף ב' חשבו את הנגזרת ב-$0$ לפי ההגדרה (גבול מנת ההפרשים). בסעיף ג' חשבו את $f'(x)$ עבור $x\neq0$ לפי כללי הגזירה ובדקו את הגבול כאשר $x\to 0$.
\end{hint}

\begin{solution}
\begin{enumerate}
  \item \textbf{כן.} לכל $x\neq 0$ מתקיים
  \[ 0\le |f(x)| = |x|^3\left|\sin\left(\tfrac{5}{x}\right)\right| \le |x|^3 \xrightarrow[x\to0]{} 0, \]
  ולכן לפי משפט הסנדוויץ' $\lim_{x\to0} f(x)=0=f(0)$. כלומר $f$ רציפה ב-$0$.
  \item \textbf{כן.} לפי הגדרת הנגזרת:
  \[ f'(0)=\lim_{h\to0}\frac{f(h)-f(0)}{h}=\lim_{h\to0}\frac{h^3\sin(5/h)}{h}=\lim_{h\to0} h^2\sin\left(\tfrac{5}{h}\right). \]
  כאן $h^2\to0$ ו-$\sin(5/h)$ חסומה, ולכן (אפסה כפול חסומה) הגבול שווה $0$. לכן $f$ גזירה ב-$0$ ו-$f'(0)=0$.
  \item \textbf{כן.} עבור $x\neq0$, לפי כלל המכפלה וכלל השרשרת:
  \[ f'(x)=3x^2\sin\left(\tfrac{5}{x}\right)+x^3\cos\left(\tfrac{5}{x}\right)\cdot\left(-\tfrac{5}{x^2}\right)=3x^2\sin\left(\tfrac{5}{x}\right)-5x\cos\left(\tfrac{5}{x}\right). \]
  שני המחוברים הם מהצורה "אפסה כפול חסומה" ($3x^2\to0$, $5x\to0$ ו-$|\sin|,|\cos|\le1$), ולכן
  \[ \lim_{x\to0}f'(x)=0=f'(0). \]
  מכאן ש-$f'$ רציפה בנקודה $x=0$.
\end{enumerate}
\end{solution}

\subsection*{שאלה 2א}

\begin{question}
מצאו את כל האסימפטוטות של הפונקציה
\[ f(x)=\frac{2x^2-3x+1}{x} \]
\end{question}

\begin{hint}
חלקו את המונה ב-$x$: $f(x)=2x-3+\frac1x$.
\end{hint}

\begin{solution}
תחום ההגדרה: $x\neq0$. לכל $x\neq 0$:
\[ f(x)=2x-3+\frac{1}{x}. \]
\textbf{אסימפטוטה אנכית:} המונה ב-$x=0$ שווה $1\neq0$, ולכן
\[ \lim_{x\to0^+}f(x)=+\infty,\qquad \lim_{x\to0^-}f(x)=-\infty, \]
כלומר $x=0$ אסימפטוטה אנכית. זו הנקודה היחידה מחוץ לתחום, ו-$f$ רציפה בכל נקודה אחרת, לכן אין אסימפטוטות אנכיות נוספות.

\textbf{אסימפטוטה משופעת/אופקית:} נחשב
\[ m=\lim_{x\to\pm\infty}\frac{f(x)}{x}=\lim_{x\to\pm\infty}\left(2-\frac3x+\frac1{x^2}\right)=2, \]
\[ n=\lim_{x\to\pm\infty}\big(f(x)-2x\big)=\lim_{x\to\pm\infty}\left(-3+\frac1x\right)=-3. \]
לכן $y=2x-3$ אסימפטוטה משופעת גם ב-$+\infty$ וגם ב-$-\infty$. כיוון ש-$f(x)\to\pm\infty$ כאשר $x\to\pm\infty$, אין אסימפטוטה אופקית.

\textbf{תשובה:} $x=0$ (אנכית) ו-$y=2x-3$ (משופעת בשני הכיוונים).
\end{solution}

\subsection*{שאלה 2ב}

\begin{question}
מצאו תחומי קמירות ונקודות פיתול של הפונקציה
\[ f(x)=x^2\cdot e^{-x} \]
\end{question}

\begin{hint}
חשבו את $f''$ ובדקו את סימנה; שימו לב ש-$e^{-x}>0$ תמיד.
\end{hint}

\begin{solution}
$f$ מוגדרת וגזירה פעמיים בכל $\R$.
\[ f'(x)=2xe^{-x}-x^2e^{-x}=(2x-x^2)e^{-x}, \]
\[ f''(x)=(2-2x)e^{-x}-(2x-x^2)e^{-x}=(x^2-4x+2)e^{-x}. \]
כיוון ש-$e^{-x}>0$, סימן $f''$ הוא סימן $x^2-4x+2$, שאפסיו $x=2\pm\sqrt2$. זהו טרינום עם מקדם מוביל חיובי, ולכן:
\begin{itemize}
  \item $f''>0$ ב-$(-\infty,\,2-\sqrt2)$ וב-$(2+\sqrt2,\,\infty)$: שם $f$ \textbf{קמורה} (קעורה כלפי מעלה).
  \item $f''<0$ ב-$(2-\sqrt2,\,2+\sqrt2)$: שם $f$ \textbf{קעורה} (קעורה כלפי מטה).
\end{itemize}
בנקודות $x=2\pm\sqrt2$ הפונקציה רציפה ו-$f''$ מחליפה סימן, ולכן אלו \textbf{נקודות פיתול}:
\[ \left(2-\sqrt2,\ (2-\sqrt2)^2e^{-(2-\sqrt2)}\right)=\left(2-\sqrt2,\ (6-4\sqrt2)e^{\sqrt2-2}\right), \]
\[ \left(2+\sqrt2,\ (6+4\sqrt2)e^{-2-\sqrt2}\right). \]
\end{solution}

\subsection*{שאלה 3}

\begin{question}
העריכו בעזרת קירוב לינארי (פולינום טיילור ממעלה ראשונה) את $\sqrt[4]{18}$ ותנו חסם לשגיאה.
\end{question}

\begin{hint}
פתחו את $f(x)=\sqrt[4]{x}$ סביב נקודה קרובה ל-$18$ שבה השורש הרביעי ידוע.
\end{hint}

\begin{hint}
את השגיאה חסמו בעזרת שארית לגרנז': $R_1(x)=\frac{f''(c)}{2}(x-a)^2$ עבור $c$ בין $a$ ל-$x$.
\end{hint}

\begin{solution}
נבחר $f(x)=x^{1/4}$ ונפתח סביב $a=16$ (כי $\sqrt[4]{16}=2$).
\[ f'(x)=\tfrac14x^{-3/4},\qquad f''(x)=-\tfrac{3}{16}x^{-7/4}. \]
לכן $f(16)=2$, $f'(16)=\frac14\cdot16^{-3/4}=\frac14\cdot\frac18=\frac1{32}$, ופולינום טיילור מסדר ראשון הוא
\[ P_1(x)=2+\frac{1}{32}(x-16). \]
הקירוב:
\[ \sqrt[4]{18}\approx P_1(18)=2+\frac{2}{32}=\frac{33}{16}=2.0625. \]

\textbf{חסם לשגיאה.} לפי משפט טיילור עם שארית לגרנז' ($f$ גזירה פעמיים ב-$(0,\infty)$), קיימת $c\in(16,18)$ כך ש-
\[ R_1(18)=f(18)-P_1(18)=\frac{f''(c)}{2!}(18-16)^2=-\frac{3}{16}c^{-7/4}\cdot\frac{4}{2}=-\frac{3}{8}c^{-7/4}. \]
הפונקציה $c^{-7/4}$ יורדת, ולכן עבור $c>16$:
\[ |R_1(18)|=\frac38c^{-7/4}<\frac38\cdot16^{-7/4}=\frac38\cdot\frac{1}{2^7}=\frac{3}{1024}\approx0.0029. \]

\textbf{תשובה:} $\sqrt[4]{18}\approx2.0625$ עם שגיאה קטנה מ-$\frac{3}{1024}\approx 0.003$. יתרה מזו, כיוון ש-$R_1<0$, הקירוב גדול מהערך האמיתי:
$2.0625-\frac{3}{1024}<\sqrt[4]{18}<2.0625$.
(לבדיקה: $\sqrt[4]{18}\approx2.05977$.)
\end{solution}

\subsection*{שאלה 4א}

\begin{question}
חשבו את הפונקציה הקדומה של
\[ f(x)=\frac{e^x}{e^{2x}+1} \]
\end{question}

\begin{hint}
הציבו $t=e^x$.
\end{hint}

\begin{solution}
נציב $t=e^x$, $dt=e^x\,dx$:
\[ \int\frac{e^x}{e^{2x}+1}\,dx=\int\frac{dt}{t^2+1}=\arctan t+C=\boxed{\arctan(e^x)+C}. \]
בדיקה: $\left(\arctan e^x\right)'=\frac{e^x}{1+e^{2x}}$.
\end{solution}

\subsection*{שאלה 4ב}

\begin{question}
מצאו את הפתרון הכללי של המשוואה הדיפרנציאלית
\[ (x^2+4)y'+2xy^2=0 \]
\end{question}

\begin{hint}
זו משוואה פרידה: העבירו את כל ה-$y$ לצד אחד ואת כל ה-$x$ לצד השני. אל תשכחו את הפתרון $y\equiv0$.
\end{hint}

\begin{solution}
המשוואה שקולה ל-
\[ y'=-\frac{2x}{x^2+4}\,y^2 \]
(כי $x^2+4>0$), וזו משוואה פרידה.

\textbf{פתרון קבוע:} $y\equiv0$ מקיים את המשוואה.

\textbf{פתרונות עם $y\neq0$:} נחלק ב-$y^2$:
\[ \frac{y'}{y^2}=-\frac{2x}{x^2+4}\ \Longrightarrow\ \int\frac{dy}{y^2}=-\int\frac{2x}{x^2+4}\,dx \ \Longrightarrow\ -\frac1y=-\ln(x^2+4)+C. \]
(השתמשנו ב-$\int\frac{2x}{x^2+4}dx=\ln(x^2+4)$, שכן המונה הוא נגזרת המכנה.) לכן
\[ \boxed{y=\frac{1}{\ln(x^2+4)+C}},\quad C\in\R, \]
בכל קטע שבו המכנה אינו מתאפס, ובנוסף הפתרון $\boxed{y\equiv0}$.

בדיקה: עבור $y=\frac1{\ln(x^2+4)+C}$ מתקיים $y'=-\frac{2x}{x^2+4}\cdot y^2$, ולכן $(x^2+4)y'+2xy^2=0$.
\end{solution}

\subsection*{שאלה 5}

\begin{question}
חשבו את האינטגרל
\[ \int\frac{6x^2+13x+8}{x^3+2x^2+2x}dx \]
\end{question}

\begin{hint}
פרקו את המכנה: $x^3+2x^2+2x=x(x^2+2x+2)$, ולגורם הריבועי אין שורשים ממשיים. השתמשו בפירוק לשברים חלקיים.
\end{hint}

\begin{hint}
את $\frac{2x+5}{x^2+2x+2}$ פצלו ל-$\frac{2x+2}{x^2+2x+2}+\frac{3}{(x+1)^2+1}$.
\end{hint}

\begin{solution}
\textbf{פירוק המכנה.} $x^3+2x^2+2x=x(x^2+2x+2)$, והדיסקרימיננטה של $x^2+2x+2$ היא $4-8<0$, ולכן הוא אי-פריק מעל $\R$. מעלת המונה קטנה ממעלת המכנה, ולכן נחפש פירוק לשברים חלקיים:
\[ \frac{6x^2+13x+8}{x(x^2+2x+2)}=\frac{A}{x}+\frac{Bx+C}{x^2+2x+2}. \]
כפל במכנה: $6x^2+13x+8=A(x^2+2x+2)+(Bx+C)x$.
\begin{itemize}
  \item הצבת $x=0$: $8=2A$, ולכן $A=4$.
  \item אז $6x^2+13x+8-4(x^2+2x+2)=2x^2+5x=Bx^2+Cx$, ולכן $B=2$, $C=5$.
\end{itemize}
\[ \frac{6x^2+13x+8}{x^3+2x^2+2x}=\frac4x+\frac{2x+5}{x^2+2x+2}=\frac4x+\frac{2x+2}{x^2+2x+2}+\frac{3}{(x+1)^2+1}. \]

\textbf{אינטגרציה.}
\begin{itemize}
  \item $\int\frac4x\,dx=4\ln|x|$.
  \item $\int\frac{2x+2}{x^2+2x+2}\,dx=\ln(x^2+2x+2)$ (המונה הוא נגזרת המכנה, והמכנה חיובי).
  \item $\int\frac{3}{(x+1)^2+1}\,dx=3\arctan(x+1)$ (הצבה $t=x+1$).
\end{itemize}

\textbf{תשובה:}
\[ \int\frac{6x^2+13x+8}{x^3+2x^2+2x}dx=\boxed{4\ln|x|+\ln(x^2+2x+2)+3\arctan(x+1)+C}. \]
(בדיקה: גזירת התוצאה מחזירה את האינטגרנד.)
\end{solution}

\subsection*{שאלה 6א}

\begin{question}
חשבו את השטח הסופי החסום בין הפונקציה $f(x)=x^2$ והישר $y=-x+2$.
\end{question}

\begin{hint}
מצאו את נקודות החיתוך, וקבעו איזו פונקציה עליונה בין שתיהן.
\end{hint}

\begin{solution}
\textbf{נקודות חיתוך:} $x^2=-x+2\iff x^2+x-2=0\iff (x+2)(x-1)=0$, כלומר $x=-2$ ו-$x=1$.
בקטע $(-2,1)$ מתקיים $-x+2-x^2=-(x+2)(x-1)>0$, ולכן הישר מעל הפרבולה. השטח:
\[ S=\int_{-2}^{1}\left(-x+2-x^2\right)dx=\left[-\frac{x^2}{2}+2x-\frac{x^3}{3}\right]_{-2}^{1}
=\left(-\frac12+2-\frac13\right)-\left(-2-4+\frac83\right)=\frac76+\frac{10}{3}=\boxed{\frac92}. \]
\end{solution}

\subsection*{שאלה 6ב}

\begin{question}
חשבו את האינטגרל הלא אמיתי $\int_0^\infty xe^{-x}dx$, או הראו שאינו מתכנס.
\end{question}

\begin{hint}
חשבו את $\int_0^R xe^{-x}dx$ באינטגרציה בחלקים, ואז השאיפו $R\to\infty$.
\end{hint}

\begin{solution}
לפי ההגדרה, $\int_0^\infty xe^{-x}dx=\lim_{R\to\infty}\int_0^R xe^{-x}dx$.
באינטגרציה בחלקים ($u=x$, $v'=e^{-x}$, כלומר $u'=1$, $v=-e^{-x}$):
\[ \int_0^R xe^{-x}dx=\Big[-xe^{-x}\Big]_0^R+\int_0^R e^{-x}dx=-Re^{-R}+\Big[-e^{-x}\Big]_0^R=-Re^{-R}-e^{-R}+1. \]
כאשר $R\to\infty$: $e^{-R}\to0$, וגם $Re^{-R}=\frac{R}{e^R}\to0$ (למשל לפי כלל לופיטל: $\lim\frac{R}{e^R}=\lim\frac1{e^R}=0$). לכן
\[ \int_0^\infty xe^{-x}dx=\lim_{R\to\infty}\left(1-Re^{-R}-e^{-R}\right)=\boxed{1}, \]
והאינטגרל מתכנס.
\end{solution}

\section*{תשפ"ב סמסטר א מועד ב}

\subsection*{שאלה 1}

\begin{question}
עבור פרמטרים כלשהם $a,b\in\R$ נגדיר פונקציה על ידי
\[ f(x)=\begin{cases} \sin(ax)\sin(\frac{1}{x}) & x<0\\ b+3 & x=0\\ x\ln(x) & x>0 \end{cases} \]
לאילו ערכים של הפרמטרים $a,b$ הפונקציה רציפה ב-$0$? נמקו.
\end{question}

\begin{hint}
חשבו בנפרד את הגבולות החד-צדדיים ב-$0$. משמאל: "אפסה כפול חסומה". מימין: כתבו $x\ln x=\frac{\ln x}{1/x}$ והשתמשו בכלל לופיטל.
\end{hint}

\begin{solution}
$f$ רציפה ב-$0$ אם ורק אם שני הגבולות החד-צדדיים קיימים ושווים ל-$f(0)=b+3$.

\textbf{גבול משמאל.} לכל $a\in\R$: $\sin(ax)\to\sin0=0$ כאשר $x\to0^-$ (רציפות הסינוס), ו-$|\sin(1/x)|\le1$. לפי "אפסה כפול חסומה":
\[ \lim_{x\to0^-}\sin(ax)\sin\left(\tfrac1x\right)=0. \]
(במפורש: $0\le|\sin(ax)\sin(1/x)|\le|\sin(ax)|\le|a||x|\to0$.) הגבול שווה $0$ \textbf{לכל} ערך של $a$.

\textbf{גבול מימין.} זהו גבול מהצורה $0\cdot(-\infty)$. נכתוב כמנה מהצורה $\frac{-\infty}{\infty}$ ונשתמש בכלל לופיטל:
\[ \lim_{x\to0^+}x\ln x=\lim_{x\to0^+}\frac{\ln x}{1/x}\overset{L}{=}\lim_{x\to0^+}\frac{1/x}{-1/x^2}=\lim_{x\to0^+}(-x)=0. \]

\textbf{מסקנה.} שני הגבולות החד-צדדיים שווים $0$, ולכן $\lim_{x\to0}f(x)=0$. הרציפות ב-$0$ שקולה ל-$b+3=0$.

\textbf{תשובה:} $f$ רציפה ב-$0$ אם ורק אם $\boxed{b=-3}$, ו-$a\in\R$ כלשהו.
\end{solution}

\subsection*{שאלה 2א}

\begin{question}
מצאו נקודות קיצון מקומי ותחומי עלייה וירידה של הפונקציה
\[ f(x)=\frac{2x}{x^2+1} \]
\end{question}

\begin{hint}
חשבו את $f'$ לפי כלל המנה ובדקו את סימנה; המכנה של $f'$ חיובי תמיד.
\end{hint}

\begin{solution}
$f$ מוגדרת וגזירה בכל $\R$ (המכנה חיובי). לפי כלל המנה:
\[ f'(x)=\frac{2(x^2+1)-2x\cdot2x}{(x^2+1)^2}=\frac{2(1-x^2)}{(x^2+1)^2}. \]
$f'(x)=0\iff x=\pm1$. סימן $f'$ הוא סימן $1-x^2$:
\begin{itemize}
  \item $f'<0$ ב-$(-\infty,-1)$: $f$ \textbf{יורדת}.
  \item $f'>0$ ב-$(-1,1)$: $f$ \textbf{עולה}.
  \item $f'<0$ ב-$(1,\infty)$: $f$ \textbf{יורדת}.
\end{itemize}
לכן ב-$x=-1$ יש \textbf{מינימום מקומי} $f(-1)=-1$, וב-$x=1$ יש \textbf{מקסימום מקומי} $f(1)=1$.
(למעשה אלו גם קיצון מוחלט: $|2x|\le x^2+1$ כי $(|x|-1)^2\ge0$, ולכן $-1\le f\le 1$.)
\end{solution}

\subsection*{שאלה 2ב}

\begin{question}
חשבו ערכי מינימום ומקסימום מוחלטים של הפונקציה הבאה בתחום הגדרתה:
\[ f(x)=\frac{\sqrt{x^2-1}}{x^2} \]
או הסבירו מדוע אינם קיימים.
\end{question}

\begin{hint}
מצאו קודם את תחום ההגדרה, ושימו לב שהפונקציה זוגית ואי-שלילית. בדקו גם את הגבול באינסוף.
\end{hint}

\begin{solution}
\textbf{תחום הגדרה:} $x^2-1\ge0$ ו-$x\neq0$, כלומר $D=(-\infty,-1]\cup[1,\infty)$.
$f$ זוגית, ולכן מספיק לחקור את $[1,\infty)$.

\textbf{מינימום מוחלט:} $f(x)\ge0$ לכל $x\in D$, ו-$f(\pm1)=0$. לכן המינימום המוחלט הוא $\boxed{0}$, המתקבל ב-$x=\pm1$.

\textbf{מקסימום מוחלט:} עבור $x>1$:
\[ f'(x)=\frac{\frac{x}{\sqrt{x^2-1}}\cdot x^2-2x\sqrt{x^2-1}}{x^4}=\frac{x^2-2(x^2-1)}{x^3\sqrt{x^2-1}}=\frac{2-x^2}{x^3\sqrt{x^2-1}}. \]
לכן $f'>0$ ב-$(1,\sqrt2)$ ו-$f'<0$ ב-$(\sqrt2,\infty)$: $f$ עולה ב-$[1,\sqrt2]$ ויורדת ב-$[\sqrt2,\infty)$ (רציפה בקצה $1$). מכאן שלכל $x\ge1$ מתקיים $f(x)\le f(\sqrt2)$, ו-
\[ f(\sqrt2)=\frac{\sqrt{2-1}}{2}=\frac12. \]
(נשים לב גם ש-$\lim_{x\to\infty}f(x)=\lim\frac{\sqrt{1-1/x^2}}{x}=0$, כך שאין "בריחה" באינסוף.) מזוגיות, אותו דבר נכון ב-$(-\infty,-1]$.

\textbf{תשובה:} המקסימום המוחלט הוא $\boxed{\tfrac12}$, המתקבל ב-$x=\pm\sqrt2$; המינימום המוחלט הוא $\boxed{0}$, המתקבל ב-$x=\pm1$.
\end{solution}

\subsection*{שאלה 3}

\begin{question}
העריכו בקירוב את הערך של $\sqrt[4]{e}$ עם שארית הקטנה מ-$\frac{1}{1000}$. עליכם לנמק היטב את תשובתכם.
\end{question}

\begin{hint}
$\sqrt[4]{e}=e^{1/4}$. השתמשו בפולינום מקלורן של $e^x$ ובשארית לגרנז', ומצאו את הסדר $n$ הקטן שמבטיח שגיאה קטנה מ-$\frac{1}{1000}$.
\end{hint}

\begin{hint}
כדי לחסום את $e^c$ עבור $0<c<\frac14$ אפשר להשתמש ב-$e^c<e<3$.
\end{hint}

\begin{solution}
נכתוב $\sqrt[4]{e}=e^{1/4}$ ונשתמש בפולינום מקלורן של $f(x)=e^x$. כל הנגזרות הן $f^{(k)}(x)=e^x$, ולכן $f^{(k)}(0)=1$ ו-
\[ P_n(x)=\sum_{k=0}^n\frac{x^k}{k!}. \]
לפי משפט טיילור עם שארית לגרנז', לכל $n$ קיימת $c\in(0,\tfrac14)$ כך ש-
\[ R_n\left(\tfrac14\right)=\frac{e^c}{(n+1)!}\left(\tfrac14\right)^{n+1}. \]
כיוון ש-$e^x$ עולה, $e^c<e^{1/4}<e<3$, ולכן
\[ 0<R_n\left(\tfrac14\right)<\frac{3}{(n+1)!\,4^{n+1}}. \]
\begin{itemize}
  \item $n=2$: החסם הוא $\frac{3}{6\cdot64}=\frac{1}{128}$ — לא מספיק.
  \item $n=3$: החסם הוא $\frac{3}{24\cdot256}=\frac{1}{2048}<\frac{1}{1000}$ — מספיק.
\end{itemize}
(ואכן $n=2$ באמת אינו מספיק, ולא רק החסם שלו: $R_2=\frac{e^c}{6\cdot64}>\frac{1}{384}>\frac1{1000}$, כי $e^c>1$.)

לכן נבחר $n=3$:
\[ \sqrt[4]{e}\approx P_3\left(\tfrac14\right)=1+\frac14+\frac{1}{2\cdot16}+\frac{1}{6\cdot64}=1+\frac14+\frac1{32}+\frac1{384}=\frac{493}{384}\approx1.28385, \]
והשגיאה מקיימת $0<\sqrt[4]{e}-\frac{493}{384}<\frac{1}{2048}<\frac{1}{1000}$.

\textbf{תשובה:} $\boxed{\sqrt[4]{e}\approx\frac{493}{384}\approx1.2839}$ (הערך האמיתי $\approx1.28403$).
\end{solution}

\subsection*{שאלה 4א}

\begin{question}
חשבו את הפונקציה הקדומה של
\[ f(x)=\cos^2x\sin^3x \]
\end{question}

\begin{hint}
כתבו $\sin^3x=(1-\cos^2x)\sin x$ והציבו $t=\cos x$.
\end{hint}

\begin{solution}
נכתוב $\cos^2x\sin^3x=\cos^2x(1-\cos^2x)\sin x$ ונציב $t=\cos x$, $dt=-\sin x\,dx$:
\[ \int\cos^2x\sin^3x\,dx=-\int t^2(1-t^2)\,dt=-\frac{t^3}{3}+\frac{t^5}{5}+C=\boxed{\frac{\cos^5x}{5}-\frac{\cos^3x}{3}+C}. \]
בדיקה: $\left(\frac{\cos^5x}{5}-\frac{\cos^3x}{3}\right)'=-\cos^4x\sin x+\cos^2x\sin x=\cos^2x\sin x(1-\cos^2x)=\cos^2x\sin^3x$.
\end{solution}

\subsection*{שאלה 4ב}

\begin{question}
פתרו את המשוואה הדיפרנציאלית
\[ xe^y=(x^2+1)y' \]
\end{question}

\begin{hint}
זו משוואה פרידה: $e^{-y}y'=\frac{x}{x^2+1}$.
\end{hint}

\begin{solution}
כיוון ש-$x^2+1>0$ ו-$e^y>0$, נוכל לכתוב
\[ e^{-y}\,y'=\frac{x}{x^2+1}. \]
(אין פתרונות קבועים: עבור $y\equiv c$ נקבל $xe^c=0$ לכל $x$, וזה בלתי אפשרי.) נבצע אינטגרציה של שני האגפים:
\[ \int e^{-y}dy=\int\frac{x}{x^2+1}dx\ \Longrightarrow\ -e^{-y}=\frac12\ln(x^2+1)-C. \]
לכן $e^{-y}=C-\frac12\ln(x^2+1)$, ומכאן
\[ \boxed{y=-\ln\left(C-\tfrac12\ln(x^2+1)\right)},\qquad C\in\R, \]
בתחום שבו $C-\frac12\ln(x^2+1)>0$ (בפרט נדרש $C>0$; התחום הוא $|x|<\sqrt{e^{2C}-1}$).

בדיקה: $y'=\frac{\frac{x}{x^2+1}}{C-\frac12\ln(x^2+1)}=\frac{x}{x^2+1}e^{y}$, ולכן $(x^2+1)y'=xe^y$.
\end{solution}

\subsection*{שאלה 5}

\begin{question}
חשבו את האינטגרל
\[ \int\frac{5x^3+4x^2+8x+12}{x^4+4x^2}dx \]
\end{question}

\begin{hint}
$x^4+4x^2=x^2(x^2+4)$. הפירוק לשברים חלקיים הוא מהצורה $\frac Ax+\frac B{x^2}+\frac{Cx+D}{x^2+4}$.
\end{hint}

\begin{solution}
\textbf{פירוק לשברים חלקיים.} $x^4+4x^2=x^2(x^2+4)$, ו-$x^2+4$ אי-פריק מעל $\R$. מעלת המונה קטנה ממעלת המכנה, ולכן
\[ \frac{5x^3+4x^2+8x+12}{x^2(x^2+4)}=\frac Ax+\frac B{x^2}+\frac{Cx+D}{x^2+4}, \]
\[ 5x^3+4x^2+8x+12=Ax(x^2+4)+B(x^2+4)+(Cx+D)x^2=(A+C)x^3+(B+D)x^2+4Ax+4B. \]
השוואת מקדמים:
\[ 4B=12\Rightarrow B=3,\quad 4A=8\Rightarrow A=2,\quad A+C=5\Rightarrow C=3,\quad B+D=4\Rightarrow D=1. \]
לכן
\[ \frac{5x^3+4x^2+8x+12}{x^4+4x^2}=\frac2x+\frac3{x^2}+\frac{3x}{x^2+4}+\frac{1}{x^2+4}. \]

\textbf{אינטגרציה.}
\begin{itemize}
  \item $\int\frac2x\,dx=2\ln|x|$, \quad $\int\frac3{x^2}\,dx=-\frac3x$.
  \item $\int\frac{3x}{x^2+4}\,dx=\frac32\int\frac{2x}{x^2+4}\,dx=\frac32\ln(x^2+4)$.
  \item $\int\frac{dx}{x^2+4}=\frac14\int\frac{dx}{(x/2)^2+1}=\frac12\arctan\frac x2$ (הצבה $t=\frac x2$).
\end{itemize}

\textbf{תשובה:}
\[ \int\frac{5x^3+4x^2+8x+12}{x^4+4x^2}dx=\boxed{2\ln|x|-\frac3x+\frac32\ln(x^2+4)+\frac12\arctan\frac x2+C}. \]
\end{solution}

\subsection*{שאלה 6א}

\begin{question}
חשבו את נפח גוף הסיבוב של הפונקציה $f(x)=\sin x$ סביב ציר ה-$x$ בקטע $[0,\pi]$.
\end{question}

\begin{hint}
השתמשו בנוסחה $V=\pi\int_a^b f^2(x)\,dx$ ובזהות $\sin^2x=\frac{1-\cos2x}{2}$.
\end{hint}

\begin{solution}
נפח גוף הסיבוב סביב ציר ה-$x$:
\[ V=\pi\int_0^\pi\sin^2x\,dx=\pi\int_0^\pi\frac{1-\cos2x}{2}\,dx=\pi\left[\frac x2-\frac{\sin2x}{4}\right]_0^\pi=\pi\cdot\frac\pi2=\boxed{\frac{\pi^2}{2}}. \]
\end{solution}

\subsection*{שאלה 6ב}

\begin{question}
בדקו התכנסות של האינטגרל הלא אמיתי
\[ \int_1^\infty\frac{\sin(\frac1x)}{2+x\sqrt{x}}dx \]
\end{question}

\begin{hint}
האינטגרנד חיובי בקטע. השתמשו במבחן ההשוואה, עם $0<\sin t\le t$ עבור $0<t\le1$.
\end{hint}

\begin{solution}
עבור $x\ge1$ מתקיים $0<\frac1x\le1<\pi$, ולכן $\sin\frac1x>0$, והאינטגרנד חיובי ורציף ב-$[1,\infty)$ (כך שהבעיה היחידה היא באינסוף). מהאי-שוויון $\sin t\le t$ עבור $t\ge0$:
\[ 0<\frac{\sin(1/x)}{2+x\sqrt x}\le\frac{1/x}{x\sqrt x}=\frac{1}{x^{5/2}}. \]
האינטגרל $\int_1^\infty\frac{dx}{x^{p}}$ מתכנס עבור $p>1$, ובפרט עבור $p=\frac52$. לפי \textbf{מבחן ההשוואה} לאינטגרלים לא אמיתיים של פונקציות אי-שליליות, האינטגרל
\[ \int_1^\infty\frac{\sin(1/x)}{2+x\sqrt x}dx \]
\textbf{מתכנס}.

(דרך חלופית: מבחן ההשוואה הגבולי עם $g(x)=x^{-5/2}$: $\lim_{x\to\infty}\frac{f(x)}{g(x)}=\lim\frac{\sin(1/x)}{1/x}\cdot\frac{x^{3/2}}{2+x^{3/2}}=1$.)
\end{solution}

\section*{תשפ"ב סמסטר א מועד מיוחד}

\subsection*{שאלה 1}

\begin{question}
עבור פרמטרים כלשהם $a,b\in\R$ נגדיר פונקציה על ידי
\[ f(x)=\begin{cases} e^{\frac1x} & x<0\\ ax+b-6 & 0\le x\le 9\\ \frac{x^2-8x-9}{3-\sqrt{x}} & x>9 \end{cases} \]
לאילו ערכים של הפרמטרים $a,b$ הפונקציה רציפה ב-$\R$? נמקו.
\end{question}

\begin{hint}
בכל אחד מהתחומים הפתוחים $f$ היא הרכבה/מנה של פונקציות רציפות. נשאר לבדוק את נקודות התפר $x=0$ ו-$x=9$.
\end{hint}

\begin{hint}
ב-$x=9$: $x^2-8x-9=(x-9)(x+1)$ ו-$x-9=(\sqrt x-3)(\sqrt x+3)$.
\end{hint}

\begin{solution}
\textbf{רציפות בתוך התחומים.} ב-$(-\infty,0)$ הפונקציה $e^{1/x}$ רציפה (הרכבת רציפות). ב-$(0,9)$ $f$ פולינום. ב-$(9,\infty)$ $f$ מנה של פונקציות רציפות שהמכנה שלה $3-\sqrt x<0$ אינו מתאפס. לכן $f$ רציפה בכל $x\neq0,9$ לכל $a,b$, ויש לבדוק רק את $x=0$ ואת $x=9$.

\textbf{הנקודה $x=0$.} $f(0)=b-6$, והגבול מימין הוא $\lim_{x\to0^+}(ax+b-6)=b-6$. משמאל: כאשר $x\to0^-$ מתקיים $\frac1x\to-\infty$, ולכן
\[ \lim_{x\to0^-}e^{1/x}=0. \]
רציפות ב-$0$ שקולה ל-$b-6=0$, כלומר $b=6$.

\textbf{הנקודה $x=9$.} $f(9)=9a+b-6$, וזה גם הגבול משמאל. מימין, עבור $x>9$:
\[ \frac{x^2-8x-9}{3-\sqrt x}=\frac{(x-9)(x+1)}{3-\sqrt x}=\frac{(\sqrt x-3)(\sqrt x+3)(x+1)}{-(\sqrt x-3)}=-(\sqrt x+3)(x+1), \]
ולכן
\[ \lim_{x\to9^+}f(x)=-(3+3)(9+1)=-60. \]
רציפות ב-$9$ שקולה ל-$9a+b-6=-60$.

\textbf{פתרון המערכת.} מ-$b=6$ נקבל $9a=-60$, כלומר $a=-\frac{20}{3}$.

\textbf{תשובה:} $f$ רציפה בכל $\R$ אם ורק אם $\boxed{a=-\tfrac{20}{3},\ b=6}$.
\end{solution}

\subsection*{שאלה 2א}

\begin{question}
מצאו את כל נקודות הקיצון המקומיות והמוחלטות (גלובליות) על כל הישר של הפונקציה
\[ f(x)=x^2\cdot e^{-x} \]
\end{question}

\begin{hint}
בדקו את הגבולות ב-$\pm\infty$ כדי להכריע אם יש קיצון מוחלט.
\end{hint}

\begin{solution}
$f$ גזירה בכל $\R$, ו-
\[ f'(x)=2xe^{-x}-x^2e^{-x}=x(2-x)e^{-x}. \]
כיוון ש-$e^{-x}>0$, $f'(x)=0\iff x=0$ או $x=2$, וסימן $f'$ הוא סימן $x(2-x)$:
\begin{itemize}
  \item $f'<0$ ב-$(-\infty,0)$: $f$ יורדת.
  \item $f'>0$ ב-$(0,2)$: $f$ עולה.
  \item $f'<0$ ב-$(2,\infty)$: $f$ יורדת.
\end{itemize}
לכן ב-$x=0$ יש \textbf{מינימום מקומי} $f(0)=0$, וב-$x=2$ יש \textbf{מקסימום מקומי} $f(2)=\frac{4}{e^2}$.

\textbf{קיצון מוחלט.}
\begin{itemize}
  \item $f(x)=x^2e^{-x}\ge0=f(0)$ לכל $x$, ולכן $x=0$ היא \textbf{מינימום מוחלט}, ערכו $0$.
  \item $\lim_{x\to-\infty}x^2e^{-x}=+\infty$ (שני הגורמים שואפים ל-$+\infty$), ולכן $f$ אינה חסומה מלעיל ו\textbf{אין מקסימום מוחלט}. בפרט המקסימום המקומי ב-$x=2$ אינו מוחלט (למשל $f(-2)=4e^2>\frac4{e^2}$).
\end{itemize}
(לשלמות: $\lim_{x\to+\infty}\frac{x^2}{e^x}=0$ לפי לופיטל פעמיים.)
\end{solution}

\subsection*{שאלה 2ב}

\begin{question}
הוכיחו כי לפונקציה $f(x)=\frac15x^5+\frac23x^3+x+7$ יש שורש ממשי יחיד.
\end{question}

\begin{hint}
הקיום נובע ממשפט ערך הביניים, והיחידות — ממונוטוניות: $f'(x)=x^4+2x^2+1$ (או ממשפט רול).
\end{hint}

\begin{solution}
\textbf{קיום.} $f$ פולינום ולכן רציפה בכל $\R$. מתקיים $f(0)=7>0$ ו-
\[ f(-2)=-\frac{32}{5}-\frac{16}{3}-2+7=-\frac{101}{15}<0. \]
לפי \textbf{משפט ערך הביניים} קיים $x_0\in(-2,0)$ עם $f(x_0)=0$.

\textbf{יחידות.}
\[ f'(x)=x^4+2x^2+1=(x^2+1)^2>0\quad\text{לכל } x, \]
ולכן $f$ עולה ממש ב-$\R$, ובפרט חד-חד-ערכית — היא מקבלת את הערך $0$ לכל היותר פעם אחת.
(לחלופין, לפי משפט רול: אם היו שני שורשים $x_1<x_2$, היה קיים $c\in(x_1,x_2)$ עם $f'(c)=0$, בסתירה ל-$f'>0$.)

לכן ל-$f$ יש שורש ממשי יחיד, והוא נמצא בקטע $(-2,0)$.
\end{solution}

\subsection*{שאלה 3}

\begin{question}
העריכו בעזרת קירוב לינארי (פולינום טיילור מסדר ראשון) את $\arctan(0.2)$ ותנו חסם לשגיאה.

עליכם לנמק היטב את תשובתכם.
\end{question}

\begin{hint}
פתחו את $f(x)=\arctan x$ סביב $a=0$ וחסמו את השגיאה בעזרת שארית לגרנז' $R_1=\frac{f''(c)}{2}x^2$.
\end{hint}

\begin{solution}
נפתח את $f(x)=\arctan x$ סביב $a=0$:
\[ f'(x)=\frac{1}{1+x^2},\qquad f''(x)=-\frac{2x}{(1+x^2)^2}. \]
$f(0)=0$, $f'(0)=1$, ולכן $P_1(x)=x$ והקירוב הוא
\[ \arctan(0.2)\approx P_1(0.2)=\boxed{0.2}. \]

\textbf{חסם לשגיאה.} לפי משפט טיילור עם שארית לגרנז', קיימת $c\in(0,0.2)$ כך ש-
\[ R_1(0.2)=\frac{f''(c)}{2}(0.2)^2=-\frac{c}{(1+c^2)^2}\cdot0.04. \]
עבור $0<c<0.2$ מתקיים $(1+c^2)^2>1$, ולכן $\frac{c}{(1+c^2)^2}<c<0.2$, ומכאן
\[ |R_1(0.2)|<0.2\cdot0.04=0.008. \]

\textbf{תשובה:} $\arctan(0.2)\approx0.2$ עם שגיאה קטנה מ-$0.008$. יתרה מזו $R_1<0$, ולכן $0.192<\arctan(0.2)<0.2$.

\textbf{הערה (חסם טוב יותר).} כיוון ש-$f''(0)=0$, מתקיים $P_2=P_1$, ולכן אפשר להשתמש בשארית מסדר שני: $R_2(0.2)=\frac{f'''(c)}{6}(0.2)^3$ עם $f'''(x)=\frac{6x^2-2}{(1+x^2)^3}$. עבור $0<c<0.2$ מתקיים $|6c^2-2|\le 2$ ו-$(1+c^2)^3>1$, ולכן $|f'''(c)|<2$ ו-
\[ |R_2(0.2)|<\frac{2}{6}\cdot0.008=\frac{0.008}{3}\approx0.0027. \]
(לבדיקה: $\arctan(0.2)\approx0.19740$, והשגיאה בפועל $\approx0.0026$.)
\end{solution}

\subsection*{שאלה 4א}

\begin{question}
חשבו את הפונקציה הקדומה של
\[ f(x)=\cos(3x)\cos(7x) \]
\end{question}

\begin{hint}
השתמשו בזהות $\cos\alpha\cos\beta=\frac12\left[\cos(\alpha-\beta)+\cos(\alpha+\beta)\right]$.
\end{hint}

\begin{solution}
לפי הזהות $\cos\alpha\cos\beta=\frac12\left[\cos(\alpha-\beta)+\cos(\alpha+\beta)\right]$:
\[ \cos(3x)\cos(7x)=\frac12\left[\cos(4x)+\cos(10x)\right]. \]
לכן
\[ \int\cos(3x)\cos(7x)\,dx=\frac12\left[\frac{\sin(4x)}{4}+\frac{\sin(10x)}{10}\right]+C=\boxed{\frac{\sin(4x)}{8}+\frac{\sin(10x)}{20}+C}. \]
\end{solution}

\subsection*{שאלה 4ב}

\begin{question}
פתרו את המשוואה הדיפרנציאלית
\[ y'=e^{x-y} \]
\end{question}

\begin{hint}
$e^{x-y}=e^x\cdot e^{-y}$, ולכן זו משוואה פרידה.
\end{hint}

\begin{solution}
$y'=e^x e^{-y}$, ולכן $e^{y}y'=e^x$ (אין פתרונות קבועים, כי $e^{x-y}>0$). אינטגרציה של שני האגפים:
\[ \int e^y\,dy=\int e^x\,dx\ \Longrightarrow\ e^y=e^x+C. \]
לכן
\[ \boxed{y=\ln\left(e^x+C\right)},\qquad C\in\R, \]
בתחום שבו $e^x+C>0$ (לכל $x$ אם $C\ge0$; עבור $C<0$ — בתחום $x>\ln(-C)$).

בדיקה: $y'=\frac{e^x}{e^x+C}=\frac{e^x}{e^y}=e^{x-y}$.
\end{solution}

\subsection*{שאלה 5}

\begin{question}
חשבו את האינטגרל
\[ \int\frac{5x^2+2x+2}{x^3-1}dx \]
\end{question}

\begin{hint}
$x^3-1=(x-1)(x^2+x+1)$, והגורם הריבועי אי-פריק. פרקו לשברים חלקיים.
\end{hint}

\begin{solution}
\textbf{פירוק לשברים חלקיים.} $x^3-1=(x-1)(x^2+x+1)$, ולגורם $x^2+x+1$ דיסקרימיננטה $1-4<0$, ולכן הוא אי-פריק. מעלת המונה קטנה ממעלת המכנה:
\[ \frac{5x^2+2x+2}{(x-1)(x^2+x+1)}=\frac{A}{x-1}+\frac{Bx+C}{x^2+x+1}, \]
\[ 5x^2+2x+2=A(x^2+x+1)+(Bx+C)(x-1). \]
\begin{itemize}
  \item הצבת $x=1$: $9=3A$, ולכן $A=3$.
  \item אז $5x^2+2x+2-3(x^2+x+1)=2x^2-x-1=(x-1)(2x+1)$, ולכן $Bx+C=2x+1$, כלומר $B=2$, $C=1$.
\end{itemize}
\[ \frac{5x^2+2x+2}{x^3-1}=\frac{3}{x-1}+\frac{2x+1}{x^2+x+1}. \]

\textbf{אינטגרציה.} $\int\frac{3}{x-1}dx=3\ln|x-1|$, והמונה $2x+1$ הוא בדיוק נגזרת המכנה $x^2+x+1>0$, ולכן $\int\frac{2x+1}{x^2+x+1}dx=\ln(x^2+x+1)$.

\textbf{תשובה:}
\[ \int\frac{5x^2+2x+2}{x^3-1}dx=\boxed{3\ln|x-1|+\ln(x^2+x+1)+C}. \]
\end{solution}

\subsection*{שאלה 6א}

\begin{question}
חשבו את השטח הסופי החסום בין גרף הפונקציה $f(x)=x^3-x$ וציר ה-$x$.
\end{question}

\begin{hint}
שימו לב שהפונקציה מחליפה סימן בין נקודות החיתוך עם הציר — יש לחשב את השטח של כל חלק בנפרד (או להשתמש באי-זוגיות).
\end{hint}

\begin{solution}
\textbf{חיתוך עם ציר ה-$x$:} $x^3-x=x(x-1)(x+1)=0\iff x\in\{-1,0,1\}$.
בקטע $(-1,0)$ מתקיים $f>0$, ובקטע $(0,1)$ מתקיים $f<0$. מחוץ ל-$[-1,1]$ התחום בין הגרף לציר אינו חסום, ולכן השטח הסופי החסום הוא שני ה"עלים" שבין $-1$ ל-$1$:
\[ S=\int_{-1}^0(x^3-x)\,dx+\int_0^1(x-x^3)\,dx. \]
\[ \int_0^1(x-x^3)\,dx=\left[\frac{x^2}{2}-\frac{x^4}{4}\right]_0^1=\frac14,\qquad \int_{-1}^0(x^3-x)\,dx=\left[\frac{x^4}{4}-\frac{x^2}{2}\right]_{-1}^0=-\left(\frac14-\frac12\right)=\frac14. \]
(השוויון נובע גם מכך ש-$f$ אי-זוגית.) לכן
\[ S=\frac14+\frac14=\boxed{\frac12}. \]
(שימו לב: $\int_{-1}^1(x^3-x)\,dx=0$ — זה אינו השטח.)
\end{solution}

\subsection*{שאלה 6ב}

\begin{question}
בדקו התכנסות של האינטגרל הלא אמיתי
\[ \int_1^\infty\frac{e^x}{e^{2x}+1}dx \]
\end{question}

\begin{hint}
הפונקציה הקדומה היא $\arctan(e^x)$ (הצבה $t=e^x$).
\end{hint}

\begin{solution}
נציב $t=e^x$, $dt=e^xdx$: $\int\frac{e^x}{e^{2x}+1}dx=\int\frac{dt}{t^2+1}=\arctan(e^x)+C$. לכן
\[ \int_1^R\frac{e^x}{e^{2x}+1}dx=\arctan(e^R)-\arctan(e). \]
כאשר $R\to\infty$: $e^R\to\infty$ ו-$\lim_{t\to\infty}\arctan t=\frac\pi2$, ולכן
\[ \int_1^\infty\frac{e^x}{e^{2x}+1}dx=\lim_{R\to\infty}\left(\arctan(e^R)-\arctan e\right)=\boxed{\frac\pi2-\arctan e}\approx0.3466. \]
האינטגרל \textbf{מתכנס}.

(לחלופין, ללא חישוב: $0<\frac{e^x}{e^{2x}+1}<\frac{e^x}{e^{2x}}=e^{-x}$, ו-$\int_1^\infty e^{-x}dx=e^{-1}$ מתכנס, ולכן לפי מבחן ההשוואה האינטגרל מתכנס.)
\end{solution}

\section*{תשפ"ג סמסטר א מועד א}

\subsection*{שאלה 1א}

\begin{question}
בדקו קיום של פונקציות הפוכה לפונקציה $f(x)=|x|+x$ בתחום הגדרתה. במקרה שהיא קיימת מצאו אותה.
\end{question}

\begin{hint}
פתחו את הערך המוחלט: מה ערכה של $f$ עבור $x<0$? לפונקציה יש הפוכה רק אם היא חד-חד-ערכית.
\end{hint}

\begin{solution}
הפונקציה מוגדרת לכל $x\in\R$. נפתח את הערך המוחלט:
\[ f(x)=\begin{cases} 2x & x\ge 0\\ 0 & x<0.\end{cases} \]
לפונקציה קיימת פונקציה הפוכה (בתחום הגדרתה) אם ורק אם היא חד-חד-ערכית. אבל למשל $f(-1)=f(-2)=0$ בעוד $-1\neq -2$, ולכן $f$ אינה חד-חד-ערכית על $\R$, ו\textbf{לא קיימת לה פונקציה הפוכה} בתחום הגדרתה.

(הערה: אם מצמצמים את $f$ לקרן $[0,\infty)$ היא שם $f(x)=2x$, חד-חד-ערכית ועל $[0,\infty)$, וההפוכה של הצמצום היא $f^{-1}(y)=\frac{y}{2}$, $y\ge 0$. אך זו כבר לא $f$ על כל תחום הגדרתה.)
\end{solution}

\subsection*{שאלה 1ב}

\begin{question}
לאילו ערכים של הפרמטר $a$ יש לפונקציה הבאה גבול ב-$x=0$?
\[ f(x)=\begin{cases} \dfrac{\sin(2x)}{x} & x<0 \\[2mm] \dfrac{x+3}{4x+a} & x\ge 0 \end{cases} \]
\end{question}

\begin{hint}
חשבו בנפרד את הגבולות החד-צדדיים ודרשו שיהיו שווים. שימו לב למקרה $a=0$.
\end{hint}

\begin{solution}
לפונקציה יש גבול ב-$x=0$ אם ורק אם שני הגבולות החד-צדדיים קיימים (סופיים) ושווים.

\textbf{גבול משמאל:} לפי הגבול היסודי $\lim_{t\to 0}\frac{\sin t}{t}=1$,
\[ \lim_{x\to 0^-}\frac{\sin(2x)}{x}=\lim_{x\to 0^-}2\cdot\frac{\sin(2x)}{2x}=2. \]

\textbf{גבול מימין:} אם $a\neq 0$, המונה והמכנה רציפים ב-$0$ והמכנה שונה מ-$0$ שם, ולכן
\[ \lim_{x\to 0^+}\frac{x+3}{4x+a}=\frac{3}{a}. \]
אם $a=0$, אז $\frac{x+3}{4x}\to +\infty$ כאשר $x\to 0^+$ (המונה שואף ל-$3$ והמכנה ל-$0^+$), והגבול מימין אינו סופי, לכן אין גבול.

לכן עבור $a\ne 0$ דרוש $\frac{3}{a}=2$, כלומר $a=\frac32$. (במקרה זה המכנה $4x+\frac32>0$ לכל $x\ge 0$, כך שהפונקציה מוגדרת היטב בסביבה ימנית של $0$.)

\textbf{תשובה:} לפונקציה יש גבול ב-$x=0$ רק עבור $a=\frac{3}{2}$, וערך הגבול הוא $2$.
\end{solution}

\subsection*{שאלה 2א}

\begin{question}
השתמשו בקירוב לינארי כדי להעריך את הביטוי $\sqrt[3]{28}$.
\end{question}

\begin{hint}
בחרו נקודה קרובה ל-$28$ שבה השורש השלישי ידוע.
\end{hint}

\begin{solution}
נגדיר $g(x)=\sqrt[3]{x}=x^{1/3}$ ונקודת בסיס $x_0=27$, שבה $g(27)=3$. הנגזרת:
\[ g'(x)=\frac{1}{3}x^{-2/3}=\frac{1}{3\sqrt[3]{x^2}},\qquad g'(27)=\frac{1}{3\cdot 9}=\frac{1}{27}. \]
הקירוב הלינארי (המשיק) בסביבת $x_0$: $g(x)\approx g(x_0)+g'(x_0)(x-x_0)$. עם $x=28$, $x-x_0=1$:
\[ \sqrt[3]{28}\approx 3+\frac{1}{27}\approx 3.037. \]
(הערך האמיתי הוא $3.0366\ldots$.)
\end{solution}

\subsection*{שאלה 2ב}

\begin{question}
תהי $f(x)$ פונקציה רציפה על כל הישר, ונניח כי לכל $x$ מתקיים $|f(x)|\le 7$. הוכיחו כי למשוואה $2x+f(x)=3$ יש לפחות פתרון אחד.
\end{question}

\begin{hint}
הגדירו $g(x)=2x+f(x)-3$ ומצאו נקודה שבה $g$ שלילית ונקודה שבה $g$ חיובית, בעזרת החסם $|f(x)|\le 7$.
\end{hint}

\begin{solution}
נגדיר $g(x)=2x+f(x)-3$. הפונקציה $g$ רציפה על כל $\R$ כסכום של פונקציות רציפות. מהנתון $-7\le f(x)\le 7$ לכל $x$, ולכן:
\[ g(-10)=-20+f(-10)-3\le -23+7=-16<0,\qquad g(10)=20+f(10)-3\ge 17-7=10>0. \]
$g$ רציפה בקטע הסגור $[-10,10]$ ומקבלת בקצותיו ערכים בעלי סימנים מנוגדים, ולכן לפי \textbf{משפט ערך הביניים} (משפט בולצאנו) קיימת $c\in(-10,10)$ שבה $g(c)=0$, כלומר $2c+f(c)=3$. לכן למשוואה יש לפחות פתרון אחד. $\blacksquare$
\end{solution}

\subsection*{שאלה 3א}

\begin{question}
מצאו תחומי קמירות ונקודות פיתול עבור הפונקציה $f(x)=\ln(1+4x^2)$.
\end{question}

\begin{solution}
תחום ההגדרה: $1+4x^2>0$ לכל $x$, ולכן $D=\R$. נגזור:
\[ f'(x)=\frac{8x}{1+4x^2},\qquad
f''(x)=\frac{8(1+4x^2)-8x\cdot 8x}{(1+4x^2)^2}=\frac{8-32x^2}{(1+4x^2)^2}=\frac{8(1-2x)(1+2x)}{(1+4x^2)^2}. \]
המכנה חיובי תמיד, ולכן סימן $f''$ הוא סימן $1-4x^2$:
\begin{itemize}
  \item $f''(x)>0$ עבור $-\frac12<x<\frac12$: הפונקציה \textbf{קמורה} (קמורה כלפי מעלה, $\cup$) בקטע $\left(-\frac12,\frac12\right)$.
  \item $f''(x)<0$ עבור $|x|>\frac12$: הפונקציה \textbf{קעורה} ($\cap$) בתחומים $\left(-\infty,-\frac12\right)$ ו-$\left(\frac12,\infty\right)$.
\end{itemize}
ב-$x=\pm\frac12$ הנגזרת השנייה מחליפה סימן והפונקציה רציפה (וגזירה) שם, ולכן אלה נקודות פיתול. $f(\pm\tfrac12)=\ln(1+1)=\ln 2$.

\textbf{תשובה:} קמורה ב-$\left(-\frac12,\frac12\right)$, קעורה ב-$\left(-\infty,-\frac12\right)\cup\left(\frac12,\infty\right)$; נקודות פיתול $\left(-\frac12,\ln 2\right)$ ו-$\left(\frac12,\ln 2\right)$.
\end{solution}

\subsection*{שאלה 3ב}

\begin{question}
הוכיחו לכל $x,y$ מתקיים $|\arctan x-\arctan y|\le|x-y|$.
\end{question}

\begin{hint}
הפעילו את משפט לגרנז' על $\arctan$ בקטע שבין $x$ ל-$y$.
\end{hint}

\begin{solution}
אם $x=y$ שני האגפים שווים $0$ ואין מה להוכיח. נניח $x\neq y$, ובלי הגבלת הכלליות $x<y$. הפונקציה $g(t)=\arctan t$ רציפה ב-$[x,y]$ וגזירה ב-$(x,y)$, ולכן לפי \textbf{משפט לגרנז'} קיימת $c\in(x,y)$ כך ש-
\[ \arctan y-\arctan x=g'(c)(y-x)=\frac{1}{1+c^2}(y-x). \]
מכיוון ש-$1+c^2\ge 1$ מתקיים $0<\frac{1}{1+c^2}\le 1$, ולכן
\[ |\arctan x-\arctan y|=\frac{1}{1+c^2}\,|x-y|\le|x-y|. \qquad\blacksquare \]
\end{solution}

\subsection*{שאלה 4א}

\begin{question}
חשבו את הגבול $\displaystyle\lim_{x\to 1}\frac{x^2+\ln(x)-1}{e^x-e}$.
\end{question}

\begin{solution}
בהצבת $x=1$: המונה $1+0-1=0$ והמכנה $e-e=0$, כלומר ביטוי מהצורה $\frac00$. המונה והמכנה גזירים בסביבת $1$, ונגזרת המכנה $e^x\neq 0$. לפי \textbf{כלל לופיטל}:
\[ \lim_{x\to1}\frac{x^2+\ln x-1}{e^x-e}=\lim_{x\to1}\frac{2x+\frac1x}{e^x}=\frac{2+1}{e}=\frac{3}{e}, \]
כאשר הגבול האחרון מחושב בהצבה (רציפות), ולכן גם הגבול המקורי קיים ושווה ל-$\frac{3}{e}$.
\end{solution}

\subsection*{שאלה 4ב}

\begin{question}
חשבו את $\sin 0.1$ בשגיאה שלא תעלה על $0.001$.
\end{question}

\begin{hint}
השתמשו בפולינום מקלורן של $\sin x$ והעריכו את השארית בצורת לגרנז'.
\end{hint}

\begin{solution}
נשתמש בפולינום מקלורן של $f(x)=\sin x$. מתקיים $f(0)=0$, $f'(0)=\cos 0=1$, $f''(0)=-\sin 0=0$, ולכן
\[ P_2(x)=x. \]
לפי נוסחת טיילור עם שארית לגרנז', קיימת $c$ בין $0$ ל-$x$ כך ש-
\[ \sin x=P_2(x)+R_2(x),\qquad R_2(x)=\frac{f'''(c)}{3!}x^3=\frac{-\cos c}{6}x^3. \]
עבור $x=0.1$, ומכיוון ש-$|\cos c|\le 1$:
\[ |R_2(0.1)|\le\frac{(0.1)^3}{6}=\frac{0.001}{6}\approx 0.000167<0.001. \]
לכן $\sin 0.1\approx 0.1$, בשגיאה קטנה מ-$0.001$ כנדרש.

(לקירוב מדויק יותר אפשר לקחת $P_4(x)=x-\frac{x^3}{6}$, ואז $\sin 0.1\approx 0.0998333$ בשגיאה $\le\frac{(0.1)^5}{120}<10^{-7}$; הערך האמיתי $0.0998334\ldots$.)
\end{solution}

\subsection*{שאלה 5א}

\begin{question}
חשבו את האינטגרל $\displaystyle\int\sqrt{2x+3}\,dx$.
\end{question}

\begin{solution}
נציב $t=2x+3$, ואז $dt=2\,dx$, כלומר $dx=\frac{dt}{2}$:
\[ \int\sqrt{2x+3}\,dx=\frac12\int t^{1/2}\,dt=\frac12\cdot\frac{t^{3/2}}{3/2}+C=\frac13(2x+3)^{3/2}+C. \]
בדיקה: $\left(\frac13(2x+3)^{3/2}\right)'=\frac13\cdot\frac32(2x+3)^{1/2}\cdot 2=\sqrt{2x+3}$.
\end{solution}

\subsection*{שאלה 5ב}

\begin{question}
חשבו את האינטגרל $\displaystyle\int_{-1}^{1}\tan x\,dx$.
\end{question}

\begin{hint}
בדקו האם הפונקציה $\tan x$ זוגית או אי-זוגית, והאם היא רציפה בקטע $[-1,1]$.
\end{hint}

\begin{solution}
מכיוון ש-$1<\frac{\pi}{2}$, הפונקציה $\tan x=\frac{\sin x}{\cos x}$ רציפה בקטע $[-1,1]$ (שם $\cos x>0$), ולכן האינטגרל המסוים קיים. היא אי-זוגית: $\tan(-x)=-\tan x$. לאינטגרל של פונקציה אי-זוגית רציפה על קטע סימטרי $[-a,a]$ ערך $0$. ישירות: פונקציה קדומה היא $-\ln|\cos x|$ (כי $(-\ln\cos x)'=\frac{\sin x}{\cos x}$), ולפי \textbf{המשפט היסודי} (ניוטון-לייבניץ):
\[ \int_{-1}^{1}\tan x\,dx=\Big[-\ln(\cos x)\Big]_{-1}^{1}=-\ln(\cos 1)+\ln(\cos(-1))=0, \]
כי $\cos$ זוגית. \textbf{תשובה:} $0$.
\end{solution}

\subsection*{שאלה 6א}

\begin{question}
חשבו את השטח החסום בין הגרפים של הפונקציות $f(x)=\sin^2x$ ו-$g(x)=\frac14$ בקטע $[0,\pi]$.
\end{question}

\begin{hint}
מצאו את נקודות החיתוך $\sin^2x=\frac14$ בקטע, וחלקו את הקטע לפי איזו פונקציה גדולה יותר. השתמשו בזהות $\sin^2x=\frac{1-\cos 2x}{2}$.
\end{hint}

\begin{solution}
\textbf{נקודות חיתוך:} $\sin^2x=\frac14\iff\sin x=\pm\frac12$. בקטע $[0,\pi]$ מתקיים $\sin x\ge 0$, ולכן $\sin x=\frac12$, כלומר $x=\frac{\pi}{6}$ או $x=\frac{5\pi}{6}$.

בקטעים $[0,\frac{\pi}{6})$ ו-$(\frac{5\pi}{6},\pi]$ מתקיים $0\le\sin x<\frac12$, ולכן $\sin^2x<\frac14$; בקטע $(\frac{\pi}{6},\frac{5\pi}{6})$ מתקיים $\sin^2x>\frac14$.

\textbf{פונקציה קדומה:} לפי $\sin^2x=\frac{1-\cos2x}{2}$,
\[ \int\left(\sin^2x-\frac14\right)dx=\int\left(\frac14-\frac{\cos 2x}{2}\right)dx=\frac{x}{4}-\frac{\sin 2x}{4}=:F(x). \]
ערכים: $F(0)=0$, $F(\frac{\pi}{6})=\frac{\pi}{24}-\frac{\sqrt3}{8}$, $F(\frac{5\pi}{6})=\frac{5\pi}{24}+\frac{\sqrt3}{8}$, $F(\pi)=\frac{\pi}{4}$.

\textbf{השטח} הוא $\int_0^\pi|f-g|\,dx$:
\begin{align*}
S_1&=\int_0^{\pi/6}\left(\tfrac14-\sin^2x\right)dx=-\big(F(\tfrac{\pi}{6})-F(0)\big)=\frac{\sqrt3}{8}-\frac{\pi}{24},\\
S_2&=\int_{\pi/6}^{5\pi/6}\left(\sin^2x-\tfrac14\right)dx=F(\tfrac{5\pi}{6})-F(\tfrac{\pi}{6})=\frac{\pi}{6}+\frac{\sqrt3}{4},\\
S_3&=\int_{5\pi/6}^{\pi}\left(\tfrac14-\sin^2x\right)dx=-\big(F(\pi)-F(\tfrac{5\pi}{6})\big)=\frac{\sqrt3}{8}-\frac{\pi}{24}.
\end{align*}
סה"כ:
\[ S=S_1+S_2+S_3=\frac{\sqrt3}{2}+\frac{\pi}{12}\approx 1.128. \]
(אם מתכוונים רק לתחום הסגור שבין שתי נקודות החיתוך, שטחו $S_2=\frac{\pi}{6}+\frac{\sqrt3}{4}\approx 0.957$.)
\end{solution}

\subsection*{שאלה 6ב}

\begin{question}
חשבו את האינטגרל הלא-אמיתי $\displaystyle\int_0^{\infty}e^{-2x}\sin x\,dx$ או הראו שאינו מתכנס.
\end{question}

\begin{hint}
חשבו את האינטגרל עד $R$ בעזרת אינטגרציה בחלקים פעמיים (האינטגרל "חוזר על עצמו"), ואז העבירו $R\to\infty$.
\end{hint}

\begin{solution}
\textbf{התכנסות:} $|e^{-2x}\sin x|\le e^{-2x}$ ו-$\int_0^\infty e^{-2x}dx=\frac12$ מתכנס, ולכן לפי \textbf{מבחן ההשוואה} האינטגרל מתכנס (בהחלט).

\textbf{חישוב:} נסמן $I(R)=\int_0^R e^{-2x}\sin x\,dx$ ונמצא פונקציה קדומה $J=\int e^{-2x}\sin x\,dx$ באינטגרציה בחלקים פעמיים ($u=e^{-2x}$):
\begin{align*}
J&=-e^{-2x}\cos x-\int 2e^{-2x}\cos x\,dx
 =-e^{-2x}\cos x-2\left(e^{-2x}\sin x+\int 2e^{-2x}\sin x\,dx\right)\\
 &=-e^{-2x}\cos x-2e^{-2x}\sin x-4J.
\end{align*}
לכן $5J=-e^{-2x}(\cos x+2\sin x)$, כלומר
\[ \int e^{-2x}\sin x\,dx=-\frac{e^{-2x}(\cos x+2\sin x)}{5}+C. \]
מכאן
\[ I(R)=\frac{1}{5}-\frac{e^{-2R}(\cos R+2\sin R)}{5}. \]
כאשר $R\to\infty$: $|\cos R+2\sin R|\le 3$ ו-$e^{-2R}\to 0$, ולכן (חסומה כפול שואפת לאפס) האיבר השני שואף ל-$0$. לכן
\[ \int_0^\infty e^{-2x}\sin x\,dx=\frac15. \]
\end{solution}

\section*{תשפ"ג סמסטר א מועד ב}

\subsection*{שאלה 1א}

\begin{question}
חשבו את הגבולות החד צדדיים של הפונקציה $f(x)=2^{\frac{1}{x}}$ בנקודה הנתונה $x=0$. קבעו האם יש לה גבול $\lim_{x\to 0}f(x)$.
\end{question}

\begin{solution}
נסמן $t=\frac1x$. כאשר $x\to 0^+$ מתקיים $t\to+\infty$, וכאשר $x\to 0^-$ מתקיים $t\to-\infty$. מכיוון ש-$2>1$, $\lim_{t\to+\infty}2^t=+\infty$ ו-$\lim_{t\to-\infty}2^t=0$. לכן (גבול של הרכבה)
\[ \lim_{x\to0^+}2^{1/x}=+\infty,\qquad \lim_{x\to0^-}2^{1/x}=0. \]
הגבולות החד-צדדיים שונים (והימני אף אינו סופי), ולכן \textbf{הגבול $\lim_{x\to0}2^{1/x}$ אינו קיים}.
\end{solution}

\subsection*{שאלה 1ב}

\begin{question}
מיינו את נקודות אי הרציפות של הפונקציה
\[ f(x)=\begin{cases}\dfrac{x^3-8}{x-2} & x\neq 2\\[2mm] 1 & x=2\end{cases} \]
\end{question}

\begin{solution}
לכל $x\neq 2$ הפונקציה היא מנה של פולינומים עם מכנה שונה מאפס, ולכן רציפה שם. נבדוק את $x=2$. לפי נוסחת הפרש חזקות שלישיות $x^3-8=(x-2)(x^2+2x+4)$, ולכן עבור $x\neq2$:
\[ f(x)=x^2+2x+4\ \Longrightarrow\ \lim_{x\to2}f(x)=4+4+4=12. \]
הגבול קיים וסופי, אך $f(2)=1\neq 12$. לכן $x=2$ היא נקודת אי-רציפות \textbf{סליקה} (ממין ראשון סליקה): אם נגדיר מחדש $f(2)=12$ הפונקציה תהיה רציפה. זו נקודת אי הרציפות היחידה.
\end{solution}

\subsection*{שאלה 2א}

\begin{question}
חשבו את משוואת המשיק למעגל $x^2+y^2=25$ בנקודה $(3,4)$.
\end{question}

\begin{hint}
גזרו את המשוואה בצורה סתומה.
\end{hint}

\begin{solution}
הנקודה על המעגל: $9+16=25$. בסביבת $(3,4)$ ($y>0$) המעגל הוא גרף של פונקציה גזירה $y(x)$. נגזור את המשוואה בצורה סתומה לפי $x$:
\[ 2x+2y\,y'=0\ \Longrightarrow\ y'=-\frac{x}{y},\qquad y'(3)=-\frac34. \]
משוואת המשיק: $y-4=-\frac34(x-3)$, כלומר
\[ y=-\frac34x+\frac{25}{4}\qquad\Longleftrightarrow\qquad 3x+4y=25. \]
(בהתאם לגאומטריה: המשיק ניצב לרדיוס בכיוון $(3,4)$.)
\end{solution}

\subsection*{שאלה 2ב}

\begin{question}
השתמשו בקירוב לינארי כדי להעריך את הביטוי $\arcsin(0.54)$.
\end{question}

\begin{hint}
השתמשו בנקודה $0.5$, שבה $\arcsin$ ידוע.
\end{hint}

\begin{solution}
נגדיר $f(x)=\arcsin x$ ונקודת בסיס $x_0=0.5$: $f(0.5)=\frac{\pi}{6}$, ו-
\[ f'(x)=\frac{1}{\sqrt{1-x^2}},\qquad f'(0.5)=\frac{1}{\sqrt{3/4}}=\frac{2}{\sqrt3}. \]
קירוב לינארי: $f(x)\approx f(x_0)+f'(x_0)(x-x_0)$, עם $x-x_0=0.04$:
\[ \arcsin(0.54)\approx\frac{\pi}{6}+\frac{2}{\sqrt3}\cdot0.04=\frac{\pi}{6}+\frac{0.08}{\sqrt3}\approx 0.5236+0.0462=0.5698. \]
(הערך האמיתי $0.5704\ldots$.)
\end{solution}

\subsection*{שאלה 3א}

\begin{question}
חשבו נקודות קיצון מוחלט של הפונקציה $f(x)=\frac{x^2}{x-3}$ בקטע $[-2,2]$.
\end{question}

\begin{solution}
הנקודה $x=3$ אינה בקטע, ולכן $f$ רציפה בקטע הסגור $[-2,2]$, ולפי \textbf{משפט ויירשטראס} היא מקבלת בו מקסימום ומינימום. הם מתקבלים בנקודות קריטיות פנימיות או בקצוות.
\[ f'(x)=\frac{2x(x-3)-x^2}{(x-3)^2}=\frac{x^2-6x}{(x-3)^2}=\frac{x(x-6)}{(x-3)^2}. \]
$f'(x)=0\iff x=0$ או $x=6$; בקטע $(-2,2)$ רק $x=0$. נשווה ערכים:
\[ f(-2)=\frac{4}{-5}=-\frac45,\qquad f(0)=0,\qquad f(2)=\frac{4}{-1}=-4. \]
\textbf{תשובה:} מקסימום מוחלט $0$ בנקודה $x=0$; מינימום מוחלט $-4$ בנקודה $x=2$.
\end{solution}

\subsection*{שאלה 3ב}

\begin{question}
למשוואה $e^x=x+1$ הוכיחו כי קיים פתרון או הוכיחו כי לא קיים פתרון בתחום $x\neq0$.
\end{question}

\begin{hint}
חקרו את הפונקציה $g(x)=e^x-x-1$: מצאו את נקודת המינימום שלה.
\end{hint}

\begin{solution}
נוכיח ש\textbf{אין} פתרון עם $x\neq 0$. נגדיר $g(x)=e^x-x-1$, גזירה על $\R$, עם $g(0)=0$ ו-
\[ g'(x)=e^x-1. \]
עבור $x<0$: $g'(x)<0$, ולכן $g$ יורדת ממש ב-$(-\infty,0]$; עבור $x>0$: $g'(x)>0$, ולכן $g$ עולה ממש ב-$[0,\infty)$. (מונוטוניות ממש נובעת ממשפט לגרנז'.) לכן:
\begin{itemize}
  \item לכל $x<0$: $g(x)>g(0)=0$;
  \item לכל $x>0$: $g(x)>g(0)=0$.
\end{itemize}
כלומר $e^x>x+1$ לכל $x\neq0$, ולמשוואה $e^x=x+1$ אין פתרון בתחום $x\ne 0$ (הפתרון היחיד שלה הוא $x=0$). $\blacksquare$
\end{solution}

\subsection*{שאלה 4א}

\begin{question}
חשבו את הגבול $\displaystyle\lim_{x\to0}\frac{x-\sin(x)}{x^3}$.
\end{question}

\begin{solution}
ביטוי מהצורה $\frac00$. נפעיל את \textbf{כלל לופיטל} (התנאים מתקיימים: גזירות בסביבת $0$, נגזרת המכנה שונה מ-$0$ ליד $0$):
\[ \lim_{x\to0}\frac{x-\sin x}{x^3}=\lim_{x\to0}\frac{1-\cos x}{3x^2}=\lim_{x\to0}\frac{\sin x}{6x}=\frac16, \]
כאשר בשלב השני הפעלנו שוב לופיטל (שוב $\frac00$) ובאחרון השתמשנו ב-$\lim\frac{\sin x}{x}=1$.
(לחלופין: $\sin x=x-\frac{x^3}{6}+o(x^3)$, ולכן $\frac{x-\sin x}{x^3}=\frac16+o(1)$.) \textbf{תשובה:} $\frac16$.
\end{solution}

\subsection*{שאלה 4ב}

\begin{question}
השתמשו בפולינום טיילור מסדר שני כדי להעריך את הביטוי $\sqrt[4]{82}$ והעריכו את השגיאה.
\end{question}

\begin{hint}
פתחו את $f(x)=x^{1/4}$ סביב $x_0=81$ והשתמשו בשארית לגרנז'.
\end{hint}

\begin{solution}
$f(x)=x^{1/4}$, $x_0=81$, $f(81)=3$.
\[ f'(x)=\tfrac14x^{-3/4},\quad f''(x)=-\tfrac{3}{16}x^{-7/4},\quad f'''(x)=\tfrac{21}{64}x^{-11/4}. \]
\[ f'(81)=\frac{1}{4\cdot27}=\frac{1}{108},\qquad f''(81)=-\frac{3}{16\cdot 3^7}=-\frac{1}{16\cdot729}=-\frac{1}{11664}. \]
פולינום טיילור מסדר 2 סביב $81$, בנקודה $x=82$ ($x-x_0=1$):
\[ \sqrt[4]{82}\approx P_2(82)=3+\frac{1}{108}-\frac{1}{2\cdot11664}=3+\frac{1}{108}-\frac{1}{23328}=\frac{70199}{23328}\approx 3.0092164. \]
\textbf{הערכת שגיאה} (שארית לגרנז'): קיימת $c\in(81,82)$ כך ש-
\[ R_2=\frac{f'''(c)}{3!}(82-81)^3=\frac{21}{64\cdot 6}\,c^{-11/4}. \]
$c^{-11/4}$ יורדת, ולכן $c^{-11/4}<81^{-11/4}=3^{-11}$, ומכאן
\[ 0<R_2<\frac{21}{384\cdot 3^{11}}=\frac{21}{384\cdot177147}\approx 3.1\cdot10^{-7}. \]
כלומר $\sqrt[4]{82}\approx3.0092164$ בשגיאה קטנה מ-$3.1\cdot 10^{-7}$ (והקירוב קטן מעט מהערך האמיתי, $3.0092167\ldots$).
\end{solution}

\subsection*{שאלה 5א}

\begin{question}
חשבו את האינטגרל $\displaystyle\int\frac{x}{\sqrt{1+x^2}}\,dx$.
\end{question}

\begin{solution}
נציב $t=1+x^2$, $dt=2x\,dx$:
\[ \int\frac{x}{\sqrt{1+x^2}}\,dx=\frac12\int t^{-1/2}\,dt=t^{1/2}+C=\sqrt{1+x^2}+C. \]
בדיקה: $\left(\sqrt{1+x^2}\right)'=\frac{2x}{2\sqrt{1+x^2}}=\frac{x}{\sqrt{1+x^2}}$.
\end{solution}

\subsection*{שאלה 5ב}

\begin{question}
פתרו את המשוואה $e^x\,dx-y\,dy=0$ עם תנאי התחלה $y(0)=1$.
\end{question}

\begin{hint}
זו משוואה פרידה: העבירו אגף ובצעו אינטגרציה לכל צד בנפרד.
\end{hint}

\begin{solution}
המשוואה פרידה: $y\,dy=e^x\,dx$ (כלומר $y\,y'=e^x$). נבצע אינטגרציה לשני האגפים:
\[ \int y\,dy=\int e^x\,dx\ \Longrightarrow\ \frac{y^2}{2}=e^x+C. \]
מתנאי ההתחלה $y(0)=1$: $\frac12=1+C$, ולכן $C=-\frac12$, ו-
\[ y^2=2e^x-1. \]
מכיוון ש-$y(0)=1>0$ ופתרון רציף אינו יכול לעבור לערכים שליליים בלי לעבור ב-$0$, נבחר את הענף החיובי:
\[ y(x)=\sqrt{2e^x-1}, \]
המוגדר כאשר $2e^x-1>0$, כלומר $x>-\ln 2$.
בדיקה: $y'=\frac{e^x}{\sqrt{2e^x-1}}$, ולכן $y\,y'=e^x$ ✓, ו-$y(0)=\sqrt{1}=1$ ✓.
\end{solution}

\subsection*{שאלה 6א}

\begin{question}
חשבו את השטח הסופי החסום בין גרף הפונקציה $f(x)=x^3+x^2-2x$ וציר ה-$x$.
\end{question}

\begin{hint}
פרקו את הפולינום לגורמים כדי למצוא את נקודות החיתוך עם ציר $x$ ואת הסימן בכל קטע.
\end{hint}

\begin{solution}
פירוק: $f(x)=x(x^2+x-2)=x(x+2)(x-1)$. נקודות החיתוך עם ציר $x$: $x=-2,0,1$. סימן: ב-$(-2,0)$ $f>0$ (למשל $f(-1)=2$), וב-$(0,1)$ $f<0$ (למשל $f(\frac12)=-\frac58$). התחום הסופי החסום בין הגרף לציר הוא מעל $[-2,1]$.

פונקציה קדומה: $F(x)=\frac{x^4}{4}+\frac{x^3}{3}-x^2$. אז $F(-2)=4-\frac83-4=-\frac83$, $F(0)=0$, $F(1)=\frac14+\frac13-1=-\frac{5}{12}$.
\[ S=\int_{-2}^0 f\,dx-\int_0^1 f\,dx=\big(F(0)-F(-2)\big)-\big(F(1)-F(0)\big)=\frac83+\frac{5}{12}=\frac{37}{12}. \]
\textbf{תשובה:} $S=\frac{37}{12}$.
\end{solution}

\subsection*{שאלה 6ב}

\begin{question}
בדקו התכנסות של האינטגרל $\displaystyle\int_1^{\infty}\frac{\sqrt{x^5}\,dx}{2x^2+3x+40}$.
\end{question}

\begin{hint}
בדקו כיצד מתנהגת הפונקציה כאשר $x\to\infty$ והשוו ל-$x^{p}$ מתאים (מבחן ההשוואה הגבולי).
\end{hint}

\begin{solution}
הפונקציה $g(x)=\frac{x^{5/2}}{2x^2+3x+40}$ רציפה וחיובית ב-$[1,\infty)$. נשווה ל-$h(x)=x^{1/2}$:
\[ \lim_{x\to\infty}\frac{g(x)}{h(x)}=\lim_{x\to\infty}\frac{x^2}{2x^2+3x+40}=\frac12\in(0,\infty). \]
לפי \textbf{מבחן ההשוואה הגבולי}, $\int_1^\infty g$ ו-$\int_1^\infty x^{1/2}dx$ מתכנסים או מתבדרים יחד. $\int_1^\infty x^{p}dx$ מתכנס רק עבור $p<-1$, ו-$p=\frac12$, לכן $\int_1^\infty x^{1/2}dx=\infty$.
(אפשר גם לשים לב שאפילו $g(x)\to\infty$, ולכן בוודאי שהאינטגרל מתבדר.)

\textbf{תשובה:} האינטגרל \textbf{מתבדר} (שווה $+\infty$).
\end{solution}

\section*{תשפ"ג סמסטר א מועד ג}

\subsection*{שאלה 1א}

\begin{question}
בדקו קיום של פונקציות הפוכה לפונקציה $f(x)=|x|+2x$ בתחום הגדרתה. במקרה שהיא קיימת מצאו אותה.
\end{question}

\begin{hint}
פרקו את הערך המוחלט: רשמו את $f$ בנפרד עבור $x\ge 0$ ועבור $x<0$, ובדקו שהיא חח"ע ועל.
\end{hint}

\begin{solution}
הפונקציה מוגדרת על כל $\R$. נפרק את הערך המוחלט:
\[ f(x)=\begin{cases} 3x & x\ge 0\\ x & x<0. \end{cases} \]
\textbf{חח"ע:} בכל אחד מהתחומים $f$ עולה ממש (שיפוע $3$ ושיפוע $1$), ובנוסף $f(x)\ge 0$ עבור $x\ge0$ ו-$f(x)<0$ עבור $x<0$, כך שערכים משני התחומים אינם יכולים להתלכד. לכן $f$ עולה ממש על כל $\R$, ובפרט חח"ע.

\textbf{על:} התמונה של $[0,\infty)$ תחת $3x$ היא $[0,\infty)$, והתמונה של $(-\infty,0)$ תחת $x$ היא $(-\infty,0)$. לכן התמונה היא כל $\R$.

מכאן ש-$f:\R\to\R$ הפיכה. נמצא את ההפכית: אם $y\ge 0$ אז $y=3x$ עם $x\ge 0$, כלומר $x=y/3$; אם $y<0$ אז $y=x$. לכן
\[ \boxed{f^{-1}(y)=\begin{cases} \dfrac{y}{3} & y\ge 0\\[2mm] y & y<0. \end{cases}} \]
(בדיקה: $f^{-1}(f(x))=x$ בשני התחומים.)
\end{solution}

\subsection*{שאלה 1ב}

\begin{question}
מיינו את נקודות אי הרציפות של הפונקציה
\[ f(x)=\begin{cases} \dfrac{\sin(2x)}{x} & x\neq 0\\[2mm] 1 & x=0 \end{cases} \]
\end{question}

\begin{solution}
לכל $x\neq 0$ הפונקציה היא מנה של פונקציות רציפות ($\sin(2x)$ ו-$x$) שהמכנה שלה אינו מתאפס, ולכן רציפה. נבדוק את $x=0$: לפי הגבול היסודי $\lim_{t\to0}\frac{\sin t}{t}=1$,
\[ \lim_{x\to0}\frac{\sin(2x)}{x}=\lim_{x\to0}2\cdot\frac{\sin(2x)}{2x}=2. \]
הגבול קיים וסופי, אך $f(0)=1\neq 2$. לכן $x=0$ היא \textbf{נקודת אי-רציפות סליקה} (ניתן לתקן ע"י הגדרה מחדש $f(0)=2$), וזו נקודת אי-הרציפות היחידה.
\end{solution}

\subsection*{שאלה 2א}

\begin{question}
השתמשו בקירוב לינארי כדי להעריך את הביטוי $\log_2 10$.
\end{question}

\begin{hint}
בחרו נקודה קרובה ל-$10$ שבה $\log_2$ ידוע, למשל $x_0=8$.
\end{hint}

\begin{solution}
נגדיר $f(x)=\log_2 x=\frac{\ln x}{\ln 2}$, ואז $f'(x)=\frac{1}{x\ln 2}$. הקירוב הלינארי סביב $x_0$ הוא
\[ f(x)\approx f(x_0)+f'(x_0)(x-x_0). \]
ניקח $x_0=8$ (שם $\log_2 8=3$) ו-$x=10$:
\[ \log_2 10\approx 3+\frac{1}{8\ln 2}\cdot 2=3+\frac{1}{4\ln 2}\approx 3+0.3607=\boxed{3.3607}. \]
(הערך המדויק $\log_2 10\approx 3.3219$; הקירוב גבוה מעט מהערך האמיתי, כצפוי, כי $\log_2$ קעורה ולכן המשיק נמצא מעל הגרף.)
\end{solution}

\subsection*{שאלה 2ב}

\begin{question}
חשבו את הנגזרת בנקודה $x=0$ של הפונקציה $y(x)$ הנתונה בצורה סתומה על ידי המשוואה
\[ y+x^3y^3-3xy=x^5 \]
\end{question}

\begin{hint}
מצאו תחילה את $y(0)$ מתוך המשוואה, ואז גזרו את שני האגפים לפי $x$ (כלל השרשרת ונגזרת מכפלה).
\end{hint}

\begin{solution}
\textbf{ערך הפונקציה:} נציב $x=0$ במשוואה: $y+0-0=0$, ולכן $y(0)=0$.

\textbf{גזירה סתומה:} נגזור את שני האגפים לפי $x$ (כאשר $y=y(x)$):
\[ y'+3x^2y^3+3x^3y^2y'-3y-3xy'=5x^4. \]
נציב $x=0,\ y=0$:
\[ y'(0)+0+0-3\cdot 0-0=0 \quad\Longrightarrow\quad \boxed{y'(0)=0}. \]
(באופן כללי $y'=\dfrac{5x^4-3x^2y^3+3y}{1+3x^3y^2-3x}$, והמכנה שווה $1\neq 0$ בנקודה $(0,0)$.)
\end{solution}

\subsection*{שאלה 3א}

\begin{question}
מצאו את הנקודות על העקומה $y=2x^3-3x^2-12x-5$ שבהן המשיק מקביל לציר $x$.
\end{question}

\begin{solution}
המשיק מקביל לציר $x$ בדיוק כאשר שיפועו $0$, כלומר $y'=0$:
\[ y'=6x^2-6x-12=6(x^2-x-2)=6(x-2)(x+1)=0\iff x=2\ \text{או}\ x=-1. \]
נחשב את ערכי $y$: $y(2)=16-12-24-5=-25$ ו-$y(-1)=-2-3+12-5=2$. הנקודות הן
\[ \boxed{(-1,\,2)\quad\text{ו-}\quad(2,\,-25)}. \]
\end{solution}

\subsection*{שאלה 3ב}

\begin{question}
מצאו תחומי קמירות ונקודות פיתול עבור הפונקציה $f(x)=e^{-x^2+2x}$
\end{question}

\begin{solution}
$f$ מוגדרת וגזירה פעמיים על כל $\R$. לפי כלל השרשרת
\[ f'(x)=(2-2x)e^{-x^2+2x}, \]
\[ f''(x)=-2e^{-x^2+2x}+(2-2x)^2e^{-x^2+2x}=(4x^2-8x+2)e^{-x^2+2x}=2(2x^2-4x+1)e^{-x^2+2x}. \]
מכיוון ש-$e^{-x^2+2x}>0$, סימן $f''$ הוא סימן $2x^2-4x+1$, שאפסיו
\[ x_{1,2}=\frac{4\pm\sqrt{16-8}}{4}=1\pm\frac{\sqrt2}{2}. \]
זו פרבולה עם מקדם מוביל חיובי, ולכן:
\begin{itemize}
  \item $f''>0$ ($f$ \textbf{קמורה}) עבור $x<1-\frac{\sqrt2}{2}$ ועבור $x>1+\frac{\sqrt2}{2}$;
  \item $f''<0$ ($f$ \textbf{קעורה}) עבור $1-\frac{\sqrt2}{2}<x<1+\frac{\sqrt2}{2}$.
\end{itemize}
בשתי הנקודות $x=1\pm\frac{\sqrt2}{2}$ הנגזרת השנייה מחליפה סימן והפונקציה רציפה, ולכן אלה נקודות פיתול. מכיוון ש-$-x^2+2x=1-(x-1)^2=1-\frac12=\frac12$ בנקודות אלה, נקודות הפיתול הן
\[ \boxed{\left(1-\tfrac{\sqrt2}{2},\ \sqrt e\right),\quad \left(1+\tfrac{\sqrt2}{2},\ \sqrt e\right)}. \]
\end{solution}

\subsection*{שאלה 4א}

\begin{question}
חשבו את הגבול $\displaystyle\lim_{x\to0}\frac{e^{2x}-e^x}{x}$
\end{question}

\begin{solution}
זהו גבול מהצורה $\frac00$, והמונה והמכנה גזירים, לכן לפי כלל לופיטל
\[ \lim_{x\to0}\frac{e^{2x}-e^x}{x}=\lim_{x\to0}\frac{2e^{2x}-e^x}{1}=2-1=\boxed{1}. \]
(לחלופין: $\frac{e^{2x}-e^x}{x}=2\cdot\frac{e^{2x}-1}{2x}-\frac{e^x-1}{x}\to 2-1=1$ לפי הגבול היסודי $\lim_{t\to0}\frac{e^t-1}{t}=1$.)
\end{solution}

\subsection*{שאלה 4ב}

\begin{question}
השתמשו בפולינום מקלורן מסדר שני כדי להעריך את הביטוי $e^{0.1}$ והעריכו את השגיאה
\end{question}

\begin{hint}
השתמשו בשארית לגראנז' $R_2(x)=\frac{f'''(c)}{3!}x^3$ עם $c$ בין $0$ ל-$0.1$, וחסמו את $e^c$.
\end{hint}

\begin{solution}
עבור $f(x)=e^x$ כל הנגזרות הן $e^x$ ושוות $1$ ב-$0$, לכן פולינום מקלורן מסדר שני הוא
\[ P_2(x)=1+x+\frac{x^2}{2}, \qquad e^{0.1}\approx P_2(0.1)=1+0.1+0.005=\boxed{1.105}. \]
\textbf{הערכת השגיאה:} לפי משפט טיילור עם שארית לגראנז', קיימת $c\in(0,0.1)$ כך ש-
\[ R_2(0.1)=e^{0.1}-P_2(0.1)=\frac{e^{c}}{3!}(0.1)^3. \]
מכיוון ש-$e^x$ עולה, $e^c<e^{0.1}<e^1<3$, ולכן
\[ 0<R_2(0.1)<\frac{3}{6}\cdot 0.001=0.0005. \]
(חסם הדוק יותר: $e^{0.1}<1.2$ נותן $R_2<0.0002$.) השגיאה חיובית, כלומר הקירוב חסר. הערך האמיתי $e^{0.1}\approx 1.10517$, והשגיאה בפועל כ-$0.00017$.
\end{solution}

\subsection*{שאלה 5א}

\begin{question}
חשבו את האינטגרל $\displaystyle\int x\cdot\sqrt{1-x^2}\,dx$
\end{question}

\begin{solution}
נציב $t=1-x^2$, $dt=-2x\,dx$, כלומר $x\,dx=-\frac12dt$:
\[ \int x\sqrt{1-x^2}\,dx=-\frac12\int t^{1/2}\,dt=-\frac12\cdot\frac{2}{3}t^{3/2}+C=\boxed{-\frac13(1-x^2)^{3/2}+C}. \]
בדיקה בגזירה: $\left(-\frac13(1-x^2)^{3/2}\right)'=-\frac13\cdot\frac32(1-x^2)^{1/2}\cdot(-2x)=x\sqrt{1-x^2}$.
\end{solution}

\subsection*{שאלה 5ב}

\begin{question}
חשבו את האינטגרל $\displaystyle\int_2^4\frac{dx}{x+4}$
\end{question}

\begin{solution}
הפונקציה $\frac{1}{x+4}$ רציפה על $[2,4]$, ופונקציה קדומה שלה היא $\ln|x+4|$. לפי המשפט היסודי (ניוטון-לייבניץ):
\[ \int_2^4\frac{dx}{x+4}=\ln(x+4)\Big|_2^4=\ln8-\ln6=\boxed{\ln\frac43}\approx 0.2877. \]
\end{solution}

\subsection*{שאלה 6א}

\begin{question}
בדקו התכנסות של האינטגרל $\displaystyle\int_0^\infty e^{-x^2}dx$
\end{question}

\begin{hint}
את האינטגרל לא ניתן לחשב בפונקציות אלמנטריות. פצלו ב-$x=1$ והשוו את $e^{-x^2}$ לפונקציה פשוטה יותר עבור $x\ge1$.
\end{hint}

\begin{solution}
נפצל: $\int_0^\infty e^{-x^2}dx=\int_0^1 e^{-x^2}dx+\int_1^\infty e^{-x^2}dx$.
האינטגרל הראשון הוא אינטגרל מסוים רגיל של פונקציה רציפה על קטע סגור, ולכן קיים וסופי.

עבור $x\ge 1$ מתקיים $x^2\ge x$, ולכן $0<e^{-x^2}\le e^{-x}$. האינטגרל
\[ \int_1^\infty e^{-x}dx=\lim_{R\to\infty}\left(e^{-1}-e^{-R}\right)=e^{-1} \]
מתכנס. לפי \textbf{מבחן ההשוואה} לאינטגרלים לא-אמיתיים של פונקציות אי-שליליות, גם $\int_1^\infty e^{-x^2}dx$ מתכנס. לכן $\boxed{\int_0^\infty e^{-x^2}dx\ \text{מתכנס}}$
(וערכו, כידוע, $\frac{\sqrt\pi}{2}$, אך זה לא נדרש).
\end{solution}

\subsection*{שאלה 6ב}

\begin{question}
חשבו את האינטגרל הלא-אמיתי $\displaystyle\int_1^\infty\frac{dx}{x^3}$ או הראו שאינו מתכנס.
\end{question}

\begin{solution}
לפי ההגדרה:
\[ \int_1^\infty\frac{dx}{x^3}=\lim_{R\to\infty}\int_1^R x^{-3}dx=\lim_{R\to\infty}\left[-\frac{1}{2x^2}\right]_1^R=\lim_{R\to\infty}\left(\frac12-\frac{1}{2R^2}\right)=\boxed{\frac12}. \]
האינטגרל מתכנס (בהתאם לכך ש-$\int_1^\infty\frac{dx}{x^p}$ מתכנס עבור $p>1$).
\end{solution}

\section*{תשפ"ג סמסטר ב מועד א}

\subsection*{שאלה 1}

\begin{question}
(שאלה חובה)
\begin{enumerate}
  \item (22 נק') חקרו באופן מלא את הפונקציה: $\displaystyle f(x)=\frac{e^{2x+1}}{(x+1)^2}$ (תחום הגדרה, חיתוך עם הצירים, סימני הפונקציה, אסימפטוטות, תחומי עלייה וירידה, נקודות קיצון, תחומי קמירות ונקודות פיתול) וציירו את סקיצת הגרף של הפונקציה.
  \item (4 נק') שרטט את גרף הפונקציה $f(-|x|)$ ?
\end{enumerate}
\end{question}

\begin{hint}
בגזירה הוציאו גורם משותף: $f'(x)=\frac{e^{2x+1}\left(2(x+1)-2\right)}{(x+1)^3}$.
\end{hint}

\begin{hint}
בסעיף ב' שימו לב ש-$f(-|x|)$ היא פונקציה זוגית, ועבור $x\ge0$ היא שווה ל-$f(-x)$, כלומר שיקוף של החלק של גרף $f$ שמשמאל לציר $y$.
\end{hint}

\begin{solution}
\begin{enumerate}
  \item \textbf{תחום הגדרה:} $x\neq -1$, כלומר $D=(-\infty,-1)\cup(-1,\infty)$.

  \textbf{חיתוך עם הצירים וסימן:} $e^{2x+1}>0$ ו-$(x+1)^2>0$ בתחום, ולכן $f(x)>0$ לכל $x\in D$; אין חיתוך עם ציר $x$. חיתוך עם ציר $y$: $f(0)=e$, הנקודה $(0,e)$.

  \textbf{אסימפטוטות:}
  \begin{itemize}
    \item אנכית: $\lim_{x\to-1^\pm}f(x)=\frac{e^{-1}}{0^+}=+\infty$, ולכן $x=-1$ אסימפטוטה אנכית (משני הצדדים הפונקציה שואפת ל-$+\infty$).
    \item ב-$-\infty$: $e^{2x+1}\to 0$ ו-$(x+1)^2\to\infty$, לכן $f(x)\to0$: $y=0$ אסימפטוטה אופקית משמאל.
    \item ב-$+\infty$: לפי לופיטל (פעמיים) $\lim_{x\to\infty}\frac{e^{2x+1}}{(x+1)^2}=\lim\frac{4e^{2x+1}}{2}=\infty$, ובאופן דומה $\frac{f(x)}{x}\to\infty$, לכן אין אסימפטוטה אופקית או משופעת מימין.
  \end{itemize}

  \textbf{מונוטוניות וקיצון:} לפי כלל המנה,
  \[ f'(x)=\frac{2e^{2x+1}(x+1)^2-e^{2x+1}\cdot2(x+1)}{(x+1)^4}=\frac{2e^{2x+1}\,x}{(x+1)^3}. \]
  $f'(x)=0\iff x=0$. סימן $f'$ הוא סימן $\frac{x}{(x+1)^3}$:
  \begin{itemize}
    \item $x<-1$: מונה שלילי, מכנה שלילי $\Rightarrow f'>0$, $f$ \textbf{עולה}.
    \item $-1<x<0$: מונה שלילי, מכנה חיובי $\Rightarrow f'<0$, $f$ \textbf{יורדת}.
    \item $x>0$: $f'>0$, $f$ \textbf{עולה}.
  \end{itemize}
  לכן ב-$x=0$ יש \textbf{מינימום מקומי} $(0,e)$. אין מקסימום מקומי (ב-$x=-1$ הפונקציה אינה מוגדרת).

  \textbf{קמירות:} גזירה נוספת נותנת
  \[ f''(x)=\frac{2e^{2x+1}(2x^2+1)}{(x+1)^4}. \]
  (נגזור $f'=2e^{2x+1}\cdot x(x+1)^{-3}$: $f''=2e^{2x+1}\left[2x(x+1)^{-3}+(x+1)^{-3}-3x(x+1)^{-4}\right]=\frac{2e^{2x+1}\left[(2x+1)(x+1)-3x\right]}{(x+1)^4}$ ו-$(2x+1)(x+1)-3x=2x^2+1$.)
  מכיוון ש-$2x^2+1>0$, מתקיים $f''>0$ בכל התחום: $f$ \textbf{קמורה} על $(-\infty,-1)$ ועל $(-1,\infty)$, ו\textbf{אין נקודות פיתול}.

  \textbf{סקיצה (תיאור במילים):} משמאל הגרף יוצא מעל האסימפטוטה $y=0$ (קרוב ל-$0$ ב-$-\infty$), עולה בצורה קמורה ושואף ל-$+\infty$ כאשר $x\to-1^-$. מימין לאסימפטוטה $x=-1$ הגרף יורד מ-$+\infty$ עד המינימום $(0,e)$ ואז עולה במהירות (אקספוננציאלית) ל-$+\infty$. הגרף כולו מעל ציר $x$.

  \item נסמן $g(x)=f(-|x|)$. מתקיים $g(-x)=g(x)$, כלומר $g$ \textbf{זוגית} – הגרף סימטרי ביחס לציר $y$. עבור $x\ge0$: $g(x)=f(-x)$, כלומר הגרף של $g$ עבור $x\ge0$ הוא שיקוף ביחס לציר $y$ של גרף $f$ עבור $x\le 0$, והגרף עבור $x<0$ הוא פשוט גרף $f$ עבור $x<0$.

  בפרט: $g$ מוגדרת עבור $x\neq\pm1$, $g(0)=e$, ויש לה אסימפטוטות אנכיות $x=\pm1$ (שם $g\to+\infty$) ואסימפטוטה אופקית $y=0$ ב-$\pm\infty$. על $(-1,0)$ $g=f$ יורדת, ועל $(0,1)$ $g$ עולה; ב-$x=0$ יש מינימום מקומי $(0,e)$ (ומכיוון ש-$f'(0)=0$ אין שם "שפיץ": הנגזרות החד-צדדיות של $g$ שתיהן $0$). על $(-\infty,-1)$ $g$ עולה מ-$0$ ל-$+\infty$, ועל $(1,\infty)$ היא יורדת מ-$+\infty$ ל-$0$. $g$ קמורה בכל אחד מהקטעים.
\end{enumerate}
\end{solution}

\subsection*{שאלה 2א}

\begin{question}
(שאלה חובה)
פתרו את האינטגרל הבא:
$\displaystyle\int\frac{(x^2+1)\cdot dx}{x^3-x}$
\end{question}

\begin{hint}
פרקו לשברים חלקיים: $x^3-x=x(x-1)(x+1)$.
\end{hint}

\begin{solution}
$x^3-x=x(x-1)(x+1)$, ודרגת המונה קטנה מדרגת המכנה. נפרק:
\[ \frac{x^2+1}{x(x-1)(x+1)}=\frac{A}{x}+\frac{B}{x-1}+\frac{C}{x+1}. \]
בשיטת הכיסוי: $A=\frac{0+1}{(0-1)(0+1)}=-1$, $B=\frac{1+1}{1\cdot 2}=1$, $C=\frac{1+1}{(-1)(-2)}=1$. לכן
\[ \int\frac{x^2+1}{x^3-x}dx=-\ln|x|+\ln|x-1|+\ln|x+1|+C=\boxed{\ln\left|\frac{x^2-1}{x}\right|+C}. \]
\end{solution}

\subsection*{שאלה 2ב}

\begin{question}
(שאלה חובה)
פתרו את האינטגרל הבא:
$\displaystyle\int\sin^2 4x\cos^2 x\,dx$
\end{question}

\begin{hint}
השתמשו בנוסחאות הורדת חזקה $\sin^2\alpha=\frac{1-\cos2\alpha}{2}$, $\cos^2\alpha=\frac{1+\cos2\alpha}{2}$.
\end{hint}

\begin{solution}
לפי נוסחאות הורדת חזקה:
\[ \sin^2 4x\cos^2x=\frac{1-\cos8x}{2}\cdot\frac{1+\cos2x}{2}=\frac14\left(1+\cos2x-\cos8x-\cos8x\cos2x\right). \]
לפי נוסחת מכפלה-לסכום $\cos8x\cos2x=\frac12(\cos10x+\cos6x)$, ולכן
\[ \sin^2 4x\cos^2x=\frac14+\frac14\cos2x-\frac14\cos8x-\frac18\cos10x-\frac18\cos6x. \]
נבצע אינטגרציה איבר-איבר:
\[ \boxed{\int\sin^2 4x\cos^2x\,dx=\frac{x}{4}+\frac{\sin2x}{8}-\frac{\sin8x}{32}-\frac{\sin6x}{48}-\frac{\sin10x}{80}+C}. \]
\end{solution}

\subsection*{שאלה 2ג}

\begin{question}
(שאלה חובה)
פתרו את האינטגרל הבא:
$\displaystyle\int_{-3}^{-1}\arctan(x+2)\,dx$
\end{question}

\begin{hint}
הציבו $t=x+2$ ושימו לב לזוגיות/אי-זוגיות.
\end{hint}

\begin{solution}
נציב $t=x+2$, $dt=dx$; הגבולות: $x=-3\mapsto t=-1$, $x=-1\mapsto t=1$:
\[ \int_{-3}^{-1}\arctan(x+2)\,dx=\int_{-1}^{1}\arctan t\,dt. \]
$\arctan$ פונקציה אי-זוגית ורציפה, והקטע סימטרי סביב $0$, לכן האינטגרל $\boxed{=0}$.

(ישירות: באינטגרציה בחלקים $\int\arctan t\,dt=t\arctan t-\frac12\ln(1+t^2)+C$, ובין $-1$ ל-$1$: $\left(\frac\pi4-\frac12\ln2\right)-\left(\frac\pi4-\frac12\ln2\right)=0$.)
\end{solution}

\subsection*{שאלה 2ד}

\begin{question}
(שאלה חובה)
פתרו את האינטגרל הבא:
$\displaystyle\int_2^5\frac{2x-1}{\sqrt{x^2-x}}\,dx$
\end{question}

\begin{hint}
המונה הוא בדיוק הנגזרת של מה שמתחת לשורש.
\end{hint}

\begin{solution}
על $[2,5]$ מתקיים $x^2-x>0$, כך שהאינטגרנד רציף. נציב $u=x^2-x$, $du=(2x-1)dx$; הגבולות $u(2)=2$, $u(5)=20$:
\[ \int_2^5\frac{2x-1}{\sqrt{x^2-x}}dx=\int_2^{20}u^{-1/2}du=2\sqrt u\Big|_2^{20}=2\sqrt{20}-2\sqrt2=\boxed{4\sqrt5-2\sqrt2}\approx 6.116. \]
\end{solution}

\subsection*{שאלה 3א}

\begin{question}
חשבו את שטח התחום החסום על ידי הקווים $y=\min\left(e^x,3\right),\ x=4,\ x=0,\ y=0$
\end{question}

\begin{hint}
מצאו היכן $e^x=3$ ופצלו את האינטגרל בנקודה זו.
\end{hint}

\begin{solution}
$e^x\le 3\iff x\le\ln3$ (כי $e^x$ עולה). לכן על $[0,4]$:
\[ \min(e^x,3)=\begin{cases} e^x & 0\le x\le\ln 3\\ 3 & \ln3\le x\le 4\end{cases} \]
ושימו לב ש-$\ln3\approx1.0986\in[0,4]$. הפונקציה חיובית, ולכן השטח הוא
\[ S=\int_0^{\ln3}e^x\,dx+\int_{\ln3}^{4}3\,dx=\left(e^{\ln3}-e^0\right)+3(4-\ln3)=2+12-3\ln3=\boxed{14-3\ln3}\approx 10.70. \]
\end{solution}

\subsection*{שאלה 3ב}

\begin{question}
מצאו את הערכים של פרמטר $b$ עבורם לגרף הפונקציה $\displaystyle f(x)=\frac{1}{x^2+b}$ אין נקודות פיתול
\end{question}

\begin{hint}
חשבו את $f''$ והפרידו למקרים $b>0$, $b=0$, $b<0$. זכרו שנקודת פיתול חייבת להיות בתחום ההגדרה.
\end{hint}

\begin{solution}
נגזור: $f'(x)=-\frac{2x}{(x^2+b)^2}$, ו-
\[ f''(x)=\frac{-2(x^2+b)^2+2x\cdot2(x^2+b)\cdot2x}{(x^2+b)^4}=\frac{6x^2-2b}{(x^2+b)^3}. \]
\begin{itemize}
  \item $b>0$: $f$ מוגדרת על כל $\R$, המכנה חיובי, והמונה $6x^2-2b$ מחליף סימן ב-$x=\pm\sqrt{b/3}$. לכן יש שתי נקודות פיתול.
  \item $b=0$: $f=\frac1{x^2}$, $x\neq0$, ו-$f''=\frac{6}{x^4}>0$: אין נקודות פיתול.
  \item $b<0$: $f$ מוגדרת עבור $x\neq\pm\sqrt{-b}$. המונה $6x^2-2b=6x^2+2|b|>0$ תמיד, ולכן $f''$ לא מתאפסת. סימן $f''$ משתנה רק במעבר דרך $x=\pm\sqrt{-b}$, שאינן בתחום ההגדרה (אסימפטוטות אנכיות), ולכן אינן נקודות פיתול. אין נקודות פיתול.
\end{itemize}
\textbf{תשובה:} לגרף אין נקודות פיתול בדיוק עבור $\boxed{b\le 0}$.
\end{solution}

\subsection*{שאלה 4א}

\begin{question}
חשבו את הגבולות הבאים
\[ (1)\ \lim_{x\to0}\frac{\ln(1+2x)}{\sin2x-\sin6x}\qquad\qquad (2)\ \lim_{x\to0}\left(\sqrt{1+2x}\right)^{\frac{1}{(-x)}} \]
\end{question}

\begin{hint}
בגבול (2) רשמו את הביטוי בצורה $e^{g(x)}$ וחשבו את הגבול של המעריך.
\end{hint}

\begin{solution}
\textbf{(1)} גבול מהצורה $\frac00$; לפי כלל לופיטל:
\[ \lim_{x\to0}\frac{\ln(1+2x)}{\sin2x-\sin6x}=\lim_{x\to0}\frac{\frac{2}{1+2x}}{2\cos2x-6\cos6x}=\frac{2}{2-6}=\boxed{-\frac12}. \]
(לחלופין, בעזרת $\ln(1+2x)\sim2x$, $\sin2x\sim2x$, $\sin6x\sim6x$: $\frac{2x}{2x-6x}=-\frac12$.)

\textbf{(2)} עבור $x$ קרוב ל-$0$ ($x>-\frac12$, $x\neq0$) הבסיס חיובי, ו-
\[ \left(\sqrt{1+2x}\right)^{-1/x}=e^{-\frac{\ln(1+2x)}{2x}}. \]
לפי הגבול היסודי (או לופיטל) $\lim_{x\to0}\frac{\ln(1+2x)}{2x}=1$. מרציפות פונקציית האקספוננט,
\[ \lim_{x\to0}\left(\sqrt{1+2x}\right)^{\frac{1}{-x}}=e^{-1}=\boxed{\frac1e}. \]
\end{solution}

\subsection*{שאלה 4ב}

\begin{question}
עבור אילו ערכים של $a$ הפונקציה $f(x)=\min\left(x^2+x,a\right)$ תהיה גזירה לכל $x$ ממשי? נמקו את תשובתכם.
\end{question}

\begin{hint}
השוו את $a$ לערך המינימלי של $x^2+x$. אם הישר $y=a$ חותך את הפרבולה, מה קורה לנגזרות החד-צדדיות בנקודת החיתוך?
\end{hint}

\begin{solution}
הפרבולה $p(x)=x^2+x=\left(x+\frac12\right)^2-\frac14$ מקבלת ערך מינימלי $-\frac14$ ב-$x=-\frac12$.
\begin{itemize}
  \item אם $a\le-\frac14$: $p(x)\ge-\frac14\ge a$ לכל $x$, ולכן $f(x)\equiv a$ קבועה, וגזירה בכל $\R$.
  \item אם $a>-\frac14$: למשוואה $x^2+x=a$ שני שורשים $x_1<x_2$, $x_{1,2}=\frac{-1\mp\sqrt{1+4a}}{2}$. אז $f(x)=x^2+x$ על $[x_1,x_2]$ ו-$f(x)=a$ מחוצה לו. בנקודה $x_1$ (שבה $f$ רציפה) הנגזרת השמאלית היא $0$ (פונקציה קבועה), והימנית היא $p'(x_1)=2x_1+1=-\sqrt{1+4a}\neq0$. הנגזרות החד-צדדיות שונות ולכן $f$ אינה גזירה ב-$x_1$ (ובאופן דומה ב-$x_2$, שם הנגזרת השמאלית $\sqrt{1+4a}$ והימנית $0$).
\end{itemize}
\textbf{תשובה:} $f$ גזירה לכל $x$ ממשי אם ורק אם $\boxed{a\le-\frac14}$.
\end{solution}

\subsection*{שאלה 5א}

\begin{question}
מצאו את משוואת המשיק לגרף הפונקציה-
\[ \begin{cases} x=\ln\left(3t^2+e\cdot t-3\right)\\ y=t^3-3t\end{cases} \]
בנקודה $M_0(1,-2)$.
\end{question}

\begin{hint}
מצאו תחילה את ערך הפרמטר $t_0$ המתאים לנקודה $M_0$ (שתי המשוואות צריכות להתקיים בו-זמנית), ואז השתמשו ב-$\frac{dy}{dx}=\frac{y'(t)}{x'(t)}$.
\end{hint}

\begin{solution}
\textbf{מציאת הפרמטר:} מ-$y=-2$: $t^3-3t+2=0\iff(t-1)^2(t+2)=0$, כלומר $t=1$ או $t=-2$.
מ-$x=1$: $3t^2+et-3=e$. עבור $t=1$: $3+e-3=e$ \checkmark. עבור $t=-2$: $12-2e-3=9-2e\neq e$ (כי $e\neq3$). לכן $t_0=1$.

\textbf{הנגזרת:} לפי נוסחת הנגזרת של פונקציה הנתונה פרמטרית,
\[ x'(t)=\frac{6t+e}{3t^2+et-3},\qquad y'(t)=3t^2-3. \]
ב-$t_0=1$: $x'(1)=\frac{6+e}{e}\neq 0$ ו-$y'(1)=0$, ולכן
\[ \frac{dy}{dx}\Big|_{M_0}=\frac{y'(1)}{x'(1)}=0. \]
\textbf{משוואת המשיק:} $y-(-2)=0\cdot(x-1)$, כלומר $\boxed{y=-2}$ (משיק אופקי).
\end{solution}

\subsection*{שאלה 5ב}

\begin{question}
תנו הגדרה של פונקציה קדומה לפונקציה $f(x)$. יש להוסיף דוגמה.
\end{question}

\begin{solution}
\textbf{הגדרה:} תהי $f$ מוגדרת על קטע $I$. פונקציה $F$ נקראת \textbf{פונקציה קדומה} של $f$ על $I$ אם $F$ גזירה על $I$ ומתקיים $F'(x)=f(x)$ לכל $x\in I$.
(אם $F$ קדומה של $f$ על קטע, אז כל הקדומות הן מהצורה $F+C$, $C$ קבוע.)

\textbf{דוגמה:} $F(x)=\frac{x^3}{3}$ היא פונקציה קדומה של $f(x)=x^2$ על $\R$, כי $\left(\frac{x^3}{3}\right)'=x^2$; גם $\frac{x^3}{3}+5$ קדומה שלה. דוגמה נוספת: $-\cos x$ קדומה של $\sin x$.
\end{solution}

\subsection*{שאלה 6א}

\begin{question}
רשמו את פולינום מקלורן מסדר 2 לפונקציה $y=f(x)$ המוגדרת בצורה סתומה: $x^2+x+e^x+2y=0$.
\end{question}

\begin{solution}
נמצא $y(0)$, $y'(0)$, $y''(0)$ בגזירה סתומה.
\begin{itemize}
  \item $x=0$: $0+0+1+2y(0)=0\Rightarrow y(0)=-\frac12$.
  \item גזירה: $2x+1+e^x+2y'=0\Rightarrow y'=-\frac{2x+1+e^x}{2}$, ולכן $y'(0)=-\frac{2}{2}=-1$.
  \item גזירה נוספת: $2+e^x+2y''=0\Rightarrow y''=-\frac{2+e^x}{2}$, ולכן $y''(0)=-\frac32$.
\end{itemize}
פולינום מקלורן מסדר 2:
\[ P_2(x)=y(0)+y'(0)x+\frac{y''(0)}{2}x^2=\boxed{-\frac12-x-\frac34x^2}. \]
(בדיקה: כאן אפשר גם לבודד $y=-\frac12(x^2+x+e^x)$ ולהציב $e^x=1+x+\frac{x^2}{2}+\dots$, ומתקבל אותו פולינום.)
\end{solution}

\subsection*{שאלה 6ב}

\begin{question}
האם למשוואה $e^x-\sin x=0$ אין פתרונות ממשיים? נמק
\end{question}

\begin{hint}
בדקו מה קורה עבור $x<0$: $e^x$ קטן מאוד, ו-$\sin x$ מגיע ל-$1$. נסו את משפט ערך הביניים על קטע מתאים.
\end{hint}

\begin{solution}
\textbf{לא} – למשוואה יש פתרונות ממשיים. נגדיר $g(x)=e^x-\sin x$, רציפה על $\R$.
\[ g(-\pi)=e^{-\pi}-0>0,\qquad g\left(-\tfrac{3\pi}{2}\right)=e^{-3\pi/2}-\sin\left(-\tfrac{3\pi}{2}\right)=e^{-3\pi/2}-1<0, \]
כי $e^{-3\pi/2}<1$. לפי \textbf{משפט ערך הביניים} קיים $c\in\left(-\frac{3\pi}2,-\pi\right)$ עם $g(c)=0$, כלומר $e^c=\sin c$ (מספרית $c\approx-3.183$). באותו אופן, בכל קטע $\left(-\frac{3\pi}{2}-2\pi k,\,-\pi-2\pi k\right)$ יש פתרון, כך שיש אינסוף פתרונות (כולם שליליים: עבור $x\ge0$, $e^x\ge1\ge\sin x$ עם שוויון בלתי אפשרי בו-זמנית, ולכן אין פתרונות אי-שליליים).
\end{solution}

\section*{תשפ"ד סמסטר א מועד א}

\subsection*{שאלה 1}

\begin{question}
חשבו את הגבול הבא:
\[ \lim_{x\to\infty}\left(\frac{x^2+x+2}{x^2+x+1}\right)^{2x^2+4} \]
\end{question}

\begin{hint}
זהו גבול מהצורה $1^\infty$. כתבו את הבסיס בצורה $1+\frac{1}{x^2+x+1}$ והשתמשו בגבול $\lim_{t\to\infty}\left(1+\frac1t\right)^t=e$.
\end{hint}

\begin{solution}
כאשר $x\to\infty$ הבסיס שואף ל-$1$ והמעריך שואף ל-$\infty$, כלומר זהו ביטוי לא מוגדר מהצורה $1^\infty$.
נכתוב את הבסיס בצורה
\[ \frac{x^2+x+2}{x^2+x+1} = 1+\frac{1}{x^2+x+1}. \]
נסמן $t=x^2+x+1$; כאשר $x\to\infty$ גם $t\to\infty$. אז
\[ \left(1+\frac{1}{x^2+x+1}\right)^{2x^2+4} = \left[\left(1+\frac1t\right)^{t}\right]^{\frac{2x^2+4}{x^2+x+1}}. \]
הביטוי שבסוגריים המרובעים שואף ל-$e$ (הגבול המפורסם $\lim_{t\to\infty}(1+\frac1t)^t=e$), והמעריך מקיים
\[ \lim_{x\to\infty}\frac{2x^2+4}{x^2+x+1} = \lim_{x\to\infty}\frac{2+\frac{4}{x^2}}{1+\frac1x+\frac1{x^2}} = 2. \]
מרציפות הפונקציה $(u,v)\mapsto u^v=e^{v\ln u}$ עבור $u>0$ (או: לפי $e^{v\ln u}$ ורציפות $\ln$ ו-$\exp$) נקבל
\[ \lim_{x\to\infty}\left(\frac{x^2+x+2}{x^2+x+1}\right)^{2x^2+4} = e^{2}. \]

\textbf{דרך נוספת:} נכתוב את הביטוי כ-$e^{(2x^2+4)\ln\left(1+\frac{1}{x^2+x+1}\right)}$. מכיוון ש-$\ln(1+s)\sim s$ כאשר $s\to0$,
\[ (2x^2+4)\ln\left(1+\frac{1}{x^2+x+1}\right) \sim \frac{2x^2+4}{x^2+x+1}\xrightarrow[x\to\infty]{}2, \]
ושוב מרציפות האקספוננט מתקבל הגבול $\boxed{e^2}$.
\end{solution}

\subsection*{שאלה 2}

\begin{question}
יהיו $a,b\in\R$. נגדיר פונקציה על ידי
\[ f(x)=\begin{cases} \frac{\sin(8x)}{x} & x<0\\[2pt] ax+b & 0\le x\le 4\\[2pt] \frac{x^2-3x-4}{2-\sqrt{x}} & x>4 \end{cases} \]
לאילו ערכים של $a,b$ הפונקציה רציפה ב-$\R$? נמקו.
\end{question}

\begin{hint}
בכל אחד משלושת התחומים הפתוחים הפונקציה רציפה כהרכבה/מנה של פונקציות רציפות. נותר לבדוק את נקודות התפר $x=0$ ו-$x=4$.
\end{hint}

\begin{hint}
בגבול ב-$x=4$ פרקו את המונה $x^2-3x-4=(x-4)(x+1)$ וכתבו $x-4=(\sqrt x-2)(\sqrt x+2)$.
\end{hint}

\begin{solution}
\textbf{רציפות בתוך כל תחום.}
\begin{itemize}
  \item בתחום $x<0$: $f(x)=\frac{\sin 8x}{x}$ היא מנה של פונקציות רציפות שהמכנה שלה אינו מתאפס, ולכן רציפה.
  \item בתחום $0<x<4$: $f(x)=ax+b$ פולינום, רציף.
  \item בתחום $x>4$: $f(x)=\frac{x^2-3x-4}{2-\sqrt x}$ מנה של פונקציות רציפות, והמכנה מתאפס רק ב-$x=4$ שאינו בתחום. לכן רציפה.
\end{itemize}
לכן $f$ רציפה ב-$\R$ אם ורק אם היא רציפה בנקודות $x=0$ ו-$x=4$.

\textbf{הנקודה $x=0$.} $f(0)=b$, ומימין $\lim_{x\to0^+}f(x)=\lim_{x\to0^+}(ax+b)=b$. משמאל, לפי הגבול היסודי $\lim_{t\to0}\frac{\sin t}{t}=1$:
\[ \lim_{x\to0^-}\frac{\sin(8x)}{x} = \lim_{x\to0^-}8\cdot\frac{\sin(8x)}{8x} = 8. \]
לכן רציפות ב-$0$ מתקיימת אם ורק אם $b=8$.

\textbf{הנקודה $x=4$.} $f(4)=4a+b$, ומשמאל $\lim_{x\to4^-}(ax+b)=4a+b$. מימין, עבור $x>4$:
\[ \frac{x^2-3x-4}{2-\sqrt x} = \frac{(x-4)(x+1)}{2-\sqrt x} = \frac{(\sqrt x-2)(\sqrt x+2)(x+1)}{-(\sqrt x-2)} = -(\sqrt x+2)(x+1), \]
ולכן
\[ \lim_{x\to4^+}f(x) = -(2+2)(4+1) = -20. \]
רציפות ב-$4$ מתקיימת אם ורק אם $4a+b=-20$.

\textbf{מסקנה.} עם $b=8$ נקבל $4a=-28$, כלומר
\[ \boxed{a=-7,\quad b=8}, \]
ורק עבור ערכים אלה $f$ רציפה בכל $\R$.
\end{solution}

\subsection*{שאלה 3}

\begin{question}
חשבו תחומי עלייה וירידה וקיצון מוחלט של הפונקציה
\[ f(x)=\sqrt{2x-x^2} \]
בתחום הגדרתה.
\end{question}

\begin{solution}
\textbf{תחום הגדרה.} יש לדרוש $2x-x^2=x(2-x)\ge0$, כלומר $0\le x\le 2$. תחום ההגדרה הוא הקטע הסגור $[0,2]$.

\textbf{נגזרת.} בקטע הפתוח $(0,2)$ הביטוי שבשורש חיובי, ולכן לפי כלל השרשרת
\[ f'(x)=\frac{2-2x}{2\sqrt{2x-x^2}}=\frac{1-x}{\sqrt{2x-x^2}},\qquad 0<x<2. \]
(בנקודות $x=0,2$ הפונקציה אינה גזירה – הנגזרת החד-צדדית אינסופית – אבל הן נקודות קצה.)

\textbf{סימן הנגזרת.} המכנה חיובי, ולכן סימן $f'$ הוא סימן $1-x$:
\begin{itemize}
  \item עבור $0<x<1$: $f'(x)>0$, ולכן $f$ עולה ממש בקטע $[0,1]$.
  \item עבור $1<x<2$: $f'(x)<0$, ולכן $f$ יורדת ממש בקטע $[1,2]$.
\end{itemize}
(המעבר לקטעים הסגורים מוצדק כי $f$ רציפה בהם – מסקנה ממשפט לגרנז'.)

\textbf{קיצון מוחלט.} $f$ רציפה בקטע הסגור $[0,2]$, ולכן לפי משפט ויירשטראס היא מקבלת בו מקסימום ומינימום. הנקודות החשודות הן נקודת הקריטית $x=1$ ונקודות הקצה:
\[ f(0)=0,\qquad f(1)=\sqrt{2-1}=1,\qquad f(2)=0. \]
לכן
\[ \boxed{\max_{[0,2]}f=f(1)=1,\qquad \min_{[0,2]}f=f(0)=f(2)=0.} \]
(אכן $f\ge0$ תמיד כשורש, כך שהמינימום $0$ ברור גם ישירות. בנוסף, $y=\sqrt{2x-x^2}$ שקולה ל-$(x-1)^2+y^2=1,\ y\ge0$ – חצי המעגל העליון שמרכזו $(1,0)$ ורדיוסו $1$, מה שמאשר את התוצאות.)
\end{solution}

\subsection*{שאלה 4}

\begin{question}
העריכו בעזרת קירוב לינארי את $\sqrt[3]{30}$, ותנו חסם לשגיאה.
\end{question}

\begin{hint}
השתמשו ב-$f(x)=\sqrt[3]{x}$ סביב הנקודה $x_0=27$, שבה השורש ידוע.
\end{hint}

\begin{hint}
לחסם השגיאה השתמשו בשארית לגרנז' של פולינום טיילור מסדר 1: $R_1(x)=\frac{f''(c)}{2}(x-x_0)^2$.
\end{hint}

\begin{solution}
נבחר $f(x)=\sqrt[3]{x}=x^{1/3}$ ו-$x_0=27$ (נקודה קרובה ל-$30$ שבה $f(27)=3$).

\textbf{הקירוב הלינארי.} $f'(x)=\frac13x^{-2/3}$, ולכן $f'(27)=\frac{1}{3\cdot 9}=\frac1{27}$. הקירוב הלינארי (פולינום טיילור מסדר 1):
\[ f(x)\approx f(27)+f'(27)(x-27)=3+\frac{x-27}{27}, \]
ובפרט
\[ \sqrt[3]{30}\approx 3+\frac{3}{27}=3+\frac19=\frac{28}{9}\approx 3.1111. \]

\textbf{חסם לשגיאה.} לפי משפט טיילור עם שארית לגרנז', קיימת $c\in(27,30)$ כך ש-
\[ \sqrt[3]{30}-\frac{28}{9} = R_1(30) = \frac{f''(c)}{2!}(30-27)^2. \]
$f''(x)=-\frac29x^{-5/3}$, ולכן
\[ |R_1(30)| = \frac{2}{9}c^{-5/3}\cdot\frac{9}{2} = \frac{1}{c^{5/3}} < \frac{1}{27^{5/3}}=\frac{1}{3^5}=\frac1{243}\approx 0.0041, \]
כי $c>27$ והפונקציה $c^{-5/3}$ יורדת.

בנוסף, $f''<0$ ולכן $R_1<0$: הקירוב הוא הערכת-יתר. לסיכום
\[ \boxed{\sqrt[3]{30}\approx\frac{28}{9}\approx3.111,\qquad 0<\frac{28}{9}-\sqrt[3]{30}<\frac{1}{243}\approx0.0041.} \]
(לבדיקה: $\sqrt[3]{30}\approx3.1072$, והשגיאה בפועל כ-$0.0039$.)
\end{solution}

\subsection*{שאלה 5}

\begin{question}
חשבו את הפונקציה הקדומה של
\[ f(x)=\frac{4x^2+9x+10}{x^3+2x^2+5x} \]
\end{question}

\begin{hint}
פרקו את המכנה: $x^3+2x^2+5x=x(x^2+2x+5)$, כאשר לגורם הריבועי אין שורשים ממשיים. פרקו לשברים חלקיים מהצורה $\frac{A}{x}+\frac{Bx+C}{x^2+2x+5}$.
\end{hint}

\begin{hint}
באינטגרל של $\frac{Bx+C}{x^2+2x+5}$ פצלו למונה שהוא נגזרת המכנה ועוד קבוע, והשלימו לריבוע: $x^2+2x+5=(x+1)^2+4$.
\end{hint}

\begin{solution}
\textbf{פירוק המכנה.} $x^3+2x^2+5x=x(x^2+2x+5)$, והדיסקרימיננטה של $x^2+2x+5$ היא $4-20<0$, כך שהגורם הריבועי אי-פריק. מעלת המונה קטנה ממעלת המכנה, לכן ניתן לפרק ישירות לשברים חלקיים:
\[ \frac{4x^2+9x+10}{x(x^2+2x+5)}=\frac{A}{x}+\frac{Bx+C}{x^2+2x+5}. \]
כפל במכנה: $4x^2+9x+10=A(x^2+2x+5)+(Bx+C)x$.
\begin{itemize}
  \item הצבת $x=0$: $10=5A$, כלומר $A=2$.
  \item השוואת מקדמי $x^2$: $4=A+B$, כלומר $B=2$.
  \item השוואת מקדמי $x$: $9=2A+C$, כלומר $C=5$.
\end{itemize}
לכן
\[ f(x)=\frac{2}{x}+\frac{2x+5}{x^2+2x+5}. \]

\textbf{אינטגרציה.} $\int\frac{2}{x}\,dx=2\ln|x|$. בשבר השני נפצל $2x+5=(2x+2)+3$, כאשר $2x+2$ היא נגזרת המכנה:
\[ \int\frac{2x+5}{x^2+2x+5}\,dx=\int\frac{2x+2}{x^2+2x+5}\,dx+3\int\frac{dx}{(x+1)^2+4}. \]
האינטגרל הראשון שווה $\ln(x^2+2x+5)$ (הצבה $u=x^2+2x+5$; ללא ערך מוחלט כי הביטוי חיובי). בשני נציב $x+1=2t$:
\[ 3\int\frac{dx}{(x+1)^2+4}=\frac32\arctan\frac{x+1}{2}. \]

\textbf{תשובה.}
\[ \boxed{\int f(x)\,dx = 2\ln|x|+\ln(x^2+2x+5)+\frac32\arctan\frac{x+1}{2}+C} \]
(בכל אחד מהקטעים $x>0$ ו-$x<0$ בנפרד; ניתן לבדוק בגזירה שמתקבלת $f$.)
\end{solution}

\subsection*{שאלה 6}

\begin{question}
חשבו את האינטגרל המסויים
\[ \int_0^{\pi}x\sin^2x\,dx \]
\end{question}

\begin{hint}
השתמשו בזהות $\sin^2x=\frac{1-\cos 2x}{2}$ ואחר כך באינטגרציה בחלקים עבור $\int x\cos2x\,dx$.
\end{hint}

\begin{solution}
לפי הזהות $\sin^2x=\frac{1-\cos2x}{2}$:
\[ \int_0^{\pi}x\sin^2x\,dx=\frac12\int_0^{\pi}x\,dx-\frac12\int_0^{\pi}x\cos2x\,dx. \]
\textbf{האינטגרל הראשון:} $\frac12\int_0^\pi x\,dx=\frac12\cdot\frac{\pi^2}{2}=\frac{\pi^2}{4}$.

\textbf{האינטגרל השני} – אינטגרציה בחלקים עם $u=x$, $dv=\cos2x\,dx$, כלומר $du=dx$, $v=\frac12\sin2x$:
\[ \int_0^{\pi}x\cos2x\,dx=\left[\frac{x\sin2x}{2}\right]_0^{\pi}-\frac12\int_0^{\pi}\sin2x\,dx = 0+\left[\frac{\cos2x}{4}\right]_0^{\pi}=\frac14-\frac14=0. \]
לכן
\[ \boxed{\int_0^{\pi}x\sin^2x\,dx=\frac{\pi^2}{4}} \]

\textbf{בדיקה (דרך נוספת):} בהצבה $x=\pi-t$ מתקבל $I=\int_0^\pi(\pi-t)\sin^2t\,dt=\pi\int_0^\pi\sin^2t\,dt-I$, ולכן $I=\frac\pi2\cdot\frac\pi2=\frac{\pi^2}{4}$.
\end{solution}

\subsection*{שאלה 7}

\begin{question}
חשבו את נפח גוף הסיבוב של הפונקציה הבאה, סביב ציר ה-$x$, בקטע $[0,\pi]$:
\[ f(x)=\sqrt{\sin x}\cdot\cos^2x \]
\end{question}

\begin{hint}
נפח גוף סיבוב סביב ציר $x$ הוא $V=\pi\int_a^b f^2(x)\,dx$. כאן $f^2(x)=\sin x\cos^4x$, ומתאימה ההצבה $u=\cos x$.
\end{hint}

\begin{solution}
בקטע $[0,\pi]$ מתקיים $\sin x\ge0$, כך ש-$f$ מוגדרת ורציפה בכל הקטע. נוסחת נפח גוף הסיבוב סביב ציר ה-$x$:
\[ V=\pi\int_0^{\pi}f^2(x)\,dx=\pi\int_0^{\pi}\sin x\cos^4x\,dx. \]
נציב $u=\cos x$, $du=-\sin x\,dx$; כאשר $x=0$: $u=1$, וכאשר $x=\pi$: $u=-1$. לכן
\[ \int_0^{\pi}\sin x\cos^4x\,dx=\int_{1}^{-1}u^4\,(-du)=\int_{-1}^{1}u^4\,du=\left[\frac{u^5}{5}\right]_{-1}^{1}=\frac25. \]
ומכאן
\[ \boxed{V=\frac{2\pi}{5}} \]
\end{solution}

\section*{תשפ"ד סמסטר א מועד ב}

\subsection*{שאלה 1}

\begin{question}
חשבו את הגבול הבא:
\[ \lim_{x\to 1}\frac{(x^2+\ln x-1)\cdot\sin(x-1)}{(x-1)(e^x-e)} \]
\end{question}

\begin{hint}
פצלו את הביטוי למכפלה של $\frac{\sin(x-1)}{x-1}$ ושל $\frac{x^2+\ln x-1}{e^x-e}$, וחשבו כל גבול בנפרד.
\end{hint}

\begin{solution}
כאשר $x\to1$ המונה והמכנה שואפים שניהם ל-$0$. נכתוב את הביטוי כמכפלה (עבור $x$ קרוב ל-$1$, $x\neq1$, $x>0$):
\[ \frac{(x^2+\ln x-1)\sin(x-1)}{(x-1)(e^x-e)}=\frac{\sin(x-1)}{x-1}\cdot\frac{x^2+\ln x-1}{e^x-e}. \]
\textbf{הגורם הראשון:} נציב $t=x-1\to0$, ומהגבול היסודי $\lim_{t\to0}\frac{\sin t}{t}=1$ נקבל $\lim_{x\to1}\frac{\sin(x-1)}{x-1}=1$.

\textbf{הגורם השני:} המונה $x^2+\ln x-1\to 1+0-1=0$ והמכנה $e^x-e\to0$, כלומר זהו מקרה $\frac00$. שתי הפונקציות גזירות בסביבת $1$ ונגזרת המכנה $e^x\neq0$, ולכן לפי כלל לופיטל
\[ \lim_{x\to1}\frac{x^2+\ln x-1}{e^x-e}=\lim_{x\to1}\frac{2x+\frac1x}{e^x}=\frac{2+1}{e}=\frac3e, \]
בתנאי שהגבול מימין קיים -- והוא אכן קיים (הפונקציה $\frac{2x+1/x}{e^x}$ רציפה ב-$1$).

\textbf{סיכום:} לפי אריתמטיקה של גבולות (מכפלת שני גבולות סופיים),
\[ \lim_{x\to 1}\frac{(x^2+\ln x-1)\sin(x-1)}{(x-1)(e^x-e)}=1\cdot\frac3e=\boxed{\frac3e}. \]
\end{solution}

\subsection*{שאלה 2}

\begin{question}
נגדיר פונקציה על ידי
\[ f(x)=\begin{cases}\frac{\ln(1+x^2)}{\sin^2 x} & x<0\\ 0 & 0\le x\le 1\\ \frac{x}{x^2-1} & x>1\end{cases} \]
מצאו את כל נקודות אי הרציפות של הפונקציה, ומיינו אותן לאי רציפות מסוג סליקה, קפיצה או עיקרית. נמקו.
\end{question}

\begin{hint}
בתוך כל אחד מהתחומים הפונקציה היא הרכבה/מנה של פונקציות אלמנטריות. בדקו היכן מתאפס מכנה, ובנפרד את נקודות התפר $x=0$ ו-$x=1$.
\end{hint}

\begin{hint}
ליד $0$: $\ln(1+x^2)\approx x^2$ ו-$\sin^2x\approx x^2$.
\end{hint}

\begin{solution}
\textbf{בתוך התחומים הפתוחים.}
\begin{itemize}
  \item בקטע $(0,1)$ הפונקציה קבועה ולכן רציפה.
  \item בקטע $(1,\infty)$: $f(x)=\frac{x}{x^2-1}$ היא מנה של פולינומים שהמכנה שלה אינו מתאפס שם ($x^2-1>0$), ולכן רציפה.
  \item בקטע $(-\infty,0)$: $f(x)=\frac{\ln(1+x^2)}{\sin^2x}$ היא מנה של פונקציות רציפות, ולכן רציפה בכל נקודה שבה $\sin x\neq0$. המכנה מתאפס בנקודות $x=-k\pi$, $k=1,2,3,\dots$, ושם הפונקציה אינה מוגדרת. בנקודות אלה המונה שואף ל-$\ln(1+k^2\pi^2)>0$ והמכנה שואף ל-$0^+$ (כי $\sin^2x>0$ בסביבה מנוקבת), ולכן
  \[ \lim_{x\to-k\pi^\pm}f(x)=+\infty. \]
  הגבולות החד-צדדיים אינם סופיים, ולכן $x=-k\pi$ ($k\in\N$) הן נקודות אי-רציפות \textbf{עיקריות} (מסוג שני).
\end{itemize}

\textbf{הנקודה $x=0$.} $f(0)=0$, ומימין $f\equiv0$ ולכן $\lim_{x\to0^+}f(x)=0$. משמאל:
\[ \lim_{x\to0^-}\frac{\ln(1+x^2)}{\sin^2x}=\lim_{x\to0^-}\frac{\ln(1+x^2)}{x^2}\cdot\left(\frac{x}{\sin x}\right)^2=1\cdot1^2=1, \]
כאשר השתמשנו בגבול היסודי $\lim_{t\to0}\frac{\ln(1+t)}{t}=1$ (עם $t=x^2\to0$) וב-$\lim_{x\to0}\frac{\sin x}{x}=1$. שני הגבולות החד-צדדיים סופיים ושונים ($1\neq0$), ולכן ב-$x=0$ יש אי-רציפות מסוג \textbf{קפיצה}.

\textbf{הנקודה $x=1$.} $f(1)=0$ ו-$\lim_{x\to1^-}f(x)=0$. מימין, המונה שואף ל-$1$ והמכנה $x^2-1\to0^+$, לכן
\[ \lim_{x\to1^+}\frac{x}{x^2-1}=+\infty. \]
הגבול החד-צדדי מימין אינו סופי, ולכן ב-$x=1$ יש אי-רציפות \textbf{עיקרית}.

\textbf{תשובה:} ב-$x=0$ -- אי-רציפות מסוג קפיצה; ב-$x=1$ -- אי-רציפות עיקרית; בנקודות $x=-k\pi$, $k=1,2,\dots$ (שבהן $f$ אינה מוגדרת) -- אי-רציפות עיקרית (הגבול שם הוא $+\infty$). בכל שאר הנקודות $f$ רציפה.
\end{solution}

\subsection*{שאלה 3}

\begin{question}
חשבו תחומי קמירות כלפי מעלה וקמירות כלפי מטה (קמירות וקעירות), ונקודות פיתול של הפונקציה
\[ f(x)=\ln\left(\frac{e^x}{1+e^x}\right) \]
בתחום הגדרתה.
\end{question}

\begin{hint}
פשטו קודם: $\ln\frac{e^x}{1+e^x}=x-\ln(1+e^x)$.
\end{hint}

\begin{solution}
\textbf{תחום הגדרה.} לכל $x$ ממשי $\frac{e^x}{1+e^x}>0$, ולכן $f$ מוגדרת על כל $\R$. לפי חוקי לוגריתמים
\[ f(x)=\ln e^x-\ln(1+e^x)=x-\ln(1+e^x). \]

\textbf{נגזרת ראשונה.}
\[ f'(x)=1-\frac{e^x}{1+e^x}=\frac{1}{1+e^x}. \]

\textbf{נגזרת שנייה.} לפי כלל השרשרת,
\[ f''(x)=-\frac{e^x}{(1+e^x)^2}. \]
לכל $x\in\R$ מתקיים $e^x>0$ ו-$(1+e^x)^2>0$, ולכן $f''(x)<0$ לכל $x$.

\textbf{מסקנה.} מאחר ש-$f''<0$ על כל $\R$, הפונקציה קמורה כלפי מטה (קעורה, $\cap$) על כל הישר $(-\infty,\infty)$, ואין תחום שבו היא קמורה כלפי מעלה. נקודת פיתול היא נקודה שבה משתנה כיוון הקמירות; מאחר ש-$f''$ אינה מחליפה סימן (ואף אינה מתאפסת), \textbf{אין לפונקציה נקודות פיתול}.
\end{solution}

\subsection*{שאלה 4}

\begin{question}
העריכו בעזרת קירוב טיילור מסדר שני את $\sqrt[3]{e}$, ותנו חסם לשגיאה.
\end{question}

\begin{hint}
השתמשו בפולינום מקלורן של $e^x$ מסדר $2$ בנקודה $x=\frac13$, ובשארית בצורת לגרנז'.
\end{hint}

\begin{solution}
נשתמש בפונקציה $f(x)=e^x$ סביב $x_0=0$, ונציב $x=\frac13$, שכן $\sqrt[3]{e}=e^{1/3}$.

\textbf{פולינום טיילור.} לכל $n$, $f^{(n)}(x)=e^x$ ולכן $f^{(n)}(0)=1$. פולינום מקלורן מסדר $2$:
\[ P_2(x)=1+x+\frac{x^2}{2}. \]
לכן
\[ \sqrt[3]{e}\approx P_2\!\left(\tfrac13\right)=1+\frac13+\frac1{18}=\frac{25}{18}\approx1.3889. \]

\textbf{חסם לשגיאה.} לפי משפט טיילור עם שארית לגרנז', קיימת $c$ בין $0$ ל-$\frac13$ כך ש-
\[ R_2\!\left(\tfrac13\right)=e^{1/3}-P_2\!\left(\tfrac13\right)=\frac{f'''(c)}{3!}\left(\frac13\right)^3=\frac{e^c}{6\cdot27}=\frac{e^c}{162}. \]
מאחר ש-$e^x$ עולה ו-$0<c<\frac13$, מתקיים $e^c<e^{1/3}$. כיוון ש-$e<3<\frac{27}{8}=1.5^3$, נקבל $e^{1/3}<1.5$. לכן
\[ 0<R_2\!\left(\tfrac13\right)<\frac{1.5}{162}=\frac1{108}\approx0.0093. \]
(אפשר גם להשתמש בחסם הגס יותר $e^c<e<3$ ולקבל שגיאה קטנה מ-$\frac{3}{162}=\frac1{54}$.)

\textbf{תשובה:} $\sqrt[3]{e}\approx\frac{25}{18}\approx1.3889$, והשגיאה חיובית וקטנה מ-$\frac1{108}$. (לבדיקה: $\sqrt[3]{e}=1.39561\ldots$, והשגיאה בפועל היא כ-$0.0067$.)
\end{solution}

\subsection*{שאלה 5}

\begin{question}
חשבו את הפונקציה הקדומה של
\[ f(x)=\frac{5x^3+x^2+x-2}{x^4+x^2} \]
\end{question}

\begin{hint}
פרקו את המכנה: $x^4+x^2=x^2(x^2+1)$, והשתמשו בפירוק לשברים חלקיים מהצורה $\frac Ax+\frac B{x^2}+\frac{Cx+D}{x^2+1}$.
\end{hint}

\begin{solution}
מעלת המונה ($3$) קטנה ממעלת המכנה ($4$), ולכן אין צורך בחילוק פולינומים. נפרק את המכנה: $x^4+x^2=x^2(x^2+1)$, כאשר $x^2+1$ אי-פריק מעל $\R$. נחפש פירוק לשברים חלקיים:
\[ \frac{5x^3+x^2+x-2}{x^2(x^2+1)}=\frac Ax+\frac B{x^2}+\frac{Cx+D}{x^2+1}. \]
נכפול ב-$x^2(x^2+1)$:
\[ 5x^3+x^2+x-2=Ax(x^2+1)+B(x^2+1)+(Cx+D)x^2=(A+C)x^3+(B+D)x^2+Ax+B. \]
השוואת מקדמים: $B=-2$, $A=1$, $A+C=5\Rightarrow C=4$, $B+D=1\Rightarrow D=3$. לכן
\[ f(x)=\frac1x-\frac2{x^2}+\frac{4x+3}{x^2+1}=\frac1x-\frac2{x^2}+2\cdot\frac{2x}{x^2+1}+\frac{3}{x^2+1}. \]
נאנטגרל איבר-איבר בעזרת האינטגרלים המיידיים $\int\frac{dx}x=\ln|x|$, $\int x^{-2}dx=-x^{-1}$, $\int\frac{2x}{x^2+1}dx=\ln(x^2+1)$ (הצבה $t=x^2+1$) ו-$\int\frac{dx}{x^2+1}=\arctan x$:
\[ \boxed{\int f(x)\,dx=\ln|x|+\frac2x+2\ln(x^2+1)+3\arctan x+C} \]
(בכל אחד מהקטעים $x>0$ ו-$x<0$ בנפרד; הקבוע יכול להיות שונה בכל קטע).

\textbf{בדיקה:} גזירה נותנת $\frac1x-\frac2{x^2}+\frac{4x}{x^2+1}+\frac3{x^2+1}=f(x)$.
\end{solution}

\subsection*{שאלה 6}

\begin{question}
חשבו את האינטגרל המסויים
\[ \int_0^1 (x+1)^2\cdot e^x\,dx \]
\end{question}

\begin{hint}
אינטגרציה בחלקים פעמיים, כאשר בכל פעם גוזרים את הפולינום.
\end{hint}

\begin{solution}
נשתמש באינטגרציה בחלקים: $\int u\,v'=uv-\int u'v$.

\textbf{שלב ראשון:} $u=(x+1)^2$, $v'=e^x$, ולכן $u'=2(x+1)$, $v=e^x$:
\[ \int_0^1(x+1)^2e^x\,dx=\Big[(x+1)^2e^x\Big]_0^1-2\int_0^1(x+1)e^x\,dx=(4e-1)-2\int_0^1(x+1)e^x\,dx. \]

\textbf{שלב שני:} $u=x+1$, $v'=e^x$, ולכן $u'=1$, $v=e^x$:
\[ \int_0^1(x+1)e^x\,dx=\Big[(x+1)e^x\Big]_0^1-\int_0^1e^x\,dx=(2e-1)-(e-1)=e. \]

\textbf{סיכום:}
\[ \int_0^1(x+1)^2e^x\,dx=4e-1-2e=\boxed{2e-1}\approx 4.4366. \]
(לחלופין: פונקציה קדומה היא $F(x)=e^x\big((x+1)^2-2(x+1)+2\big)=e^x(x^2+1)$, ולפי המשפט היסודי $F(1)-F(0)=2e-1$.)
\end{solution}

\subsection*{שאלה 7}

\begin{question}
בדקו התכנסות של האיטגרל הלא-אמיתי הבא:
\[ \int_2^\infty \frac{\sqrt[5]{x}}{3x^2+x+100}\,dx \]
\end{question}

\begin{hint}
לאינסוף, הפונקציה מתנהגת כמו $\frac{x^{1/5}}{3x^2}$. השוו (במבחן ההשוואה הגבולי) ל-$\frac{1}{x^{9/5}}$.
\end{hint}

\begin{solution}
הפונקציה $f(x)=\frac{\sqrt[5]{x}}{3x^2+x+100}$ רציפה וחיובית על $[2,\infty)$ (המכנה חיובי), ולכן האינטגרל לא-אמיתי רק בגלל הגבול העליון, ומותר להשתמש במבחני ההשוואה לפונקציות אי-שליליות.

נשווה ל-$g(x)=\frac{1}{x^{9/5}}=\frac{x^{1/5}}{x^2}$, שגם היא חיובית על $[2,\infty)$:
\[ \lim_{x\to\infty}\frac{f(x)}{g(x)}=\lim_{x\to\infty}\frac{x^{1/5}\cdot x^{2}}{x^{1/5}(3x^2+x+100)}=\lim_{x\to\infty}\frac{1}{3+\frac1x+\frac{100}{x^2}}=\frac13. \]
הגבול סופי וחיובי, ולכן לפי \textbf{מבחן ההשוואה הגבולי} האינטגרלים $\int_2^\infty f$ ו-$\int_2^\infty g$ מתכנסים או מתבדרים יחד.

האינטגרל $\int_2^\infty\frac{dx}{x^p}$ מתכנס אם ורק אם $p>1$; כאן $p=\frac95>1$, ובאופן מפורש
\[ \int_2^\infty x^{-9/5}dx=\lim_{R\to\infty}\left[-\frac54x^{-4/5}\right]_2^R=\frac54\cdot2^{-4/5}<\infty. \]

\textbf{תשובה:} האינטגרל $\int_2^\infty\frac{\sqrt[5]{x}}{3x^2+x+100}\,dx$ \textbf{מתכנס}.
\end{solution}

\section*{תשפ"ד סמסטר א מועד מיוחד}

\subsection*{שאלה 1}

\begin{question}
חשבו את הגבול הבא:
\[ \lim_{x\to0^+}\frac{\ln(e^x-1)}{3\cdot\ln x} \]
\end{question}

\begin{hint}
זהו מקרה $\frac{\infty}{\infty}$; השתמשו בכלל לופיטל.
\end{hint}

\begin{solution}
כאשר $x\to0^+$: $e^x-1\to0^+$ ולכן $\ln(e^x-1)\to-\infty$, וגם $3\ln x\to-\infty$. זהו מקרה $\frac{\infty}{\infty}$. המונה והמכנה גזירים ב-$(0,\infty)$ ונגזרת המכנה $\frac3x\neq0$, לכן לפי כלל לופיטל (בתנאי שהגבול מימין קיים):
\[ \lim_{x\to0^+}\frac{\ln(e^x-1)}{3\ln x}=\lim_{x\to0^+}\frac{\frac{e^x}{e^x-1}}{\frac3x}=\lim_{x\to0^+}\frac{e^x}{3}\cdot\frac{x}{e^x-1}. \]
מתקיים $e^x\to1$ ומהגבול היסודי $\lim_{x\to0}\frac{e^x-1}{x}=1$ נקבל $\frac{x}{e^x-1}\to1$. לפי אריתמטיקה של גבולות הגבול מימין קיים ושווה $\frac13\cdot1=\frac13$, ולכן
\[ \lim_{x\to0^+}\frac{\ln(e^x-1)}{3\ln x}=\boxed{\frac13}. \]
\end{solution}

\subsection*{שאלה 2}

\begin{question}
נגדיר פונקציה על ידי
\[ f(x)=\begin{cases}2^{\frac{\sin x}{x}} & x<0\\ ax+b & 0\le x\le 1\\ \frac{x^2+x-2}{\sqrt{x}-1} & x>1\end{cases} \]
לאילו ערכים של הפרמטרים $a,b$ הפונקציה רציפה בכל $\R$? נמקו.
\end{question}

\begin{hint}
בתוך כל תחום הפונקציה רציפה; נשאר לדרוש רציפות בנקודות התפר $x=0$ ו-$x=1$.
\end{hint}

\begin{hint}
עבור $x>1$: $x^2+x-2=(x-1)(x+2)$ ו-$x-1=(\sqrt x-1)(\sqrt x+1)$.
\end{hint}

\begin{solution}
\textbf{בתוך התחומים.} ב-$(-\infty,0)$ הפונקציה $\frac{\sin x}{x}$ רציפה (מנה של רציפות, מכנה שונה מ-$0$), ו-$2^t$ רציפה, ולכן ההרכבה רציפה. ב-$(0,1)$ זו פונקציה לינארית, רציפה. ב-$(1,\infty)$ זו מנה של פונקציות רציפות שהמכנה שלה $\sqrt x-1>0$, ולכן רציפה. לכן לכל $a,b$ הפונקציה רציפה בכל נקודה פרט אולי ל-$x=0$ ו-$x=1$.

\textbf{רציפות ב-$x=0$.} $f(0)=b$, ומימין $\lim_{x\to0^+}(ax+b)=b$. משמאל, מאחר ש-$\lim_{x\to0}\frac{\sin x}{x}=1$ ו-$2^t$ רציפה ב-$t=1$, לפי משפט הגבול של פונקציה מורכבת
\[ \lim_{x\to0^-}2^{\frac{\sin x}{x}}=2^1=2. \]
לכן $f$ רציפה ב-$0$ אם ורק אם $b=2$.

\textbf{רציפות ב-$x=1$.} $f(1)=a+b$ ו-$\lim_{x\to1^-}f(x)=a+b$. מימין, עבור $x>1$:
\[ \frac{x^2+x-2}{\sqrt x-1}=\frac{(x-1)(x+2)}{\sqrt x-1}=\frac{(\sqrt x-1)(\sqrt x+1)(x+2)}{\sqrt x-1}=(\sqrt x+1)(x+2), \]
ולכן $\lim_{x\to1^+}f(x)=(1+1)(1+2)=6$. לכן $f$ רציפה ב-$1$ אם ורק אם $a+b=6$.

\textbf{תשובה:} $f$ רציפה בכל $\R$ אם ורק אם $\boxed{b=2,\ a=4}$.
\end{solution}

\subsection*{שאלה 3}

\begin{question}
חשבו תחומי מינימום ומקסימום מוחלטים של הפונקציה
\[ f(x)=(x+2)\cdot e^x \]
בתחום הגדרתה, או הסבירו מדוע אינם קיימים.
\end{question}

\begin{hint}
מצאו את תחומי העלייה והירידה בעזרת $f'$, ובדקו את התנהגות הפונקציה כאשר $x\to\pm\infty$.
\end{hint}

\begin{solution}
\textbf{תחום הגדרה:} כל $\R$; $f$ גזירה בכל נקודה.

\textbf{נגזרת:} לפי כלל המכפלה
\[ f'(x)=e^x+(x+2)e^x=(x+3)e^x. \]
כיוון ש-$e^x>0$, סימן $f'$ הוא סימן $x+3$: $f'<0$ ב-$(-\infty,-3)$ ו-$f'>0$ ב-$(-3,\infty)$. לכן (לפי מסקנה ממשפט לגרנז') $f$ יורדת ממש ב-$(-\infty,-3]$ ועולה ממש ב-$[-3,\infty)$.

\textbf{מינימום מוחלט:} מהמונוטוניות, לכל $x\le-3$ מתקיים $f(x)\ge f(-3)$, ולכל $x\ge-3$ מתקיים $f(x)\ge f(-3)$. לכן $x=-3$ היא נקודת מינימום מוחלט, וערך המינימום הוא
\[ f(-3)=(-3+2)e^{-3}=-\frac1{e^3}\approx-0.0498. \]

\textbf{מקסימום מוחלט:} $\lim_{x\to\infty}(x+2)e^x=+\infty$, ולכן הפונקציה אינה חסומה מלעיל ואין לה מקסימום מוחלט. (לשלמות: $\lim_{x\to-\infty}(x+2)e^x=\lim_{x\to-\infty}\frac{x+2}{e^{-x}}=0$ לפי לופיטל, כך שהפונקציה שואפת ל-$0^-$ משמאל, אך מימין אינה חסומה.)

\textbf{תשובה:} מינימום מוחלט $-e^{-3}$ המתקבל ב-$x=-3$; מקסימום מוחלט אינו קיים כי $f(x)\to+\infty$ כאשר $x\to\infty$.
\end{solution}

\subsection*{שאלה 4}

\begin{question}
הוכיחו כי למשוואה הבאה יש פתרון ממשי חיובי יחיד:
\[ 2x^3+\arctan x+\sqrt{x}=7 \]
\end{question}

\begin{hint}
קיום: משפט ערך הביניים עבור $g(x)=2x^3+\arctan x+\sqrt x-7$ בקטע מתאים. יחידות: הראו ש-$g$ עולה ממש.
\end{hint}

\begin{solution}
נגדיר $g(x)=2x^3+\arctan x+\sqrt x-7$ על $[0,\infty)$. $g$ רציפה שם כסכום של פונקציות רציפות, וגזירה ב-$(0,\infty)$.

\textbf{קיום.} $g(0)=-7<0$, ו-
\[ g(2)=16+\arctan2+\sqrt2-7=9+\arctan 2+\sqrt2>0. \]
לפי משפט ערך הביניים קיימת $x_0\in(0,2)$ עם $g(x_0)=0$, כלומר פתרון חיובי של המשוואה.

\textbf{יחידות.} לכל $x>0$:
\[ g'(x)=6x^2+\frac1{1+x^2}+\frac1{2\sqrt x}>0. \]
לכן $g$ עולה ממש ב-$(0,\infty)$ (מסקנה ממשפט לגרנז': אם $0<x_1<x_2$ אז $g(x_2)-g(x_1)=g'(c)(x_2-x_1)>0$). פונקציה עולה ממש מקבלת כל ערך לכל היותר פעם אחת, ולכן יש לה לכל היותר אפס אחד ב-$(0,\infty)$.

(לחלופין: אם היו שני פתרונות חיוביים $x_1<x_2$, אז לפי משפט רול הייתה $c\in(x_1,x_2)$ עם $g'(c)=0$, בסתירה ל-$g'>0$.)

\textbf{מסקנה:} למשוואה יש פתרון ממשי חיובי יחיד (מספרית $x_0\approx1.3486$).
\end{solution}

\subsection*{שאלה 5}

\begin{question}
חשבו את הפונקציה הקדומה של
\[ f(x)=\sin^2x\cdot\cos^5x \]
\end{question}

\begin{hint}
החזקה של $\cos x$ אי-זוגית: הפרידו גורם אחד $\cos x$, כתבו $\cos^4x=(1-\sin^2x)^2$ והציבו $t=\sin x$.
\end{hint}

\begin{solution}
נכתוב $\cos^5x=\cos^4x\cdot\cos x=(1-\sin^2x)^2\cos x$. לכן
\[ \int\sin^2x\cos^5x\,dx=\int\sin^2x(1-\sin^2x)^2\cos x\,dx. \]
נציב $t=\sin x$, $dt=\cos x\,dx$:
\[ \int t^2(1-t^2)^2\,dt=\int\left(t^2-2t^4+t^6\right)dt=\frac{t^3}3-\frac{2t^5}5+\frac{t^7}7+C. \]
נחזור למשתנה $x$:
\[ \boxed{\int\sin^2x\cos^5x\,dx=\frac{\sin^3x}{3}-\frac{2\sin^5x}{5}+\frac{\sin^7x}{7}+C}. \]
\textbf{בדיקה:} הנגזרת היא $\cos x\left(\sin^2x-2\sin^4x+\sin^6x\right)=\cos x\sin^2x(1-\sin^2x)^2=\sin^2x\cos^5x$.
\end{solution}

\subsection*{שאלה 6}

\begin{question}
חשבו את נפח גוף הסיבוב סביב ציר ה-$x$ של הפונקציה
\[ f(x)=\frac{1}{\sqrt{x}(x-1)} \]
בקטע $[4,9]$.
\end{question}

\begin{hint}
נפח גוף סיבוב: $V=\pi\int_a^b f^2(x)\,dx$. כאן $f^2(x)=\frac1{x(x-1)^2}$ -- השתמשו בשברים חלקיים.
\end{hint}

\begin{solution}
$f$ רציפה על $[4,9]$, ונפח גוף הסיבוב סביב ציר ה-$x$ הוא
\[ V=\pi\int_4^9f^2(x)\,dx=\pi\int_4^9\frac{dx}{x(x-1)^2}. \]
\textbf{שברים חלקיים:}
\[ \frac1{x(x-1)^2}=\frac Ax+\frac B{x-1}+\frac C{(x-1)^2}\ \Longrightarrow\ 1=A(x-1)^2+Bx(x-1)+Cx. \]
הצבת $x=0$ נותנת $A=1$; הצבת $x=1$ נותנת $C=1$; השוואת מקדם $x^2$: $A+B=0$, ולכן $B=-1$:
\[ \frac1{x(x-1)^2}=\frac1x-\frac1{x-1}+\frac1{(x-1)^2}. \]
\textbf{אינטגרציה:} פונקציה קדומה (ל-$x>1$) היא $F(x)=\ln x-\ln(x-1)-\frac1{x-1}$, ולפי המשפט היסודי
\begin{align*}
\int_4^9\frac{dx}{x(x-1)^2}&=\left[\ln\frac{x}{x-1}-\frac1{x-1}\right]_4^9=\left(\ln\frac98-\frac18\right)-\left(\ln\frac43-\frac13\right)\\
&=\ln\left(\frac98\cdot\frac34\right)+\frac13-\frac18=\ln\frac{27}{32}+\frac5{24}.
\end{align*}
\textbf{תשובה:}
\[ V=\pi\left(\frac5{24}+\ln\frac{27}{32}\right)=\pi\left(\frac5{24}-\ln\frac{32}{27}\right)\approx0.1207. \]
\end{solution}

\subsection*{שאלה 7}

\begin{question}
חשבו את האיטגרל הלא-אמיתי הבא, או הראו התבדרות שלו:
\[ \int_2^\infty x\cdot e^{-x}\,dx \]
\end{question}

\begin{hint}
אינטגרציה בחלקים עם $u=x$, $v'=e^{-x}$, ואחר כך גבול כאשר הגבול העליון שואף לאינסוף.
\end{hint}

\begin{solution}
לפי ההגדרה, $\int_2^\infty xe^{-x}\,dx=\lim_{R\to\infty}\int_2^Rxe^{-x}\,dx$, אם הגבול קיים.

\textbf{פונקציה קדומה:} אינטגרציה בחלקים עם $u=x$, $v'=e^{-x}$ ($u'=1$, $v=-e^{-x}$):
\[ \int xe^{-x}\,dx=-xe^{-x}+\int e^{-x}\,dx=-xe^{-x}-e^{-x}+C=-(x+1)e^{-x}+C. \]
לכן
\[ \int_2^Rxe^{-x}\,dx=\Big[-(x+1)e^{-x}\Big]_2^R=3e^{-2}-\frac{R+1}{e^R}. \]
\textbf{מעבר לגבול:} $\lim_{R\to\infty}\frac{R+1}{e^R}$ הוא מקרה $\frac\infty\infty$, ולפי כלל לופיטל $\lim_{R\to\infty}\frac{1}{e^R}=0$. לכן
\[ \int_2^\infty xe^{-x}\,dx=\boxed{\frac{3}{e^2}}\approx0.406, \]
והאינטגרל מתכנס.
\end{solution}

\section*{תשפ"ד סמסטר ב מועד א}

\subsection*{שאלה 1א}

\begin{question}
חשבו
\[ \lim_{x\to\infty} x\left(\ln(x+1)-\ln(x)\right) \]
\end{question}

\begin{hint}
כתבו $\ln(x+1)-\ln x=\ln\left(1+\frac1x\right)$ והציבו $t=\frac1x$.
\end{hint}

\begin{solution}
לפי חוקי הלוגריתם, $\ln(x+1)-\ln(x)=\ln\frac{x+1}{x}=\ln\left(1+\frac1x\right)$. נציב $t=\frac1x$; כאשר $x\to\infty$ מתקיים $t\to0^+$, ולכן
\[ \lim_{x\to\infty} x\ln\left(1+\frac1x\right)=\lim_{t\to0^+}\frac{\ln(1+t)}{t}. \]
זהו גבול מהצורה $\frac00$; לפי כלל לופיטל (או הגבול היסודי $\frac{\ln(1+t)}{t}\to1$):
\[ \lim_{t\to0^+}\frac{\ln(1+t)}{t}=\lim_{t\to0^+}\frac{1/(1+t)}{1}=1. \]
התשובה: $\boxed{1}$.
\end{solution}

\subsection*{שאלה 1ב}

\begin{question}
חשבו
\[ \lim_{x\to 0} \frac{x^2+\sin(x)\cos(x)}{1-e^x} \]
\end{question}

\begin{solution}
כאשר $x\to0$ המונה שואף ל-$0+\sin0\cos0=0$ והמכנה ל-$1-e^0=0$, כלומר גבול מהצורה $\frac00$. המונה והמכנה גזירים, ונגזרת המכנה $-e^x\neq0$ בסביבת $0$, לכן לפי כלל לופיטל:
\[ \lim_{x\to0}\frac{x^2+\sin x\cos x}{1-e^x}=\lim_{x\to0}\frac{2x+\cos^2x-\sin^2x}{-e^x}=\frac{0+1-0}{-1}=-1, \]
כאשר השתמשנו ברציפות הביטוי שהתקבל ב-$0$. (לחלופין: $\sin x\cos x=\frac12\sin2x$, ומחלקים מונה ומכנה ב-$x$: $\frac{x+\frac{\sin 2x}{2x}}{\frac{1-e^x}{x}}\to\frac{0+1}{-1}$.)
התשובה: $\boxed{-1}$.
\end{solution}

\subsection*{שאלה 2}

\begin{question}
נגדיר פונקציה באופן הבא:
\[ f(x)=\begin{cases} x^2\sin\left(\frac{3}{x}\right) & x\neq0\\ 0 & x=0\end{cases} \]
\begin{enumerate}
  \item האם $f(x)$ רציפה בנקודה $x=0$?
  \item האם $f(x)$ גזירה בנקודה $x=0$?
\end{enumerate}
\end{question}

\begin{hint}
$\sin\left(\frac3x\right)$ חסומה; השתמשו במשפט הסנדוויץ' (או "אפסה כפול חסומה").
\end{hint}

\begin{solution}
\begin{enumerate}
  \item לכל $x\neq0$ מתקיים $\left|\sin\frac3x\right|\le1$, ולכן
  \[ 0\le |f(x)|=x^2\left|\sin\frac3x\right|\le x^2. \]
  מכיוון ש-$x^2\to0$ כאשר $x\to0$, לפי משפט הסנדוויץ' $|f(x)|\to0$, כלומר $\lim_{x\to0}f(x)=0=f(0)$. לכן $f$ \textbf{רציפה} ב-$x=0$.
  \item נחשב את הנגזרת לפי ההגדרה:
  \[ f'(0)=\lim_{h\to0}\frac{f(h)-f(0)}{h}=\lim_{h\to0}\frac{h^2\sin\frac3h}{h}=\lim_{h\to0}h\sin\frac3h. \]
  שוב $0\le\left|h\sin\frac3h\right|\le|h|\to0$, ולפי משפט הסנדוויץ' הגבול קיים ושווה $0$ (פונקציה השואפת לאפס כפול פונקציה חסומה). לכן $f$ \textbf{גזירה} ב-$x=0$ ו-$f'(0)=0$.
\end{enumerate}
\end{solution}

\subsection*{שאלה 3}

\begin{question}
חקרו את הפונקציה הבאה (נק' קיצון מקומיות, תחומי קמירות/קעירות, אסימפטוטות):
\[ f(x)=\frac{2-x-3x^2}{x^2-1} \]
שרטטו את הגרף שלה.
\end{question}

\begin{hint}
פרקו את המונה והמכנה לגורמים — יש גורם משותף.
\end{hint}

\begin{hint}
אחרי הצמצום ניתן לכתוב $f(x)=-3-\frac{1}{x-1}$ (עבור $x\neq\pm1$).
\end{hint}

\begin{solution}
\textbf{תחום הגדרה:} $x^2-1\neq0$, כלומר $x\neq\pm1$.

\textbf{פישוט:} $2-x-3x^2=-(3x^2+x-2)=-(3x-2)(x+1)$ ו-$x^2-1=(x-1)(x+1)$. לכן לכל $x\neq\pm1$:
\[ f(x)=\frac{-(3x-2)(x+1)}{(x-1)(x+1)}=\frac{2-3x}{x-1}=-3-\frac{1}{x-1}. \]
(בדיקה: $-3(x-1)-1=2-3x$.)

\textbf{הנקודה $x=-1$:} $\lim_{x\to-1}f(x)=\frac{2+3}{-2}=-\frac52$ — גבול סופי, ולכן זו נקודת אי-רציפות סליקה ("חור" בגרף בנקודה $(-1,-\frac52)$), ואין בה אסימפטוטה אנכית.

\textbf{חיתוך עם הצירים:} $f(0)=-2$; $f(x)=0\iff x=\frac23$. הנקודות $(0,-2)$, $(\frac23,0)$.

\textbf{נגזרת ראשונה:} $f'(x)=\frac{1}{(x-1)^2}>0$ לכל $x$ בתחום. לכן $f$ עולה בכל אחד מהקטעים $(-\infty,-1)$, $(-1,1)$, $(1,\infty)$, ו\textbf{אין נקודות קיצון מקומיות} (הנגזרת לא מתאפסת בשום מקום, ולפי משפט פרמה אין קיצון פנימי).

\textbf{נגזרת שנייה:} $f''(x)=-\frac{2}{(x-1)^3}$.
\begin{itemize}
  \item עבור $x<1$ ($x\neq-1$): $(x-1)^3<0$, לכן $f''>0$ — $f$ \textbf{קמורה} ($\cup$) ב-$(-\infty,-1)$ וב-$(-1,1)$.
  \item עבור $x>1$: $f''<0$ — $f$ \textbf{קעורה} ($\cap$) ב-$(1,\infty)$.
\end{itemize}
אין נקודות פיתול: שינוי הקמירות קורה רק ב-$x=1$, שאינה בתחום ההגדרה.

\textbf{אסימפטוטות:}
\begin{itemize}
  \item אנכית: $x=1$. $\lim_{x\to1^-}f(x)=-3-\frac{1}{0^-}=+\infty$, $\lim_{x\to1^+}f(x)=-\infty$.
  \item אופקית: $\lim_{x\to\pm\infty}f(x)=-3$ (המעלות במונה ובמכנה שוות, יחס המקדמים המובילים $\frac{-3}{1}$). לכן $y=-3$ אסימפטוטה אופקית בשני הכיוונים; אין אסימפטוטה משופעת.
\end{itemize}

\textbf{הגרף:} היפרבולה $y=-3-\frac1{x-1}$ (הזזה של $-\frac1x$) עם חור בנקודה $(-1,-\frac52)$. משמאל ל-$x=1$ הגרף עולה מהאסימפטוטה $y=-3$ (מעליה) דרך החור $(-1,-2.5)$, $(0,-2)$, $(\frac23,0)$ ושואף ל-$+\infty$ כאשר $x\to1^-$, והוא קמור. מימין ל-$x=1$ הגרף עולה מ-$-\infty$ ומתקרב מלמטה ל-$y=-3$, והוא קעור.
\end{solution}

\subsection*{שאלה 4א}

\begin{question}
כתבו את פולינום מקלורן מסדר 3 של $f(x)=\ln(1+x)$.
\end{question}

\begin{solution}
נחשב נגזרות ב-$0$:
\[ f(x)=\ln(1+x),\quad f'(x)=\frac1{1+x},\quad f''(x)=-\frac1{(1+x)^2},\quad f'''(x)=\frac{2}{(1+x)^3}, \]
ולכן $f(0)=0,\ f'(0)=1,\ f''(0)=-1,\ f'''(0)=2$. פולינום מקלורן:
\[ P_3(x)=f(0)+f'(0)x+\frac{f''(0)}{2!}x^2+\frac{f'''(0)}{3!}x^3=\boxed{x-\frac{x^2}{2}+\frac{x^3}{3}}. \]
\end{solution}

\subsection*{שאלה 4ב}

\begin{question}
מצאו קירוב לינארי של $\sqrt{15}$
\end{question}

\begin{hint}
השתמשו במשיק לגרף $g(x)=\sqrt x$ בנקודה קרובה ל-$15$ שבה השורש ידוע.
\end{hint}

\begin{solution}
נשתמש בפונקציה $g(x)=\sqrt x$ סביב הנקודה $a=16$ (הקרובה ל-$15$, שבה השורש ידוע). $g(16)=4$, $g'(x)=\frac{1}{2\sqrt x}$, $g'(16)=\frac18$. הקירוב הלינארי (משוואת המשיק):
\[ g(x)\approx g(16)+g'(16)(x-16)=4+\frac{x-16}{8}. \]
עבור $x=15$: $\sqrt{15}\approx4-\frac18=\boxed{3.875}$.
(לשם השוואה $\sqrt{15}\approx3.873$. מאחר ש-$g''<0$ — $g$ קעורה — המשיק נמצא מעל הגרף, ולכן זהו קירוב מלמעלה; לפי שארית לגרנז' השגיאה היא $\frac{|g''(c)|}{2}\cdot1^2=\frac{1}{8c^{3/2}}<\frac{1}{8\cdot 15^{3/2}}<0.003$.)
\end{solution}

\subsection*{שאלה 5א}

\begin{question}
חשבו
\[ \int\frac{x^3+2x^2-3x+5}{x^2+x-6}\,dx \]
\end{question}

\begin{hint}
חילוק פולינומים ואז פירוק לשברים חלקיים ($x^2+x-6=(x+3)(x-2)$).
\end{hint}

\begin{solution}
מעלת המונה גדולה ממעלת המכנה, לכן נחלק פולינומים:
\[ x^3+2x^2-3x+5=(x+1)(x^2+x-6)+(2x+11), \]
(בדיקה: $(x+1)(x^2+x-6)=x^3+2x^2-5x-6$, והשארית היא $x^3+2x^2-3x+5-(x^3+2x^2-5x-6)=2x+11$.) לכן
\[ \frac{x^3+2x^2-3x+5}{x^2+x-6}=x+1+\frac{2x+11}{(x+3)(x-2)}. \]
פירוק לשברים חלקיים: $\frac{2x+11}{(x+3)(x-2)}=\frac{A}{x+3}+\frac{B}{x-2}$, כלומר $2x+11=A(x-2)+B(x+3)$. הצבת $x=2$: $15=5B\Rightarrow B=3$; הצבת $x=-3$: $5=-5A\Rightarrow A=-1$. לכן
\[ \int\frac{x^3+2x^2-3x+5}{x^2+x-6}\,dx=\int\left(x+1-\frac1{x+3}+\frac3{x-2}\right)dx=\boxed{\frac{x^2}{2}+x-\ln|x+3|+3\ln|x-2|+C}. \]
\end{solution}

\subsection*{שאלה 5ב}

\begin{question}
חשבו
\[ \int_1^e(\ln(x))^2\,dx \]
\end{question}

\begin{hint}
אינטגרציה בחלקים פעמיים, עם $u=(\ln x)^2,\ v'=1$.
\end{hint}

\begin{solution}
אינטגרציה בחלקים עם $u=(\ln x)^2$, $v'=1$, כלומר $u'=\frac{2\ln x}{x}$, $v=x$:
\[ \int_1^e(\ln x)^2dx=\Big[x(\ln x)^2\Big]_1^e-\int_1^e 2\ln x\,dx=e-2\int_1^e\ln x\,dx. \]
שוב בחלקים ($u=\ln x,\ v'=1$): $\int_1^e\ln x\,dx=\big[x\ln x\big]_1^e-\int_1^e1\,dx=e-(e-1)=1$. לכן
\[ \int_1^e(\ln x)^2dx=\boxed{e-2}\approx0.718. \]
\end{solution}

\section*{תשפ"ד סמסטר ב מועד ב}

\subsection*{שאלה 1א}

\begin{question}
חשבו
\[ \lim_{x\to0}\sqrt[x]{1-2x} \]
\end{question}

\begin{hint}
כתבו $\sqrt[x]{1-2x}=(1-2x)^{1/x}=e^{\frac{\ln(1-2x)}{x}}$.
\end{hint}

\begin{solution}
לכל $x\neq0$ קרוב ל-$0$ מתקיים $1-2x>0$, ולכן
\[ \sqrt[x]{1-2x}=(1-2x)^{1/x}=e^{\frac{\ln(1-2x)}{x}}. \]
נחשב את גבול המעריך. זהו גבול מהצורה $\frac00$, ולפי כלל לופיטל:
\[ \lim_{x\to0}\frac{\ln(1-2x)}{x}=\lim_{x\to0}\frac{\frac{-2}{1-2x}}{1}=-2. \]
מרציפות פונקציית האקספוננט נקבל
\[ \lim_{x\to0}\sqrt[x]{1-2x}=e^{-2}=\boxed{\frac1{e^2}}. \]
(זו גם צורה של הגבול היסודי $\lim_{t\to0}(1+t)^{1/t}=e$ עם $t=-2x$.)
\end{solution}

\subsection*{שאלה 1ב}

\begin{question}
חשבו
\[ \lim_{x\to0}\frac{\tan(x)-\sin(x)}{\sin^3(x)} \]
\end{question}

\begin{hint}
הוציאו $\tan x$ כגורם משותף במונה: $\tan x-\sin x=\tan x\,(1-\cos x)$.
\end{hint}

\begin{solution}
בסביבה מנוקבת של $0$ מתקיים $\sin x\neq0$, $\cos x\neq0$, ו-
\[ \tan x-\sin x=\frac{\sin x}{\cos x}-\sin x=\frac{\sin x\,(1-\cos x)}{\cos x}. \]
לכן, בעזרת $1-\cos^2x=\sin^2x$:
\[ \frac{\tan x-\sin x}{\sin^3x}=\frac{1-\cos x}{\cos x\,\sin^2x}=\frac{1-\cos x}{\cos x\,(1-\cos x)(1+\cos x)}=\frac{1}{\cos x\,(1+\cos x)}. \]
(הצמצום מותר כי $\cos x\neq1$ עבור $0<|x|<2\pi$.) הביטוי שהתקבל רציף ב-$0$, ולכן
\[ \lim_{x\to0}\frac{\tan x-\sin x}{\sin^3x}=\frac{1}{1\cdot2}=\boxed{\frac12}. \]
\end{solution}

\subsection*{שאלה 2}

\begin{question}
לכל $a,b\in\R$ נגדיר את הפונקציה הבאה:
\[ f(x)=\begin{cases} \sin(ax) & x<0\\ b-2 & x=0\\ x^2\ln(x) & 0>x\end{cases} \]
\begin{enumerate}
  \item עבור אילו $a,b\in\R$ הפונקציה $f(x)$ רציפה בנקודה $x=0$?
  \item האם קיימים $a,b\in\R$ עבורם הפונקציה $f(x)$ גזירה בנקודה $x=0$?
\end{enumerate}
\end{question}

\begin{hint}
חשבו את הגבולות החד-צדדיים ב-$0$. לגבול $\lim_{x\to0^+}x^2\ln x$ (וגם $x\ln x$) השתמשו בכלל לופיטל אחרי כתיבה כמנה.
\end{hint}

\begin{hint}
גזירות גוררת רציפות. חשבו את הנגזרות החד-צדדיות לפי ההגדרה.
\end{hint}

\begin{solution}
\textbf{הערה:} בתנאי של השורה השלישית נכתב "$0>x$", וזו כמובן טעות דפוס — השורה הראשונה כבר מכסה את $x<0$, וב-$x<0$ הביטוי $\ln x$ אינו מוגדר. הכוונה היא $0<x$, ונפתור בהתאם.
\begin{enumerate}
  \item \textbf{גבול משמאל:} $\lim_{x\to0^-}\sin(ax)=\sin0=0$ (רציפות הסינוס), לכל $a$.

  \textbf{גבול מימין:} $\lim_{x\to0^+}x^2\ln x$ הוא מהצורה $0\cdot(-\infty)$. נכתוב כמנה ונשתמש בכלל לופיטל:
  \[ \lim_{x\to0^+}\frac{\ln x}{x^{-2}}=\lim_{x\to0^+}\frac{1/x}{-2x^{-3}}=\lim_{x\to0^+}\left(-\frac{x^2}{2}\right)=0. \]
  לכן הגבול של $f$ ב-$0$ קיים ושווה $0$ לכל $a$, ו-$f$ רציפה ב-$0$ אם ורק אם $f(0)=b-2=0$.
  \textbf{תשובה:} $f$ רציפה ב-$0$ עבור $b=2$ ו-$a\in\R$ כלשהו.
  \item אם $f$ גזירה ב-$0$ היא בפרט רציפה שם, ולכן חייב להתקיים $b=2$, כלומר $f(0)=0$. נחשב את הנגזרות החד-צדדיות לפי ההגדרה:
  \[ f'_-(0)=\lim_{h\to0^-}\frac{\sin(ah)-0}{h}=a\lim_{h\to0^-}\frac{\sin(ah)}{ah}=a \]
  (עבור $a=0$ המנה זהותית $0$, וגם אז הגבול $0=a$), ו-
  \[ f'_+(0)=\lim_{h\to0^+}\frac{h^2\ln h-0}{h}=\lim_{h\to0^+}h\ln h=\lim_{h\to0^+}\frac{\ln h}{1/h}\overset{\text{לופיטל}}{=}\lim_{h\to0^+}\frac{1/h}{-1/h^2}=\lim_{h\to0^+}(-h)=0. \]
  $f$ גזירה ב-$0$ אם ורק אם שתי הנגזרות החד-צדדיות שוות, כלומר $a=0$.
  \textbf{תשובה:} כן — $f$ גזירה ב-$0$ בדיוק עבור $a=0,\ b=2$, ואז $f'(0)=0$.
\end{enumerate}
\end{solution}

\subsection*{שאלה 3}

\begin{question}
חקרו את הפונקציה הבאה (נק' קיצון מקומיות, תחומי עליה/ירידה, אסימפטוטות):
\[ f(x)=\frac{e^x}{x^2-3} \]
שרטטו את הגרף שלה.
\end{question}

\begin{hint}
בנגזרת, המונה הוא $e^x(x^2-2x-3)=e^x(x-3)(x+1)$. שימו לב שנקודות הקיצון חייבות להיות בתחום ההגדרה, וש-$-\sqrt3<-1<\sqrt3<3$.
\end{hint}

\begin{solution}
\textbf{תחום הגדרה:} $x^2\neq3$, כלומר $x\neq\pm\sqrt3$.

\textbf{סימן וחיתוך עם הצירים:} $e^x>0$, לכן אין חיתוך עם ציר $x$, ו-$f>0$ עבור $|x|>\sqrt3$, $f<0$ עבור $|x|<\sqrt3$. חיתוך עם ציר $y$: $f(0)=-\frac13$.

\textbf{נגזרת:} לפי כלל המנה
\[ f'(x)=\frac{e^x(x^2-3)-e^x\cdot2x}{(x^2-3)^2}=\frac{e^x(x^2-2x-3)}{(x^2-3)^2}=\frac{e^x(x-3)(x+1)}{(x^2-3)^2}. \]
המכנה ו-$e^x$ חיוביים, ולכן סימן $f'$ הוא סימן $(x-3)(x+1)$. מתקיים $-\sqrt3\approx-1.73<-1<\sqrt3\approx1.73<3$:
\begin{itemize}
  \item $(-\infty,-\sqrt3)$ ו-$(-\sqrt3,-1)$: $f'>0$, $f$ \textbf{עולה}.
  \item $(-1,\sqrt3)$ ו-$(\sqrt3,3)$: $f'<0$, $f$ \textbf{יורדת}.
  \item $(3,\infty)$: $f'>0$, $f$ \textbf{עולה}.
\end{itemize}

\textbf{נקודות קיצון:} $f'$ מחליפה סימן מ-$+$ ל-$-$ ב-$x=-1$, לכן זו נקודת \textbf{מקסימום מקומי}: $f(-1)=\frac{e^{-1}}{-2}=-\frac1{2e}\approx-0.18$. ב-$x=3$ הסימן מתחלף מ-$-$ ל-$+$, לכן \textbf{מינימום מקומי}: $f(3)=\frac{e^3}{6}\approx3.35$.

\textbf{אסימפטוטות:}
\begin{itemize}
  \item אנכיות $x=\pm\sqrt3$ (המונה $e^{\pm\sqrt3}\neq0$ והמכנה מתאפס):
  $\lim_{x\to-\sqrt3^-}f=+\infty$, $\lim_{x\to-\sqrt3^+}f=-\infty$, $\lim_{x\to\sqrt3^-}f=-\infty$, $\lim_{x\to\sqrt3^+}f=+\infty$ (לפי סימן $x^2-3$ משני הצדדים).
  \item ב-$-\infty$: $e^x\to0$ ו-$x^2-3\to\infty$, לכן $f\to0$: אסימפטוטה אופקית $y=0$ (הגרף מעליה).
  \item ב-$+\infty$: לפי כלל לופיטל (פעמיים) $\lim_{x\to\infty}\frac{e^x}{x^2-3}=\lim\frac{e^x}{2}=\infty$, וגם $\frac{f(x)}{x}=\frac{e^x}{x^3-3x}\to\infty$; לכן אין אסימפטוטה אופקית או משופעת ב-$+\infty$.
\end{itemize}

\textbf{הגרף:} משמאל הגרף יוצא מעל ציר $x$ קרוב ל-$y=0$, עולה ל-$+\infty$ כאשר $x\to-\sqrt3^-$. בין $-\sqrt3$ ל-$\sqrt3$ הגרף שלילי: עולה מ-$-\infty$ עד המקסימום המקומי $(-1,-\frac1{2e})$, יורד דרך $(0,-\frac13)$ ל-$-\infty$ כאשר $x\to\sqrt3^-$. מימין ל-$\sqrt3$ הגרף יורד מ-$+\infty$ עד המינימום המקומי $(3,\frac{e^3}{6})$ ואז עולה ל-$+\infty$.
\end{solution}

\subsection*{שאלה 4א}

\begin{question}
כתבו את פולינום מקלורן מסדר 4 של $f(x)=\sin(x)\cos(x)$
\end{question}

\begin{hint}
$\sin x\cos x=\frac12\sin(2x)$.
\end{hint}

\begin{solution}
לפי זהות הזווית הכפולה $f(x)=\sin x\cos x=\frac12\sin(2x)$. פולינום מקלורן של $\sin t$ מסדר 4 הוא $t-\frac{t^3}{6}$ (המקדם של $t^4$ הוא $0$). נציב $t=2x$:
\[ P_4(x)=\frac12\left(2x-\frac{(2x)^3}{6}\right)=\boxed{x-\frac{2}{3}x^3}. \]
(ישירות: $f'(x)=\cos2x$, $f''=-2\sin2x$, $f'''=-4\cos2x$, $f^{(4)}=8\sin2x$; בנקודה $0$: $f=0,\ f'=1,\ f''=0,\ f'''=-4,\ f^{(4)}=0$, ולכן $P_4=x-\frac{4}{6}x^3$.)
\end{solution}

\subsection*{שאלה 4ב}

\begin{question}
מצאו קירוב לינארי של $5^{\frac{2}{3}}$
\end{question}

\begin{hint}
בחרו $g(x)=x^{2/3}$ ונקודה קרובה ל-$5$ שבה הערך ידוע.
\end{hint}

\begin{solution}
נשתמש ב-$g(x)=x^{2/3}$ סביב $a=8$ (שם $8^{2/3}=4$). $g'(x)=\frac23x^{-1/3}$, ולכן $g'(8)=\frac23\cdot\frac12=\frac13$. הקירוב הלינארי:
\[ g(x)\approx g(8)+g'(8)(x-8)=4+\frac{x-8}{3}. \]
עבור $x=5$: $5^{2/3}\approx4+\frac{-3}{3}=\boxed{3}$.
(הערך האמיתי $\approx2.924$; השגיאה גדולה יחסית כי $5$ רחוקה מ-$8$. בחירה אחרת, למשל $a=1$, נותנת $1+\frac23\cdot4\approx3.67$ — פחות מדויק.)
\end{solution}

\subsection*{שאלה 5א}

\begin{question}
חשבו
\[ \int\frac{2x^3+7x^2-5x-19}{x^2+x-6}\,dx \]
\end{question}

\begin{hint}
חילוק פולינומים ואז שברים חלקיים, $x^2+x-6=(x+3)(x-2)$.
\end{hint}

\begin{solution}
חילוק פולינומים:
\[ 2x^3+7x^2-5x-19=(2x+5)(x^2+x-6)+(2x+11), \]
(בדיקה: $(2x+5)(x^2+x-6)=2x^3+7x^2-7x-30$, והשארית $(-5x-19)-(-7x-30)=2x+11$.) שברים חלקיים: $\frac{2x+11}{(x+3)(x-2)}=\frac{A}{x+3}+\frac{B}{x-2}$, כלומר $2x+11=A(x-2)+B(x+3)$. עבור $x=2$: $15=5B$, $B=3$. עבור $x=-3$: $5=-5A$, $A=-1$. לכן
\[ \int\frac{2x^3+7x^2-5x-19}{x^2+x-6}dx=\int\left(2x+5-\frac1{x+3}+\frac3{x-2}\right)dx=\boxed{x^2+5x-\ln|x+3|+3\ln|x-2|+C}. \]
\end{solution}

\subsection*{שאלה 5ב}

\begin{question}
חשבו
\[ \int_0^{\pi}e^x\sin(x)\,dx \]
\end{question}

\begin{hint}
אינטגרציה בחלקים פעמיים; האינטגרל המקורי יופיע שוב באגף ימין, ופותרים משוואה עבורו.
\end{hint}

\begin{solution}
נסמן $I=\int_0^\pi e^x\sin x\,dx$. בחלקים עם $u=\sin x$, $v'=e^x$:
\[ I=\big[e^x\sin x\big]_0^\pi-\int_0^\pi e^x\cos x\,dx=0-\int_0^\pi e^x\cos x\,dx. \]
שוב בחלקים עם $u=\cos x$, $v'=e^x$:
\[ \int_0^\pi e^x\cos x\,dx=\big[e^x\cos x\big]_0^\pi+\int_0^\pi e^x\sin x\,dx=(-e^\pi-1)+I. \]
לכן $I=e^\pi+1-I$, כלומר
\[ I=\boxed{\frac{e^\pi+1}{2}}\approx12.07. \]
\end{solution}

\section*{תשפ"ד סמסטר ב מועד ג}

\subsection*{שאלה 1א}

\begin{question}
חשבו
\[ \lim_{x\longrightarrow\infty}\left(\frac{x+3}{x-4}\right)^x \]
\end{question}

\begin{hint}
כתבו $\frac{x+3}{x-4}=1+\frac{7}{x-4}$ והשתמשו בגבול $\lim_{t\to\infty}\left(1+\frac1t\right)^t=e$.
\end{hint}

\begin{solution}
לכל $x>4$: $\frac{x+3}{x-4}=1+\frac{7}{x-4}$. נסמן $t=\frac{x-4}{7}\to\infty$, אז $x=7t+4$ ו-
\[ \left(\frac{x+3}{x-4}\right)^x=\left(1+\frac1t\right)^{7t+4}=\left[\left(1+\frac1t\right)^{t}\right]^7\cdot\left(1+\frac1t\right)^4\xrightarrow[t\to\infty]{}e^7\cdot1. \]
(לחלופין: $x\ln\frac{x+3}{x-4}=\frac{\ln(x+3)-\ln(x-4)}{1/x}$, ולפי לופיטל הגבול הוא $\lim\frac{\frac1{x+3}-\frac1{x-4}}{-1/x^2}=\lim\frac{7x^2}{(x+3)(x-4)}=7$.)
התשובה: $\boxed{e^7}$.
\end{solution}

\subsection*{שאלה 1ב}

\begin{question}
חשבו
\[ \lim_{x\longrightarrow-8}\frac{\sqrt{1+x}-3}{2+\sqrt[3]{x}} \]
\end{question}

\begin{hint}
בדקו קודם את תחום ההגדרה של $\sqrt{1+x}$ ליד $x=-8$.
\end{hint}

\begin{solution}
\textbf{כפי שהשאלה מודפסת}, הביטוי $\sqrt{1+x}$ מוגדר (בממשיים) רק עבור $x\ge-1$, ולכן הפונקציה אינה מוגדרת באף סביבה מנוקבת של $x=-8$, והגבול \textbf{אינו מוגדר}. (גם ההצבה $x=-8$ נותנת $\sqrt{-7}$.)

ככל הנראה נפלה טעות דפוס, והכוונה ל-$\frac{\sqrt{1-x}-3}{2+\sqrt[3]{x}}$ — אז בנקודה $x=-8$ המונה $\sqrt9-3=0$ והמכנה $2+(-2)=0$, גבול מהצורה $\frac00$. נפתור גבול זה. נכפיל ב"צמודים": במונה $\sqrt{1-x}-3=\frac{(1-x)-9}{\sqrt{1-x}+3}=\frac{-(x+8)}{\sqrt{1-x}+3}$, ובמכנה, לפי $a^3+b^3=(a+b)(a^2-ab+b^2)$ עם $a=2,\ b=\sqrt[3]x$:
\[ 2+\sqrt[3]{x}=\frac{8+x}{4-2\sqrt[3]{x}+\sqrt[3]{x^2}}. \]
לכן עבור $x\neq-8$:
\[ \frac{\sqrt{1-x}-3}{2+\sqrt[3]{x}}=-\frac{4-2\sqrt[3]{x}+\sqrt[3]{x^2}}{\sqrt{1-x}+3}\xrightarrow[x\to-8]{}-\frac{4+4+4}{3+3}=-2, \]
לפי רציפות הביטוי האחרון ב-$x=-8$. (אותה תוצאה מתקבלת מכלל לופיטל: $\frac{-\frac{1}{2\sqrt{1-x}}}{\frac{1}{3\sqrt[3]{x^2}}}\to\frac{-1/6}{1/12}=-2$.)
התשובה: כפי שמודפס — הגבול אינו מוגדר; בגרסה המתוקנת $\sqrt{1-x}$ — הגבול הוא $\boxed{-2}$.
\end{solution}

\subsection*{שאלה 2}

\begin{question}
לכל $a,b\in\R$ נגדיר את הפונקציה הבאה:
\[ f(x)=\begin{cases} ax^2+bx+1 & x<0\\ b-2 & x=0\\ bx^2+ax+1 & 0<x\end{cases} \]
\begin{enumerate}
  \item עבור אילו $a,b\in\R$ הפונקציה $f(x)$ רציפה בנקודה $x=0$?
  \item האם קיימים $a,b\in\R$ עבורם הפונקציה $f(x)$ גזירה בנקודה $x=0$?
\end{enumerate}
\end{question}

\begin{hint}
גזירות גוררת רציפות, אז תנאי סעיף א' חייב להתקיים. חשבו את הנגזרות החד-צדדיות לפי ההגדרה.
\end{hint}

\begin{solution}
\begin{enumerate}
  \item הפולינומים רציפים, לכן
  \[ \lim_{x\to0^-}f(x)=a\cdot0+b\cdot0+1=1,\qquad \lim_{x\to0^+}f(x)=1. \]
  הגבול ב-$0$ קיים ושווה $1$ לכל $a,b$, ו-$f$ רציפה ב-$0$ אם ורק אם $f(0)=b-2=1$.
  \textbf{תשובה:} $b=3$ ו-$a\in\R$ כלשהו.
  \item אם $f$ גזירה ב-$0$ אז היא רציפה שם, לכן $b=3$ ו-$f(0)=1$. נחשב נגזרות חד-צדדיות לפי ההגדרה:
  \[ f'_-(0)=\lim_{h\to0^-}\frac{ah^2+bh+1-1}{h}=\lim_{h\to0^-}(ah+b)=b, \qquad f'_+(0)=\lim_{h\to0^+}\frac{bh^2+ah+1-1}{h}=\lim_{h\to0^+}(bh+a)=a. \]
  $f$ גזירה ב-$0$ אם ורק אם $f'_-(0)=f'_+(0)$, כלומר $a=b$. יחד עם $b=3$:
  \textbf{תשובה:} כן, בדיוק עבור $a=b=3$, ואז $f'(0)=3$.
\end{enumerate}
\end{solution}

\subsection*{שאלה 3}

\begin{question}
חקרו את הפונקציה הבאה (נק' קיצון מקומיות, תחומי עליה/ירידה, אסימפטוטות):
\[ f(x)=e^{-(x^3-3x)} \]
שרטטו את הגרף שלה.
\end{question}

\begin{hint}
$f=e^{g}$ עם $g(x)=3x-x^3$; מכיוון ש-$e^t$ עולה, $f$ עולה/יורדת בדיוק היכן ש-$g$ עולה/יורדת.
\end{hint}

\begin{solution}
\textbf{תחום הגדרה:} כל $\R$; $f$ רציפה וגזירה בכל מקום, ו-$f(x)>0$ לכל $x$ (אין חיתוך עם ציר $x$). חיתוך עם ציר $y$: $f(0)=e^0=1$.

\textbf{נגזרת:} לפי כלל השרשרת
\[ f'(x)=e^{-(x^3-3x)}\cdot(-(3x^2-3))=3(1-x^2)\,e^{-(x^3-3x)}=3(1-x)(1+x)\,e^{3x-x^3}. \]
האקספוננט חיובי, לכן סימן $f'$ הוא סימן $1-x^2$:
\begin{itemize}
  \item $(-\infty,-1)$: $f'<0$, $f$ \textbf{יורדת}.
  \item $(-1,1)$: $f'>0$, $f$ \textbf{עולה}.
  \item $(1,\infty)$: $f'<0$, $f$ \textbf{יורדת}.
\end{itemize}

\textbf{נקודות קיצון:} ב-$x=-1$ הנגזרת עוברת מ-$-$ ל-$+$: \textbf{מינימום מקומי} $f(-1)=e^{-(-1+3)}=e^{-2}\approx0.135$. ב-$x=1$ הנגזרת עוברת מ-$+$ ל-$-$: \textbf{מקסימום מקומי} $f(1)=e^{-(1-3)}=e^{2}\approx7.39$.

\textbf{אסימפטוטות:}
\begin{itemize}
  \item אנכיות: אין, $f$ רציפה על כל הישר.
  \item $x\to+\infty$: $3x-x^3=-x^3\left(1-\frac{3}{x^2}\right)\to-\infty$, לכן $f(x)\to0$: אסימפטוטה אופקית $y=0$ (הגרף מעליה).
  \item $x\to-\infty$: $3x-x^3\to+\infty$, לכן $f(x)\to+\infty$; גם $\frac{f(x)}{x}\to-\infty$ (מונה $\to+\infty$, מכנה $\to-\infty$), ולכן אין אסימפטוטה אופקית או משופעת ב-$-\infty$.
\end{itemize}

\textbf{הגרף:} מגיע מ-$+\infty$ בצד שמאל, יורד עד המינימום המקומי $(-1,e^{-2})$, עולה דרך $(0,1)$ עד המקסימום המקומי $(1,e^2)$, ואז יורד ושואף ל-$0$ (מעל ציר $x$) כאשר $x\to\infty$. הגרף כולו מעל ציר $x$.
\end{solution}

\subsection*{שאלה 4א}

\begin{question}
כתבו את פולינום מקלורן מסדר 3 של $f(x)=\sin(\cos(x))$
\end{question}

\begin{hint}
$f$ זוגית, ולכן מקדם $x$ ומקדם $x^3$ מתאפסים; נשאר לחשב את $f(0)$ ו-$f''(0)$.
\end{hint}

\begin{solution}
נחשב נגזרות לפי כלל השרשרת:
\[ f'(x)=\cos(\cos x)\cdot(-\sin x)=-\sin x\cos(\cos x), \]
\[ f''(x)=-\cos x\cos(\cos x)-\sin x\cdot\sin(\cos x)\cdot\sin x=-\cos x\cos(\cos x)-\sin^2x\,\sin(\cos x). \]
בנקודה $0$: $f(0)=\sin1$, $f'(0)=0$, $f''(0)=-\cos1$. מכיוון ש-$f(-x)=\sin(\cos(-x))=f(x)$, הפונקציה זוגית, ולכן כל נגזרותיה מסדר אי-זוגי מתאפסות ב-$0$ (הנגזרת של פונקציה זוגית היא אי-זוגית), ובפרט $f'''(0)=0$. לכן
\[ P_3(x)=f(0)+f'(0)x+\frac{f''(0)}{2}x^2+\frac{f'''(0)}{6}x^3=\boxed{\sin1-\frac{\cos1}{2}x^2}. \]
(בקירוב: $0.8415-0.2702x^2$.)
\end{solution}

\subsection*{שאלה 4ב}

\begin{question}
מצאו קירוב לינארי של $\sqrt[3]{7^2}$ בעזרת בפונקציה $f(x)=\sqrt[3]{x^2}$
\end{question}

\begin{hint}
השתמשו במשיק לגרף $f$ בנקודה קרובה ל-$7$ שבה הערך ידוע.
\end{hint}

\begin{solution}
$f(x)=\sqrt[3]{x^2}=x^{2/3}$ סביב $a=8$, שם $f(8)=4$. $f'(x)=\frac23x^{-1/3}$, $f'(8)=\frac23\cdot\frac12=\frac13$. הקירוב הלינארי:
\[ f(x)\approx4+\frac13(x-8), \qquad \sqrt[3]{7^2}=f(7)\approx4-\frac13=\boxed{\frac{11}{3}\approx3.667}. \]
(הערך האמיתי $\sqrt[3]{49}\approx3.659$.)
\end{solution}

\subsection*{שאלה 5א}

\begin{question}
חשבו
\[ \int\frac{x^3}{x^2+x-2}\,dx \]
\end{question}

\begin{hint}
חילוק פולינומים ושברים חלקיים, $x^2+x-2=(x+2)(x-1)$.
\end{hint}

\begin{solution}
חילוק פולינומים:
\[ x^3=(x-1)(x^2+x-2)+(3x-2), \]
(בדיקה: $(x-1)(x^2+x-2)=x^3-3x+2$, והשארית $x^3-(x^3-3x+2)=3x-2$.) שברים חלקיים: $\frac{3x-2}{(x+2)(x-1)}=\frac{A}{x+2}+\frac{B}{x-1}$, כלומר $3x-2=A(x-1)+B(x+2)$. עבור $x=1$: $1=3B$, $B=\frac13$. עבור $x=-2$: $-8=-3A$, $A=\frac83$. לכן
\[ \int\frac{x^3}{x^2+x-2}dx=\int\left(x-1+\frac{8/3}{x+2}+\frac{1/3}{x-1}\right)dx=\boxed{\frac{x^2}{2}-x+\frac83\ln|x+2|+\frac13\ln|x-1|+C}. \]
\end{solution}

\subsection*{שאלה 5ב}

\begin{question}
חשבו
\[ \int_0^{\ln2}xe^{-x}\,dx \]
\end{question}

\begin{hint}
אינטגרציה בחלקים עם $u=x$.
\end{hint}

\begin{solution}
אינטגרציה בחלקים עם $u=x$, $v'=e^{-x}$, כלומר $u'=1$, $v=-e^{-x}$:
\[ \int_0^{\ln2}xe^{-x}dx=\Big[-xe^{-x}\Big]_0^{\ln2}+\int_0^{\ln2}e^{-x}dx=-\frac{\ln2}{2}+\Big[-e^{-x}\Big]_0^{\ln2}=-\frac{\ln2}{2}+\left(1-\frac12\right), \]
כי $e^{-\ln2}=\frac12$. לכן
\[ \int_0^{\ln2}xe^{-x}dx=\boxed{\frac{1-\ln2}{2}}\approx0.153. \]
\end{solution}

\section*{תשפ"ה סמסטר א בוחן לדוגמה}

\subsection*{חלק א', שאלה 1}

\begin{question}
חשבו את הגבול הבא:
\[
\lim_{x\to\infty}\left(\frac{x^2+x+2}{x^2+x+1}\right)^{2x^2+4}
\]
\end{question}

\begin{hint}
זהו גבול מהצורה $1^\infty$. כתבו את הבסיס בצורה $1+\frac{1}{x^2+x+1}$.
\end{hint}

\begin{solution}
הבסיס שואף ל-1 והמעריך ל-$\infty$ — גבול מהצורה $1^\infty$. נכתוב את הבסיס כ-
\[
\frac{x^2+x+2}{x^2+x+1}=\frac{x^2+x+1+1}{x^2+x+1}=1+\frac{1}{x^2+x+1},
\]
ונכפול ונחלק את המעריך ב-$x^2+x+1$:
\[
\left(1+\frac{1}{x^2+x+1}\right)^{2x^2+4}
=\left[\left(1+\frac{1}{x^2+x+1}\right)^{x^2+x+1}\right]^{\frac{2x^2+4}{x^2+x+1}} .
\]
\textbf{הבסיס הפנימי:} נציב $y=x^2+x+1$; כאשר $x\to\infty$ גם $y\to\infty$, ולפי הגבול היסודי
\[
\lim_{y\to\infty}\left(1+\frac1y\right)^y=e
\]
(ומשפט הגבול של הרכבה) הבסיס הפנימי שואף ל-$e$.

\textbf{המעריך:} חלוקה ב-$x^2$ נותנת
\[
\lim_{x\to\infty}\frac{2x^2+4}{x^2+x+1}=\lim_{x\to\infty}\frac{2+\frac4{x^2}}{1+\frac1x+\frac1{x^2}}=2 .
\]
\textbf{סיכום:} אם $u(x)\to e>0$ ו-$v(x)\to2$ אז $u^v=e^{v\ln u}\to e^{2\ln e}=e^2$, מרציפות $\ln$ ו-$\exp$ (גבול של מנה/מכפלה ורציפות). לכן
\[
\lim_{x\to\infty}\left(\frac{x^2+x+2}{x^2+x+1}\right)^{2x^2+4}=\boxed{e^2}.
\]
\end{solution}

\subsection*{חלק א', שאלה 2}

\begin{question}
יהיו $a,b\in\R$. נגדיר פונקציה על ידי
\[
f(x)=\begin{cases}
\frac{\sin(8x)}{x} & x<0\\
ax+b & 0\le x\le 4\\
\frac{x^2-3x-4}{2-\sqrt{x}} & x>4
\end{cases}
\]
לאילו ערכים של $a,b$ הפונקציה רציפה ב-$\R$ ? נמקו.
\end{question}

\begin{hint}
בכל נקודה $x\neq0,4$ הפונקציה רציפה בלי קשר ל-$a,b$. בדקו רק את נקודות התפר $x=0$ ו-$x=4$.
\end{hint}

\begin{hint}
ב-$x=4$: $x^2-3x-4=(x-4)(x+1)$, וכפלו וחלקו ב-$2+\sqrt x$.
\end{hint}

\begin{solution}
\textbf{נקודות שאינן נקודות תפר.} בכל נקודה $x<0$ הפונקציה היא מנה של פונקציות רציפות עם מכנה שונה מ-0; בכל נקודה $0<x<4$ היא פולינום; ובכל נקודה $x>4$ היא מנה של פונקציות רציפות עם מכנה $2-\sqrt x\neq0$ (כי $\sqrt x>2$). מכיוון שבכל אחת מהנקודות האלה $f$ מתלכדת עם הנוסחה בסביבה שלמה, $f$ רציפה בהן לכל $a,b$. נותר לבדוק את $x=0$ ואת $x=4$.

\textbf{הנקודה $x=0$.} $f(0)=b$. משמאל, לפי $\lim_{y\to0}\frac{\sin y}{y}=1$ עם $y=8x$:
\[
\lim_{x\to0^-}\frac{\sin(8x)}{x}=8\lim_{x\to0^-}\frac{\sin(8x)}{8x}=8\lim_{y\to0^-}\frac{\sin y}{y}=8\cdot1=8 .
\]
מימין: $\lim_{x\to0^+}(ax+b)=b=f(0)$ (רציפות פולינום). לכן $f$ רציפה ב-0 אם ורק אם $\boxed{b=8}$.

\textbf{הנקודה $x=4$.} $f(4)=4a+b=4a+8$, ומשמאל $\lim_{x\to4^-}(ax+8)=4a+8=f(4)$. מימין, גבול מהצורה $\frac00$. נכפול בצמוד:
\[
\frac{x^2-3x-4}{2-\sqrt x}\cdot\frac{2+\sqrt x}{2+\sqrt x}=\frac{(x-4)(x+1)(2+\sqrt x)}{4-x}=-(x+1)(2+\sqrt x)\qquad(x\neq4),
\]
ולכן, מרציפות הביטוי האחרון ב-4,
\[
\lim_{x\to4^+}f(x)=-(4+1)(2+\sqrt4)=-5\cdot4=-20 .
\]
$f$ רציפה ב-4 אם ורק אם $4a+8=-20$, כלומר $\boxed{a=-7}$.

\textbf{תשובה:} $f$ רציפה ב-$\R$ אם ורק אם $a=-7,\ b=8$.

(הערה: בפתרון בכתב יד המצורף הסימן "$-$" הושמט בשורה $-\lim(x+1)(2+\sqrt x)=5\cdot4=20$, אבל בשורה הבאה נכתב נכון $4a+8=-20$ והתשובה $a=-7$ נכונה.)
\end{solution}

\subsection*{חלק א', שאלה 3}

\begin{question}
הוכיחו כי למשוואה הבאה יש שורש ממשי ומצאו קטע שאורכו לכל היותר 1 המכיל את השורש:
\[
x^2+\ln(x)=0
\]
\end{question}

\begin{hint}
הגדירו $f(x)=x^2+\ln x$ על $(0,\infty)$ וחפשו נקודות שבהן קל לחשב את $\ln$, כמו $x=1$ ו-$x=\frac1e$.
\end{hint}

\begin{solution}
נגדיר $f(x)=x^2+\ln x$ על $(0,\infty)$. $f$ רציפה שם כסכום של פולינום ושל פונקציית הלוגריתם, הרציפה בתחום הגדרתה.

נחשב:
\[
f\!\left(\tfrac1e\right)=\frac1{e^2}+\ln\frac1e=\frac{1}{e^2}-1<0,\qquad f(1)=1^2+\ln1=1>0 ,
\]
כי $e^2>1$.

$f$ רציפה בקטע הסגור $\left[\frac1e,1\right]\subset(0,\infty)$ ומחליפה סימן בקצותיו. לפי משפט ערך הביניים (בולצאנו) קיים $c\in\left(\frac1e,1\right)$ עם $f(c)=0$, כלומר שורש ממשי של המשוואה.

\textbf{הקטע המבוקש:} $\left[\frac1e,1\right]$, שאורכו $1-\frac1e<1$.

(הערה: השורש יחיד, כי $f$ עולה ממש כסכום של $x^2$ ו-$\ln x$, שתיהן עולות ממש על $(0,\infty)$. מספרית $c\approx0.653$. אפשר גם לבחור קטע שלא מצריך את $e$, למשל $\left[\frac12,1\right]$: $f(\frac12)=\frac14-\ln2<0$.)
\end{solution}

\subsection*{חלק ב', שאלה 4}

\begin{question}
תהי $f(x)$ פונקציה זוגית. אילו מהטענות הבאות נכונה?
\begin{enumerate}
  \item הפונקציה $g(x)=x^2f(x)-\frac{1}{|x|+1}$ היא פונקציה זוגית.
  \item הפונקציה $g(x)=x^3f(x)-x|x|$ היא פונקציה זוגית.
  \item הפונקציה $g(x)=x^3f(x)+1$ היא פונקציה זוגית.
  \item הפונקציה $g(x)=\frac{f(x)}{x^2+x}$ היא פונקציה אי-זוגית.
\end{enumerate}
\end{question}

\begin{hint}
מכפלה של זוגית בזוגית היא זוגית; מכפלה של אי-זוגית בזוגית היא אי-זוגית. להפרכה מספיקה דוגמה אחת, למשל $f\equiv1$.
\end{hint}

\begin{solution}
\textbf{התשובה הנכונה: (א)}.

\textbf{(א) נכונה:} תחום ההגדרה של $g$ הוא תחום ההגדרה של $f$, שהוא סימטרי (כי $f$ זוגית). לכל $x$ בתחום, מכיוון ש-$f(-x)=f(x)$, $(-x)^2=x^2$ ו-$|-x|=|x|$:
\[
g(-x)=(-x)^2f(-x)-\frac{1}{|-x|+1}=x^2f(x)-\frac{1}{|x|+1}=g(x).
\]

\textbf{למה האחרות שגויות} (הטענות צריכות להתקיים לכל $f$ זוגית, ולכן מספיקה דוגמה נגדית אחת; ניקח $f(x)\equiv1$, שהיא זוגית):
\begin{itemize}
  \item (ב) — $g(x)=x^3-x|x|$: $g(2)=8-4=4$ ואילו $g(-2)=-8+4=-4\neq g(2)$. (באופן כללי $g$ זו היא אי-זוגית, כהפרש של שתי פונקציות אי-זוגיות.)
  \item (ג) — $g(x)=x^3+1$: $g(1)=2$ ואילו $g(-1)=0$.
  \item (ד) — $g(x)=\frac{1}{x^2+x}=\frac{1}{x(x+1)}$ מוגדרת ב-$x=1$ אבל לא ב-$x=-1$, ולכן תחום ההגדרה אינו סימטרי ו-$g$ אינה אי-זוגית. (גם ללא שיקול התחום: $g(2)=\frac16$ ו-$g(-2)=\frac{1}{4-2}=\frac12\neq-\frac16=-g(2)$.)
\end{itemize}
\end{solution}

\subsection*{חלק ב', שאלה 5}

\begin{question}
תהי
\[
f(x)=\begin{cases}
3x-1 & x<0\\
x^2 & x>0
\end{cases}
\]
אילו מהטענות הבאות נכונה?
\begin{enumerate}
  \item הפונקציה אינה מונוטונית.
  \item הפונקציה מונוטונית יורדת ממש.
  \item הפונקציה מונוטונית עולה ממש.
  \item הפונקציה מונוטונית אך לא ממש.
\end{enumerate}
\end{question}

\begin{hint}
בדקו כל חלק בנפרד, ואז השוו ערך כלשהו משמאל ל-0 לערך כלשהו מימין ל-0.
\end{hint}

\begin{solution}
\textbf{התשובה הנכונה: (ג)}.

(שימו לב: כפי שהפונקציה כתובה, היא אינה מוגדרת ב-$x=0$; בפתרון המצורף נכתב $x\le0$ בחלק הראשון. התשובה זהה בשני המקרים.)

יהיו $x_1<x_2$ בתחום ההגדרה. נבחין בשלושה מקרים:
\begin{itemize}
  \item $x_1<x_2<0$: $f(x_2)-f(x_1)=3(x_2-x_1)>0$.
  \item $0<x_1<x_2$: $f(x_2)-f(x_1)=x_2^2-x_1^2=(x_2-x_1)(x_2+x_1)>0$, כי שני הגורמים חיוביים.
  \item $x_1<0<x_2$: $f(x_1)=3x_1-1<-1<0<x_2^2=f(x_2)$.
\end{itemize}
(אם מגדירים $f(0)=3\cdot0-1=-1$, גם המקרים $x_1<0=x_2$ ו-$x_1=0<x_2$ מקיימים $f(x_1)<f(x_2)$ באותו אופן.)
בכל המקרים $f(x_1)<f(x_2)$, ולכן $f$ עולה ממש.

\textbf{למה האחרות שגויות:} (א) — הראינו ש-$f$ מונוטונית. (ב) — למשל $f(1)=1>f(-1)=-4$, ולכן אינה יורדת. (ד) — "מונוטונית אך לא ממש" אומר שקיימים $x_1<x_2$ עם $f(x_1)=f(x_2)$, וזה לא קורה כי $f$ עולה ממש.
\end{solution}

\subsection*{חלק ב', שאלה 6}

\begin{question}
תהי $f(x)$ הפונקציה שהוגדרה בשאלה קודמת. אילו מהטענות הבאות נכונה?
\begin{enumerate}
  \item הפונקציה חח"ע ועל $\R$.
  \item הפונקציה חח"ע אך אינה על $\R$.
  \item הפונקציה על $\R$ אך אינה חח"ע.
  \item הפונקציה אינה חח"ע ואינה על $\R$.
\end{enumerate}
\end{question}

\begin{solution}
\textbf{התשובה הנכונה: (ב)}.

(הפונקציה מהשאלה הקודמת: $f(x)=3x-1$ עבור $x<0$ ו-$f(x)=x^2$ עבור $x>0$.)

\textbf{חח"ע:} בשאלה הקודמת הראינו ש-$f$ עולה ממש, ופונקציה עולה ממש היא חח"ע: אם $x_1\neq x_2$, נניח $x_1<x_2$, אז $f(x_1)<f(x_2)$ ובפרט $f(x_1)\neq f(x_2)$.

\textbf{אינה על $\R$:} עבור $x<0$ מתקיים $f(x)=3x-1<-1$, ועבור $x>0$ מתקיים $f(x)=x^2>0$. לכן אף ערך בקטע $[-1,0]$ אינו מתקבל; למשל אין $x$ עם $f(x)=-\frac12$. (התמונה היא $(-\infty,-1)\cup(0,\infty)$; אם מגדירים $f(0)=-1$ כמו בפתרון המצורף, התמונה היא $(-\infty,-1]\cup(0,\infty)$, ועדיין $-\frac12$ אינו מתקבל.)

\textbf{למה האחרות שגויות:} (א) ו-(ג) טוענות ש-$f$ על $\R$ — הפרכנו. (ג) ו-(ד) טוענות ש-$f$ אינה חח"ע — הפרכנו.
\end{solution}

\subsection*{חלק ב', שאלה 7}

\begin{question}
בדקו רציפות של הפונקציה הבאה בנקודה 0:
\[
f(x)=\begin{cases}
\sin(x)\cos(\frac{1}{x}) & x<0\\
\frac{1}{1+x^2} & x\ge 0
\end{cases}
\]
אילו מהטענות הבאות נכונה?
\begin{enumerate}
  \item הפונקציה רציפה ב-0.
  \item הפונקציה בעלת אי-רציפות סליקה ב-0.
  \item הפונקציה בעלת אי-רציפות מסוג קפיצה ב-0.
  \item הפונקציה בעלת אי-רציפות עיקרית ב-0.
\end{enumerate}
\end{question}

\begin{hint}
משמאל: פונקציה השואפת ל-0 כפול פונקציה חסומה.
\end{hint}

\begin{solution}
\textbf{התשובה הנכונה: (ג)}.

\textbf{משמאל:} $\lim_{x\to0^-}\sin x=0$ (רציפות $\sin$) ו-$\left|\cos\frac1x\right|\le1$ לכל $x\neq0$. לפי המשפט "אפסה כפול חסומה שואפת ל-0" (או בסנדוויץ' $-|\sin x|\le\sin x\cos\frac1x\le|\sin x|$):
\[
\lim_{x\to0^-}f(x)=\lim_{x\to0^-}\sin(x)\cos\left(\tfrac1x\right)=0 .
\]
\textbf{מימין:} $\frac{1}{1+x^2}$ רציפה, ולכן $\lim_{x\to0^+}f(x)=\frac{1}{1+0}=1=f(0)$.

שני הגבולות החד-צדדיים קיימים וסופיים אך שונים ($0\neq1$), ולכן יש ב-0 אי-רציפות מסוג קפיצה.

\textbf{למה האחרות שגויות:} (א) ו-(ב) — שתיהן דורשות שהגבול הדו-צדדי ב-0 יהיה קיים, והוא אינו קיים. (ד) — אף על פי ש-$\cos\frac1x$ לבדה אינה בעלת גבול ב-0, המכפלה עם $\sin x$ כן שואפת ל-0, ולכן שני הגבולות החד-צדדיים קיימים וסופיים; אי-רציפות עיקרית דורשת שאחד מהם לא יהיה קיים כגבול סופי.
\end{solution}

\section*{תשפ"ה סמסטר א בוחן}

\subsection*{חלק א', שאלה 1}

\begin{question}
חשבו את הגבול הבא:
\[
\lim_{x\to\infty}\left(\sqrt{x^2+5x}-\sqrt{x^2-5x}\right)
\]
\end{question}

\begin{hint}
זהו גבול מהצורה $\infty-\infty$. כפלו וחלקו בצמוד $\sqrt{x^2+5x}+\sqrt{x^2-5x}$.
\end{hint}

\begin{solution}
עבור $x>5$ שני השורשים מוגדרים וחיוביים. זהו גבול מהצורה $\infty-\infty$, ולכן נכפול ונחלק בביטוי הצמוד (החיובי):
\[
\sqrt{x^2+5x}-\sqrt{x^2-5x}
=\frac{(x^2+5x)-(x^2-5x)}{\sqrt{x^2+5x}+\sqrt{x^2-5x}}
=\frac{10x}{\sqrt{x^2+5x}+\sqrt{x^2-5x}} .
\]
נחלק מונה ומכנה ב-$x$. מכיוון ש-$x>0$ מתקיים $x=\sqrt{x^2}$, ולכן
\[
=\frac{10}{\sqrt{1+\frac5x}+\sqrt{1-\frac5x}} .
\]
כאשר $x\to\infty$: $\frac5x\to 0$, ומרציפות פונקציית השורש ב-1 נקבל $\sqrt{1\pm\frac5x}\to 1$. לפי אריתמטיקה של גבולות (המכנה שואף ל-$2\neq 0$):
\[
\lim_{x\to\infty}\left(\sqrt{x^2+5x}-\sqrt{x^2-5x}\right)=\frac{10}{1+1}=\boxed{5}.
\]
\end{solution}

\subsection*{חלק א', שאלה 2}

\begin{question}
יהיו $a,b\in\R$. נגדיר פונקציה על ידי
\[
f(x)=\begin{cases}
\frac{x\sin(ax)}{1-\cos x} & x<0\\
b & x=0\\
\frac{e^x-1}{xe^x} & x>0
\end{cases}
\]
לאילו ערכים של $a,b$ הפונקציה רציפה ב-$0$ ? נמקו.
\end{question}

\begin{hint}
$f$ רציפה ב-0 אם ורק אם שני הגבולות החד-צדדיים ב-0 קיימים ושווים ל-$f(0)=b$.
\end{hint}

\begin{hint}
השתמשו בגבולות היסודיים $\frac{\sin t}{t}\to1$, $\frac{1-\cos x}{x^2}\to\frac12$, $\frac{e^x-1}{x}\to1$ (אסור להשתמש בכלל לופיטל בבוחן).
\end{hint}

\begin{solution}
הפונקציה רציפה ב-0 אם ורק אם
\[
\lim_{x\to0^-}f(x)=\lim_{x\to0^+}f(x)=f(0)=b .
\]

\textbf{הגבול משמאל.} נשתמש בגבול היסודי
\[
\lim_{x\to0}\frac{1-\cos x}{x^2}=\frac12 ,
\]
הנובע מהזהות $1-\cos x=2\sin^2\frac x2$: $\frac{1-\cos x}{x^2}=\frac12\left(\frac{\sin(x/2)}{x/2}\right)^2\to\frac12$.
אם $a\neq0$ נכתוב, עבור $x<0$ קרוב ל-0:
\[
\frac{x\sin(ax)}{1-\cos x}=a\cdot\frac{\sin(ax)}{ax}\cdot\frac{x^2}{1-\cos x}\xrightarrow[x\to0^-]{}a\cdot1\cdot2=2a ,
\]
לפי $\lim_{t\to0}\frac{\sin t}{t}=1$ (הצבה $t=ax\to0$) ואריתמטיקה של גבולות.
אם $a=0$ אז $f(x)=0$ לכל $x<0$ והגבול משמאל הוא $0=2a$. בכל מקרה
\[
\lim_{x\to0^-}f(x)=2a .
\]

\textbf{הגבול מימין.}
\[
\frac{e^x-1}{xe^x}=\frac{e^x-1}{x}\cdot\frac{1}{e^x}\xrightarrow[x\to0^+]{}1\cdot\frac{1}{e^0}=1 ,
\]
לפי הגבול היסודי $\lim_{x\to0}\frac{e^x-1}{x}=\ln e=1$ (מדף הנוסחאות) ורציפות $e^x$ ב-0.

\textbf{מסקנה.} $f$ רציפה ב-0 אם ורק אם $2a=1=b$, כלומר
\[
\boxed{a=\tfrac12,\quad b=1.}
\]
\end{solution}

\subsection*{חלק א', שאלה 3}

\begin{question}
חשבו את הפונקציה ההופכית לפונקציה $f(x)=\ln(1+\sqrt{x})$ ורשמו את תחום ההגדרה של הפונקציה ההפוכה.
\end{question}

\begin{hint}
תחום ההגדרה של הפונקציה ההפוכה הוא התמונה של $f$. מצאו את התמונה, ואז חלצו את $x$ מהמשוואה $y=\ln(1+\sqrt x)$.
\end{hint}

\begin{solution}
\textbf{תחום ההגדרה והתמונה של $f$.} $\sqrt x$ מוגדר עבור $x\ge0$, ואז $1+\sqrt x\ge1>0$ ולכן $\ln(1+\sqrt x)$ מוגדר. לכן $D_f=[0,\infty)$.

$f$ עולה ממש: $\sqrt x$ עולה ממש על $[0,\infty)$ ו-$\ln$ עולה ממש, והרכבה של פונקציות עולות ממש היא עולה ממש. בפרט $f$ חד-חד-ערכית ולכן הפיכה על התמונה שלה.

התמונה: $f(0)=\ln 1=0$ ולכן $f(x)\ge0$ לכל $x\ge0$. לכל $y\ge0$ נמצא בהמשך $x\ge0$ עם $f(x)=y$, ולכן $\operatorname{Im}f=[0,\infty)$.

\textbf{חישוב ההופכית.} יהי $y\ge0$. נפתור $y=\ln(1+\sqrt x)$:
\[
e^y=1+\sqrt x\iff \sqrt x=e^y-1 .
\]
מכיוון ש-$y\ge0$ מתקיים $e^y-1\ge0$, ולכן יש פתרון יחיד $x=(e^y-1)^2\ge0$. (עבור $y<0$ היה מתקבל $\sqrt x<0$ — אין פתרון, כצפוי מכך ש-$y$ אינו בתמונה.)

לכן
\[
\boxed{f^{-1}(x)=\left(e^x-1\right)^2,\qquad D_{f^{-1}}=[0,\infty).}
\]
\textbf{בדיקה:} עבור $x\ge0$: $f(f^{-1}(x))=\ln\left(1+\sqrt{(e^x-1)^2}\right)=\ln\left(1+|e^x-1|\right)=\ln(e^x)=x$, כי $e^x-1\ge0$. שימו לב שהגבלת התחום ל-$[0,\infty)$ הכרחית: הנוסחה $(e^x-1)^2$ מוגדרת לכל $x$, אבל עבור $x<0$ היא אינה הפכית של $f$.
\end{solution}

\subsection*{חלק ב', שאלה 4}

\begin{question}
יהיו
\[
f(x)=e^x,\qquad g(x)=\cos(2x)
\]
אילו מהטענות הבאות נכונה?
\begin{enumerate}
  \item הפונקציה $f\circ g$ מחזורית.
  \item הפונקציה $f\circ g$ זוגית.
  \item הפונקציה $f\circ g$ חסומה ב-$\R$.
  \item כל התשובות נכונות
\end{enumerate}
\end{question}

\begin{solution}
\textbf{התשובה הנכונה: (ד)}.

$(f\circ g)(x)=f(g(x))=e^{\cos 2x}$, מוגדרת לכל $x\in\R$. נבדוק כל טענה:
\begin{itemize}
  \item \textbf{(א) נכונה:} $\cos$ מחזורית עם מחזור $2\pi$, ולכן $\cos(2(x+\pi))=\cos(2x+2\pi)=\cos 2x$, ומכאן $(f\circ g)(x+\pi)=e^{\cos 2x}=(f\circ g)(x)$ לכל $x$. כלומר $f\circ g$ מחזורית (עם מחזור $\pi$).
  \item \textbf{(ב) נכונה:} $\cos$ זוגית, ולכן $(f\circ g)(-x)=e^{\cos(-2x)}=e^{\cos 2x}=(f\circ g)(x)$ לכל $x$.
  \item \textbf{(ג) נכונה:} $-1\le\cos 2x\le 1$ ו-$e^t$ עולה ממש, ולכן $e^{-1}\le e^{\cos 2x}\le e$ לכל $x\in\R$.
\end{itemize}
מכיוון ששלוש הטענות נכונות, התשובה היא (ד). (כל אחת מ-(א), (ב), (ג) לבדה אינה התשובה המלאה, שכן גם האחרות נכונות.)
\end{solution}

\subsection*{חלק ב', שאלה 5}

\begin{question}
תהי $f(x)=3^{-3x+3}$. אילו מהטענות הבאות נכונה?
\begin{enumerate}
  \item הפונקציה אינה מונוטונית.
  \item הפונקציה מונוטונית יורדת ממש.
  \item הפונקציה מונוטונית עולה ממש.
  \item הפונקציה מונוטונית אך לא ממש.
\end{enumerate}
\end{question}

\begin{solution}
\textbf{התשובה הנכונה: (ב)}.

$f$ היא הרכבה $f=h\circ u$ של $u(x)=-3x+3$ ו-$h(t)=3^t$. הפונקציה $u$ יורדת ממש (פונקציה לינארית עם שיפוע $-3<0$) והפונקציה $h$ עולה ממש (פונקציה מעריכית עם בסיס $3>1$).

לכן אם $x_1<x_2$ אז $u(x_1)>u(x_2)$, ומכאן $3^{u(x_1)}>3^{u(x_2)}$, כלומר $f(x_1)>f(x_2)$: $f$ יורדת ממש על $\R$.
(לחלופין: $f(x)=27\cdot\left(\frac{1}{27}\right)^{x}$, פונקציה מעריכית עם בסיס $\frac1{27}<1$ כפול קבוע חיובי.)

\textbf{למה האחרות שגויות:} (א) — הראינו שהיא מונוטונית. (ג) — למשל $f(0)=27>f(1)=1$, ולכן אינה עולה. (ד) — "מונוטונית אך לא ממש" אומר שיש $x_1<x_2$ עם $f(x_1)=f(x_2)$, אבל $f$ יורדת ממש ולכן חד-חד-ערכית.
\end{solution}

\subsection*{חלק ב', שאלה 6}

\begin{question}
תהי $f(x)=x^3+x^2$. אילו מהטענות הבאות נכונה?
\begin{enumerate}
  \item הפונקציה חח"ע ועל $\R$.
  \item הפונקציה חח"ע אך אינה על $\R$.
  \item הפונקציה על $\R$ אך אינה חח"ע.
  \item הפונקציה אינה חח"ע ואינה על $\R$.
\end{enumerate}
\end{question}

\begin{hint}
פרקו: $x^3+x^2=x^2(x+1)$. כמה שורשים יש?
\end{hint}

\begin{solution}
\textbf{התשובה הנכונה: (ג)}.

\textbf{אינה חח"ע:} $f(x)=x^2(x+1)$, ולכן $f(0)=0=f(-1)$ עם $0\neq-1$.

\textbf{על $\R$:} $f$ פולינום ולכן רציפה על $\R$. כמו כן
\[
\lim_{x\to\infty}x^3\left(1+\tfrac1x\right)=\infty,\qquad \lim_{x\to-\infty}x^3\left(1+\tfrac1x\right)=-\infty .
\]
יהי $y\in\R$. לפי הגדרת הגבולות האינסופיים קיימים $x_1<x_2$ עם $f(x_1)<y<f(x_2)$. לפי משפט ערך הביניים ($f$ רציפה על $[x_1,x_2]$) קיים $c\in(x_1,x_2)$ עם $f(c)=y$. לכן $f$ על $\R$.

\textbf{למה האחרות שגויות:} (א) ו-(ב) טוענות ש-$f$ חח"ע — הפרכנו. (ד) טוענת ש-$f$ אינה על — הוכחנו שהיא על.
\end{solution}

\subsection*{חלק ב', שאלה 7}

\begin{question}
בדקו רציפות של הפונקציה הבאה בנקודה 2:
\[
f(x)=\begin{cases}
\frac{1}{1+4^{\frac{1}{x-2}}} & x\neq 2\\
0 & x=2
\end{cases}
\]
אילו מהטענות הבאות נכונה?
\begin{enumerate}
  \item הפונקציה רציפה ב-2.
  \item הפונקציה בעלת אי-רציפות סליקה ב-2.
  \item הפונקציה בעלת אי-רציפות מסוג קפיצה ב-2.
  \item הפונקציה בעלת אי-רציפות עיקרית ב-2.
\end{enumerate}
\end{question}

\begin{hint}
חשבו בנפרד את הגבולות החד-צדדיים: לאן שואף $\frac{1}{x-2}$ כאשר $x\to2^+$ וכאשר $x\to2^-$?
\end{hint}

\begin{solution}
\textbf{התשובה הנכונה: (ג)}.

נחשב את הגבולות החד-צדדיים ב-2.

\textbf{מימין:} כאשר $x\to2^+$, $x-2\to0^+$ ולכן $\frac1{x-2}\to+\infty$. מכיוון ש-$4>1$, $\lim_{t\to+\infty}4^t=\infty$, ולכן $1+4^{\frac{1}{x-2}}\to\infty$ ו-
\[
\lim_{x\to2^+}f(x)=0 .
\]
\textbf{משמאל:} כאשר $x\to2^-$, $\frac1{x-2}\to-\infty$ ו-$\lim_{t\to-\infty}4^t=0$, ולכן
\[
\lim_{x\to2^-}f(x)=\frac{1}{1+0}=1 .
\]
שני הגבולות החד-צדדיים קיימים (סופיים) אך שונים, ולכן הגבול $\lim_{x\to2}f(x)$ אינו קיים, וזו בדיוק אי-רציפות מסוג קפיצה (מסוג ראשון לא סליקה).

\textbf{למה האחרות שגויות:} (א) — הגבול ב-2 לא קיים, ולכן $f$ אינה רציפה ב-2 (העובדה ש-$f(2)=0$ שווה לגבול מימין נותנת רק רציפות מימין). (ב) — אי-רציפות סליקה דורשת שהגבול הדו-צדדי יהיה קיים. (ד) — אי-רציפות עיקרית (מסוג שני) פירושה שלפחות אחד הגבולות החד-צדדיים אינו קיים כגבול סופי, וכאן שניהם קיימים.
\end{solution}

\section*{תשפ"ה סמסטר א מבחן לדוגמה}

\subsection*{שאלה 1א}

\begin{question}
נגדיר פונקציה על ידי
\[ f(x) = \begin{cases} x^3 \sin(\frac{5}{x}) & x \ne 0 \\ 0 & x = 0 \end{cases} \]
האם הפונקציה רציפה וגזירה בנקודה $x = 0$? נמקו.
\end{question}

\begin{hint}
השתמשו בכך ש-$\sin$ חסומה: "אפסה כפול חסומה" שואפת לאפס. לגזירות חשבו את גבול מנת ההפרשים לפי ההגדרה.
\end{hint}

\begin{solution}
\textbf{רציפות.} לכל $x \ne 0$ מתקיים $\left|\sin\frac{5}{x}\right| \le 1$, ולכן
\[ 0 \le |f(x)| = |x|^3\left|\sin\tfrac{5}{x}\right| \le |x|^3 \xrightarrow[x\to0]{} 0. \]
לפי כלל הסנדוויץ' $\lim_{x\to0} f(x) = 0 = f(0)$, ולכן $f$ רציפה ב-$x = 0$.

\textbf{גזירות.} לפי הגדרת הנגזרת:
\[ \frac{f(x) - f(0)}{x - 0} = \frac{x^3\sin\frac5x}{x} = x^2\sin\frac5x, \qquad 0 \le \left| x^2 \sin\tfrac5x \right| \le x^2 \xrightarrow[x\to0]{} 0. \]
שוב לפי כלל הסנדוויץ' (אפסה כפול חסומה), הגבול קיים ושווה $0$. לכן $f$ גזירה ב-$x = 0$ ו-$f'(0) = 0$.

\textbf{תשובה:} $f$ רציפה וגזירה ב-$x = 0$, ו-$f'(0) = 0$. (גזירות גוררת רציפות, כך שהרציפות נובעת גם מהגזירות.)
\end{solution}

\subsection*{שאלה 1ב}

\begin{question}
חשבו את הגבול
\[ \lim_{x\to\frac{\pi}{4}} (\tan x)^{\tan(2x)} \]
\end{question}

\begin{hint}
זה ביטוי $1^\infty$. כתבו $(\tan x)^{\tan 2x} = e^{\tan(2x)\ln(\tan x)}$ וחשבו את גבול המעריך (למשל בהצבה $t = \tan x$ ובכלל לופיטל).
\end{hint}

\begin{solution}
בסביבה מנוקבת של $\frac{\pi}{4}$ מתקיים $\tan x > 0$, ולכן הביטוי מוגדר ו-
\[ (\tan x)^{\tan 2x} = e^{\tan(2x)\cdot\ln(\tan x)}. \]
כאשר $x \to \frac{\pi}{4}$: $\tan x \to 1$ ו-$|\tan 2x| \to \infty$ (ביטוי $1^\infty$). נחשב את גבול המעריך. לפי נוסחת הזווית הכפולה $\tan 2x = \frac{2\tan x}{1 - \tan^2 x}$; נציב $t = \tan x$, ואז $t \to 1$ כאשר $x \to \frac{\pi}{4}$ ($\tan$ רציפה), ו-$t \ne 1$ בסביבה מנוקבת:
\[ \lim_{x\to\frac{\pi}{4}} \tan(2x)\ln(\tan x) = \lim_{t\to1} \frac{2t\ln t}{1 - t^2}. \]
זה ביטוי $\frac00$; לפי כלל לופיטל:
\[ \lim_{t\to1} \frac{2t\ln t}{1 - t^2} = \lim_{t\to1} \frac{2\ln t + 2}{-2t} = \frac{0 + 2}{-2} = -1. \]
(לחלופין: $\frac{2t\ln t}{1 - t^2} = -\frac{2t}{1 + t}\cdot\frac{\ln t}{t - 1} \to -1\cdot 1$.)
מרציפות פונקציית האקספוננט:
\[ \boxed{\lim_{x\to\frac{\pi}{4}} (\tan x)^{\tan(2x)} = e^{-1} = \frac1e}. \]
שימו לב שהגבול הוא דו-צדדי: אף ש-$\tan 2x \to +\infty$ משמאל ו-$\tan 2x \to -\infty$ מימין, המכפלה $\tan(2x)\ln(\tan x)$ שואפת ל-$-1$ משני הצדדים.
\end{solution}

\subsection*{שאלה 2א}

\begin{question}
מצאו את כל האסימפטוטות של הפונקציה
\[ f(x) = \frac{2x^2 - 3x + \ln(x)}{x} \]
נמקו.
\end{question}

\begin{hint}
קבעו את תחום ההגדרה, ובדקו אסימפטוטה אנכית ב-$x = 0^+$ ואסימפטוטה משופעת $y = ax + b$ כאשר $x \to +\infty$.
\end{hint}

\begin{solution}
נכתוב
\[ f(x) = \frac{2x^2 - 3x + \ln x}{x} = 2x - 3 + \frac{\ln x}{x}. \]
\textbf{תחום הגדרה:} $x > 0$ (בגלל $\ln x$; ממילא $x \ne 0$). $f$ רציפה ב-$(0, \infty)$.

\textbf{אסימפטוטה אנכית.} הנקודה היחידה החשודה היא קצה התחום $x = 0$. כאשר $x \to 0^+$: $\ln x \to -\infty$ ו-$\frac1x \to +\infty$, ולכן $\frac{\ln x}{x} = \ln x\cdot\frac1x \to -\infty$, ואילו $2x - 3 \to -3$. לכן
\[ \lim_{x\to0^+} f(x) = -\infty, \]
ו-$x = 0$ אסימפטוטה אנכית (מימין). בכל נקודה אחרת $f$ רציפה ולכן אין שם אסימפטוטה אנכית.

\textbf{אסימפטוטה משופעת.} התחום אינסופי רק ב-$+\infty$. נחשב
\[ a = \lim_{x\to\infty}\frac{f(x)}{x} = \lim_{x\to\infty}\left( 2 - \frac3x + \frac{\ln x}{x^2} \right) = 2, \]
\[ b = \lim_{x\to\infty}\bigl(f(x) - 2x\bigr) = \lim_{x\to\infty}\left( -3 + \frac{\ln x}{x} \right) = -3, \]
כאשר השתמשנו ב-$\lim_{x\to\infty}\frac{\ln x}{x} = 0$ (לפי לופיטל: $\lim \frac{1/x}{1} = 0$), ומכאן גם $\frac{\ln x}{x^2} \to 0$. לכן $y = 2x - 3$ אסימפטוטה משופעת ב-$+\infty$, ואין אסימפטוטה אופקית.

\textbf{תשובה:} אסימפטוטה אנכית $x = 0$, ואסימפטוטה משופעת $y = 2x - 3$ כאשר $x \to +\infty$.
\end{solution}

\subsection*{שאלה 2ב}

\begin{question}
מצאו תחומי קמירות ונקודות פיתול של הפונקציה
\[ f(x) = x^2 \cdot e^{-x} \]
נמקו.
\end{question}

\begin{hint}
$f''(x) = e^{-x}(x^2 - 4x + 2)$; בדקו את סימן הגורם הריבועי.
\end{hint}

\begin{solution}
$f(x) = x^2 e^{-x}$ מוגדרת וגזירה אינסוף פעמים בכל $\R$.
\[ f'(x) = 2xe^{-x} - x^2e^{-x} = e^{-x}(2x - x^2), \]
\[ f''(x) = -e^{-x}(2x - x^2) + e^{-x}(2 - 2x) = e^{-x}(x^2 - 4x + 2). \]
מכיוון ש-$e^{-x} > 0$, הסימן של $f''$ הוא הסימן של $x^2 - 4x + 2$, ששורשיו
\[ x_{1,2} = \frac{4 \pm \sqrt{16 - 8}}{2} = 2 \pm \sqrt2. \]
זו פרבולה "מחייכת", ולכן:
\begin{itemize}
  \item $f'' > 0$ ב-$(-\infty, 2 - \sqrt2)$ וב-$(2 + \sqrt2, \infty)$: שם $f$ \textbf{קמורה};
  \item $f'' < 0$ ב-$(2 - \sqrt2, 2 + \sqrt2)$: שם $f$ \textbf{קעורה}.
\end{itemize}
בנקודות $x = 2 \pm \sqrt2$ הפונקציה גזירה ו-$f''$ מחליפה סימן, ולכן אלה נקודות פיתול. הערכים:
\[ f(2 \pm \sqrt2) = (2 \pm \sqrt2)^2 e^{-(2 \pm \sqrt2)} = (6 \pm 4\sqrt2)\,e^{-2 \mp \sqrt2}. \]

\textbf{תשובה:} $f$ קמורה ב-$(-\infty, 2 - \sqrt2)$ וב-$(2 + \sqrt2, \infty)$, קעורה ב-$(2 - \sqrt2, 2 + \sqrt2)$, ונקודות הפיתול הן $x = 2 - \sqrt2$ ו-$x = 2 + \sqrt2$.
\end{solution}

\subsection*{שאלה 3}

\begin{question}
העריכו בעזרת קירוב לינארי (פולינום טיילור ממעלה ראשונה) את $\sqrt[4]{18}$ ותנו חסם לשגיאה.
\end{question}

\begin{hint}
פתחו את $f(x) = \sqrt[4]{x}$ סביב $a = 16$, כי $\sqrt[4]{16} = 2$, וחסמו את שארית לגרנז' $R_1$.
\end{hint}

\begin{solution}
נגדיר $f(x) = \sqrt[4]{x} = x^{1/4}$ ונפתח סביב $a = 16$, הנקודה הקרובה ל-$18$ שבה השורש הרביעי ידוע. עבור $x > 0$:
\[ f'(x) = \frac14 x^{-3/4}, \qquad f''(x) = -\frac{3}{16} x^{-7/4}, \]
ולכן $f(16) = 2$ ו-$f'(16) = \frac{1}{4\cdot 16^{3/4}} = \frac{1}{4\cdot 8} = \frac{1}{32}$.

\textbf{קירוב לינארי:}
\[ P_1(x) = f(16) + f'(16)(x - 16) = 2 + \frac{x - 16}{32}, \qquad
   \sqrt[4]{18} \approx P_1(18) = 2 + \frac{2}{32} = 2 + \frac{1}{16} = \frac{33}{16} = 2.0625. \]

\textbf{חסם לשגיאה.} לפי משפט טיילור עם שארית לגרנז', קיימת $c \in (16, 18)$ כך ש-
\[ R_1(18) = \frac{f''(c)}{2!}(18 - 16)^2 = \frac{-\frac{3}{16}c^{-7/4}}{2}\cdot 4 = -\frac{3}{8\,c^{7/4}}. \]
מכיוון ש-$c > 16$, מתקיים $c^{7/4} > 16^{7/4} = 2^7 = 128$, ולכן
\[ |R_1(18)| = \frac{3}{8c^{7/4}} < \frac{3}{8\cdot128} = \frac{3}{1024} \approx 0.0029. \]

\textbf{תשובה:} $\sqrt[4]{18} \approx \frac{33}{16} = 2.0625$, עם שגיאה קטנה מ-$\frac{3}{1024}$. מכיוון ש-$R_1 < 0$, הקירוב גדול מהערך האמיתי: $\frac{33}{16} - \frac{3}{1024} < \sqrt[4]{18} < \frac{33}{16}$.

(בדיקה: $\sqrt[4]{18} \approx 2.05977$, והשגיאה בפועל כ-$0.00273 < \frac{3}{1024}$.)
\end{solution}

\subsection*{שאלה 4א}

\begin{question}
חשבו את האינטגרל
\[ \int \frac{6x^2 + 13x + 8}{x^3 + 2x^2 + 2x}\,dx \]
\end{question}

\begin{hint}
$x^3 + 2x^2 + 2x = x(x^2 + 2x + 2)$ והגורם הריבועי אי-פריק. פרקו ל-$\frac{A}{x} + \frac{Bx + C}{x^2 + 2x + 2}$.
\end{hint}

\begin{solution}
\textbf{פירוק המכנה:} $x^3 + 2x^2 + 2x = x(x^2 + 2x + 2)$, והדיסקרימיננטה של $x^2 + 2x + 2$ היא $4 - 8 < 0$, לכן הגורם הריבועי אי-פריק ($x^2 + 2x + 2 = (x + 1)^2 + 1 > 0$). מעלת המונה קטנה ממעלת המכנה.

\textbf{שברים חלקיים:}
\[ \frac{6x^2 + 13x + 8}{x(x^2 + 2x + 2)} = \frac{A}{x} + \frac{Bx + C}{x^2 + 2x + 2}, \]
כלומר $6x^2 + 13x + 8 = A(x^2 + 2x + 2) + (Bx + C)x$.
\begin{itemize}
  \item $x = 0$: $8 = 2A$, ולכן $A = 4$.
  \item מקדמי $x^2$: $6 = A + B$, ולכן $B = 2$.
  \item מקדמי $x$: $13 = 2A + C$, ולכן $C = 5$.
\end{itemize}
לכן
\[ \int \frac{6x^2 + 13x + 8}{x^3 + 2x^2 + 2x}\,dx = \int \left( \frac4x + \frac{2x + 2}{x^2 + 2x + 2} + \frac{3}{(x + 1)^2 + 1} \right) dx, \]
כאשר כתבנו $2x + 5 = (2x + 2) + 3$.
\begin{itemize}
  \item $\int \frac4x\,dx = 4\ln|x| + C$.
  \item בהצבה $t = x^2 + 2x + 2$, $dt = (2x + 2)\,dx$: $\int \frac{2x + 2}{x^2 + 2x + 2}\,dx = \ln(x^2 + 2x + 2) + C$.
  \item בהצבה $u = x + 1$: $\int \frac{3\,dx}{(x + 1)^2 + 1} = 3\arctan(x + 1) + C$.
\end{itemize}
\textbf{תשובה:}
\[ \boxed{\int \frac{6x^2 + 13x + 8}{x^3 + 2x^2 + 2x}\,dx = 4\ln|x| + \ln(x^2 + 2x + 2) + 3\arctan(x + 1) + C.} \]
\end{solution}

\subsection*{שאלה 4ב}

\begin{question}
מצאו את הפתרון הכללי של המשוואה הדיפרנציאלית
\[ (x^2 + 4)y' + 2xy^2 = 0 \]
\end{question}

\begin{hint}
המשוואה פרידה: $\frac{dy}{y^2} = -\frac{2x}{x^2 + 4}\,dx$. אל תשכחו את הפתרון $y \equiv 0$.
\end{hint}

\begin{solution}
נכתוב את המשוואה בצורה ($x^2 + 4 > 0$ לכל $x$):
\[ \frac{dy}{dx} = -\frac{2x}{x^2 + 4}\, y^2, \]
וזו משוואה פרידה.

\textbf{הפתרון הקבוע:} $y \equiv 0$ מקיים את המשוואה, ולכן הוא פתרון.

\textbf{עבור $y \ne 0$:}
\[ \int \frac{dy}{y^2} = -\int \frac{2x}{x^2 + 4}\,dx. \]
באגף שמאל $-\frac1y$. באגף ימין, בהצבה $t = x^2 + 4$, $dt = 2x\,dx$: $\int \frac{2x}{x^2 + 4}\,dx = \ln(x^2 + 4) + C$. לכן
\[ -\frac{1}{y} = -\ln(x^2 + 4) + C \quad\Longrightarrow\quad y = \frac{1}{\ln(x^2 + 4) - C}. \]

\textbf{תשובה:} הפתרון הכללי הוא
\[ \boxed{y = \frac{1}{\ln(x^2 + 4) + C}}, \quad C \in \R \]
(שינינו את שם הקבוע; בכל קטע שבו המכנה אינו מתאפס), ובנוסף הפתרון $y \equiv 0$.
\end{solution}

\subsection*{שאלה 5א}

\begin{question}
חשבו את נפח גוף הסיבוב של הפונקציה $f(x) = \sqrt{\cos x} \cdot \sin^2 x$ בקטע $[0, \frac{\pi}{2}]$.
\end{question}

\begin{hint}
בסיבוב סביב ציר ה-$x$: $V = \pi\int_0^{\pi/2} f^2(x)\,dx = \pi\int_0^{\pi/2}\cos x\,\sin^4 x\,dx$; הציבו $t = \sin x$.
\end{hint}

\begin{solution}
נבין את השאלה כסיבוב גרף $f$ סביב ציר ה-$x$. $f$ מוגדרת ורציפה ב-$[0, \frac{\pi}{2}]$ (שם $\cos x \ge 0$), ונפח גוף הסיבוב הוא
\[ V = \pi\int_0^{\pi/2} f^2(x)\,dx = \pi\int_0^{\pi/2} \cos x\,\sin^4 x\,dx. \]
נציב $t = \sin x$, $dt = \cos x\,dx$; כאשר $x = 0$ מתקיים $t = 0$, וכאשר $x = \frac{\pi}{2}$ מתקיים $t = 1$:
\[ V = \pi\int_0^1 t^4\,dt = \pi\left[\frac{t^5}{5}\right]_0^1 = \frac{\pi}{5}. \]
\textbf{תשובה:} $\boxed{V = \frac{\pi}{5}}$.
\end{solution}

\subsection*{שאלה 5ב}

\begin{question}
חשבו את האינטגרל הלא אמיתי $\int_0^\infty xe^{-x}\,dx$, או הראו שאינו מתכנס.
\end{question}

\begin{hint}
מצאו פונקציה קדומה באינטגרציה בחלקים, ואז חשבו את הגבול כאשר הגבול העליון שואף לאינסוף.
\end{hint}

\begin{solution}
האינטגרנד רציף ב-$[0, \infty)$, ולפי ההגדרה
\[ \int_0^\infty xe^{-x}\,dx = \lim_{A\to\infty}\int_0^A xe^{-x}\,dx. \]
\textbf{פונקציה קדומה} באינטגרציה בחלקים, $u = x$, $v' = e^{-x}$ ($u' = 1$, $v = -e^{-x}$):
\[ \int xe^{-x}\,dx = -xe^{-x} + \int e^{-x}\,dx = -xe^{-x} - e^{-x} + C = -(x + 1)e^{-x} + C. \]
לכן, לפי ניוטון-לייבניץ,
\[ \int_0^A xe^{-x}\,dx = \Bigl[-(x + 1)e^{-x}\Bigr]_0^A = 1 - \frac{A + 1}{e^A}. \]
כאשר $A \to \infty$, $\frac{A + 1}{e^A} \to 0$ (ביטוי $\frac{\infty}{\infty}$; לפי לופיטל $\lim \frac{1}{e^A} = 0$).

\textbf{תשובה:} האינטגרל מתכנס ו-$\boxed{\int_0^\infty xe^{-x}\,dx = 1}$.
\end{solution}

\section*{תשפ"ה סמסטר א מועד א}

\subsection*{שאלה 1א}

\begin{question}
נגדיר פונקציה על ידי
\[ f(x) = \begin{cases} ax^2 + b & x \le 0 \\ \dfrac{\ln(e^x - 1)}{2\ln x} & x > 0 \end{cases} \]
לאילו ערכים של הפרמטרים $a, b$ הפונקציה רציפה בנקודה $x=0$? נמקו.
\end{question}

\begin{hint}
כתבו $e^x - 1 = x \cdot \frac{e^x - 1}{x}$ והשתמשו בגבול הידוע $\lim_{x\to 0}\frac{e^x-1}{x} = 1$.
\end{hint}

\begin{solution}
לכל $a,b$ מתקיים $f(0) = a\cdot 0 + b = b$, ולשמאל של $0$ הפונקציה היא הפולינום $ax^2+b$, ולכן
$\lim_{x\to 0^-} f(x) = b = f(0)$. לכן $f$ רציפה ב-$0$ אם ורק אם $\lim_{x\to 0^+} f(x) = b$.

נחשב את הגבול מימין. עבור $x>0$ מתקיים $e^x - 1 > 0$ ולכן
\[ \ln(e^x - 1) = \ln\left( x \cdot \frac{e^x-1}{x} \right) = \ln x + \ln\frac{e^x - 1}{x}. \]
לכן
\[ f(x) = \frac{\ln x + \ln\frac{e^x-1}{x}}{2\ln x} = \frac12 + \frac{\ln\frac{e^x-1}{x}}{2\ln x}. \]
מכיוון ש-$\lim_{x\to 0}\frac{e^x - 1}{x} = 1$ (זו הנגזרת של $e^x$ ב-$0$) ו-$\ln$ רציפה ב-$1$, המונה של השבר השני שואף ל-$\ln 1 = 0$,
בעוד ש-$2\ln x \to -\infty$. לכן השבר השני שואף ל-$0$ ומתקבל
\[ \lim_{x\to 0^+} f(x) = \frac12. \]
(אפשר גם להשתמש בכלל לופיטל במקרה $\frac{\infty}{\infty}$:
$\lim_{x\to0^+}\frac{\ln(e^x-1)}{2\ln x} = \lim_{x\to0^+}\frac{e^x/(e^x-1)}{2/x} = \lim_{x\to 0^+}\frac{e^x}{2}\cdot\frac{x}{e^x-1} = \frac12$.)

\textbf{תשובה:} $f$ רציפה ב-$x=0$ אם ורק אם $b = \frac12$, ו-$a$ כלשהו ($a \in \R$).
\end{solution}

\subsection*{שאלה 1ב}

\begin{question}
כמה פתרונות ממשיים יש למשוואה הבאה?
\[ x + \sqrt{x} = 17 \]
נמקו היטב.
\end{question}

\begin{hint}
הגדירו $g(x) = x + \sqrt{x} - 17$ על $[0,\infty)$: קיום פתרון נובע ממשפט ערך הביניים, ויחידותו ממונוטוניות.
\end{hint}

\begin{solution}
נגדיר $g(x) = x + \sqrt{x} - 17$. המשוואה מוגדרת רק עבור $x \ge 0$, ו-$g$ רציפה על $[0,\infty)$.

\textbf{קיום:} $g(0) = -17 < 0$ ו-$g(17) = \sqrt{17} > 0$. לפי משפט ערך הביניים קיימת $c \in (0,17)$ עם $g(c) = 0$.

\textbf{יחידות:} עבור $x > 0$ מתקיים $g'(x) = 1 + \frac{1}{2\sqrt{x}} > 0$, ולכן (מסקנה ממשפט לגרנז') $g$ עולה ממש על $[0,\infty)$
(למעשה זה ברור גם ישירות: סכום של שתי פונקציות עולות ממש). פונקציה עולה ממש מקבלת כל ערך לכל היותר פעם אחת, ולכן יש לכל היותר פתרון אחד.

\textbf{תשובה:} למשוואה פתרון ממשי \textbf{אחד בדיוק}. (למעשה, בהצבה $t = \sqrt{x}\ge 0$ מקבלים $t^2 + t - 17 = 0$, כלומר $t = \frac{-1+\sqrt{69}}{2}$, ו-$x = \left(\frac{-1+\sqrt{69}}{2}\right)^2 \approx 13.35$.)
\end{solution}

\subsection*{שאלה 2א}

\begin{question}
מצאו את כל האסימפטוטות של הפונקציה
\[ f(x) = \frac{x^3 - 3x^2 + 3x + \sqrt{x}}{x^2 + x - 2} \]
נמקו.
\end{question}

\begin{hint}
שימו לב לתחום ההגדרה: בגלל $\sqrt{x}$ הפונקציה מוגדרת רק עבור $x \ge 0$. אילו אפסים של המכנה נמצאים בתחום?
\end{hint}

\begin{solution}
\textbf{תחום הגדרה.} $\sqrt{x}$ מוגדר עבור $x\ge0$, והמכנה מתפרק: $x^2 + x - 2 = (x-1)(x+2)$. לכן התחום הוא $[0,1)\cup(1,\infty)$
(הנקודה $x=-2$ אינה בתחום בכלל).

\textbf{אסימפטוטות אנכיות.} $f$ רציפה בכל נקודה בתחומה (מנה של פונקציות רציפות), ולכן אסימפטוטה אנכית יכולה להיות רק ב-$x=1$.
שם המונה שואף ל-$1 - 3 + 3 + 1 = 2 \neq 0$ והמכנה שואף ל-$0$, כאשר $x^2+x-2 < 0$ עבור $0\le x<1$ ו-$>0$ עבור $x>1$. לכן
\[ \lim_{x\to1^-} f(x) = -\infty, \qquad \lim_{x\to1^+} f(x) = +\infty, \]
ו-$x = 1$ אסימפטוטה אנכית. ב-$x=0$ הפונקציה מוגדרת ($f(0)=0$) ורציפה מימין, כך שאין שם אסימפטוטה.

\textbf{אסימפטוטות משופעות/אופקיות.} יש לבדוק רק $x\to+\infty$ (הפונקציה אינה מוגדרת עבור $x<0$). נחפש $y = mx + n$:
\[ m = \lim_{x\to\infty}\frac{f(x)}{x} = \lim_{x\to\infty}\frac{x^3 - 3x^2 + 3x + \sqrt{x}}{x^3 + x^2 - 2x}
= \lim_{x\to\infty}\frac{1 - \frac3x + \frac3{x^2} + x^{-5/2}}{1 + \frac1x - \frac2{x^2}} = 1. \]
בחילוק פולינומים $x^3 - 3x^2 + 3x = (x^2+x-2)(x-4) + (9x - 8)$, ולכן
\[ f(x) - x = -4 + \frac{9x - 8 + \sqrt{x}}{x^2 + x - 2} \xrightarrow[x\to\infty]{} -4 + 0 = -4, \]
כי במונה של השבר החזקה הגבוהה היא $x$ ובמכנה $x^2$. לכן $n=-4$.

\textbf{תשובה:} אסימפטוטה אנכית $x = 1$, ואסימפטוטה משופעת $y = x - 4$ כאשר $x \to +\infty$. אין אסימפטוטות נוספות.
\end{solution}

\subsection*{שאלה 2ב}

\begin{question}
מצאו מינימום ומקסימום מוחלטים של הפונקציה
\[ f(x) = e^x \cdot (x-1)^2 \]
או הסבירו מדוע אינם קיימים. נמקו.
\end{question}

\begin{hint}
שימו לב ש-$f(x) \ge 0$ לכל $x$. מה קורה כאשר $x \to \infty$?
\end{hint}

\begin{solution}
$f$ מוגדרת וגזירה על כל $\R$.
\[ f'(x) = e^x(x-1)^2 + 2e^x(x-1) = e^x(x-1)(x+1). \]
נקודות קריטיות: $x = \pm 1$. $f' > 0$ על $(-\infty,-1)$, $f'<0$ על $(-1,1)$, $f'>0$ על $(1,\infty)$.
לכן $x=-1$ מקסימום מקומי ($f(-1) = 4/e$) ו-$x=1$ מינימום מקומי ($f(1)=0$).

\textbf{מינימום מוחלט:} לכל $x$ מתקיים $e^x>0$ ו-$(x-1)^2 \ge 0$, ולכן $f(x) \ge 0 = f(1)$. לכן המינימום המוחלט הוא $0$, ומתקבל ב-$x=1$.

\textbf{מקסימום מוחלט:} $\lim_{x\to\infty} e^x(x-1)^2 = \infty$ (מכפלה של שני גורמים השואפים ל-$\infty$), ולכן $f$ אינה חסומה מלעיל ואין לה מקסימום מוחלט.
(המקסימום המקומי $4/e$ ב-$x=-1$ אינו מוחלט, למשל $f(3) = 4e^3 > 4/e$.)

\textbf{תשובה:} מינימום מוחלט $0$ בנקודה $x=1$; מקסימום מוחלט אינו קיים כי $f(x)\to\infty$ כאשר $x\to\infty$.
\end{solution}

\subsection*{שאלה 3}

\begin{question}
חשבו בקירוב את $\sqrt[3]{2}$, עם שגיאה הקטנה מ-$\frac{1}{10}$. נמקו.
\end{question}

\begin{hint}
השתמשו בפולינום טיילור של $f(x) = \sqrt[3]{x}$ סביב $x_0 = 1$, והעריכו את השגיאה בעזרת שארית לגרנז'.
\end{hint}

\begin{hint}
בדקו: האם קירוב לינארי (סדר 1) מספיק כדי להבטיח שגיאה קטנה מ-$\frac1{10}$? אם לא, עברו לסדר 2.
\end{hint}

\begin{solution}
נגדיר $f(x) = x^{1/3}$ ונפתח סביב $x_0 = 1$ (שם הערכים ידועים). הנגזרות:
\[ f'(x) = \tfrac13 x^{-2/3},\quad f''(x) = -\tfrac29 x^{-5/3},\quad f'''(x) = \tfrac{10}{27}x^{-8/3}, \]
ולכן $f(1) = 1$, $f'(1) = \frac13$, $f''(1) = -\frac29$.

\textbf{ניסיון בסדר 1.} $P_1(2) = 1 + \frac13 = \frac43$, ולפי שארית לגרנז' קיימת $c\in(1,2)$ עם
$R_1 = \frac{f''(c)}{2!}(2-1)^2 = -\frac19 c^{-5/3}$. מכאן רק $|R_1| < \frac19$, וזה \emph{אינו} מבטיח שגיאה קטנה מ-$\frac1{10}$. לכן נעבור לסדר 2.

\textbf{סדר 2.} פולינום טיילור מסדר 2 סביב $1$:
\[ P_2(x) = 1 + \tfrac13(x-1) - \tfrac19(x-1)^2, \qquad P_2(2) = 1 + \tfrac13 - \tfrac19 = \tfrac{11}{9}. \]
לפי משפט טיילור עם שארית לגרנז', קיימת $c \in (1,2)$ כך ש-
\[ \sqrt[3]{2} - \tfrac{11}{9} = R_2(2) = \frac{f'''(c)}{3!}(2-1)^3 = \frac{10}{27\cdot 6}\,c^{-8/3} = \frac{5}{81}\,c^{-8/3}. \]
מכיוון ש-$c > 1$ מתקיים $c^{-8/3} < 1$, ולכן
\[ 0 < R_2(2) < \frac{5}{81} < \frac{1}{10} \qquad (\text{כי } 50 < 81). \]

\textbf{תשובה:} $\sqrt[3]{2} \approx \frac{11}{9} \approx 1.222$, עם שגיאה קטנה מ-$\frac{5}{81}<\frac1{10}$ (ובפרט $\frac{11}{9} < \sqrt[3]{2} < \frac{11}{9} + \frac{5}{81}$).
לבדיקה: $\sqrt[3]{2} \approx 1.2599$, והשגיאה בפועל היא כ-$0.038$.

\textbf{דרך חלופית (קירוב לינארי בנקודה קרובה יותר).} $\left(\frac54\right)^3 = \frac{125}{64}$, שקרוב ל-$2$. קירוב לינארי סביב $x_0 = \frac{125}{64}$:
$f(2) \approx \frac54 + \frac13\cdot\frac{16}{25}\cdot\frac{3}{64} = \frac54 + \frac1{100} = 1.26$, והשגיאה חסומה ע"י $\frac19 c^{-5/3}\left(\frac3{64}\right)^2 < \frac{1}{9}\cdot\frac{9}{4096}<\frac1{10}$ (כי $c>1$).
\end{solution}

\subsection*{שאלה 4א}

\begin{question}
חשבו את האינטגרל
\[ \int \frac{2x^2 + 8}{x^3 - 4x^2 + 8x}\,dx \]
\end{question}

\begin{hint}
פרקו את המכנה: $x^3 - 4x^2 + 8x = x(x^2 - 4x + 8)$, כאשר לגורם הריבועי אין שורשים ממשיים, ופרקו לשברים חלקיים.
\end{hint}

\begin{solution}
\textbf{פירוק המכנה.} $x^3 - 4x^2 + 8x = x(x^2 - 4x + 8)$, ול-$x^2 - 4x + 8 = (x-2)^2 + 4$ אין שורשים ממשיים (הדיסקרימיננטה $16 - 32 < 0$).

\textbf{שברים חלקיים.} נחפש
\[ \frac{2x^2 + 8}{x(x^2 - 4x + 8)} = \frac{A}{x} + \frac{Bx + C}{x^2 - 4x + 8}, \]
כלומר $2x^2 + 8 = A(x^2 - 4x + 8) + x(Bx + C)$. הצבת $x = 0$ נותנת $8 = 8A$, כלומר $A = 1$. אז
$2x^2 + 8 = x^2 - 4x + 8 + Bx^2 + Cx$, ומהשוואת מקדמים $B = 1$, $C = 4$.

\textbf{אינטגרציה.}
\[ \int\frac{dx}{x} = \ln|x|. \]
לגורם השני נכתוב $x + 4 = \frac12(2x - 4) + 6$, כאשר $(x^2 - 4x + 8)' = 2x - 4$:
\[ \int\frac{x+4}{x^2-4x+8}\,dx = \frac12\int\frac{2x-4}{x^2-4x+8}\,dx + 6\int\frac{dx}{(x-2)^2 + 4}
= \frac12\ln(x^2-4x+8) + 6\cdot\frac12\arctan\frac{x-2}{2}, \]
כאשר באינטגרל האחרון השתמשנו ב-$\int\frac{dt}{t^2 + a^2} = \frac1a\arctan\frac ta$ עם $t = x-2$, $a = 2$.

\textbf{תשובה:}
\[ \int \frac{2x^2 + 8}{x^3 - 4x^2 + 8x}\,dx = \ln|x| + \frac12\ln(x^2 - 4x + 8) + 3\arctan\frac{x-2}{2} + C. \]
(בדיקה: גזירת התוצאה מחזירה את האינטגרנד.)
\end{solution}

\subsection*{שאלה 4ב}

\begin{question}
מצאו את הפתרון הכללי של המשוואה הדיפרנציאלית
\[ y' = e^{x+y} \]
\end{question}

\begin{hint}
$e^{x+y} = e^x e^y$, ולכן זו משוואה פרידה (מפרידת משתנים).
\end{hint}

\begin{solution}
$y' = e^x e^y$. מכיוון ש-$e^y > 0$ תמיד, אין פתרונות קבועים (עבור $y\equiv c$ היינו מקבלים $0 = e^{x+c}$, סתירה), ומותר לחלק ב-$e^y$:
\[ e^{-y}\,y' = e^x \quad\Longrightarrow\quad \int e^{-y}\,dy = \int e^x\,dx \quad\Longrightarrow\quad -e^{-y} = e^x + c. \]
נסמן $C = -c$: $e^{-y} = C - e^x$. הצד השמאלי חיובי, ולכן הפתרון מוגדר רק עבור $x$ שמקיימים $e^x < C$ (בפרט נדרש $C > 0$), ושם
\[ y = -\ln\left(C - e^x\right), \qquad C > 0,\ x < \ln C. \]
\textbf{בדיקה:} $y' = \frac{e^x}{C - e^x}$, ומצד שני $e^{x+y} = e^x\cdot\frac{1}{C - e^x}$. שווה.

\textbf{תשובה:} הפתרון הכללי $y(x) = -\ln(C - e^x)$ (באופן שקול: $e^{-y} + e^x = C$), עם קבוע $C>0$, בתחום $x < \ln C$.
\end{solution}

\subsection*{שאלה 5א}

\begin{question}
חשבו את השטח החסום בין הפונקציה
\[ f(x) = \sin^2 x \cdot \cos^3 x \]
וציר ה-$x$ בקטע $[0, \frac{\pi}{2}]$.
\end{question}

\begin{hint}
כתבו $\cos^3 x = (1 - \sin^2 x)\cos x$ והציבו $t = \sin x$.
\end{hint}

\begin{solution}
בקטע $[0,\frac\pi2]$ מתקיים $\sin^2 x \ge 0$ ו-$\cos x\ge 0$, ולכן $f(x) \ge 0$ והשטח הוא פשוט $\int_0^{\pi/2} f(x)\,dx$.
נכתוב $\sin^2 x\cos^3 x = \sin^2 x(1 - \sin^2 x)\cos x$ ונציב $t = \sin x$, $dt = \cos x\,dx$; כאשר $x$ עובר מ-$0$ ל-$\frac\pi2$, $t$ עובר מ-$0$ ל-$1$:
\[ S = \int_0^{\pi/2}\sin^2 x\cos^3 x\,dx = \int_0^1 (t^2 - t^4)\,dt = \left[\frac{t^3}{3} - \frac{t^5}{5}\right]_0^1 = \frac13 - \frac15 = \frac{2}{15}. \]
\textbf{תשובה:} $S = \frac{2}{15}$.
\end{solution}

\subsection*{שאלה 5ב}

\begin{question}
חשבו את האינטגרל הלא אמיתי הבא, או הראו שאינו מתכנס:
\[ \int_1^\infty \frac{1}{(1+x)\cdot\sqrt{x}}\,dx \]
\end{question}

\begin{hint}
הציבו $t = \sqrt{x}$.
\end{hint}

\begin{solution}
לפי ההגדרה, $\int_1^\infty g = \lim_{R\to\infty}\int_1^R g$. נחשב את האינטגרל על $[1,R]$ בהצבה $t = \sqrt{x}$, כלומר $x = t^2$, $dx = 2t\,dt$; כאשר $x$ עובר מ-$1$ ל-$R$, $t$ עובר מ-$1$ ל-$\sqrt R$:
\[ \int_1^R \frac{dx}{(1+x)\sqrt{x}} = \int_1^{\sqrt R}\frac{2t\,dt}{(1+t^2)\,t} = 2\int_1^{\sqrt R}\frac{dt}{1+t^2} = 2\left(\arctan\sqrt R - \arctan 1\right). \]
כאשר $R \to \infty$ גם $\sqrt R \to \infty$, ו-$\arctan\sqrt R \to \frac\pi2$. לכן
\[ \int_1^\infty \frac{dx}{(1+x)\sqrt{x}} = 2\left(\frac\pi2 - \frac\pi4\right) = \frac{\pi}{2}. \]
(התכנסות אפשר גם לראות מראש במבחן ההשוואה: $0 < \frac{1}{(1+x)\sqrt x} < \frac{1}{x^{3/2}}$, ו-$\int_1^\infty x^{-3/2}dx$ מתכנס.)

\textbf{תשובה:} האינטגרל מתכנס וערכו $\frac{\pi}{2}$.
\end{solution}

\section*{תשפ"ה סמסטר א מועד ב}

\subsection*{שאלה 1א}

\begin{question}
עבור מספר ממשי $a$ נגדיר פונקציה על ידי
\[ f(x) = \begin{cases} \dfrac{\sqrt{x+4} - 2}{x^2 + 2x} & x < 0 \\ a\cdot|x - 1| & x \ge 0 \end{cases} \]
(11 נק') לאיזה ערך של $a$ הפונקציה רציפה בנקודה $x = 0$? נמקו.
\end{question}

\begin{hint}
הכפילו מונה ומכנה בצמוד $\sqrt{x+4} + 2$.
\end{hint}

\begin{solution}
$f(0) = a\cdot|0 - 1| = a$, ומימין ל-$0$ (בסביבה ימנית) $f(x) = a|x-1|$ רציפה, ולכן $\lim_{x\to0^+}f(x) = a = f(0)$.
לכן $f$ רציפה ב-$0$ אם ורק אם $\lim_{x\to 0^-} f(x) = a$.

עבור $-2 < x < 0$ נכפיל בצמוד:
\[ \frac{\sqrt{x+4} - 2}{x^2 + 2x} = \frac{(x + 4) - 4}{x(x+2)\left(\sqrt{x+4} + 2\right)} = \frac{1}{(x+2)\left(\sqrt{x+4}+2\right)}. \]
הביטוי האחרון רציף ב-$x=0$, ולכן
\[ \lim_{x\to 0^-} f(x) = \frac{1}{2\cdot(2 + 2)} = \frac18. \]
\textbf{תשובה:} $f$ רציפה ב-$0$ אם ורק אם $a = \frac18$.
\end{solution}

\subsection*{שאלה 1ב}

\begin{question}
עבור מספר ממשי $a$ נגדיר פונקציה על ידי
\[ f(x) = \begin{cases} \dfrac{\sqrt{x+4} - 2}{x^2 + 2x} & x < 0 \\ a\cdot|x - 1| & x \ge 0 \end{cases} \]
(10 נק') לאיזה ערך של $a$ הפונקציה גזירה בנקודה $x = 1$? נמקו.
\end{question}

\begin{hint}
חשבו את הנגזרות החד-צדדיות של $a|x-1|$ ב-$x=1$.
\end{hint}

\begin{solution}
בסביבה של $x = 1$ (למשל עבור $x>0$) $f(x) = a|x - 1|$, ו-$f(1) = 0$. נחשב נגזרות חד-צדדיות:
\[ f'_+(1) = \lim_{h\to0^+}\frac{a|h| - 0}{h} = a, \qquad f'_-(1) = \lim_{h\to0^-}\frac{a|h|}{h} = \lim_{h\to 0^-}\frac{-ah}{h} = -a. \]
$f$ גזירה ב-$1$ אם ורק אם שתי הנגזרות החד-צדדיות קיימות ושוות: $a = -a$, כלומר $a = 0$.
במקרה זה $f \equiv 0$ ב-$[0,\infty)$ ו-$f'(1) = 0$.

\textbf{תשובה:} $f$ גזירה ב-$x=1$ אם ורק אם $a = 0$.
\end{solution}

\subsection*{שאלה 2}

\begin{question}
נגדיר פונקציה
\[ f(x) = \frac{x^2}{e^x} \]
\begin{enumerate}
  \item (11 נק') חשבו נקודות קיצון מקומי, ותחומי עלייה וירידה של הפונקציה. נמקו.
  \item (10 נק') חשבו נקודות פיתול, ותחומי קמירות של הפונקציה. נמקו.
\end{enumerate}
\end{question}

\begin{solution}
$f(x) = x^2e^{-x}$ מוגדרת וגזירה אינסוף פעמים על כל $\R$.
\begin{enumerate}
  \item
  \[ f'(x) = 2xe^{-x} - x^2e^{-x} = x(2 - x)e^{-x}. \]
  מכיוון ש-$e^{-x} > 0$, הסימן של $f'$ הוא הסימן של $x(2-x)$:
  \begin{itemize}
    \item $x < 0$: $f' < 0$, $f$ יורדת;
    \item $0 < x < 2$: $f' > 0$, $f$ עולה;
    \item $x > 2$: $f' < 0$, $f$ יורדת.
  \end{itemize}
  \textbf{תשובה:} $f$ יורדת ב-$(-\infty,0]$, עולה ב-$[0,2]$ ויורדת ב-$[2,\infty)$.
  לפי מבחן הנגזרת הראשונה (החלפת סימן של $f'$): $x = 0$ נקודת מינימום מקומי, $f(0) = 0$; $x = 2$ נקודת מקסימום מקומי, $f(2) = \frac{4}{e^2}$.

  \item
  \[ f''(x) = (2 - 2x)e^{-x} - (2x - x^2)e^{-x} = (x^2 - 4x + 2)e^{-x}. \]
  $f''(x) = 0 \iff x^2 - 4x + 2 = 0 \iff x = 2 \pm \sqrt{2}$. הסימן של $f''$ הוא סימן הפרבולה $x^2 - 4x + 2$ (הפותחת כלפי מעלה):
  \begin{itemize}
    \item $x < 2 - \sqrt2$ או $x > 2+\sqrt2$: $f'' > 0$, $f$ קמורה (כלפי מעלה, $\cup$);
    \item $2-\sqrt2 < x < 2 + \sqrt 2$: $f'' < 0$, $f$ קעורה (כלפי מטה, $\cap$).
  \end{itemize}
  בשתי הנקודות $x = 2\pm\sqrt2$ $f''$ מחליפה סימן, ולכן הן נקודות פיתול.

  \textbf{תשובה:} $f$ קמורה ב-$(-\infty, 2 - \sqrt2]$ וב-$[2+\sqrt2, \infty)$, וקעורה ב-$[2-\sqrt2, 2+\sqrt2]$.
  נקודות הפיתול: $x = 2 \pm \sqrt2$, עם ערכים $f(2\pm\sqrt2) = (6 \pm 4\sqrt2)\,e^{-(2\pm\sqrt2)}$
  (כלומר בערך $(0.586,\ 0.191)$ ו-$(3.414,\ 0.384)$).
\end{enumerate}
\end{solution}

\subsection*{שאלה 3}

\begin{question}
חשבו בעזרת קירוב לינארי את $\ln(1.2)$, והעריכו את השגיאה. נמקו.
\end{question}

\begin{hint}
השתמשו בפולינום טיילור מסדר 1 של $\ln x$ סביב $x_0 = 1$, והעריכו את השגיאה בעזרת שארית לגרנז'.
\end{hint}

\begin{solution}
נגדיר $f(x) = \ln x$ ונבחר $x_0 = 1$ (נקודה קרובה ל-$1.2$ שבה הערכים ידועים). $f(1) = 0$, $f'(x) = \frac1x$, $f'(1) = 1$, $f''(x) = -\frac{1}{x^2}$.

\textbf{קירוב לינארי:} $P_1(x) = f(1) + f'(1)(x - 1) = x - 1$, ולכן
\[ \ln(1.2) \approx P_1(1.2) = 0.2. \]

\textbf{הערכת השגיאה:} לפי משפט טיילור עם שארית לגרנז', קיימת $c\in(1, 1.2)$ כך ש-
\[ \ln(1.2) - 0.2 = R_1(1.2) = \frac{f''(c)}{2!}(1.2 - 1)^2 = -\frac{1}{2c^2}\cdot 0.04 = -\frac{0.02}{c^2}. \]
מכיוון ש-$1 < c < 1.2$ מתקיים $1 < c^2 < 1.44$, ולכן
\[ \frac{0.02}{1.44} < |R_1| < 0.02, \qquad\text{כלומר}\qquad 0.0138 < |R_1| < 0.02, \]
והשגיאה שלילית, כלומר הקירוב $0.2$ גדול מהערך האמיתי.

\textbf{תשובה:} $\ln(1.2) \approx 0.2$, עם שגיאה קטנה מ-$0.02$ (ובפרט $0.18 < \ln(1.2) < 0.2$).
לבדיקה: $\ln(1.2) \approx 0.1823$, והשגיאה בפועל כ-$0.0177$.
\end{solution}

\subsection*{שאלה 4א}

\begin{question}
חשבו את האינטגרל
\[ \int_0^{\frac{\sqrt2}{2}} \frac{x^2}{\sqrt{1 - x^2}}\,dx \]
תזכורת: $\sin(\frac\pi4) = \cos(\frac\pi4) = \frac{\sqrt2}{2}$.
\end{question}

\begin{hint}
הציבו $x = \sin t$, ואחר כך השתמשו בזהות $\sin^2 t = \frac{1 - \cos 2t}{2}$.
\end{hint}

\begin{solution}
נציב $x = \sin t$ עם $t\in[0,\frac\pi4]$ (פונקציה גזירה ועולה, $\sin 0 = 0$, $\sin\frac\pi4 = \frac{\sqrt2}2$), $dx = \cos t\,dt$.
בתחום זה $\cos t > 0$, ולכן $\sqrt{1 - \sin^2 t} = \cos t$:
\[ \int_0^{\frac{\sqrt2}{2}}\frac{x^2}{\sqrt{1-x^2}}\,dx = \int_0^{\pi/4}\frac{\sin^2 t}{\cos t}\cos t\,dt = \int_0^{\pi/4}\sin^2 t\,dt
= \int_0^{\pi/4}\frac{1 - \cos 2t}{2}\,dt = \left[\frac t2 - \frac{\sin 2t}{4}\right]_0^{\pi/4}. \]
לכן
\[ \int_0^{\frac{\sqrt2}{2}}\frac{x^2}{\sqrt{1-x^2}}\,dx = \frac\pi8 - \frac{\sin\frac\pi2}{4} = \frac{\pi}{8} - \frac14. \]
(האינטגרנד רציף על הקטע הסגור, כי $1 - x^2 \ge \frac12$ שם, כך שזה אינטגרל רגיל.)

\textbf{תשובה:} $\frac\pi8 - \frac14 \approx 0.1427$.
\end{solution}

\subsection*{שאלה 4ב}

\begin{question}
מצאו את הפתרון הכללי של המשוואה הדיפרנציאלית
\[ (x^4 + 4)yy' = x^3 \]
\end{question}

\begin{hint}
זו משוואה פרידה: $y\,dy = \frac{x^3}{x^4 + 4}\,dx$.
\end{hint}

\begin{solution}
$x^4 + 4 > 0$ לכל $x$, ולכן אפשר לחלק: $y y' = \frac{x^3}{x^4 + 4}$. זו משוואה פרידה. נבצע אינטגרציה לפי $x$ של שני האגפים:
\[ \int y\,dy = \int\frac{x^3}{x^4 + 4}\,dx \quad\Longrightarrow\quad \frac{y^2}{2} = \frac14\ln(x^4 + 4) + c, \]
כאשר באגף ימין הצבנו $u = x^4 + 4$, $du = 4x^3dx$. כפל ב-$2$ וסימון $C = 2c$:
\[ y^2 = \frac12\ln(x^4 + 4) + C. \]
(שימו לב ש-$y\equiv0$ אינו פתרון, כי אז נקבל $0 = x^3$, שאינו מתקיים לכל $x$. כמו כן, לא נדרשה חלוקה ב-$y$: השתמשנו רק בכך ש-$yy' = \left(\frac{y^2}{2}\right)'$.)

\textbf{תשובה:} הפתרון הכללי נתון בצורה סתומה ע"י $y^2 = \frac12\ln(x^4+4) + C$, או במפורש
\[ y = \pm\sqrt{\tfrac12\ln(x^4 + 4) + C}, \]
בתחום שבו הביטוי בתוך השורש חיובי. \textbf{בדיקה:} גזירת $y^2 = \frac12\ln(x^4+4)+C$ נותנת $2yy' = \frac{2x^3}{x^4+4}$, כלומר $(x^4+4)yy' = x^3$.
\end{solution}

\subsection*{שאלה 5א}

\begin{question}
חשבו את נפח גוף הסיבוב סביב ציר ה-$x$ בקטע $[1,4]$ של הפונקציה
\[ f(x) = \frac{\sqrt{2x+1}}{x^2 + x} \]
\end{question}

\begin{hint}
נפח גוף הסיבוב הוא $V = \pi\int_a^b f^2(x)\,dx$. שימו לב ש-$(x^2 + x)' = 2x + 1$.
\end{hint}

\begin{solution}
נפח גוף הסיבוב של גרף $f$ סביב ציר ה-$x$ מעל $[a,b]$ הוא $V = \pi\int_a^b f(x)^2\,dx$. כאן
$f(x)^2 = \frac{2x+1}{(x^2 + x)^2}$, רציפה על $[1,4]$. נציב $u = x^2 + x$, $du = (2x + 1)\,dx$; כאשר $x$ עובר מ-$1$ ל-$4$, $u$ עובר מ-$2$ ל-$20$:
\[ V = \pi\int_1^4\frac{2x+1}{(x^2+x)^2}\,dx = \pi\int_2^{20}\frac{du}{u^2} = \pi\left[-\frac1u\right]_2^{20} = \pi\left(\frac12 - \frac1{20}\right) = \frac{9\pi}{20}. \]
\textbf{תשובה:} $V = \frac{9\pi}{20}$.
\end{solution}

\subsection*{שאלה 5ב}

\begin{question}
חשבו את האינטגרל הלא אמיתי הבא, או הראו שאינו מתכנס:
\[ \int_1^\infty \frac{\ln x}{x}\,dx \]
\end{question}

\begin{hint}
מצאו פונקציה קדומה בהצבה $u = \ln x$.
\end{hint}

\begin{solution}
לפי ההגדרה $\int_1^\infty\frac{\ln x}{x}dx = \lim_{R\to\infty}\int_1^R\frac{\ln x}{x}dx$. בהצבה $u = \ln x$, $du = \frac{dx}{x}$:
\[ \int_1^R\frac{\ln x}{x}\,dx = \int_0^{\ln R}u\,du = \frac{(\ln R)^2}{2} \xrightarrow[R\to\infty]{} \infty. \]
(לחלופין, במבחן ההשוואה: עבור $x \ge e$ מתקיים $\frac{\ln x}{x} \ge \frac1x \ge 0$, ו-$\int_e^\infty\frac{dx}x$ מתבדר.)

\textbf{תשובה:} האינטגרל אינו מתכנס (מתבדר ל-$\infty$).
\end{solution}

\section*{תשפ"ה סמסטר א מועד מיוחד}

\subsection*{שאלה 1א}

\begin{question}
חשבו את הגבול
\[ \lim_{x\to\infty}\quad x\cdot\left(\sqrt{x^2 + 4} - \sqrt{x^2 - 4}\right) \]
\end{question}

\begin{hint}
הכפילו וחלקו בצמוד $\sqrt{x^2+4} + \sqrt{x^2-4}$.
\end{hint}

\begin{solution}
עבור $x \ge 2$ נכפיל ונחלק בצמוד (שהוא חיובי):
\[ x\left(\sqrt{x^2+4} - \sqrt{x^2-4}\right) = x\cdot\frac{(x^2 + 4) - (x^2 - 4)}{\sqrt{x^2+4} + \sqrt{x^2-4}} = \frac{8x}{\sqrt{x^2+4} + \sqrt{x^2-4}}. \]
נחלק מונה ומכנה ב-$x$ (עבור $x>0$, $\sqrt{x^2 + 4} = x\sqrt{1 + 4/x^2}$):
\[ = \frac{8}{\sqrt{1 + \frac4{x^2}} + \sqrt{1 - \frac4{x^2}}} \xrightarrow[x\to\infty]{} \frac{8}{1 + 1} = 4, \]
לפי אריתמטיקה של גבולות ורציפות השורש.

\textbf{תשובה:} הגבול שווה ל-$4$.
\end{solution}

\subsection*{שאלה 1ב}

\begin{question}
נגדיר פונקציה על ידי
\[ f(x) = \begin{cases} x^3\cdot\cos(\frac1x) & x \neq 0 \\ 0 & x = 0 \end{cases} \]
האם הפונקציה גזירה ב-$x = 0$? \\
אם כן חשבו את הנגזרת שלה בנקודה זו. אם לא נמקו.
\end{question}

\begin{hint}
השתמשו בהגדרת הנגזרת ובכך ש-$|\cos(\frac1x)| \le 1$.
\end{hint}

\begin{solution}
לפי הגדרת הנגזרת, עבור $h \neq 0$:
\[ \frac{f(h) - f(0)}{h} = \frac{h^3\cos(\frac1h)}{h} = h^2\cos\left(\tfrac1h\right). \]
מכיוון ש-$|\cos(\frac1h)| \le 1$ מתקיים $0 \le \left|h^2\cos(\frac1h)\right| \le h^2 \to 0$, ולכן לפי משפט הסנדוויץ' (סדרה/פונקציה חסומה כפול שואפת לאפס)
\[ \lim_{h\to 0}\frac{f(h) - f(0)}{h} = 0. \]
\textbf{תשובה:} כן, $f$ גזירה ב-$x=0$ ו-$f'(0) = 0$.
\end{solution}

\subsection*{שאלה 2}

\begin{question}
נגדיר פונקציה
\[ f(x) = \frac{x}{x^4 + 3} \]
\begin{enumerate}
  \item (11 נק') חשבו נקודות קיצון מקומי, ותחומי עלייה וירידה של הפונקציה בתחום הגדרתה. נמקו.
  \item (10 נק') חשבו את כל האסימפטוטות של הפונקציה. נמקו.
\end{enumerate}
\end{question}

\begin{solution}
תחום ההגדרה: $x^4 + 3 \ge 3 > 0$ לכל $x$, ולכן $f$ מוגדרת, רציפה וגזירה על כל $\R$.
\begin{enumerate}
  \item לפי כלל המנה:
  \[ f'(x) = \frac{(x^4 + 3) - x\cdot 4x^3}{(x^4+3)^2} = \frac{3 - 3x^4}{(x^4 + 3)^2} = \frac{3(1 - x^2)(1 + x^2)}{(x^4+3)^2}. \]
  המכנה ו-$1 + x^2$ חיוביים, ולכן סימן $f'$ הוא סימן $1 - x^2$:
  \begin{itemize}
    \item $|x| < 1$: $f' > 0$, $f$ עולה ב-$[-1, 1]$;
    \item $|x| > 1$: $f' < 0$, $f$ יורדת ב-$(-\infty, -1]$ וב-$[1, \infty)$.
  \end{itemize}
  לפי מבחן הנגזרת הראשונה: $x=-1$ נקודת מינימום מקומי, $f(-1) = -\frac14$; $x = 1$ נקודת מקסימום מקומי, $f(1) = \frac14$.
  (מכיוון ש-$f\to0$ באינסוף, אלה גם הקיצונים המוחלטים.)

  \item \textbf{אנכיות:} $f$ רציפה על כל $\R$, ולכן בכל נקודה $x_0$ מתקיים $\lim_{x\to x_0}f(x) = f(x_0)$ סופי, ואין אסימפטוטות אנכיות.

  \textbf{אופקיות/משופעות:}
  \[ \lim_{x\to\pm\infty}\frac{x}{x^4 + 3} = \lim_{x\to\pm\infty}\frac{1/x^3}{1 + 3/x^4} = 0. \]
  לכן $y = 0$ אסימפטוטה אופקית גם ב-$+\infty$ וגם ב-$-\infty$ (ולכן אין אסימפטוטה משופעת נוספת: $m = \lim \frac{f(x)}{x} = 0$, $n = \lim f(x) = 0$).

  \textbf{תשובה:} האסימפטוטה היחידה היא $y = 0$ (אופקית, בשני הכיוונים); אין אסימפטוטות אנכיות.
\end{enumerate}
\end{solution}

\subsection*{שאלה 3}

\begin{question}
חשבו בעזרת קירוב לינארי את $\sqrt[3]{10}$, והעריכו את השגיאה. נמקו.
\end{question}

\begin{hint}
בחרו נקודה קרובה ל-$10$ שבה השורש השלישי ידוע: $x_0 = 8$. העריכו את השגיאה בעזרת שארית לגרנז'.
\end{hint}

\begin{solution}
נגדיר $f(x) = x^{1/3}$ ונבחר $x_0 = 8$ (כי $\sqrt[3]{8} = 2$). הנגזרות:
\[ f'(x) = \tfrac13x^{-2/3}, \qquad f''(x) = -\tfrac29x^{-5/3}, \]
ולכן $f(8) = 2$, $f'(8) = \frac13\cdot\frac14 = \frac1{12}$.

\textbf{קירוב לינארי:} $P_1(x) = 2 + \frac{1}{12}(x - 8)$, ולכן
\[ \sqrt[3]{10} \approx P_1(10) = 2 + \frac{2}{12} = \frac{13}{6} \approx 2.1667. \]

\textbf{הערכת השגיאה:} לפי משפט טיילור עם שארית לגרנז' קיימת $c \in (8, 10)$ כך ש-
\[ \sqrt[3]{10} - \frac{13}{6} = R_1(10) = \frac{f''(c)}{2!}(10 - 8)^2 = -\frac19c^{-5/3}\cdot 4 = -\frac{4}{9}c^{-5/3}. \]
מכיוון ש-$c > 8$ מתקיים $c^{-5/3} < 8^{-5/3} = \frac1{32}$, ולכן
\[ |R_1(10)| < \frac49\cdot\frac{1}{32} = \frac{1}{72} \approx 0.0139, \]
והשגיאה שלילית, כלומר הקירוב גדול מהערך האמיתי.

\textbf{תשובה:} $\sqrt[3]{10} \approx \frac{13}{6} \approx 2.1667$, עם שגיאה קטנה מ-$\frac{1}{72}$; בפרט $\frac{13}{6} - \frac1{72} < \sqrt[3]{10} < \frac{13}{6}$.
לבדיקה: $\sqrt[3]{10} \approx 2.1544$, והשגיאה בפועל כ-$0.0122$.
\end{solution}

\subsection*{שאלה 4א}

\begin{question}
חשבו את האינטגרל
\[ \int\frac{7x^2 + 6}{x^3 + 2x^2 + 2x}\,dx \]
\end{question}

\begin{hint}
$x^3 + 2x^2 + 2x = x\left((x+1)^2 + 1\right)$. פרקו לשברים חלקיים.
\end{hint}

\begin{solution}
\textbf{פירוק המכנה.} $x^3 + 2x^2 + 2x = x(x^2 + 2x + 2)$, ול-$x^2 + 2x + 2 = (x+1)^2 + 1$ אין שורשים ממשיים.

\textbf{שברים חלקיים.}
\[ \frac{7x^2 + 6}{x(x^2 + 2x + 2)} = \frac{A}{x} + \frac{Bx + C}{x^2 + 2x + 2}, \qquad 7x^2 + 6 = A(x^2 + 2x + 2) + x(Bx + C). \]
הצבת $x = 0$: $6 = 2A$, כלומר $A = 3$. אז $7x^2 + 6 = 3x^2 + 6x + 6 + Bx^2 + Cx$, ומהשוואת מקדמים $B = 4$, $C = -6$.

\textbf{אינטגרציה.} נכתוב $4x - 6 = 2(2x + 2) - 10$, כאשר $(x^2 + 2x + 2)' = 2x + 2$:
\[ \int\frac{4x - 6}{x^2 + 2x + 2}\,dx = 2\int\frac{2x+2}{x^2+2x+2}\,dx - 10\int\frac{dx}{(x+1)^2 + 1} = 2\ln(x^2 + 2x + 2) - 10\arctan(x + 1). \]

\textbf{תשובה:}
\[ \int\frac{7x^2 + 6}{x^3 + 2x^2 + 2x}\,dx = 3\ln|x| + 2\ln(x^2 + 2x + 2) - 10\arctan(x+1) + C. \]
(בדיקה: גזירת התוצאה מחזירה את האינטגרנד.)
\end{solution}

\subsection*{שאלה 4ב}

\begin{question}
מצאו את הפתרון הכללי של המשוואה הדיפרנציאלית
\[ y' - xy^2 = 4x \]
\end{question}

\begin{hint}
העבירו אגף: $y' = x(y^2 + 4)$ — זו משוואה פרידה.
\end{hint}

\begin{solution}
נעביר אגף: $y' = xy^2 + 4x = x(y^2 + 4)$. זו משוואה פרידה. מכיוון ש-$y^2 + 4 > 0$ תמיד, אין פתרונות קבועים (עבור $y\equiv c$: $0 = x(c^2+4)$ לא מתקיים לכל $x$), ומותר לחלק:
\[ \frac{y'}{y^2 + 4} = x \quad\Longrightarrow\quad \int\frac{dy}{y^2 + 4} = \int x\,dx \quad\Longrightarrow\quad \frac12\arctan\frac y2 = \frac{x^2}{2} + c. \]
נכפול ב-$2$ ונסמן $C = 2c$: $\arctan\frac y2 = x^2 + C$, ולכן
\[ y = 2\tan\left(x^2 + C\right), \]
בתחום שבו $x^2 + C \in (-\frac\pi2, \frac\pi2)$ (כי $\arctan$ מקבלת ערכים רק בקטע זה).

\textbf{בדיקה:} $y' = 2\cdot 2x\left(1 + \tan^2(x^2 + C)\right) = 4x + x\cdot 4\tan^2(x^2+C) = 4x + xy^2$. מתקיים.

\textbf{תשובה:} $y(x) = 2\tan(x^2 + C)$, $C \in \R$.
\end{solution}

\subsection*{שאלה 5א}

\begin{question}
חשבו את השטח הסופי החסום בין גרף הפונקציה $f(x) = xe^x$ והישר $y = e\cdot x$.
\end{question}

\begin{hint}
מצאו תחילה את נקודות החיתוך: $xe^x = ex \iff x(e^x - e) = 0$. את $\int xe^x\,dx$ חשבו באינטגרציה בחלקים.
\end{hint}

\begin{solution}
\textbf{נקודות חיתוך:} $xe^x = ex \iff x(e^x - e) = 0 \iff x = 0$ או $e^x = e$, כלומר $x = 0$ או $x = 1$.

\textbf{איזה גרף מעל:} $ex - xe^x = x(e - e^x)$. עבור $0 < x < 1$: $x > 0$ ו-$e^x < e$, ולכן $ex > xe^x$ — הישר מעל הגרף.
(עבור $x < 0$ שתי הפונקציות לא נחתכות שוב, והתחום ביניהן אינסופי ושטחו אינסופי, כי $ex \to -\infty$ בעוד $xe^x \to 0$; לכן "השטח הסופי" הוא זה שבין $0$ ל-$1$.)

\textbf{חישוב:} באינטגרציה בחלקים ($u = x$, $dv = e^xdx$): $\int xe^x\,dx = xe^x - e^x + C$. לכן
\[ S = \int_0^1\left(ex - xe^x\right)dx = \left[\frac{ex^2}{2}\right]_0^1 - \left[xe^x - e^x\right]_0^1 = \frac e2 - \left((e - e) - (0 - 1)\right) = \frac e2 - 1. \]
\textbf{תשובה:} $S = \frac e2 - 1 \approx 0.359$.
\end{solution}

\subsection*{שאלה 5ב}

\begin{question}
חשבו את האינטגרל הלא אמיתי הבא, או הראו שאינו מתכנס:
\[ \int_e^\infty\frac{1}{x\ln^2 x}\,dx \]
\end{question}

\begin{hint}
הציבו $u = \ln x$.
\end{hint}

\begin{solution}
לפי ההגדרה, $\int_e^\infty = \lim_{R\to\infty}\int_e^R$. בהצבה $u = \ln x$, $du = \frac{dx}{x}$ ($x = e \mapsto u = 1$, $x = R\mapsto u = \ln R$):
\[ \int_e^R\frac{dx}{x\ln^2x} = \int_1^{\ln R}\frac{du}{u^2} = \left[-\frac1u\right]_1^{\ln R} = 1 - \frac{1}{\ln R} \xrightarrow[R\to\infty]{} 1, \]
כי $\ln R \to \infty$.

\textbf{תשובה:} האינטגרל מתכנס וערכו $1$.
\end{solution}

\section*{תשפ"ו סמסטר א בוחן}

\subsection*{חלק א', שאלה 1}

\begin{question}
נגדיר פונקציה על ידי
\[
f(x)=\begin{cases}
x^3 & x<1\\
2x+1 & x\ge 1
\end{cases}
\]
\begin{enumerate}
  \item (5 נק') שרטטו את גרף הפונקציה.
  \item (15 נק') האם פונקציה זו הפיכה כפונקציה מ-$\R$ ל-$\R$?
  אם כן - חשבו את הפונקציה ההופכית שלה $f^{-1}:\R\to\R$.
  אם לא - נמקו.
\end{enumerate}
\end{question}

\begin{hint}
פונקציה הפיכה מ-$\R$ ל-$\R$ חייבת להיות גם חח"ע וגם על $\R$. מהי התמונה של כל אחד משני החלקים?
\end{hint}

\begin{solution}
\begin{enumerate}
  \item \textbf{הגרף:} עבור $x<1$ זהו גרף הפרבולה הקובית $y=x^3$: עולה ממש, עובר בנקודות $(-1,-1)$, $(0,0)$ (נקודת פיתול עם משיק אופקי), ומתקרב לנקודה $(1,1)$ — נקודה זו \emph{אינה} שייכת לגרף (עיגול ריק). עבור $x\ge1$ זו קרן ישרה עם שיפוע 2 שמתחילה בנקודה $(1,3)$ (עיגול מלא) ועוברת ב-$(2,5)$. בנקודה $x=1$ יש "קפיצה" מגובה 1 לגובה 3.
  \item \textbf{הפונקציה אינה הפיכה כפונקציה מ-$\R$ ל-$\R$}, כי היא אינה על $\R$.

  נמצא את התמונה. עבור $x<1$: $x^3$ עולה ממש, ולכן $x^3<1^3=1$; כלומר ערכי החלק הראשון קטנים מ-1 (ולמעשה מכסים את $(-\infty,1)$). עבור $x\ge1$: $2x+1\ge3$; כלומר ערכי החלק השני גדולים או שווים ל-3 (ומכסים את $[3,\infty)$). לכן
  \[
  \operatorname{Im}f=(-\infty,1)\cup[3,\infty).
  \]
  בפרט, לא קיים $x\in\R$ עם $f(x)=2$: אם $x<1$ אז $f(x)<1$, ואם $x\ge1$ אז $f(x)\ge3$. לכן $f$ אינה על $\R$, ולכן אין לה הופכית $f^{-1}:\R\to\R$.

  (הערה: $f$ כן חח"ע — היא עולה ממש, שכן כל חלק עולה ממש וכל ערך של החלק הראשון קטן מכל ערך של החלק השני. לכן $f$ הפיכה כפונקציה מ-$\R$ אל התמונה שלה $(-\infty,1)\cup[3,\infty)$, עם $f^{-1}(y)=\sqrt[3]{y}$ עבור $y<1$ ו-$f^{-1}(y)=\frac{y-1}{2}$ עבור $y\ge3$; אבל לא כפונקציה מ-$\R$ ל-$\R$.)
\end{enumerate}
\end{solution}

\subsection*{חלק א', שאלה 2}

\begin{question}
יהיו $a,b\in\R$. נגדיר פונקציה על ידי
\[
f(x)=\begin{cases}
\frac{x^2-4x+3}{x-1} & x<1\\
a & x=1\\
2b\cdot\frac{\ln x}{x-1} & x>1
\end{cases}
\]
לאילו ערכים של $a,b$ הפונקציה רציפה ב-$1$ ? נמקו.
\end{question}

\begin{hint}
בגבול משמאל פרקו את המונה לגורמים. בגבול מימין הציבו $y=x-1$ והשתמשו בגבול היסודי $\frac{\ln(1+y)}{y}\to1$.
\end{hint}

\begin{solution}
הפונקציה רציפה ב-1 אם ורק אם הגבולות החד-צדדיים ב-1 קיימים ושווים ל-$f(1)=a$.

\textbf{הגבול משמאל.} $x^2-4x+3=(x-1)(x-3)$, ועבור $x\neq1$ מותר לצמצם:
\[
\lim_{x\to1^-}f(x)=\lim_{x\to1^-}\frac{(x-1)(x-3)}{x-1}=\lim_{x\to1^-}(x-3)=1-3=-2 ,
\]
כאשר בשוויון האחרון השתמשנו ברציפות הפולינום $x-3$.

\textbf{הגבול מימין.} לפי אריתמטיקה של גבולות (הקבוע $2b$ יוצא מהגבול), ועם ההצבה $y=x-1$ ($y\to0^+$ כאשר $x\to1^+$):
\[
\lim_{x\to1^+}f(x)=2b\cdot\lim_{x\to1^+}\frac{\ln x}{x-1}=2b\cdot\lim_{y\to0^+}\frac{\ln(1+y)}{y}=2b\cdot1=2b ,
\]
לפי הגבול היסודי $\lim_{y\to0}\frac{\ln(1+y)}{y}=1$ (מדף הנוסחאות: $\lim_{y\to0}\frac{\log_a(1+y)}{y}=\frac{1}{\ln a}$ עם $a=e$).

\textbf{מסקנה.} $f$ רציפה ב-1 אם ורק אם $-2=a=2b$, כלומר
\[
\boxed{a=-2,\quad b=-1.}
\]
\end{solution}

\subsection*{חלק א', שאלה 3}

\begin{question}
הוכיחו שלמשוואה
\[
x^5+e^x=0
\]
יש פתרון ממשי, ומצאו קטע שאורכו לכל היותר 1 המכיל את הפתרון. נמקו.
\end{question}

\begin{hint}
הגדירו $f(x)=x^5+e^x$ ובדקו את הסימן שלה בנקודות שלמות קרובות ל-0.
\end{hint}

\begin{solution}
נגדיר $f(x)=x^5+e^x$. $f$ רציפה על $\R$ כסכום של פולינום ושל פונקציה מעריכית, שתיהן רציפות.

נחשב:
\[
f(0)=0^5+e^0=1>0,\qquad f(-1)=(-1)^5+e^{-1}=-1+\frac1e<0 ,
\]
כי $e>1$ ולכן $\frac1e<1$.

$f$ רציפה בקטע הסגור $[-1,0]$ ומחליפה בו סימן בקצוות. לפי משפט ערך הביניים (משפט בולצאנו) קיים $c\in(-1,0)$ כך ש-$f(c)=0$, כלומר $c^5+e^c=0$, ו-$c$ הוא פתרון ממשי של המשוואה.

\textbf{הקטע המבוקש:} $[-1,0]$, שאורכו 1.

(הערה: הפתרון יחיד — $f$ עולה ממש כסכום של שתי פונקציות עולות ממש $x^5$ ו-$e^x$ — ולכן "הפתרון" מוגדר היטב. מספרית $c\approx-0.845$.)
\end{solution}

\subsection*{חלק ב', שאלה 4}

\begin{question}
יהיו
\[
f(x)=\sqrt[3]{x},\qquad g(x)=\cos(x)-\sin(x)
\]
אילו מהטענות הבאות נכונה?
\begin{enumerate}
  \item הפונקציה $f\circ g$ זוגית.
  \item הפונקציה $f\circ g$ אי זוגית.
  \item הפונקציה $f\circ g$ מונוטונית.
  \item הפונקציה $f\circ g$ מחזורית.
\end{enumerate}
\end{question}

\begin{solution}
\textbf{התשובה הנכונה: (ד)}.

$(f\circ g)(x)=f(g(x))=\sqrt[3]{\cos x-\sin x}$, מוגדרת לכל $x\in\R$ (שורש שלישי מוגדר לכל מספר ממשי).

\textbf{(ד) נכונה:} $\sin$ ו-$\cos$ מחזוריות עם מחזור $2\pi$, ולכן לכל $x$:
\[
(f\circ g)(x+2\pi)=\sqrt[3]{\cos(x+2\pi)-\sin(x+2\pi)}=\sqrt[3]{\cos x-\sin x}=(f\circ g)(x).
\]

\textbf{למה האחרות שגויות:} נסמן $h=f\circ g$.
\begin{itemize}
  \item (א) — $h(-x)=\sqrt[3]{\cos x+\sin x}$. למשל $h\!\left(\frac\pi4\right)=\sqrt[3]{0}=0$ אבל $h\!\left(-\frac\pi4\right)=\sqrt[3]{\sqrt2}\neq0$, ולכן $h$ אינה זוגית.
  \item (ב) — פונקציה אי-זוגית המוגדרת ב-0 מקיימת $h(0)=-h(0)$, כלומר $h(0)=0$; אבל $h(0)=\sqrt[3]{1}=1\neq0$.
  \item (ג) — פונקציה מחזורית שאינה קבועה אינה מונוטונית: למשל $h(0)=1$, $h\!\left(\frac\pi4\right)=0$, $h(2\pi)=1$, כלומר $h(0)>h(\frac\pi4)<h(2\pi)$ — הפונקציה יורדת ואחר כך עולה.
\end{itemize}
\end{solution}

\subsection*{חלק ב', שאלה 5}

\begin{question}
חשבו את הגבול
\[
\lim_{x\to\infty}\left(\frac{x^2+2x}{3x^2+7x}\right)^{x^2}
\]
\begin{enumerate}
  \item הגבול קיים במובן הרחב ושווה $\infty$.
  \item הגבול קיים ושווה $e$.
  \item הגבול קיים ושווה $0$.
  \item הגבול קיים ושווה $1$.
\end{enumerate}
\end{question}

\begin{hint}
לאן שואף הבסיס? זה \emph{לא} גבול מהצורה $1^\infty$.
\end{hint}

\begin{solution}
\textbf{התשובה הנכונה: (ג)}.

הבסיס שואף ל-$\frac13$ (חלוקה ב-$x^2$), ולא ל-1, ולכן זה אינו גבול מהצורה $1^\infty$. נשתמש בכלל הסנדוויץ'. עבור $x>0$:
\[
0\le\frac{x^2+2x}{3x^2+7x}\le\frac{x^2+2x}{3x^2+6x}=\frac{x(x+2)}{3x(x+2)}=\frac13 ,
\]
כי המונה חיובי והקטנת המכנה החיובי ($7x>6x$) מגדילה את השבר. מכיוון שהפונקציה $t\mapsto t^{x^2}$ עולה על $[0,\infty)$ (עבור $x^2>0$):
\[
0\le\left(\frac{x^2+2x}{3x^2+7x}\right)^{x^2}\le\left(\frac13\right)^{x^2}.
\]
כאשר $x\to\infty$ גם $x^2\to\infty$, ומכיוון ש-$0<\frac13<1$ מתקיים $\left(\frac13\right)^{x^2}\to0$. לפי כלל הסנדוויץ'
\[
\lim_{x\to\infty}\left(\frac{x^2+2x}{3x^2+7x}\right)^{x^2}=\boxed{0}.
\]

\textbf{למה האחרות שגויות:} (א) — הביטוי חסום ע"י $\left(\frac13\right)^{x^2}\le1$, ולכן אינו שואף ל-$\infty$. (ב) ו-(ד) — ערכים אלה מתקבלים מטיפול שגוי בביטוי כאילו היה מהצורה $1^\infty$; כאן הבסיס שואף ל-$\frac13<1$ והגבול הוא 0.
\end{solution}

\subsection*{חלק ב', שאלה 6}

\begin{question}
תהי $f(x)=x+e^x$. אילו מהטענות הבאות נכונה?
\begin{enumerate}
  \item הפונקציה חח"ע ועל $\R$.
  \item הפונקציה חח"ע אך אינה על $\R$.
  \item הפונקציה על $\R$ אך אינה חח"ע.
  \item הפונקציה אינה חח"ע ואינה על $\R$.
\end{enumerate}
\end{question}

\begin{solution}
\textbf{התשובה הנכונה: (א)}.

\textbf{חח"ע:} יהיו $x_1<x_2$. מכיוון ש-$e^x$ עולה ממש, $e^{x_1}<e^{x_2}$. נחבר את שני אי-השוויונות ונקבל $x_1+e^{x_1}<x_2+e^{x_2}$, כלומר $f$ עולה ממש ובפרט חח"ע.

\textbf{על $\R$:} $f$ רציפה על $\R$ כסכום של פונקציות רציפות, ו-
\[
\lim_{x\to-\infty}(x+e^x)=-\infty+0=-\infty,\qquad \lim_{x\to\infty}(x+e^x)=\infty+\infty=\infty .
\]
יהי $y\in\R$. לפי הגבולות האלה קיימים $x_1<x_2$ עם $f(x_1)<y<f(x_2)$, ולפי משפט ערך הביניים על $[x_1,x_2]$ קיים $c$ עם $f(c)=y$. לכן $f$ על $\R$.

\textbf{למה האחרות שגויות:} (ב) ו-(ד) טוענות ש-$f$ אינה על, (ג) ו-(ד) טוענות שאינה חח"ע — שתי הטענות הופרכו לעיל.
\end{solution}

\subsection*{חלק ב', שאלה 7}

\begin{question}
בדקו את רציפות הפונקציה הבאה בנקודה 3:
\[
f(x)=\begin{cases}
\frac{\sqrt{x+7}-\sqrt{10}}{\sqrt{x+4}-\sqrt{7}} & x>3\\
\sqrt{x^2+1} & x\le 3
\end{cases}
\]
אילו מהטענות הבאות נכונה?
\begin{enumerate}
  \item הפונקציה רציפה ב-3.
  \item הפונקציה בעלת אי-רציפות סליקה ב-3.
  \item הפונקציה בעלת אי-רציפות מסוג קפיצה ב-3.
  \item הפונקציה בעלת אי-רציפות עיקרית ב-3.
\end{enumerate}
\end{question}

\begin{hint}
בגבול מימין כפלו מונה ומכנה בצמודים של שני הביטויים: $\sqrt{x+7}+\sqrt{10}$ ו-$\sqrt{x+4}+\sqrt7$.
\end{hint}

\begin{solution}
\textbf{התשובה הנכונה: (ג)}.

\textbf{משמאל:} $\sqrt{x^2+1}$ רציפה (הרכבה של פולינום ושורש), ולכן
\[
\lim_{x\to3^-}f(x)=\sqrt{3^2+1}=\sqrt{10}=f(3).
\]
\textbf{מימין:} גבול מהצורה $\frac00$. נכפול בצמודים:
\[
\frac{\sqrt{x+7}-\sqrt{10}}{\sqrt{x+4}-\sqrt7}\cdot\frac{\sqrt{x+4}+\sqrt7}{\sqrt{x+4}+\sqrt7}\cdot\frac{\sqrt{x+7}+\sqrt{10}}{\sqrt{x+7}+\sqrt{10}}
=\frac{\bigl((x+7)-10\bigr)\left(\sqrt{x+4}+\sqrt7\right)}{\bigl((x+4)-7\bigr)\left(\sqrt{x+7}+\sqrt{10}\right)}
=\frac{\sqrt{x+4}+\sqrt7}{\sqrt{x+7}+\sqrt{10}},
\]
כאשר צמצמנו את $x-3\neq0$. הביטוי האחרון רציף ב-3, ולכן
\[
\lim_{x\to3^+}f(x)=\frac{\sqrt7+\sqrt7}{\sqrt{10}+\sqrt{10}}=\frac{2\sqrt7}{2\sqrt{10}}=\frac{\sqrt7}{\sqrt{10}} .
\]
שני הגבולות החד-צדדיים קיימים וסופיים, אך $\frac{\sqrt7}{\sqrt{10}}\neq\sqrt{10}$. לכן הגבול ב-3 אינו קיים, ויש ב-3 אי-רציפות מסוג קפיצה.

\textbf{למה האחרות שגויות:} (א) — הגבול ב-3 אינו קיים, ולכן אין רציפות (יש רק רציפות משמאל). (ב) — אי-רציפות סליקה דורשת קיום הגבול הדו-צדדי. (ד) — אי-רציפות עיקרית דורשת שלפחות אחד הגבולות החד-צדדיים לא יהיה קיים כגבול סופי, וכאן שניהם קיימים.
\end{solution}

\section*{תשפ"ו סמסטר א מועד א}

\subsection*{שאלה 1א}

\begin{question}
כמה פתרונות ממשיים יש למשוואה הבאה?
\[ x^4 + 2x^2 = 1 \]
נמקו היטב.
\end{question}

\begin{hint}
הציבו $t = x^2$ וקבלו משוואה ריבועית ב-$t$; זכרו ש-$t = x^2$ חייב להיות אי-שלילי.
\end{hint}

\begin{solution}
\textbf{דרך 1 (הצבה).} נעביר אגפים: $x^4 + 2x^2 - 1 = 0$. נציב $t = x^2 \ge 0$ ונקבל
\[ t^2 + 2t - 1 = 0 \quad\Longrightarrow\quad t_{1,2} = \frac{-2 \pm \sqrt{4 + 4}}{2} = -1 \pm \sqrt{2}. \]
הפתרון $t = -1 - \sqrt{2} < 0$ נפסל, כי $x^2 \ge 0$ לכל $x$ ממשי, ולכן אין לו פתרון ממשי.
הפתרון $t = -1 + \sqrt{2} > 0$ (כי $\sqrt 2 > 1$) נותן
\[ x^2 = \sqrt{2} - 1 \quad\Longrightarrow\quad x_{1,2} = \pm\sqrt{\sqrt{2} - 1}, \]
שני מספרים ממשיים שונים.

\textbf{דרך 2 (חקירה).} נגדיר $f(x) = x^4 + 2x^2 - 1$, פונקציה רציפה וגזירה בכל $\R$. מתקיים
\[ f'(x) = 4x^3 + 4x = 4x(x^2 + 1), \]
ומכיוון ש-$x^2 + 1 > 0$, הסימן של $f'$ הוא הסימן של $x$: $f' < 0$ ב-$(-\infty, 0)$ ו-$f' > 0$ ב-$(0, \infty)$.
לכן $f$ יורדת ממש ב-$(-\infty, 0]$ ועולה ממש ב-$[0, \infty)$, ובכל אחד מהקטעים האלה יש לה לכל היותר אפס אחד.
כעת $f(0) = -1 < 0$ ו-$f(1) = f(-1) = 2 > 0$. לפי משפט ערך הביניים יש ל-$f$ אפס ב-$(-1, 0)$ ואפס ב-$(0, 1)$, ולפי המונוטוניות אלה האפסים היחידים.

\textbf{תשובה:} למשוואה יש בדיוק \textbf{שני} פתרונות ממשיים, $x = \pm\sqrt{\sqrt 2 - 1}$.
\end{solution}

\subsection*{שאלה 1ב}

\begin{question}
חשבו את הגבול הבא
\[ \lim_{x\to 0} \frac{x - \sin x}{\cos x - 1} \]
\end{question}

\begin{hint}
זהו את סוג הביטוי ($\frac{0}{0}$) והשתמשו בכלל לופיטל.
\end{hint}

\begin{solution}
כאשר $x \to 0$ המונה $x - \sin x \to 0$ והמכנה $\cos x - 1 \to 0$, כלומר ביטוי מהצורה $\frac{0}{0}$.
המונה והמכנה גזירים בסביבת $0$, ונגזרת המכנה $-\sin x \ne 0$ בסביבה מנוקבת של $0$. לפי כלל לופיטל (בתנאי שהגבול החדש קיים):
\[ \lim_{x\to 0} \frac{x - \sin x}{\cos x - 1} = \lim_{x\to 0} \frac{1 - \cos x}{-\sin x}. \]
גם זה ביטוי $\frac{0}{0}$. נחשב אותו בעזרת $1 - \cos x = 2\sin^2\frac{x}{2}$:
\[ \frac{1 - \cos x}{-\sin x} = \frac{2\sin^2\frac{x}{2}}{-\sin x}
   = -\frac{2\left(\dfrac{\sin\frac{x}{2}}{\frac{x}{2}}\right)^2 \cdot \dfrac{x^2}{4}}{\dfrac{\sin x}{x} \cdot x}
   = -\frac{\left(\dfrac{\sin\frac{x}{2}}{\frac{x}{2}}\right)^2}{\dfrac{\sin x}{x}} \cdot \frac{x}{2}
   \xrightarrow[x\to 0]{} -\frac{1^2}{1}\cdot 0 = 0, \]
כאשר השתמשנו בגבול היסודי $\lim_{u\to 0}\frac{\sin u}{u} = 1$ ובאריתמטיקה של גבולות.
(לחלופין, הפעלה נוספת של לופיטל: $\lim_{x\to0}\frac{\sin x}{-\cos x} = \frac{0}{-1} = 0$.)
מכיוון שהגבול אחרי לופיטל קיים, גם הגבול המקורי קיים ושווה לו:
\[ \boxed{\lim_{x\to 0} \frac{x - \sin x}{\cos x - 1} = 0}. \]
\end{solution}

\subsection*{שאלה 2א}

\begin{question}
חשבו את כל האסימפטוטות של הפונקציה
\[ f(x) = \frac{x^2 - 2\ln(x)}{x - 1} \]
\end{question}

\begin{hint}
התחילו מתחום ההגדרה: הנקודות ה"חשודות" לאסימפטוטה אנכית הן קצוות התחום ונקודות שבהן המכנה מתאפס. אסימפטוטה משופעת $y = ax + b$ יש לבדוק רק בכיוון שבו התחום אינסופי.
\end{hint}

\begin{solution}
\textbf{תחום הגדרה:} $\ln x$ מוגדר עבור $x > 0$, והמכנה מתאפס ב-$x = 1$. לכן התחום הוא $(0,1)\cup(1,\infty)$, ו-$f$ רציפה בו.

\textbf{אסימפטוטות אנכיות.} הנקודות החשודות הן $x = 0$ (קצה התחום) ו-$x = 1$.
\begin{itemize}
  \item $x \to 0^+$: המונה $x^2 - 2\ln x \to 0 - (-\infty) = +\infty$ והמכנה $x - 1 \to -1$, ולכן
  \[ \lim_{x\to 0^+} f(x) = -\infty. \]
  לכן $x = 0$ אסימפטוטה אנכית (מימין).
  \item $x \to 1$: המונה $\to 1 - 2\ln 1 = 1 > 0$ והמכנה $\to 0$, כאשר $x - 1 < 0$ משמאל ו-$x - 1 > 0$ מימין. לכן
  \[ \lim_{x\to 1^-} f(x) = -\infty, \qquad \lim_{x\to 1^+} f(x) = +\infty, \]
  ו-$x = 1$ אסימפטוטה אנכית (משני הצדדים).
\end{itemize}

\textbf{אסימפטוטות אופקיות/משופעות.} התחום אינסופי רק בכיוון $+\infty$, לכן בודקים רק $x \to +\infty$. נחפש $y = ax + b$:
\[ a = \lim_{x\to\infty} \frac{f(x)}{x} = \lim_{x\to\infty} \frac{x^2 - 2\ln x}{x^2 - x}
     = \lim_{x\to\infty} \frac{x^2\left(1 - \frac{2\ln x}{x^2}\right)}{x^2\left(1 - \frac{1}{x}\right)} = \frac{1 - 0}{1 - 0} = 1, \]
כי $\frac{\ln x}{x^2} \to 0$ (למשל לפי לופיטל: $\lim \frac{1/x}{2x} = 0$).
\[ b = \lim_{x\to\infty} \bigl(f(x) - x\bigr) = \lim_{x\to\infty} \frac{x^2 - 2\ln x - x^2 + x}{x - 1}
     = \lim_{x\to\infty} \frac{x\left(1 - \frac{2\ln x}{x}\right)}{x\left(1 - \frac{1}{x}\right)} = 1, \]
כי $\frac{\ln x}{x} \to 0$. מכיוון ששני הגבולות סופיים, הישר $y = x + 1$ הוא אסימפטוטה משופעת ב-$+\infty$ (ואין אסימפטוטה אופקית).

\textbf{תשובה:} אסימפטוטות אנכיות $x = 0$ ו-$x = 1$; אסימפטוטה משופעת $y = x + 1$ כאשר $x \to +\infty$.
\end{solution}

\subsection*{שאלה 2ב}

\begin{question}
מצאו מינימום ומקסימום מוחלטים של הפונקציה
\[ f(x) = \frac{\sqrt{x}}{e^x} \]
בתחום הגדרתה, או הסבירו מדוע אינם קיימים.
\end{question}

\begin{hint}
מצאו את הנקודות הקריטיות, קבעו תחומי עלייה וירידה, ובדקו את התנהגות $f$ בקצוות התחום ($x = 0$ ו-$x \to \infty$).
\end{hint}

\begin{solution}
\textbf{תחום הגדרה:} $x \ge 0$ (בגלל $\sqrt{x}$). $f$ רציפה ב-$[0,\infty)$, $f(x) \ge 0$ לכל $x$, ו-$f(0) = 0$.

\textbf{נגזרת} (עבור $x > 0$):
\[ f'(x) = \frac{\frac{1}{2\sqrt{x}}\, e^x - \sqrt{x}\, e^x}{e^{2x}} = \frac{e^x(1 - 2x)}{2\sqrt{x}\, e^{2x}} = \frac{1 - 2x}{2\sqrt{x}\, e^x}. \]
המכנה חיובי, ולכן $f'(x) = 0 \iff x = \frac12$, $f' > 0$ ב-$(0, \frac12)$ ו-$f' < 0$ ב-$(\frac12, \infty)$.
כלומר $f$ עולה ממש ב-$[0, \frac12]$ ויורדת ממש ב-$[\frac12, \infty)$ (כולל הקצה, כי $f$ רציפה שם).

\textbf{מקסימום מוחלט:} מהמונוטוניות, $f(x) \le f(\frac12)$ לכל $x \ge 0$, ולכן
\[ \max f = f\left(\tfrac12\right) = \frac{\sqrt{1/2}}{e^{1/2}} = \frac{1}{\sqrt{2e}}, \]
המתקבל ב-$x = \frac12$.

\textbf{מינימום מוחלט:} $f(x) \ge 0 = f(0)$ לכל $x$ בתחום (מונה אי-שלילי ומכנה חיובי), ולכן
\[ \min f = f(0) = 0, \]
המתקבל ב-$x = 0$. (שימו לב: אמנם $\lim_{x\to\infty} f(x) = 0$, אבל הערך $0$ מתקבל בפועל ב-$x = 0$, ולכן המינימום קיים.)

\textbf{תשובה:} מינימום מוחלט $0$ ב-$x = 0$; מקסימום מוחלט $\frac{1}{\sqrt{2e}}$ ב-$x = \frac12$.
\end{solution}

\subsection*{שאלה 3}

\begin{question}
חשבו בעזרת קירוב טיילור מסדר שני את $\ln\left(\frac{6}{5}\right) = \ln(1.2)$, והעריכו את השגיאה.

שימו לב שהקירוב וחסם השגיאה צריכים להיות מספרים רציונליים.
\end{question}

\begin{hint}
פתחו את $f(x) = \ln x$ סביב $x_0 = 1$ (שם הערכים של $f$ ונגזרותיה רציונליים) והשתמשו בשארית לגרנז' מסדר שני.
\end{hint}

\begin{solution}
נגדיר $f(x) = \ln x$ ונפתח סביב $x_0 = 1$, נקודה קרובה ל-$1.2$ שבה כל הערכים ידועים. $f$ גזירה אינסוף פעמים ב-$(0,\infty)$:
\[ f(x) = \ln x,\quad f'(x) = \frac{1}{x},\quad f''(x) = -\frac{1}{x^2},\quad f'''(x) = \frac{2}{x^3}, \]
ולכן $f(1) = 0$, $f'(1) = 1$, $f''(1) = -1$.

\textbf{פולינום טיילור מסדר שני סביב $1$:}
\[ P_2(x) = f(1) + f'(1)(x - 1) + \frac{f''(1)}{2!}(x - 1)^2 = (x - 1) - \frac{(x - 1)^2}{2}. \]
לכן
\[ \ln(1.2) \approx P_2(1.2) = 0.2 - \frac{0.04}{2} = 0.18 = \frac{9}{50}. \]

\textbf{הערכת השגיאה.} לפי משפט טיילור עם שארית לגרנז', קיימת נקודה $c$ בין $1$ ל-$1.2$ כך ש-
\[ R_2(1.2) = \ln(1.2) - P_2(1.2) = \frac{f'''(c)}{3!}(1.2 - 1)^3 = \frac{2}{6c^3}\cdot (0.2)^3 = \frac{0.008}{3c^3}. \]
מכיוון ש-$1 < c < 1.2$, מתקיים $c^3 > 1$, ולכן $\frac{1}{c^3} < 1$ ו-
\[ 0 < R_2(1.2) < \frac{0.008}{3} = \frac{8}{3000} = \frac{1}{375}. \]

\textbf{תשובה:} $\ln(1.2) \approx 0.18 = \frac{9}{50}$, והשגיאה קטנה מ-$\frac{1}{375}$ (בפרט, מכיוון שהשארית חיובית, $0.18 < \ln(1.2) < 0.18 + \frac{1}{375}$).

(בדיקה: $\ln 1.2 \approx 0.18232$, והשגיאה בפועל היא כ-$0.0023 < \frac{1}{375} \approx 0.00267$.)
\end{solution}

\subsection*{שאלה 4א}

\begin{question}
חשבו את האינטגרל
\[ \int \frac{-3x - 7}{x^3 + x^2 - 2}\,dx \]
רמז: שימו לב שהמכנה מתאפס ב-$x = 1$.
\end{question}

\begin{hint}
חלקו את המכנה ב-$(x - 1)$, בדקו שהגורם הריבועי שנותר אי-פריק, ופרקו לשברים חלקיים מהצורה $\frac{A}{x-1} + \frac{Bx + C}{x^2 + 2x + 2}$.
\end{hint}

\begin{solution}
\textbf{פירוק המכנה.} $1^3 + 1^2 - 2 = 0$, ולכן $(x - 1)$ מחלק את המכנה. חילוק פולינומים נותן
\[ x^3 + x^2 - 2 = (x - 1)(x^2 + 2x + 2). \]
הדיסקרימיננטה של $x^2 + 2x + 2$ היא $4 - 8 < 0$, לכן הגורם הריבועי אי-פריק מעל $\R$ (ובפרט $x^2 + 2x + 2 = (x+1)^2 + 1 > 0$).

\textbf{שברים חלקיים.} דרגת המונה קטנה מדרגת המכנה, לכן נחפש
\[ \frac{-3x - 7}{(x - 1)(x^2 + 2x + 2)} = \frac{A}{x - 1} + \frac{Bx + C}{x^2 + 2x + 2}, \]
כלומר $-3x - 7 = A(x^2 + 2x + 2) + (Bx + C)(x - 1)$ לכל $x$.
\begin{itemize}
  \item $x = 1$: $-10 = 5A$, ולכן $A = -2$.
  \item $x = 0$: $-7 = 2A - C$, ולכן $C = 2A + 7 = 3$.
  \item השוואת מקדמי $x^2$: $0 = A + B$, ולכן $B = 2$.
\end{itemize}
לכן
\[ \int \frac{-3x - 7}{x^3 + x^2 - 2}\,dx = \int \left( \frac{-2}{x - 1} + \frac{2x + 3}{x^2 + 2x + 2} \right) dx
   = \int \left( \frac{-2}{x - 1} + \frac{2x + 2}{x^2 + 2x + 2} + \frac{1}{(x + 1)^2 + 1} \right) dx. \]
\begin{itemize}
  \item $\int \frac{-2}{x - 1}\,dx = -2\ln|x - 1| + C$.
  \item במחובר השני המונה הוא נגזרת המכנה: בהצבה $t = x^2 + 2x + 2$, $dt = (2x + 2)\,dx$,
  \[ \int \frac{2x + 2}{x^2 + 2x + 2}\,dx = \int \frac{dt}{t} = \ln|t| + C = \ln(x^2 + 2x + 2) + C. \]
  \item בהשלמה לריבוע ובהצבה $u = x + 1$: $\int \frac{dx}{(x + 1)^2 + 1} = \arctan(x + 1) + C$.
\end{itemize}
\textbf{תשובה:}
\[ \boxed{\int \frac{-3x - 7}{x^3 + x^2 - 2}\,dx = -2\ln|x - 1| + \ln(x^2 + 2x + 2) + \arctan(x + 1) + C.} \]
(בדיקה: גזירת התוצאה מחזירה את האינטגרנד.)
\end{solution}

\subsection*{שאלה 4ב}

\begin{question}
מצאו את הפתרון הכללי של המשוואה הדיפרנציאלית
\[ y' = y^2 x \sin x \]
\end{question}

\begin{hint}
זו משוואה פרידה: העבירו את $y^2$ לאגף של $dy$, ואת $\int x\sin x\,dx$ חשבו באינטגרציה בחלקים. אל תשכחו את הפתרון $y \equiv 0$.
\end{hint}

\begin{solution}
זו משוואה פרידה: $\frac{dy}{dx} = y^2 \cdot x\sin x$.

\textbf{הפתרון הקבוע:} $y \equiv 0$ מקיים $y' = 0 = 0^2 \cdot x\sin x$, ולכן הוא פתרון.

\textbf{עבור $y \ne 0$:} נפריד משתנים ונבצע אינטגרציה:
\[ \int \frac{dy}{y^2} = \int x\sin x\,dx. \]
באגף שמאל $\int y^{-2}\,dy = -\frac{1}{y}$. באגף ימין אינטגרציה בחלקים עם $u = x$, $dv = \sin x\,dx$ (כלומר $du = dx$, $v = -\cos x$):
\[ \int x\sin x\,dx = -x\cos x + \int \cos x\,dx = -x\cos x + \sin x + C. \]
לכן
\[ -\frac{1}{y} = -x\cos x + \sin x + C \quad\Longrightarrow\quad \boxed{y = \frac{1}{x\cos x - \sin x - C}}, \]
כאשר $C$ קבוע ממשי כלשהו (בכל קטע שבו המכנה אינו מתאפס).

\textbf{תשובה:} הפתרון הכללי הוא $y = \dfrac{1}{x\cos x - \sin x + C}$ ($C \in \R$, שינינו את שם הקבוע), ובנוסף הפתרון $y \equiv 0$.
\end{solution}

\subsection*{שאלה 5א}

\begin{question}
חשבו את נפח גוף הסיבוב סביב ציר ה-$x$ של הפונקציה
\[ f(x) = \sin x + \cos x \]
בקטע $[0, \frac{\pi}{2}]$.
\end{question}

\begin{hint}
פתחו את $(\sin x + \cos x)^2$ והשתמשו בזהויות $\sin^2 x + \cos^2 x = 1$ ו-$2\sin x\cos x = \sin 2x$.
\end{hint}

\begin{solution}
נפח גוף הסיבוב של גרף $f$ סביב ציר ה-$x$ בקטע $[a,b]$ הוא $V = \pi\int_a^b f^2(x)\,dx$. כאן
\[ f^2(x) = \sin^2 x + 2\sin x\cos x + \cos^2 x = 1 + \sin 2x, \]
ולכן, לפי המשפט היסודי של החשבון האינטגרלי (ניוטון-לייבניץ),
\[ V = \pi\int_0^{\pi/2} (1 + \sin 2x)\,dx = \pi\left[ x - \frac{\cos 2x}{2} \right]_0^{\pi/2}
     = \pi\left[ \left(\frac{\pi}{2} - \frac{\cos\pi}{2}\right) - \left(0 - \frac{\cos 0}{2}\right) \right]
     = \pi\left[ \frac{\pi}{2} + \frac12 + \frac12 \right]. \]
\textbf{תשובה:} $\boxed{V = \pi\left(\frac{\pi}{2} + 1\right) = \frac{\pi^2}{2} + \pi}$.
\end{solution}

\subsection*{שאלה 5ב}

\begin{question}
חשבו את האינטגרל המוכלל הבא, או הראו שהוא מתבדר:
\[ \int_e^\infty \frac{dx}{x \cdot \ln^3 x} \]
\end{question}

\begin{hint}
הציבו $t = \ln x$.
\end{hint}

\begin{solution}
האינטגרנד $\frac{1}{x\ln^3 x}$ רציף וחיובי ב-$[e, \infty)$, ולכן האינטגרל מוכלל רק בגלל הגבול העליון האינסופי. לפי ההגדרה
\[ \int_e^\infty \frac{dx}{x\ln^3 x} = \lim_{A\to\infty} \int_e^A \frac{dx}{x\ln^3 x}. \]
נמצא פונקציה קדומה בהצבה $t = \ln x$, $dt = \frac{dx}{x}$:
\[ \int \frac{dx}{x\ln^3 x} = \int \frac{dt}{t^3} = -\frac{1}{2t^2} + C = -\frac{1}{2\ln^2 x} + C. \]
לכן, עבור $A > e$:
\[ \int_e^A \frac{dx}{x\ln^3 x} = \left[ -\frac{1}{2\ln^2 x} \right]_e^A = -\frac{1}{2\ln^2 A} + \frac{1}{2\ln^2 e} = \frac12 - \frac{1}{2\ln^2 A}. \]
כאשר $A \to \infty$ מתקיים $\ln A \to \infty$, ולכן $\frac{1}{2\ln^2 A} \to 0$.

\textbf{תשובה:} האינטגרל מתכנס ו-$\boxed{\int_e^\infty \frac{dx}{x\ln^3 x} = \frac12}$.
\end{solution}

\section*{תשפ"ו סמסטר א מועד ב}

\subsection*{שאלה 1א}

\begin{question}
חשבו את התמונה של הפונקציה
\[ f(x) = x^3 - 12x + 8 \]
על הקטע $[-3, 5]$. כלומר, חשבו את הקבוצה
\[ \{ f(x) : -3 \le x \le 5 \} \]
נמקו היטב.
\end{question}

\begin{hint}
פונקציה רציפה על קטע סגור מקבלת מינימום ומקסימום (ויירשטראס) וכל ערך ביניהם (ערך הביניים). נשאר למצוא את המינימום והמקסימום.
\end{hint}

\begin{solution}
$f$ פולינום, ולכן רציפה בקטע הסגור $[-3, 5]$. לפי משפט ויירשטראס השני $f$ מקבלת בקטע מינימום $c$ ומקסימום $d$, ולפי משפט ערך הביניים היא מקבלת כל ערך בין $c$ ל-$d$. לכן התמונה היא בדיוק הקטע הסגור $[c, d]$, ונותר לחשב את $c$ ו-$d$.

נקודות קיצון מוחלט של פונקציה גזירה בקטע סגור נמצאות בנקודות קריטיות פנימיות או בקצוות.
\[ f'(x) = 3x^2 - 12 = 3(x^2 - 4) = 0 \iff x = \pm 2, \]
ושתי הנקודות נמצאות בתוך הקטע. נחשב את ערכי $f$ בנקודות החשודות:
\[ f(-3) = -27 + 36 + 8 = 17,\quad f(-2) = -8 + 24 + 8 = 24,\quad f(2) = 8 - 24 + 8 = -8,\quad f(5) = 125 - 60 + 8 = 73. \]
לכן המינימום הוא $c = -8$ (ב-$x = 2$) והמקסימום הוא $d = 73$ (ב-$x = 5$).

\textbf{תשובה:} $\{ f(x) : -3 \le x \le 5 \} = [-8, 73]$.
\end{solution}

\subsection*{שאלה 1ב}

\begin{question}
חשבו את הגבול הבא
\[ \lim_{x\to\infty} \left( \frac{x^3 + 1}{x^3 - 2} \right)^{x^3 + 3} \]
\end{question}

\begin{hint}
זה ביטוי מהצורה $1^\infty$. כתבו $\frac{x^3 + 1}{x^3 - 2} = 1 + \frac{3}{x^3 - 2}$ והשתמשו בגבול $\left(1 + \frac{1}{u}\right)^u \to e$.
\end{hint}

\begin{solution}
כאשר $x \to \infty$ הבסיס שואף ל-$1$ והמעריך ל-$\infty$ (ביטוי $1^\infty$). נכתוב
\[ \frac{x^3 + 1}{x^3 - 2} = \frac{(x^3 - 2) + 3}{x^3 - 2} = 1 + \frac{3}{x^3 - 2}, \qquad x^3 + 3 = (x^3 - 2) + 5, \]
ולכן (עבור $x$ גדול, כך ש-$x^3 - 2 > 0$)
\[ \left( \frac{x^3 + 1}{x^3 - 2} \right)^{x^3 + 3} = \left( 1 + \frac{3}{x^3 - 2} \right)^{x^3 - 2} \cdot \left( 1 + \frac{3}{x^3 - 2} \right)^{5}. \]
נסמן $u = \frac{x^3 - 2}{3} \to \infty$. אז הגורם הראשון שווה $\left[\left(1 + \frac{1}{u}\right)^{u}\right]^3 \to e^3$, לפי הגבול היסודי $\lim_{u\to\infty}\left(1 + \frac1u\right)^u = e$ ורציפות $t \mapsto t^3$.
הגורם השני שואף ל-$1^5 = 1$, כי $\frac{3}{x^3 - 2} \to 0$.
לפי אריתמטיקה של גבולות:
\[ \boxed{\lim_{x\to\infty} \left( \frac{x^3 + 1}{x^3 - 2} \right)^{x^3 + 3} = e^3 \cdot 1 = e^3}. \]
(דרך חלופית: $\ln$ של הביטוי הוא $(x^3 + 3)\ln\left(1 + \frac{3}{x^3 - 2}\right)$, ומכיוון ש-$\ln(1 + t) \sim t$ כאשר $t \to 0$, הוא שואף ל-$\lim \frac{3(x^3 + 3)}{x^3 - 2} = 3$.)
\end{solution}

\subsection*{שאלה 2א}

\begin{question}
חשבו תחומי קמירות ונקודות פיתול של הפונקציה
\[ f(x) = \sqrt[3]{6x + 12} \]
\end{question}

\begin{hint}
כתבו $f(x) = (6x + 12)^{1/3}$, גזרו פעמיים ובדקו את סימן $f''$. שימו לב לנקודה שבה $f''$ אינה מוגדרת.
\end{hint}

\begin{solution}
\textbf{תחום הגדרה:} שורש שלישי מוגדר לכל מספר ממשי, לכן $f$ מוגדרת ורציפה בכל $\R$. נכתוב $f(x) = (6x + 12)^{1/3}$. עבור $x \ne -2$ (כלומר $6x + 12 \ne 0$), לפי כלל השרשרת:
\[ f'(x) = \frac13 (6x + 12)^{-2/3} \cdot 6 = 2(6x + 12)^{-2/3}, \]
\[ f''(x) = 2\cdot\left(-\frac23\right)(6x + 12)^{-5/3} \cdot 6 = -8(6x + 12)^{-5/3} = \frac{-8}{(6x + 12)^{5/3}}. \]
ב-$x = -2$ הפונקציה אינה גזירה: $\frac{f(x) - f(-2)}{x + 2} = \frac{(6(x+2))^{1/3}}{x + 2} = \frac{6^{1/3}}{(x + 2)^{2/3}} \to +\infty$ (משיק אנכי).

\textbf{סימן $f''$:} החזקה $(6x + 12)^{5/3}$ היא בעלת אותו סימן כמו $6x + 12$ (חזקה אי-זוגית של שורש שלישי). לכן
\begin{itemize}
  \item $x < -2 \iff 6x + 12 < 0 \iff f''(x) > 0$, ולכן $f$ קמורה ב-$(-\infty, -2]$;
  \item $x > -2 \iff 6x + 12 > 0 \iff f''(x) < 0$, ולכן $f$ קעורה ב-$[-2, \infty)$.
\end{itemize}

\textbf{נקודת פיתול:} ב-$x = -2$ הפונקציה רציפה, יש לגרף משיק (אנכי) והקמירות מתחלפת משני צדי הנקודה. לכן $(-2, f(-2)) = (-2, 0)$ היא נקודת פיתול, אף על פי ש-$f$ אינה גזירה שם. אין נקודות פיתול נוספות, כי $f'' \ne 0$ בכל נקודה אחרת.

\textbf{תשובה:} $f$ קמורה ב-$(-\infty, -2)$, קעורה ב-$(-2, \infty)$, ונקודת הפיתול היא $(-2, 0)$.
\end{solution}

\subsection*{שאלה 2ב}

\begin{question}
הוכיחו כי לכל $x \ge 0$ מתקיים
\[ 0 \le \frac{x^2}{e^x} \le 1 \]
\end{question}

\begin{hint}
חקרו את $g(x) = \frac{x^2}{e^x}$ ב-$[0, \infty)$ ומצאו את המקסימום שלה.
\end{hint}

\begin{solution}
נגדיר $g(x) = \frac{x^2}{e^x} = x^2 e^{-x}$ ב-$[0, \infty)$; $g$ רציפה וגזירה שם.

\textbf{אי-השוויון השמאלי:} $x^2 \ge 0$ ו-$e^x > 0$, ולכן $g(x) \ge 0$ לכל $x$.

\textbf{אי-השוויון הימני:} נחשב
\[ g'(x) = \frac{2x e^x - x^2 e^x}{e^{2x}} = \frac{x(2 - x)}{e^x}. \]
ב-$(0, 2)$ מתקיים $g' > 0$ וב-$(2, \infty)$ מתקיים $g' < 0$. לכן $g$ עולה ב-$[0, 2]$ ויורדת ב-$[2, \infty)$, ו-$x = 2$ היא נקודת מקסימום מוחלט של $g$ ב-$[0, \infty)$. לכן לכל $x \ge 0$:
\[ g(x) \le g(2) = \frac{4}{e^2} = \left( \frac{2}{e} \right)^2 < 1, \]
כי $0 < 2 < e$ ולכן $0 < \frac{2}{e} < 1$.

מכאן, לכל $x \ge 0$: $0 \le \frac{x^2}{e^x} \le \frac{4}{e^2} < 1$, ובפרט $0 \le \frac{x^2}{e^x} \le 1$. $\blacksquare$
\end{solution}

\subsection*{שאלה 3}

\begin{question}
חשבו בעזרת קירוב טיילור את $\sqrt[3]{11}$, כך שהשגיאה תהיה קטנה מ-$\frac{1}{10}$. נמקו.

שימו לב שהקירוב וחסם השגיאה צריכים להיות מספרים רציונליים.
\end{question}

\begin{hint}
פתחו את $f(x) = \sqrt[3]{x}$ סביב $a = 8$ (כי $\sqrt[3]{8} = 2$). התחילו מפולינום מסדר ראשון ובדקו אם שארית לגרנז' כבר קטנה מ-$\frac{1}{10}$.
\end{hint}

\begin{solution}
נגדיר $f(x) = \sqrt[3]{x} = x^{1/3}$ ונפתח סביב $a = 8$, הנקודה הקרובה ל-$11$ שבה השורש השלישי ידוע ($\sqrt[3]{8} = 2$). נחפש את $n$ המינימלי כך ש-$|R_n(11)| < \frac{1}{10}$.

\textbf{נגזרות} (עבור $x > 0$):
\[ f'(x) = \frac13 x^{-2/3}, \qquad f''(x) = -\frac29 x^{-5/3}, \]
ולכן $f(8) = 2$ ו-$f'(8) = \frac{1}{3\cdot 8^{2/3}} = \frac{1}{3\cdot 4} = \frac{1}{12}$.

\textbf{שארית לגרנז' עבור $n = 1$:} לפי משפט טיילור, קיימת $c$ בין $8$ ל-$11$ כך ש-
\[ R_1(11) = \frac{f''(c)}{2!}(11 - 8)^2 = \frac{-\frac29 c^{-5/3}}{2}\cdot 9 = -\frac{1}{c^{5/3}}. \]
מכיוון ש-$8 < c < 11$, מתקיים $c^{5/3} > 8^{5/3} = 2^5 = 32$, ולכן
\[ |R_1(11)| = \frac{1}{c^{5/3}} < \frac{1}{32} < \frac{1}{10}. \]
מכאן שפולינום טיילור מסדר ראשון מספיק ($n = 1$).

\textbf{הקירוב:}
\[ P_1(x) = f(8) + f'(8)(x - 8) = 2 + \frac{1}{12}(x - 8), \]
\[ \sqrt[3]{11} = f(11) \approx P_1(11) = 2 + \frac{3}{12} = 2\frac14 = \frac94. \]

\textbf{תשובה:} $\sqrt[3]{11} \approx \frac{9}{4} = 2.25$, עם שגיאה קטנה מ-$\frac{1}{32} < \frac{1}{10}$. (מכיוון ש-$R_1(11) < 0$, הקירוב גדול מהערך האמיתי: $\frac94 - \frac{1}{32} < \sqrt[3]{11} < \frac94$.)

(בדיקה: $\sqrt[3]{11} \approx 2.22398$, והשגיאה בפועל כ-$0.026 < \frac{1}{32} \approx 0.031$.)
\end{solution}

\subsection*{שאלה 4א}

\begin{question}
חשבו את האינטגרל
\[ \int \frac{6x^3 - 12x + 24}{x^4 - 4x^3 + 8x^2}\,dx \]
\end{question}

\begin{hint}
פרקו את המכנה: $x^4 - 4x^3 + 8x^2 = x^2(x^2 - 4x + 8)$, והגורם הריבועי אי-פריק. הפירוק לשברים חלקיים הוא מהצורה $\frac{A}{x} + \frac{B}{x^2} + \frac{Cx + D}{x^2 - 4x + 8}$.
\end{hint}

\begin{solution}
זו פונקציה רציונלית, ונפעל לפי השלבים הרגילים.
\begin{itemize}
  \item \textbf{מעלות:} מעלת המונה ($3$) קטנה ממעלת המכנה ($4$), ולכן אין צורך בחילוק פולינומים.
  \item \textbf{פירוק המכנה:} $x^4 - 4x^3 + 8x^2 = x^2(x^2 - 4x + 8)$, והדיסקרימיננטה של $x^2 - 4x + 8$ היא $16 - 32 < 0$, לכן הגורם הריבועי אי-פריק; למעשה $x^2 - 4x + 8 = (x - 2)^2 + 4$.
  \item \textbf{שברים חלקיים:}
  \[ \frac{6x^3 - 12x + 24}{x^2(x^2 - 4x + 8)} = \frac{A}{x} + \frac{B}{x^2} + \frac{Cx + D}{x^2 - 4x + 8}, \]
  כלומר $6x^3 - 12x + 24 = Ax(x^2 - 4x + 8) + B(x^2 - 4x + 8) + (Cx + D)x^2$. השוואת מקדמים:
  \begin{align*}
    x^0 &: \ 8B = 24 \ \Rightarrow\ B = 3, \\
    x^1 &: \ 8A - 4B = -12 \ \Rightarrow\ 8A = 0 \ \Rightarrow\ A = 0, \\
    x^2 &: \ -4A + B + D = 0 \ \Rightarrow\ D = -3, \\
    x^3 &: \ A + C = 6 \ \Rightarrow\ C = 6.
  \end{align*}
  לכן
  \[ \frac{6x^3 - 12x + 24}{x^4 - 4x^3 + 8x^2} = \frac{3}{x^2} + \frac{6x - 3}{x^2 - 4x + 8}. \]
\end{itemize}
\textbf{אינטגרציה.} נכתוב $6x - 3 = 3(2x - 4) + 9$, כך שבחלק הראשון המונה הוא נגזרת המכנה, ובחלק השני נשלים לריבוע:
\[ \int \left( \frac{3}{x^2} + 3\cdot\frac{2x - 4}{x^2 - 4x + 8} + \frac{9}{(x - 2)^2 + 4} \right) dx. \]
\begin{itemize}
  \item $\int \frac{3}{x^2}\,dx = -\frac{3}{x} + C$.
  \item בהצבה $t = x^2 - 4x + 8$, $dt = (2x - 4)\,dx$: $\int \frac{2x - 4}{x^2 - 4x + 8}\,dx = \ln|t| = \ln(x^2 - 4x + 8) + C$ (המכנה חיובי תמיד).
  \item בהצבה $x - 2 = 2u$, $dx = 2\,du$: $\int \frac{9\,dx}{(x - 2)^2 + 4} = \int \frac{18\,du}{4u^2 + 4} = \frac92\arctan u + C = \frac92 \arctan\frac{x - 2}{2} + C$.
\end{itemize}
\textbf{תשובה:}
\[ \boxed{\int \frac{6x^3 - 12x + 24}{x^4 - 4x^3 + 8x^2}\,dx = -\frac{3}{x} + 3\ln(x^2 - 4x + 8) + \frac92\arctan\left(\frac{x - 2}{2}\right) + C.} \]
\end{solution}

\subsection*{שאלה 4ב}

\begin{question}
מצאו את הפתרון הכללי של המשוואה הדיפרנציאלית
\[ \sqrt{1 - x^2} \cdot y' = y^2 \arcsin x \]
\end{question}

\begin{hint}
המשוואה פרידה. אחרי ההפרדה, באגף של $x$ הציבו $t = \arcsin x$.
\end{hint}

\begin{solution}
המשוואה מוגדרת עבור $-1 < x < 1$, ושם $\sqrt{1 - x^2} > 0$. נכתוב אותה בצורה
\[ \frac{dy}{dx} = y^2 \cdot \frac{\arcsin x}{\sqrt{1 - x^2}}, \]
וזו משוואה פרידה.

\textbf{הפתרון הקבוע:} $y \equiv 0$ מקיים את המשוואה ($0 = 0$), ולכן הוא פתרון.

\textbf{עבור $y \ne 0$:} נפריד משתנים:
\[ \int \frac{dy}{y^2} = \int \frac{\arcsin x}{\sqrt{1 - x^2}}\,dx. \]
באגף שמאל מתקבל $-\frac{1}{y}$. באגף ימין נציב $t = \arcsin x$, $dt = \frac{dx}{\sqrt{1 - x^2}}$:
\[ \int \frac{\arcsin x}{\sqrt{1 - x^2}}\,dx = \int t\,dt = \frac{t^2}{2} + C = \frac{\arcsin^2 x}{2} + C. \]
לכן
\[ -\frac{1}{y} = \frac{\arcsin^2 x}{2} + C \quad\Longrightarrow\quad y = \frac{-1}{\frac{\arcsin^2 x}{2} + C} = \frac{-2}{\arcsin^2 x + 2C}. \]

\textbf{תשובה:} הפתרון הכללי הוא
\[ \boxed{y = \frac{-2}{\arcsin^2 x + C}}, \quad C \in \R \]
(בכל קטע שבו המכנה אינו מתאפס), ובנוסף הפתרון $y \equiv 0$.
\end{solution}

\subsection*{שאלה 5א}

\begin{question}
חשבו את השטח החסום בין ציר ה-$x$ ובין גרף הפונקציה
\[ f(x) = xe^x \]
בקטע $[-1, 1]$.
\end{question}

\begin{hint}
שימו לב ש-$f$ מחליפה סימן ב-$x = 0$: השטח הוא $\int_{-1}^1 |f(x)|\,dx$. את הפונקציה הקדומה מצאו באינטגרציה בחלקים.
\end{hint}

\begin{solution}
מכיוון ש-$e^x > 0$, הסימן של $f(x) = xe^x$ הוא הסימן של $x$: $f \le 0$ ב-$[-1, 0]$ ו-$f \ge 0$ ב-$[0, 1]$. לכן השטח הוא
\[ S = \int_{-1}^1 |f(x)|\,dx = -\int_{-1}^0 f(x)\,dx + \int_0^1 f(x)\,dx. \]
\textbf{פונקציה קדומה} באינטגרציה בחלקים, עם $u = x$, $v' = e^x$ (כלומר $u' = 1$, $v = e^x$):
\[ F(x) = \int xe^x\,dx = xe^x - \int e^x\,dx = xe^x - e^x + C = (x - 1)e^x + C. \]
לפי ניוטון-לייבניץ:
\[ S = -\bigl(F(0) - F(-1)\bigr) + \bigl(F(1) - F(0)\bigr) = F(1) - 2F(0) + F(-1). \]
כאן $F(1) = 0$, $F(0) = -1$, $F(-1) = -2e^{-1}$, ולכן
\[ S = 0 - 2\cdot(-1) - \frac{2}{e} = 2 - \frac{2}{e}. \]
\textbf{תשובה:} $\boxed{S = 2\left(1 - \frac1e\right)}$.
\end{solution}

\subsection*{שאלה 5ב}

\begin{question}
בדקו התכנסות או התבדרות של האינטגרל המוכלל הבא:
\[ \int_1^\infty \frac{\cos x}{x^2 - \ln x} \]
הערה: אתם יכולים להניח שהמכנה חיובי בתחום האינטגרציה.
\end{question}

\begin{hint}
האינטגרנד אינו בעל סימן קבוע. בדקו התכנסות בהחלט: חסמו $|f(x)| \le \frac{1}{x^2 - \ln x}$ והשוו (במבחן ההשוואה הגבולי) ל-$\frac{1}{x^2}$.
\end{hint}

\begin{solution}
נסמן $f(x) = \frac{\cos x}{x^2 - \ln x}$. המכנה חיובי ב-$[1, \infty)$ (כפי שמותר להניח; למעשה $\ln x \le x - 1 < x \le x^2$ עבור $x \ge 1$), ולכן $f$ רציפה שם והאינטגרל מוכלל רק בגלל הגבול האינסופי. $f$ מחליפה סימן אינסוף פעמים, לכן נבדוק \textbf{התכנסות בהחלט}.

\textbf{חסימה:} לכל $x \ge 1$,
\[ 0 \le |f(x)| = \frac{|\cos x|}{x^2 - \ln x} \le \frac{1}{x^2 - \ln x} =: h(x). \]

\textbf{התכנסות $\int_1^\infty h$:} נשווה ל-$g(x) = \frac{1}{x^2}$. שתי הפונקציות חיוביות, ו-
\[ \lim_{x\to\infty} \frac{g(x)}{h(x)} = \lim_{x\to\infty} \frac{x^2 - \ln x}{x^2} = \lim_{x\to\infty} \left( 1 - \frac{\ln x}{x^2} \right) = 1, \]
כי $\frac{\ln x}{x^2} \to 0$ (לופיטל: $\frac{1/x}{2x} \to 0$). הגבול סופי וחיובי, ולכן לפי מבחן ההשוואה הגבולי $\int_1^\infty h$ ו-$\int_1^\infty g$ מתכנסים או מתבדרים יחד. מכיוון ש-$\int_1^\infty \frac{dx}{x^\alpha}$ מתכנס עבור $\alpha > 1$, ובפרט $\int_1^\infty \frac{dx}{x^2} = 1$, גם $\int_1^\infty h$ מתכנס.

\textbf{מסקנה:} לפי מבחן ההשוואה (כי $0 \le |f| \le h$), $\int_1^\infty |f|$ מתכנס. אינטגרל מוכלל שמתכנס בהחלט מתכנס, ולכן
\[ \int_1^\infty \frac{\cos x}{x^2 - \ln x}\,dx \ \text{מתכנס (ואף בהחלט).} \]
\end{solution}

\end{document}
